\documentclass[a4paper,11pt]{memoir}

% changemarks without draft mode {{{
% https://tex.stackexchange.com/questions/33559/how-to-use-memoir-change-markers-without-the-draft-option
\makeatletter
\renewcommand{\v@rid}[2]{%
  \@bsphack
  \ifchangemarks
     \marginpar[#1]{#2}%
  \fi
  \@esphack}
\makeatother
%\changemarks
% }}}

\usepackage[utf8]{inputenc}
\usepackage[T1]{fontenc}
\usepackage[svgnames]{xcolor}
\usepackage{libertine}
\usepackage[frenchb]{babel}
% TODO \usepackage[french]{babelbib}
\usepackage{mathpartir}
\usepackage{syntax}
\usepackage{amssymb}
\usepackage{tikz}
\usepackage{amsmath}
\usepackage{float}
\usepackage{url}
\usepackage{xspace}
\usepackage{graphicx}
\usepackage[shortlabels]{enumitem}
\usepackage[backref=page]{hyperref}
\usepackage{links}
\usepackage[a4paper]{geometry}
\usepackage{amsthm}
\usepackage{stmaryrd}
\usepackage[normalem]{ulem}
\usepackage{wasysym}
\usepackage{multirow}
\usepackage{thmtools}
\usepackage{gramster}
\usepackage{switch}
\usepackage{pelagiccod}
\usepackage{conteq}

\newcommand{\irule}[3]{ \inferrule*[right=(#1)]{#2}{#3} }
\def\iaxiom#1#2{\irule{#1}{ }{#2}}

% http://tex.stackexchange.com/questions/3964/mathematical-symbol-for-disjoint-set-union
\makeatletter
\def\moverlay{\mathpalette\mov@rlay}
\def\mov@rlay#1#2{\leavevmode\vtop{%
   \baselineskip\z@skip \lineskiplimit-\maxdimen
   \ialign{\hfil$\m@th#1##$\hfil\cr#2\crcr}}}
\newcommand{\charfusion}[3][\mathord]{
    #1{\ifx#1\mathop\vphantom{#2}\fi
        \mathpalette\mov@rlay{#2\cr#3}
      }
    \ifx#1\mathop\expandafter\displaylimits\fi}
\makeatother

\newcommand{\dcup}{\charfusion[\mathbin]{\cup}{+}}
\newcommand{\bigdcup}{\charfusion[\mathbin]{\bigcup}{+}}

\def\opbin{\boxplus\xspace}
\def\opun{\boxminus\xspace}

% Types
\newcommand{\tInt}{\textsc{Int}\xspace}
\newcommand{\tFloat}{\textsc{Float}\xspace}
\newcommand{\tUnit}{\textsc{Unit}\xspace}

\def\structRec#1#2{\textsc{rec}(\&#1) #2} % TODO

\newcommand{\tFV}[1]{\textrm{FreeVars}(#1)}

\def\opxor{\mathbin{\char`\^}}

% Constant
\newcommand{\cNil}{\textsc{Nil}}

% Program & language names
\newcommand{\ctonewspeak}{\textsc{c2newspeak}}
\newcommand{\cpp}{\textsc{cpp}}
\newcommand{\newspeak}{\textsc{Newspeak}}

% Newspeak
\newcommand{\npkIf}[3]{\textsc{If}(#1) \{ #2 \} \textsc{Else} \{ #3 \}}
\newcommand{\npkDoWith}[2]{\textsc{Do}\{ #1 \}\textsc{With}~#2:}
\newcommand{\npkWhile}[1]{\textsc{While}(1)\{#1\}}
\newcommand{\npkGoto}[1]{\textsc{Goto}~#1}
\newcommand{\npkReturn}[1]{\textsc{Return}~#1}

% cML
\def\iPass{\textsc{Pass}\xspace}
\def\iIf#1#2#3{\npkIf{#1}{#2}{#3}}
\def\iWhile#1#2{\textsc{While}(#1)\{#2\}}
\def\iReturn#1{\textsc{Return}(#1)}
\def\iDecl#1#2#3{\textsc{Decl}~#1=#2~\textsc{in}\{#3\}}

\def\iThen#1#2{\textsc{If}(#1) \{ #2 \}}

\def\eNull{\textsc{Null}\xspace}
\def\eUnit{(\,)\xspace}

\def\vRet{\underline{R}}

% sets pour semantique
\def\sId{\textsc{Id}\xspace}
\def\sMem{\textsc{Mem}\xspace}
\def\sAddr{\textsc{Addr}\xspace}
\def\sInt{\textsc{Int}\xspace}
\def\sFloat{\textsc{Float}\xspace}
\def\sPhi{\Phi}

% Memory state (instructions)
% #1 : memory
% #2 : instruction
\def\msi#1#2{ \langle #2, #1 \rangle }

% Reduction (expressions)
\def\mm#1#2#3#4{ \msi{#1}{#2} → \msi{#3}{#4} }
\def\mms#1#2{\mm{m}{#1}{m}{#2}}

% Reduction (instructions)
\def\mmi#1#2#3#4{ \mm{#1}{#2}{#3}{#4} }

% Reduction (left-values)
\def\mmlv#1#2#3#4{ \mm{#1}{#2}{#3}{#4} }

% Typage sém
\def\semtyp#1#2#3{#1 ⊧ #2 : #3}

% Compatibilité Γ / m
\def\mcomp#1#2{#1 ⊧ #2}

% Reduction contexts
\def\ctxOp#1#2{#1~\opbin~#2}
\def\ctxEmpty{\bullet\xspace}
\def\ctxUnOp#1{\opun~#1}
\def\ctxLvDeref#1{*#1}
\def\ctxLvField#1#2{#1.#2}
\def\ctxLvIndex#1#2{#1[#2]}
\def\ctxSet#1#2{#1 ← #2}

% Context substitution
\def\ctxSub#1#2{#1\llparenthesis#2\rrparenthesis}
\def\ctxSubL#1#2{\ctxSub{#1}{#2}_L}
\def\ctxSubI#1#2{\ctxSub{#1}{#2}_I}

% Qualifs
\def\qKernel{\textsc{Kernel}\xspace}
\def\qUser{\textsc{User}\xspace}
\def\vTainted#1{\widehat{\diamondsuit}~#1}
\def\eTaint#1{\diamondsuit~#1}

\def\ptrK#1{#1~*}
\def\ptrU#1{#1~@}

% Eval de phrase
\def\ph#1#2#3{#1 ⊩ #2 → #3}

\def\phstar#1#2#3{#1 ⊩ #2 →^* #3}

\def\phprog#1#2{⊩ #1 →^* #2}

% typage
\def\ty#1#2#3{#1 ⊢ #2 : #3}
\def\tyi#1#2{#1 ⊢ #2}
\def\typh#1#2#3{#1 ⊢ #2 → #3}
\def\typhstar#1#2#3{#1 ⊢ #2 →^* #3}

\def\sVal{\textsc{Val}}
\def\sByte{\textsc{Byte}}

\def\powerset#1{\mathcal P \left({#1}\right)}

% Functions with placeholders
\def\phx#1{\ensuremath{#1{\cdot}}}
\def\phxx#1{\ensuremath{#1{\cdot}{\cdot}}}

% Progs bien formés
\def\WF#1{\textsc{wf}(#1)}
\def\NORET#1{\textsc{NoRet}(#1)}

\newcommand{\insertcode}[1]{\input{gen/#1.pyg.tex}}

% \chapter* + add to toc
%\newcommand{\chapterstar}[1]{\chapter*{#1} \addcontentsline{toc}{chapter}{#1}}

% Register names
\def\defreg#1{\expandafter\def\csname#1\endcsname{\textsc{#1}\xspace}}

\defreg{eax}
\defreg{ebx}
\defreg{ecx}
\defreg{edx}
\defreg{esi}
\defreg{edi}
\defreg{ebp}
\defreg{esp}
\defreg{eip}
\defreg{eflags}

\def\crtrois{\textsc{CR3}}

\newcommand{\ring}[1]{\emph{ring}~#1}

\newcommand{\asminstr}[1]{\texttt{#1}}

\newcommand{\wip}[0]{%
\begin{center}
    \includegraphics[width=5cm]{gen/travaux.pdf}
\end{center}
%\textbf{(wip)}
}

\def\listcons{\texttt{(::)}\xspace}

\tikzstyle{bigbrace} = [decorate,decoration={brace,amplitude=3mm}]

% Debug tikz bounding box
\def\drawbb{\draw[thick] (current bounding box.south west) rectangle (current bounding box.north east);}

% SubFloat + verbatim
\makeatletter
\newbox\sf@box
\newenvironment{SubFloat}[2][]%
{\def\sf@one{#1}%
\def\sf@two{#2}%
\setbox\sf@box\hbox
\bgroup}%
{ \egroup
\ifx\@empty\sf@two\@empty\relax
\def\sf@two{\@empty}
\fi
\ifx\@empty\sf@one\@empty\relax
\subfloat[\sf@two]{\box\sf@box}%
\else
\subfloat[\sf@one][\sf@two]{\box\sf@box}%
\fi}
\makeatother

% comes from the memoir manual {{{

\newcommand*{\addtotextheight}[1]{%
\twocolumn
\addtolength{\textheight}{#1}%
\onecolumn}

%% for main body, bottom of text at 1in, footer below
%% top of header at 1in, first text line double spaced
%% below base of header
\newlength{\linespace}
\setlength{\linespace}{\baselineskip} % current equivalent of \onelineskip
\setlength{\headheight}{\onelineskip}
\setlength{\headsep}{\linespace}
\addtolength{\headsep}{-\topskip}
\setlength{\uppermargin}{1in}
\addtolength{\uppermargin}{\headheight}
\addtolength{\uppermargin}{\headsep}

%% and for the bottom
\setlength{\lowermargin}{1in}
\setlength{\textheight}{\paperheight}

\addtolength{\textheight}{-\uppermargin}
\addtolength{\textheight}{-\lowermargin}
%% footnote settings
\setlength{\footskip}{\onelineskip}
\setlength{\footnotesep}{\onelineskip}


%% the fiddle lengths (..ta.. for title/approval page, others for prelims)
\newlength{\toptafiddle} \setlength{\toptafiddle}{2\linespace}
\newlength{\bottafiddle} \setlength{\bottafiddle}{0pt}
\newlength{\topfiddle}
\setlength{\topfiddle}{\toptafiddle}
\newlength{\botfiddle}
\setlength{\botfiddle}{\onelineskip}

%%% general macro for Abstract, etc., headings
\newcommand*{\pretoctitle}[1]{{\clearpage\centering \vspace*{-\topfiddle}#1\par}}

%%% Start the ACKNOWLEDGEMENTS
\newcommand{\asuacknowledgements}{\pretoctitle{REMERCIEMENTS}}

%%% Start the ABSTRACT
\newcommand{\asuabstract}{ \addtotextheight{-\botfiddle} \pretoctitle{ABSTRACT}}

\addtodef{\mainmatter}{\addtotextheight{\botfiddle}{}}

%% make it easy to center any dedication
\newcommand{\asudedication}[1]{%
{\clearpage\mbox{}\vfill\centering #1 \par\vfill\clearpage}}

% }}} /memoir

% {{{ subtitle command
\newcommand\subtitle[1]{%
  \newcommand\thesubtitle{#1}}
% }}}

% Language name
\def\langname{\textsc{Safespeak}\xspace}

% Center vertically + remove header & footer
\newenvironment{headingpage}{
\thispagestyle{simple}
\vspace*{\fill}
}{
\vspace*{\fill}
}

% Hoare triple
\def\hoare#1#2#3{\ensuremath{\{#1\}~#2~\{#3\}}}

\def\eqdef{\overset{\mathrm{def}}{=\joinrel=}\xspace}

% Lens
\def\mkLens#1#2{\langle #1 | #2 \rangle}
\def\setLens#1#2{\textsc{Lens}_{#1, #2}}
\def\lensGetX#1{\ensuremath{\mathrm{get}_{#1}}}
\def\lensPutX#1{\ensuremath{\mathrm{put}_{#1}}}
\def\lensGet#1#2{\lensGetX{#1}(#2)}
\def\lensPut#1#2#3{\lensPutX{#1}(#2, #3)}
\def\lensNodeBig#1#2{
\vcenter{\hbox{
\begin{tikzpicture}[scale=0.7]
\draw[fill=#1] (-0.5,-0.5) rectangle ++ (1,1);
\draw[fill=#2] (0,0) circle (0.3);
\end{tikzpicture}
}}
}
\def\lensInner#1{
\vcenter{\hbox{
\begin{tikzpicture}[scale=0.7]
\draw[fill=#1] (0, 0) circle (0.3);
\end{tikzpicture}
}}
}
\def\lensNodeBigIdx#1#2#3#4#5{
\vcenter{\hbox{
\begin{tikzpicture}[scale=0.7]
  \draw[fill=#1] (-1.5,-1.5) rectangle ++ (3,3);
  \node[fill=#2, draw, star, minimum size=6mm] at (-0.75, 0.75) {};
  \node[fill=#3, draw, star, minimum size=6mm] at ( 0.75, 0.75) {};
  \node[fill=#4, draw, circle, minimum size=6mm] at (-0.75,-0.75) {};
  \node[fill=#5, draw, regular polygon, regular polygon sides=5, minimum size=6mm] at ( 0.75,-0.75) {};
  \node at (-2,  0.75) {a};
  \node at ( 2,  0.75) {b};
  \node at (-2, -0.75) {c};
  \node at ( 2, -0.75) {d};
\end{tikzpicture}
}}
}
\def\lensInnerStar#1{
\vcenter{\hbox{
\begin{tikzpicture}[scale=0.7]
\node[draw, fill=#1, star] {};
\end{tikzpicture}
}}
}

% Erreurs

\def\serr#1{
  \switch{#1}
    {taint}{Ω_{taint}}
    {array}{Ω_{array}}
    {ptr}{Ω_{ptr}}
    {div}{Ω_{div}}
    {field}{Ω_{field}}
    {var}{Ω_{var}}
    {END}
    % important : étendre la description dans le chapitre 4 et la figure de
    % liste
}

%\def\uGet#1#2{#1 ⇐_U #2}
%\def\uPut#1#2{#1 ⇒_U #2}
\def\cfuF{\textrm{copy\_from\_user}\xspace}
\def\ctuF{\textrm{copy\_to\_user}\xspace}
\def\uGet#1#2{\cfuF(#1, #2)}
\def\uPut#1#2{\ctuF(#1, #2)}

% Litteral notation

\def\eStruct#1{\{#1\}}
\def\eArray#1{[#1]}
\def\tStruct#1{\eStruct{#1}}
\def\eFun#1#2{\mathrm{fun} (#1) \{ #2 \}}

\def\gramorx{\hspace{3mm}\gramor\xspace}
\def\stFun#1{\textsc{Fun}_n}
\def\tComp#1#2{\mathrm{Comp}(#1, #2)}

% Box on top on rules fig
\def\ruleheader#1{
  \begin{minipage}{\linewidth}
    \raggedleft
    \fbox{#1}
  \end{minipage}
}

\def\termFst#1{\ensuremath{\textrm{fst}~#1}}
\def\termSnd#1{\ensuremath{\textrm{snd}~#1}}

\def\ptrtype{\texttt{ptrtype}\xspace}

\def\addrInt{\textrm{addr}^\natural \xspace}
\def\interpPtr#1{\textrm{interp}_{(#1)}}

\tikzset{>=latex}

\def\partccl#1{
  \clearpage~\thispagestyle{cleared}\cleartorecto
  \chapter*{Conclusion de la partie~\ref{#1}}
  \addcontentsline{toc}{chapter}{Conclusion de la partie~\ref{#1}}
}

\def\hexa#1{0x#1}
\def\floatop#1{#1.}

\def\cLookup#1#2{\mathrm{Lookup} (#1, #2)}
\def\cExtend#1#2#3{\mathrm{Extend}(#1, #2 ↦ #3)}
\def\cCleanup#1{\mathrm{Cleanup}(#1)}

\newcommand{\motif}[1]{\cleaders\hbox{#1}\hfil}
\newcommand{\chain}[2]{\leavevmode\hbox to #2{\motif{#1}}}
\newcommand{\diamonds}{$\mkern-.6mu \diamond \mkern-.6mu$}

\def\jolibreak{
\fancybreak{\chain{\diamonds}{0.25\textwidth}}
}

\def\cDom#1{\mathrm{dom}(#1)}

\def\cVarname#1{\mathrm{Varname}(#1)}

\def\cVisible#1{\mathrm{Visible}(#1)}

\def\cRepr#1{\mathrm{Repr}(#1)}

\deflink{cmm}{The C~-~- language}{http://www.cminusminus.org/}
\deflink{ocaml}{OCaml -- Home}{http://ocaml.org/}
\deflink{penjili}{Penjili project}{http://www.penjili.org/}
\deflink{python}{Python Programming Language -- Official Website}{http://www.python.org/}
\deflink{rust}{The Rust Programming Language}{http://www.rust-lang.org/}
\deflink{sparse}{Sparse - a Semantic Parser for {C}}{https://sparse.wiki.kernel.org/index.php/Main_Page}

\hyphenation{con-sis-te}
\hyphenation{con-trô-lé}


\usepackage{color,calc,graphicx,soul,fourier}
\definecolor{nicered}{rgb}{.647,.129,.149}
\makeatletter
\newlength\dlf@normtxtw
\setlength\dlf@normtxtw{\textwidth}
\def\myhelvetfont{\def\sfdefault{mdput}}
\newsavebox{\feline@chapter}
\newcommand\feline@chapter@marker[1][4cm]{%
\sbox\feline@chapter{%
\resizebox{!}{#1}{\fboxsep=1pt%
\colorbox{nicered}{\color{white}\bfseries\sffamily\thechapter}%
}}%
\rotatebox{90}{%
\resizebox{%
\heightof{\usebox{\feline@chapter}}+\depthof{\usebox{\feline@chapter}}}%
{!}{\scshape\so\@chapapp}}\quad%
\raisebox{\depthof{\usebox{\feline@chapter}}}{\usebox{\feline@chapter}}%
}
\newcommand\feline@chm[1][4cm]{%
\sbox\feline@chapter{\feline@chapter@marker[#1]}%
\makebox[0pt][l]{% aka \rlap
\makebox[1cm][r]{\usebox\feline@chapter}%
}}
\makechapterstyle{daleif1}{
\renewcommand\chapnamefont{\normalfont\Large\scshape\raggedleft\so}
\renewcommand\chaptitlefont{\normalfont\huge\bfseries\scshape\color{nicered}}
\renewcommand\chapternamenum{}
\renewcommand\printchaptername{}
\renewcommand\printchapternum{\null\hfill\feline@chm[2.5cm]\par}
\renewcommand\afterchapternum{\par\vskip\midchapskip}
\renewcommand\printchaptertitle[1]{\chaptitlefont\raggedleft ##1\par}
}
\makeatother
\chapterstyle{daleif1}

\setlength\beforechapskip{0pt}


\def\opbinintset{\{+,-,\times,/,\&,|,\opxor,\&\&,||,\ll,\gg,≤,≥,<,> \}}
\def\opbinfloatset{\{+_.,-_.,\times_.,/_.,≤_.,≥_.,<_.,>_.\}}

\def\disprule#1{
  \switch{#1}
    {Cst-Int}{ \irule{Cst-Int}
                 { }
                 { Γ ⊢ n : \tInt}
    }
    {Cst-Float}{ \irule{Cst-Float}
                   { }
                   { Γ ⊢ d : \tFloat}
    }
    {Cst-Null}{ \irule{Cst-Null}
                  { }
                  { Γ ⊢ \eNull : t*}
    }
    {Cst-Unit}{ \irule{Cst-Unit}
                  { }
                  { Γ ⊢ \eUnit : \tUnit }
    }
    {Lv-Var}{ \irule{Lv-Var}
                { x : t ∈ Γ }
                { Γ ⊢ x : t }
    }
    {Lv-Deref}{ \irule{Lv-Deref}
                  { Γ ⊢ e : t* }
                  { Γ ⊢ *e : t }
    }
    {Lv-Index}{ \irule{Lv-Index}
                  { Γ ⊢ e : \tInt \\
                    Γ ⊢ lv : t[~]
                  }
                  { Γ ⊢ lv[e] : t }
    }
    {Lv-Field}{ \irule{Lv-Field}
                  { (l, t) ∈ S \\
                    Γ ⊢ lv : S
                  }
                  { Γ ⊢ lv.l_S : t }
    }
    {Op-Int}{ \irule{Op-Int}
                { \opbin ∈ \opbinintset{} \\
                  Γ ⊢ e_1 : \tInt \\
                  Γ ⊢ e_2 : \tInt
                }{
                  Γ ⊢ e_1~\opbin~e_2 : \tInt
                }
    }
    {Op-Float}{ \irule{Op-Float}
                { \opbin ∈ \opbinfloatset{} \\
                  Γ ⊢ e_1 : \tFloat \\
                  Γ ⊢ e_2 : \tFloat
                }{
                  Γ ⊢ e_1~\opbin~e_2 : \tFloat
                }
    }
    {Eq-Num}{ \irule{Eq-Num}
                {t ∈ \{ \tInt, \tFloat \}}
                {\textsc{Eq}(t)}
    }
    {Eq-Ptr}{ \irule{Eq-Ptr}
                { }
                {\textsc{Eq}(t*)}
    }
    {Eq-Array}{ \irule{Eq-Array}
                  {\textsc{Eq}(t)}
                  {\textsc{Eq}(t[~])}
    }
    {Eq-Struct}{ \irule{Eq-Struct}
                  {∀ i ∈ [1;n]. \textsc{Eq}(t_i)}
                  {\textsc{Eq}(\tStruct{l_1 : t_1 ; … l_n : t_n})}
    }
    {Op-Eq}{ \irule{Op-Eq}
             { \opbin ∈ \{=,≠\} \\
               Γ ⊢ e_1 : t \\
               Γ ⊢ e_2 : t \\
               \textsc{Eq}(t)
             }{
               Γ ⊢ e_1~\opbin~e_2 : \tInt
             }
    }
    {Unop-Plus-Int}{ \irule{Unop-Plus-Int}
                       { Γ ⊢ e : \tInt }
                       { Γ ⊢ + e : \tInt }
    }
    {Unop-Plus-Float}{ \irule{Unop-Plus-Float}
                         { Γ ⊢ e : \tFloat }
                         { Γ ⊢ +. e : \tFloat }
    }
    {Unop-Minus-Int}{ \irule{Unop-Minus-Int}
                        { Γ ⊢ e : \tInt }
                        { Γ ⊢ - e : \tInt }
    }
    {Unop-Minus-Float}{ \irule{Unop-Minus-Float}
                          { Γ ⊢ e : \tFloat }
                          { Γ ⊢ -. e : \tFloat }
    }
    {Unop-Not}{ \irule{Unop-Not}
                { \opun ∈ \{ \sim, ! \} \\
                  Γ ⊢ e : \tInt
                }{
                  Γ ⊢ \opun~e : \tInt
                }
              }
    {Ptr-Arith}{ \irule{Ptr-Arith}
                   { \opbin ∈ \{ +_p, -_p \} \\
                     Γ ⊢ e_1 : t* \\
                     Γ ⊢ e_2 : \tInt \\
                   }
                   { Γ ⊢ e_1~\opbin~e_2 : t* }
               }
    {Addr}{ \irule{Addr}
              { Γ ⊢ lv : t }
              { Γ ⊢ \&lv : t* }
          }
    {Set}{ \irule{Set}
             { Γ ⊢ lv : t \\
               Γ ⊢ e : t
             }
             { Γ ⊢ lv ← e : t }
         }
    {Array}{ \irule{Array}
               { ∀ i ∈ [1;n] , Γ ⊢ e_i : t }
               { Γ ⊢ \eArray{ e_1 ; … ; e_n }
                       : t[~]
               }
           }
    {Struct}{ \irule{Struct}
                { ∀ i ∈ [1;n], Γ ⊢ e_i : t_i
                }
                { Γ ⊢ \eStruct{ l_1 : e_1 ; … ; l_n : e_n }
                    : \tStruct{ l_1 : t_1 ; … ; l_n : t_n }
                }
            }
    {Call}{ \irule{Call}
              { Γ ⊢ e : ( t_1 ,
                         … , t_n ) → t \\
                ∀ i ∈ [1;n], Γ ⊢ e_i : t_i
              }
              { Γ ⊢ e ( e_1
                         , …
                         , e_n
                         )
                      : t
              }
          }
    {Pass}{ \irule{Pass}
              { }
              {Γ ⊢ \iPass}
          }
    {Seq}{ \irule{Seq}
             { Γ ⊢ i_1 \\
               Γ ⊢ i_2
             }{
               Γ ⊢ i_1;i_2
             }
         }
    {Exp}{ \irule{Exp}
             { Γ ⊢ e : t }
             { Γ ⊢ e }
         }
    {Decl}{ \irule{Decl}
              { Γ ⊢ e : t
             \\ Γ, x:t ⊢ i
              }
              { Γ ⊢ \iDecl{x}{e}{i} }
          }
    {If}{ \irule{If}
            { Γ ⊢ e : \tInt \\
              Γ ⊢ i_1 \\
              Γ ⊢ i_2
            }
            { Γ ⊢ \iIf{e}{i_1}{i_2} }
        }
    {While}{ \irule{While}
               { Γ ⊢ e : \tInt \\
                 Γ ⊢ i
               }
               { Γ ⊢ \iWhile{e}{i} }
           }
    {Return}{ \irule{Return}
                { Γ ⊢ \vRet ← e }
                { Γ ⊢ \iReturn{e} }
            }
    {Fun}{ \irule{Fun}
             { Γ' = (Γ - \vRet),
                     a_1 : t_1,
                     …,
                     a_n : t_n,
                     \vRet : t \\
               Γ' ⊢ i
             }
             { Γ ⊢ \eFun{a_1, …, a_n}{i}
                    : (t_1, …, t_n) → t
             }
         }
    {T-Exp}{ \irule{T-Exp}
                { Γ ⊢ e : t }
                { \typh{Γ}{e}{Γ} }
            }
    {T-Var}{ \irule{T-Var}
                { Γ ⊢ e : t \\
                  Γ' = (x, t), Γ
                }
                { \typh{Γ}{x=e}{Γ'} }
            }
    {S-Int}{ \irule{S-Int}
               { }
               {\semtyp{m}{n}{\tInt}}
           }
    {S-Float}{ \irule{S-Float}
                 { }
                 {\semtyp{m}{d}{\tFloat}}
    }
    {S-Unit}{ \irule{S-Unit}
                { }
                {\semtyp{m}{\eUnit}{\tUnit}}
    }
    {S-Null}{ \irule{S-Null}
                { }
                {\semtyp{m}{\eNull}{τ~*}}
    }
    {S-Ptr}{ \irule{S-Ptr}
               {\semtypphi{m}{φ}{τ}}
               {\semtyp{m}{\widehat{\&}~φ}{τ~*}}
    }
    {S-Array}{ \irule{S-Array}
                 { ∀i ∈ [1;n]. \semtyp{m}{v_i}{τ}}
                 { \semtyp{m}{\widehat{ \eArray{ v_1;…;v_n } }}{τ[~]} }
    }
    {S-Struct}{ \irule{S-Struct}
                  { ∀i ∈ [1;n]. \semtyp{m}{v_i}{τ_i} }
                  { \semtyp{m}{ \widehat{ \eStruct{ l_1:v_1;…;l_n:v_n } } }
                                        { \tStruct{ l_1:τ_1;…;l_n:τ_n } }
                  }
    }
    {S-Fun}{ \irule{S-Fun}
               { }
               {\semtyp{m}{\eFun{x_1, …, x_n}{i}}{\stFun{n}}}
    }
    {Taint}{ \irule{Taint}
                 { Γ ⊢ lv : t }
                 { Γ ⊢ \eTaint{lv} : \ptrU{t} }
    }
    {GetU}{ \irule{GetU}
              { Γ ⊢ e_d : \ptrK{t}
             \\ Γ ⊢ e_s : \ptrU{t}
              }
              { Γ ⊢ \uGet{e_d}{e_s} : \tInt }
    }
    {PutU}{ \irule{PutU}
              { Γ ⊢ e_d : \ptrU{t}
             \\ Γ ⊢ e_s : \ptrK{t}
              }
              { Γ ⊢ \uPut{e_d}{e_s} : \tInt }
    }
    {AddrOf-Kernel}{ \irule{AddrOf-Kernel}
                       { Γ ⊢ lv : t }
                       { Γ ⊢ \& lv : \ptrK{t} }
    }
    {Ptr-Arith-User}{ \irule{Ptr-Arith-User}
                         { \opbin ∈ \{ +_p, -_p \} \\
                           Γ ⊢ e_1 : \ptrU{t} \\
                           Γ ⊢ e_2 : \tInt \\
                         }
                         { Γ ⊢ e_1~\opbin~e_2 : \ptrU{t} }
    }
    {END}
}

\def\semrule#1{
  \switch{#1}
    {Ctx}{ \irule{Ctx}
             { \mm{m}{i}{m'}{i'} }
             { \mm{m}{\ctxSub{C}{i}}{m'}{\ctxSub{C}{i'}} }
    }
    {Phi-Var}{ \irule{Phi-Var}
                 { a = \cLookup{x}{m} }
                 {
                   \mmlv{m}{x}{m}{a}
                 }
    }
    {Phi-Deref}{ \irule{Phi-Deref}
                   { v = \widehat{\&}~φ }
                   { \mms{*~v }{φ} }
    }
    {Phi-Struct}{ \irule{Phi-Struct}
                    {  }
                    { \mms{φ.l_S}{φ\widehat{.}l} }
    }
    {Phi-Array}{ \irule{Phi-Array}
                   { }
                   { \mms{φ[n]}{φ\widehat{[n]}} }
    }
    {Exp-Cst}{ \irule{Exp-Cst}
                 { }
                 {\mms{c}{\widehat{c}}}
             }
    {Exp-Fun}{ \irule{Exp-Fun}
                 { }
                 {\mms{f}{\widehat{f}}}
             }
    {Exp-Lv}{ \irule{Exp-Lv}
              { }
              {\mm{m}{φ}{m}{m[φ]_Φ}}
    }
    {Exp-UnOp}{ \irule{Exp-UnOp}
                { }
                {\mms{\opun~v}{\widehat{\opun}~v}}
    }
    {Exp-BinOp}{ \irule{Exp-BinOp}
                   { }
                   {\mms{v_1~\opbin~v_2}{v_1~\widehat{\opbin}~v_2}}
    }
    {Exp-AddrOf}{ \irule{Exp-AddrOf}
                    { }
                    {\mm{m}{\&~φ}{m}{\widehat{\&}~φ}}
    }
    {Exp-Set}{ \irule{Exp-Set}
                 { }
                 {\mm{m}{φ ← v}{m[φ ← v]_Φ}{v}}
    }
    {Exp-Struct}{ \irule{Exp-Struct}
                   { }
                   { \mms{          \eStruct{ l_1 : v_1 ; … ; l_n : v_n } }
                         { \widehat{\eStruct{ l_1 : v_1 ; … ; l_n : v_n } }}
                   }
    }
    {Exp-Array}{ \irule{Exp-Array}
                 { }
                 {
                   \mms{      \eArray{ v_1, … , v_n } }
                   { \widehat{\eArray{ v_1, … , v_n } }}
                 }
    }
    {Exp-Call-Return}{ \irule{Exp-Call-Return}
                          { m' = \cCleanup{m}
                         \\ v' = \cV{|m|}{v}
                          }
                          { \mm{m}{\eFun{a_1, …, a_n}{\iReturn{v}} (v_1, …, v_n)}
                               {m'}{v'}
                          }
                     }
    {Exp-Call-Ctx}{ \irule{Exp-Call-Ctx}
                        { m_1 = \cPush{m_0}{(a_1 ↦ v_1, … a_n ↦ v_n)}
                       \\ \mm{m_1}{i}{m_2}{i'}
                       \\ ∀ i ∈ [1;n], v'_i = m_2[(|m_2|, a_i)]_A
                       \\ m_3 = \cPop{m_2}
                        }
                        { \mm{m_0}{\eFun{a_1, …, a_n}{i} (v_1, …, v_n) }
                          { m_3 }{\eFun{a_1, …, a_n}{i'} (v'_1, …, v'_n) }
                        }
                  }
    {Exp-Call-Err}{ \irule{Exp-Call-Err}
                     { m' = \cPush{m}{(a_1 ↦ v_1, …, a_n ↦ v_n)}
                    \\ \msi{m'}{i} → Ω
                     }
                     { \msi{m}{\eFun{a_1, …, a_n}{i} (v_1, …, v_n) } → Ω }
                  }
    {Seq}{ \irule{Seq}
              { }
              { \mmi{m}{(\iPass;i)}{m}{i} }
    }
    {Exp}{ \irule{Exp}
             { }
             { \mmi{m}{v}{m}{\iPass} }
    }
    {Decl}{ \irule{Decl}
            { m' = \cExtend{m}{x}{v}
           \\ \mmi{m'}{i}{m''}{\iPass}
           \\ m''' = \cCleanVar{m''}{(|m|, x)}
            }
            {\mmi{m}{\iDecl{x}{v}{i}}{m'''}{\iPass}}
    }
    {Decl-Pass}{ \irule{Decl-Pass}
                     { m' = \cCleanVar{m - x}{(|m|, x)} }
                     { \mm{m}{\iDecl{x}{v}{\iPass}}
                         {m'}{\iPass}
                     }
    }
    {Decl-Return}{ \irule{Decl-Return}
                        { m' = \cCleanVar{m - x}{(|m|, x)}
                       \\ v'' = \cVV{v'}{(|m|, x)}
                        }
                        { \mm{m}{ \iDecl{x}{v}{\iReturn{v'}} }
                            {m'}{\iReturn{v''}}
                        }
    }
    {Decl-Ctx}{ \irule{Decl-Ctx}
                   { m' = \cExtend{m}{x}{v}
                   \\ \mm{m'}{i}{m''}{i'}
                   \\ v' = m''[(|m''|, x)]_A
                   \\ m''' = m'' - x
                   }
                   { \mm{m}{\iDecl{x}{v}{i}}
                       {m'''}{\iDecl{x}{v'}{i'}}
                   }
    }
    {Decl-Err}{ \irule{Decl-Err}
                    { \msi{m}{i} → Ω }
                    { \msi{m}{\iDecl{x}{v}{i}} → Ω }
    }
    {If-False}{ \irule{If-False}
                  { }
                  { \mmi{m}{\iIf{0}{i_t}{i_f}}{m}{i_f} }
    }
    {If-True}{ \irule{If-True}
                 { v ≠ 0 }
                 { \mmi{m}{\iIf{v}{i_t}{i_f}}{m}{i_t} }
    }
    {While}{ \irule{While}
               { }
               { \mmi{m}{\iWhile{e}{i}}{m}{\iThen{e}{i;\iWhile{e}{i}}}
               }
    }
    {Return}{ \irule{Return}
              { }
              { \mmi{m}{\iReturn{v};i}{m}{\iReturn{v}} }
    }
    {T-Exp}{ \irule{T-Exp}
                {\mm{m}{e}{m'}{v}}
                {\ph{m}{e}{m'}}
            }
    {T-Var}{ \irule{T-Var}
                {\mm{m}{e}{m'}{v}}
                {\ph{(s,g)}{x = e}{(s, (x ↦ v)::g)}}
    }
    {T*-Nil}{ \irule{T*-Nil}
                 { }
                 {\phstar{m}{[~]}{m}}
    }
    {T*-Cons}{ \irule{T*-Cons}
                  {\ph{m}{p}{m'} \\
                   \phstar{m'}{ps}{m''}}
                  {\phstar{m}{p::ps}{m''}}
    }
    {Exp-Err}{ \irule{Exp-Err}
                 { }
                 { \msi{m}{Ω} → Ω }
    }
    {Eval-Err}{ \irule{Eval-Err}
                  { \msi{m}{e} → Ω }
                  { \msi{m}{\ctxSub{C}{e}} → Ω }
% TODO préciser que ça s'étend à i dans progres
% return
    }
    {Prog}{ \irule{Prog}
              {\phstar{([~], [~])}{P}{m}}
              {\phprog{P}{m}}
    }
    {Expr-Tainted}{ \irule{Expr-Tainted}
                      { }
                      { \mm{m}{\eTaint{φ}}{m}{\vTainted{φ}} }
    }
    {User-Get-OK}{ \irule{User-Get-OK}
                     { \widehat{\&}~φ = m[φ_d]_Φ
                    \\ v  = m[φ_s]_Φ
                    \\ m' = m[φ ← v]_Φ
                     }
                     { \mm{m}{\uGet{φ_d}{\vTainted{φ_s}}}
                       {m'}{0}
                     }
    }
    {User-Get-Err}{ \irule{User-Get-Err}
                      {
                        ∄~φ_s. φ = \vTainted{φ_s}
                      }
                      { \mm{m}{\uGet{φ_d}{φ}}
                           {m}{-1}
                      }
    }
    {User-Put-OK}{ \irule{User-Put-OK}
                     { \widehat{\&}~φ = m[φ_d]_Φ
                    \\ v = m[φ_s]_Φ
                    \\ m' = m[φ ← v]_Φ
                     }
                     { \mm{m}{ \uPut{\vTainted{φ_d}}{φ_s} }
                         {m'}{0}
                     }
    }
    {User-Put-Err}{ \irule{User-Put-Err}
                      {
                        ∄~φ_d. φ = \vTainted{φ_d}
                      }
                      { \mm{m}{\uPut{φ}{φ_s}}
                           {m}{-1}
                      }
    }
    {END}
}

\def\figctx{
  \begin{align*}
  \gramdef{Contextes}{C}
    { \ctxEmpty }{}
    { \ctxOp{C}{e} }{}
    { \ctxOp{v}{C} }{}
    { \ctxUnOp{C} }{}
    { \&~C }{}
    { \ctxSet{C}{e} }{}
    { \ctxSet{φ}{C} }{}
    { \eStruct{ l_1:v_1 ; … ; l_i:C ; … ; l_n:e_n } }{}
    { \eArray{ v_1 ; … ; C ; … ; e_n } }{}
    { C (e_1, …, e_n) }{}
    { f (v_1, …, C, …, e_n) }{}
    { \ctxLvField{C}{l_S} }{}
    { \ctxLvIndex{C}{e} }{}
    { \ctxLvIndex{φ}{C} }{}
    { *~C }{}
    { C;i }{}
    { \iIf{C}{i_1}{i_2} }{}
    { \iReturn{C} }{}
    { \iDecl{x}{C}{i} }{}
    {END}
  \end{align*}
}

\def\figstxdata{
  \begin{align*}
  \gramdef{Constantes}{c}
                 { n      }{ Entier }
                 { d      }{ Flottant }
                 { \eNull }{ Pointeur nul }
                 { \eUnit }{ Valeur unité }
                 {END}
  \\
  \\
  \gramdef{Expressions}{e}
                 { c               }{ Constante }
                 { lv              }{ Accès mémoire }
                 { \opun~e         }{ Opération unaire }
                 { e~\opbin~e      }{ Opération binaire }
                 { \& lv           }{ Pointeur }
                 { lv ← e          }{ Affectation }
                 { \eStruct{
                    l_1 : e_1
                    ; …
                    ; l_n : e_n }  }{ Structure }
                 { \eArray{e_1 ;…; e_n} }{ Tableau }
                 { f               }{ Fonction }
                 { e (e_1, …, e_n) }{ Appel de fonction }
                 {END}
  \\
  \\
  \gramdef{Left-values}{lv}
                  { x      }{ Variable }
                  { *e     }{ Déréférencement }
                  { lv.l_S }{ Accès à un champ }
                  { lv[e]  }{ Accès à un élément }
                  {END}
  \\
  \\
  \gramdef{Fonctions}{f}
                 { \eFun{x_1, …, x_n}{i} }{Arguments, corps}
                 {END}
  \end{align*}
}

\def\figstxinstr{
  \begin{align*}
  \gramdef{Instructions}{i}
                 { \iPass          }{Instruction vide}
                 { i;i             }{Séquence}
                 { e               }{Expression}
                 { \iDecl{x}{e}{i} }{Déclaration de variable}
                 { \iIf{e}{i}{i}   }{Alternative}
                 { \iWhile{e}{i}   }{Boucle}
                 { \iReturn{e}     }{Retour de fonction}
                 {END}
  \\
  \\
  \gramdef{Phrases}{p}
                 { x = e }{Variable globale}
                 { e     }{Évaluation d'expression}
                 {END}
  \\
  \\
  \gramdef{Programme}{P}
                 { (p_1, …, p_n)}{Phrases}
                 {END}
  \end{align*}
}

\def\figstxops{
  \begin{align*}
  \gramdef{Opérateurs \linebreak binaires}{\opbin}
      { +,-,\times,/,\% }{Arithmétique entière}
      { +.,-.,\times.,/.}{Arithmétique flottante}
      { +_p, -_p        }{Arithmétique de pointeurs}
      { =,≠,≤,≥,<,>     }{Comparaisons}
      { \&,|,\opxor     }{Opérateurs bit à bit}
      { \&\&,||         }{Opérateurs logiques}
      { \ll, \gg        }{Décalages}
      {END}
  \\
  \\
  \gramdef{Opérateurs unaires}{\opun}
      { +,-   }{Arithmétique entière}
      { +.,-. }{Arithmétique flottante}
      { \sim  }{Négation bit à bit}
      { !     }{Négation logique}
      {END}
  \end{align*}
}


\usetikzlibrary{calc}
\usetikzlibrary{shapes}
\usetikzlibrary{arrows}
\usetikzlibrary{decorations.pathreplacing}
\usetikzlibrary{patterns}

\usepackage{fancyvrb}

\makeatletter
\def\PY@reset{\let\PY@it=\relax \let\PY@bf=\relax%
    \let\PY@ul=\relax \let\PY@tc=\relax%
    \let\PY@bc=\relax \let\PY@ff=\relax}
\def\PY@tok#1{\csname PY@tok@#1\endcsname}
\def\PY@toks#1+{\ifx\relax#1\empty\else%
    \PY@tok{#1}\expandafter\PY@toks\fi}
\def\PY@do#1{\PY@bc{\PY@tc{\PY@ul{%
    \PY@it{\PY@bf{\PY@ff{#1}}}}}}}
\def\PY#1#2{\PY@reset\PY@toks#1+\relax+\PY@do{#2}}

\expandafter\def\csname PY@tok@gd\endcsname{\def\PY@tc##1{\textcolor[rgb]{0.67,0.00,0.00}{##1}}}
\expandafter\def\csname PY@tok@gu\endcsname{\let\PY@bf=\textbf\def\PY@tc##1{\textcolor[rgb]{0.50,0.00,0.50}{##1}}}
\expandafter\def\csname PY@tok@gt\endcsname{\def\PY@tc##1{\textcolor[rgb]{0.67,0.00,0.00}{##1}}}
\expandafter\def\csname PY@tok@gs\endcsname{\let\PY@bf=\textbf}
\expandafter\def\csname PY@tok@gr\endcsname{\def\PY@tc##1{\textcolor[rgb]{0.67,0.00,0.00}{##1}}}
\expandafter\def\csname PY@tok@cm\endcsname{\let\PY@it=\textit\def\PY@tc##1{\textcolor[rgb]{0.67,0.67,0.67}{##1}}}
\expandafter\def\csname PY@tok@vg\endcsname{\def\PY@tc##1{\textcolor[rgb]{0.67,0.00,0.00}{##1}}}
\expandafter\def\csname PY@tok@m\endcsname{\def\PY@tc##1{\textcolor[rgb]{0.00,0.60,0.60}{##1}}}
\expandafter\def\csname PY@tok@mh\endcsname{\def\PY@tc##1{\textcolor[rgb]{0.00,0.60,0.60}{##1}}}
\expandafter\def\csname PY@tok@cs\endcsname{\let\PY@it=\textit\def\PY@tc##1{\textcolor[rgb]{0.00,0.00,0.67}{##1}}}
\expandafter\def\csname PY@tok@ge\endcsname{\let\PY@it=\textit}
\expandafter\def\csname PY@tok@vc\endcsname{\def\PY@tc##1{\textcolor[rgb]{0.67,0.00,0.00}{##1}}}
\expandafter\def\csname PY@tok@il\endcsname{\def\PY@tc##1{\textcolor[rgb]{0.00,0.60,0.60}{##1}}}
\expandafter\def\csname PY@tok@go\endcsname{\def\PY@tc##1{\textcolor[rgb]{0.53,0.53,0.53}{##1}}}
\expandafter\def\csname PY@tok@cp\endcsname{\def\PY@tc##1{\textcolor[rgb]{0.30,0.51,0.09}{##1}}}
\expandafter\def\csname PY@tok@gi\endcsname{\def\PY@tc##1{\textcolor[rgb]{0.00,0.67,0.00}{##1}}}
\expandafter\def\csname PY@tok@gh\endcsname{\let\PY@bf=\textbf\def\PY@tc##1{\textcolor[rgb]{0.00,0.00,0.50}{##1}}}
\expandafter\def\csname PY@tok@ni\endcsname{\let\PY@bf=\textbf\def\PY@tc##1{\textcolor[rgb]{0.50,0.00,0.00}{##1}}}
\expandafter\def\csname PY@tok@nn\endcsname{\let\PY@ul=\underline\def\PY@tc##1{\textcolor[rgb]{0.00,0.67,0.67}{##1}}}
\expandafter\def\csname PY@tok@no\endcsname{\def\PY@tc##1{\textcolor[rgb]{0.67,0.00,0.00}{##1}}}
\expandafter\def\csname PY@tok@na\endcsname{\def\PY@tc##1{\textcolor[rgb]{0.12,0.56,1.00}{##1}}}
\expandafter\def\csname PY@tok@nb\endcsname{\def\PY@tc##1{\textcolor[rgb]{0.00,0.67,0.67}{##1}}}
\expandafter\def\csname PY@tok@nc\endcsname{\let\PY@ul=\underline\def\PY@tc##1{\textcolor[rgb]{0.00,0.67,0.00}{##1}}}
\expandafter\def\csname PY@tok@nd\endcsname{\def\PY@tc##1{\textcolor[rgb]{0.53,0.53,0.53}{##1}}}
\expandafter\def\csname PY@tok@si\endcsname{\def\PY@tc##1{\textcolor[rgb]{0.67,0.33,0.00}{##1}}}
\expandafter\def\csname PY@tok@nf\endcsname{\def\PY@tc##1{\textcolor[rgb]{0.00,0.67,0.00}{##1}}}
\expandafter\def\csname PY@tok@s2\endcsname{\def\PY@tc##1{\textcolor[rgb]{0.67,0.33,0.00}{##1}}}
\expandafter\def\csname PY@tok@vi\endcsname{\def\PY@tc##1{\textcolor[rgb]{0.67,0.00,0.00}{##1}}}
\expandafter\def\csname PY@tok@nt\endcsname{\let\PY@bf=\textbf\def\PY@tc##1{\textcolor[rgb]{0.12,0.56,1.00}{##1}}}
\expandafter\def\csname PY@tok@nv\endcsname{\def\PY@tc##1{\textcolor[rgb]{0.67,0.00,0.00}{##1}}}
\expandafter\def\csname PY@tok@s1\endcsname{\def\PY@tc##1{\textcolor[rgb]{0.67,0.33,0.00}{##1}}}
\expandafter\def\csname PY@tok@gp\endcsname{\def\PY@tc##1{\textcolor[rgb]{0.33,0.33,0.33}{##1}}}
\expandafter\def\csname PY@tok@sh\endcsname{\def\PY@tc##1{\textcolor[rgb]{0.67,0.33,0.00}{##1}}}
\expandafter\def\csname PY@tok@ow\endcsname{\def\PY@tc##1{\textcolor[rgb]{0.00,0.00,0.67}{##1}}}
\expandafter\def\csname PY@tok@sx\endcsname{\def\PY@tc##1{\textcolor[rgb]{0.67,0.33,0.00}{##1}}}
\expandafter\def\csname PY@tok@bp\endcsname{\def\PY@tc##1{\textcolor[rgb]{0.00,0.67,0.67}{##1}}}
\expandafter\def\csname PY@tok@c1\endcsname{\let\PY@it=\textit\def\PY@tc##1{\textcolor[rgb]{0.67,0.67,0.67}{##1}}}
\expandafter\def\csname PY@tok@kc\endcsname{\def\PY@tc##1{\textcolor[rgb]{0.00,0.00,0.67}{##1}}}
\expandafter\def\csname PY@tok@c\endcsname{\let\PY@it=\textit\def\PY@tc##1{\textcolor[rgb]{0.67,0.67,0.67}{##1}}}
\expandafter\def\csname PY@tok@mf\endcsname{\def\PY@tc##1{\textcolor[rgb]{0.00,0.60,0.60}{##1}}}
\expandafter\def\csname PY@tok@err\endcsname{\def\PY@tc##1{\textcolor[rgb]{0.94,0.00,0.00}{##1}}\def\PY@bc##1{\setlength{\fboxsep}{0pt}\colorbox[rgb]{0.94,0.63,0.63}{\strut ##1}}}
\expandafter\def\csname PY@tok@kd\endcsname{\def\PY@tc##1{\textcolor[rgb]{0.00,0.00,0.67}{##1}}}
\expandafter\def\csname PY@tok@ss\endcsname{\def\PY@tc##1{\textcolor[rgb]{0.00,0.00,0.67}{##1}}}
\expandafter\def\csname PY@tok@sr\endcsname{\def\PY@tc##1{\textcolor[rgb]{0.00,0.60,0.60}{##1}}}
\expandafter\def\csname PY@tok@mo\endcsname{\def\PY@tc##1{\textcolor[rgb]{0.00,0.60,0.60}{##1}}}
\expandafter\def\csname PY@tok@mi\endcsname{\def\PY@tc##1{\textcolor[rgb]{0.00,0.60,0.60}{##1}}}
\expandafter\def\csname PY@tok@kn\endcsname{\def\PY@tc##1{\textcolor[rgb]{0.00,0.00,0.67}{##1}}}
\expandafter\def\csname PY@tok@kr\endcsname{\def\PY@tc##1{\textcolor[rgb]{0.00,0.00,0.67}{##1}}}
\expandafter\def\csname PY@tok@s\endcsname{\def\PY@tc##1{\textcolor[rgb]{0.67,0.33,0.00}{##1}}}
\expandafter\def\csname PY@tok@kp\endcsname{\def\PY@tc##1{\textcolor[rgb]{0.00,0.00,0.67}{##1}}}
\expandafter\def\csname PY@tok@w\endcsname{\def\PY@tc##1{\textcolor[rgb]{0.73,0.73,0.73}{##1}}}
\expandafter\def\csname PY@tok@kt\endcsname{\def\PY@tc##1{\textcolor[rgb]{0.00,0.67,0.67}{##1}}}
\expandafter\def\csname PY@tok@sc\endcsname{\def\PY@tc##1{\textcolor[rgb]{0.67,0.33,0.00}{##1}}}
\expandafter\def\csname PY@tok@sb\endcsname{\def\PY@tc##1{\textcolor[rgb]{0.67,0.33,0.00}{##1}}}
\expandafter\def\csname PY@tok@k\endcsname{\def\PY@tc##1{\textcolor[rgb]{0.00,0.00,0.67}{##1}}}
\expandafter\def\csname PY@tok@se\endcsname{\def\PY@tc##1{\textcolor[rgb]{0.67,0.33,0.00}{##1}}}
\expandafter\def\csname PY@tok@sd\endcsname{\def\PY@tc##1{\textcolor[rgb]{0.67,0.33,0.00}{##1}}}

\def\PYZbs{\char`\\}
\def\PYZus{\char`\_}
\def\PYZob{\char`\{}
\def\PYZcb{\char`\}}
\def\PYZca{\char`\^}
\def\PYZam{\char`\&}
\def\PYZlt{\char`\<}
\def\PYZgt{\char`\>}
\def\PYZsh{\char`\#}
\def\PYZpc{\char`\%}
\def\PYZdl{\char`\$}
\def\PYZti{\char`\~}
% for compatibility with earlier versions
\def\PYZat{@}
\def\PYZlb{[}
\def\PYZrb{]}
\makeatother



\floatstyle{boxed}
\restylefloat{figure}
\newsubfloat{figure}

% Epigraphs
\epigraphfontsize{\itshape}
\setlength\epigraphwidth{8cm}
\setlength\epigraphrule{0pt}

% conditional for xetex or luatex
\newif\ifxetexorluatex
\ifxetex
  \xetexorluatextrue
\else
  \ifluatex
    \xetexorluatextrue
  \else
    \xetexorluatexfalse
  \fi
\fi
%
\ifxetexorluatex%
  \usepackage{fontspec}
  \usepackage{libertine} % or use \setmainfont to choose any font on your system
  \newfontfamily\quotefont[Ligatures=TeX]{Linux Libertine O} % selects Libertine as the quote font
\else
  \usepackage[utf8]{inputenc}
  \usepackage[T1]{fontenc}
  %\usepackage{libertine} % or any other font package
  \newcommand*\quotefont{\fontfamily{LinuxLibertineT-LF}} % selects Libertine as the quote font
\fi

\newcommand*\quotesize{60} % if quote size changes, need a way to make shifts relative
% Make commands for the quotes
\newcommand*{\openquote}
   {\tikz[remember picture,overlay,xshift=-4ex,yshift=-2.5ex]
   \node (OQ) {\quotefont\fontsize{\quotesize}{\quotesize}\selectfont``};\kern0pt}

\newcommand*{\closequote}[1]
  {\tikz[remember picture,overlay,xshift=4ex,yshift={#1}]
   \node (CQ) {\quotefont\fontsize{\quotesize}{\quotesize}\selectfont''};}

% select a colour for the shading
\colorlet{shadecolor}{white}

\newcommand*\shadedauthorformat{\emph} % define format for the author argument

% Now a command to allow left, right and centre alignment of the author
\newcommand*\authoralign[1]{%
  \if#1l
    \def\authorfill{}\def\quotefill{\hfill}
  \else
    \if#1r
      \def\authorfill{\hfill}\def\quotefill{}
    \else
      \if#1c
        \gdef\authorfill{\hfill}\def\quotefill{\hfill}
      \else\typeout{Invalid option}
      \fi
    \fi
  \fi}
% wrap everything in its own environment which takes one argument (author) and one optional argument
% specifying the alignment [l, r or c]
%
\newenvironment{shadequote}[2][l]%
{\authoralign{#1}
\ifblank{#2}
   {\def\shadequoteauthor{}\def\yshift{-2ex}\def\quotefill{\hfill}}
   {\def\shadequoteauthor{\par\authorfill\shadedauthorformat{#2}}\def\yshift{2ex}}
\begin{snugshade}\begin{quote}\openquote}
{\shadequoteauthor\quotefill\closequote{\yshift}\end{quote}\end{snugshade}}



\setlist[itemize]{label=$\bullet$}

\maxtocdepth{subsection}    % What appears in toc
\maxsecnumdepth{subsection} % What get numbered

% Part header
\usepackage{adjustbox}
\newlength{\currentparindent}
\definecolor{lightgrey}{gray}{0.7}
\def\partintro{}
\renewcommand{\afterpartskip}{
\ifdefempty{\partintro}{}{
\vfil
\noindent
\setlength{\currentparindent}{\parindent}
{%
\setlength{\fboxsep}{20pt}%
\adjustbox{minipage=\textwidth-2\fboxsep-2\fboxrule,cfbox=lightgrey,center}%
          {\setlength{\parindent}{\currentparindent}%
           \partintro}%
}%
}
\vfil
\newpage}

%\theoremstyle{definition}
\newtheorem{definition}{Définition}[chapter]
\newtheorem{theorem}{Théorème}[chapter]
\newtheorem{lemma}{Lemme}[chapter]
\newtheorem{example}{Exemple}[chapter]

\setlength{\headsep}{5mm}
\addtolength{\textheight}{-1cm}



\title{Analyse de sécurité de logiciels système par typage statique}
\subtitle{Application au noyau Linux}
\author{Étienne Millon}
\date{xx-yy-2013}

\usepackage{url}
\ifpdf
\usepackage{pdfcolmk}
\fi
%% check if using xelatex rather than pdflatex
\ifxetex
\usepackage{fontspec}
\fi
\usepackage{graphicx}
%%\usepackage{hyperref}
%% drawing package
\usepackage{tikz}
%% for dingbats
\usepackage{pifont}
\providecommand{\HUGE}{\Huge}% if not using memoir
\newlength{\drop}% for my convenience
%% specify the Webomints family
\newcommand*{\wb}[2]{\fontsize{#1}{#2}\usefont{U}{webo}{xl}{n}}
%% select a (FontSite) font by its font family ID
\newcommand*{\FSfont}[1]{\fontencoding{T1}\fontfamily{#1}\selectfont}
%% if you don’t have the FontSite fonts either \renewcommand*{\FSfont}[1]{}
%% or use your own choice of family.
%% select a (TeX Font) font by its font family ID
\newcommand*{\TXfont}[1]{\fontencoding{T1}\fontfamily{#1}\selectfont}
%% Generic publisher’s logo
\newcommand*{\plogo}{\fbox{$\mathcal{PL}$}}

%% Some shades
\definecolor{Dark}{gray}{0.2}
\definecolor{MedDark}{gray}{0.4}
\definecolor{Medium}{gray}{0.6}
\definecolor{Light}{gray}{0.8}
%%%% Additional font series macros
\makeatletter
%%%% light series
%% e.g., kernel doc, section s: line 12 or thereabouts
\DeclareRobustCommand\ltseries
{\not@math@alphabet\ltseries\relax
\fontseries\ltdefault\selectfont}
%% e.g., kernel doc, section t: line 32 or thereabouts
\newcommand{\ltdefault}{l}
%% e.g., kernel doc, section v: line 19 or thereabouts
\DeclareTextFontCommand{\textlt}{\ltseries}
% heavy(bold) series
\DeclareRobustCommand\hbseries
{\not@math@alphabet\hbseries\relax
\fontseries\hbdefault\selectfont}
\newcommand{\hbdefault}{hb}
\DeclareTextFontCommand{\texthb}{\hbseries}
\makeatother


\usetikzlibrary{fadings}

\newcommand*{\titleUL}{
  \newgeometry{left=2cm, right=2cm, top=2cm, bottom=2cm}%
  \begingroup% University of Liege
\drop=0.1\textheight
\begin{center}
{\textsc{Université Pierre et Marie Curie\\
École Doctorale Informatique, Télécommunications et Électronique
}}\\
\vspace{1cm}
\begin{tikzpicture}
  \begin{scope}
    \node (A) {\includegraphics[height=\textwidth]{curies.pdf}};
    \path[opacity=0.2, fill=white]
      (A.north west) rectangle (A.south east);
    %\path[opacity=0.8, fill=white, path fading=north, fading transform={rotate=-45}]
    %(A.north west) rectangle (A.south east);
    %\clip (A.center) circle(8.5cm);
    %\shade[opacity=0.2, shading=axis,bottom color=black!10,top color=black,shading angle=-45]
    %(A.south west) rectangle (A.north east);
  \end{scope}
\end{tikzpicture}
\vspace{-13cm}
\\[\drop]
\begin{minipage}{13cm}
  \begin{center}
\rule{\textwidth}{1pt}\par
\vspace{0.5\baselineskip}
{\huge\bfseries \thetitle\\
\vspace{5mm}
\large --- \thesubtitle ---
}\\[0.5\baselineskip]
\rule{\textwidth}{1pt}
  \end{center}
\end{minipage}
\par
\vspace{1cm}
{{\Large\textsc{\theauthor}} \\
sous la direction d'Emmanuel Chailloux et de Sarah Zennou
}
\vspace{16mm}\\
{{\LARGE\textsc{Thèse}}\\
pour obtenir le titre de\\
{\Large Docteur en Sciences}\\
{\large mention Informatique}
}
\vspace{5mm}\\
{
  Version du \today{} \\
}
{\large Soutenue le \thedate{} devant un jury composé de\\

\textbf{Président}\\
---\\
\textbf{Rapporteurs}\\
---\\
---\\
%Sandrine Blazy\\
%Pierre Jouvelot\\
\textbf{Directeurs}\\
Emmanuel Chailloux\\
Sarah Zennou\\
\textbf{Examinateurs}\\
---\\
---
}
\end{center}
\endgroup
\restoregeometry
}


\begin{document}

\frontmatter

\begin{titlingpage}
\titleUL{}
\end{titlingpage}

% TODx decom
%\cleartorecto %{{{

%\begin{headingpage}
%\hfill%
%\begin{minipage}{0.6\textwidth}
%\shorthandoff{;}
%\begin{shadequote}[r]{Peter van den Linden}
%Many C programmers believe that \enquote{strong typing}
%just means pounding extra hard on the keyboard.
%\end{shadequote}

%%The fundamental problem addressed by a type theory is to ensure that programs
%%have meaning. The fundamental problem caused by a type theory is that meaningful
%%programs may not have meanings ascribed to them. The quest for richer type
%%systems results from this tension.
 %%Mark Manasse

%\shorthandon{;}
%\end{minipage}
%\end{headingpage}%}}}

%{{{
\asuacknowledgements

Remerciements

%}}}

\clearpage

%\asudedication{Dédicace}

\tableofcontents

\mainmatter

\chapter{Introduction}

\label{cha:intro}
Communication,
audiovisuel,
transports,
médecine:
tout ces domaines se sont transformés dans les dernières décennies,
en particulier grâce à la révolution numérique.
En effet le plus petit appareil électrique contient maintenant des composants
matériels programmables.

% TODO faire une intro "la vie est numérique" avec du Berry @ CdF

En 2013, on pense bien sûr aux téléphones portables dont la fonctionnalité et la
complexité les rapprochent des ordinateurs de bureau. Par exemple, le système
d'exploitation Android de Google est basé sur le noyau Linux, destiné à la base
aux micro-ordinateurs.

Le noyau d'un système d'exploitation est chargé de faire l'intermédiaire entre
le matériel (processeur, mémoire, périphériques, …) et les applications
exécutées sur celui-ci (par exemple un navigateur web, une calculatrice ou un
carnet d'adresses). Il doit aussi garantir la sécurité de celles-ci.

En tant qu'intermédiaire de confiance, le noyau a un certain nombre de
responsabilités et est le seul à avoir accès à certaines informations sensibles.
Il est capital de s'assurer qu'il est bien le seul à pouvoir y accéder. En
particulier, il faut pouvoir vérifier que les requêtes faites par l'utilisateur
au noyau ne peuvent pas volontairement ou involontairement détourner ce dernier
et lui faire fuiter des informations confidentielles.

Le problème est que comme tous les logiciels, les noyaux de système
d'exploitation sont écrits par des humains, qui ne sont pas parfaits. Loin de
là: on estime qu'avant relecture, 1000 lignes de code contiennent entre 5 et 100
erreurs de programmation en moyenne.
% TODO sourcer
Les activités de relecture et de
débogage ont beau prendre la majeure partie du temps de développement, il est
facile de laisser passer des défauts de programmation.

Une technique efficace est de réaliser des tests, c'est-à-dire exécuter le
programme sous un environnement contrôlé. On peut alors détecter des
comportements non désirés. Mais même avec des tests les plus exhaustifs il n'est
pas possible de couvrir tous les cas d'utilisation.

Une autre approche est d'analyser le code source du programme avant de
l'exécuter et de refuser de lancer les programmes qui contiennent des erreurs,
quitte à limiter la possibilité d'en écrire certains. C'est l'analyse statique
de programmes.

Une des techniques d'analyse statique les plus répandues est le typage statique,
qui consiste à associer à chaque morceau de programme, une étiquette
décrivant quel genre de valeur sera produite par son évaluation. Par exemple, si
$n$ est le nom d'une variable entière, alors $n + 2$ produira toujours une
valeur entière.

Pour garantir l'isolation d'un noyau de système d'exploitation, un des points
cruciaux est de restreindre la manière dont sont traitées les informations
provenant des programmes utilisateur. Le but de cet thèse est de montrer que le
typage statique peut être utilisé pour détecter et interdire ces manipulations
dangereuses.

\section{Rôle d'un système d'exploitation}

Un ordinateur est constitué de nombreux composants matériels: microprocesseur,
mémoire, et divers périphériques. Et au niveau de l'utilisateur, des dizaines de
logiciels permettent d'effectuer toutes sortes de calculs et de communications.
Le système d'exploitation permet de faire l'interface entre ces niveaux
d'abstraction.

Au cours de l'histoire des systèmes informatiques, la manière de les programmer
a beaucoup évolué. Au départ, les programmeurs avaient accès au matériel dans
son intégralité: toute la mémoire pouvait être accédée, toutes les instructions
pouvaient être utilisées.

Néanmoins c'est un peu restrictif, puisque cela ne permet qu'à une personne
d'interagir avec le système. Dans la seconde moitié des années 60, sont apparus
les premiers systèmes ``à temps partagé'', permettant à plusieurs utilisateurs
de travailler en même temps.

Permettre l'exécution de plusieurs programmes en même temps est une idée
révolutionnaire, mais elle n'est pas sans difficultés techniques: en effet les
ressources de la machine doivent être aussi partagées entre les utilisateurs et
les programmes. Par exemple, plusieurs programmes vont utiliser le processeur
les uns à la suite des autres; et chaque programme aura à sa disposition une
partie de la mémoire principale, ou du disque dur.

Si plusieurs programmes s'exécutent de manière concurrente sur le même matériel,
il faut s'assurer que l'un ne puisse pas écrire dans la mémoire de l'autre, ou
que les deux n'utilisent pas la carte réseau en même temps. Ce sont des rôles du
système d'exploitation.

Ainsi, au lieu d'accéder directement au matériel via des instructions de bas
niveau, les programmes communiquent avec le noyau, qui centralise donc les
appels au matériel, et abstrait certaines opérations.

Par exemple, comparons ce qui se passe concrètement lors de la copie de données
depuis un cédérom ou une clef USB.\@

\begin{itemize}

  \item Dans le cas du cédérom, il faut interroger le bus SATA, interroger le
    lecteur sur la présence d'un disque dans le lecteur, activer le moteur,
    calculer le numéro de trame des données sur le disque, demander la lecture,
    puis déclencher une copie de la mémoire.

  \item Avec une clef, il faut interroger le bus USB, rechercher le bon numéro
    de périphérique, le bon numéro de canal dans celui-ci, lui appliquer une
    commande de lecture au bon numéro de bloc, puis copier la mémoire.

\end{itemize}

Ces deux opérations, bien qu'elles aient la même intention (copier de la mémoire
depuis un périphérique amovible), ne sont pas effectuées en extension de la même
manière. C'est pourquoi le système d'exploitation fournit les notions de
fichier, lecteur, etc: le programmeur n'a plus qu'à utiliser des commandes de
haut niveau (``monter un lecteur'', ``ouvrir un fichier'', ``lire dans un
fichier'') et selon le type de lecteur, le système d'exploitation effectuera les
actions appropriées.

En résumé, un système d'exploitation est l'intermédiaire entre le logiciel et le
matériel, et en particulier est responsable de la gestion de la mémoire, des
périphériques et des processus. Les détails d'implantation ne sont pas présentés
à l'utilisateur; à la place il manipule des abstractions, comme la notion de
fichier. Pour une explication détaillée du concept de système d'exploitation
ainsi que des cas d'étude, on pourra se référer à~\cite{tanenbaum}.

\section{Séparation entre noyau et espace utilisateur}

Puisque le noyau est garant du bon fonctionnement du système, il ne doit pas
pouvoir être manipulé directement par l'utilisateur ou les programmes exécutés.
Ainsi, il est nécessaire de mettre en place des protections entre les espaces
noyau et utilisateur.

Au niveau matériel, on utilise la notion de \emph{niveaux de privilèges} pour
déterminer s'il est possible d'exécuter une instruction.

D'une part, le processeur contient un niveau de privilège intrinsèque. D'autre
part, chaque zone mémoire contenant du code ou des données possède également un
niveau de privilège minimum nécessaire. L'exécution d'une instruction est alors
possible si et seulement si le niveau de privilège du processeur est supérieur à
celui de l'instruction et des opérandes mémoires qui y sont présentes
\footnote{
  Ici ``supérieur'' est synonyme de ``plus privilégié''. Dans l'implantation
  d'Intel présentée dans le chapitre~\ref{cha:os}, les niveaux sont numérotés de
  0 à 3 où le niveau 0 est le plus privilégié.
}.

Par exemple, supposons qu'un programme utilisateur contienne l'instruction
\texttt{mov \%eax, \hexa{54c3b2f7}}. Celle-ci consiste à déplacer le contenu du
registre \eax{} vers l'adresse mémoire \texttt{\hexa{54c3b2f7}}. Puisque cette
adresse est dans l'espace utilisateur (la raison est décrite dans le
chapitre~\ref{cha:os}), aucune erreur de protection mémoire n'est déclenchée.

Ainsi, pour une instruction manipulant des données en mémoire, les accès
possibles sont décrits dans la figure~\ref{fig:erreursec}. En cas
d'impossibilité (signalée par un \Square), une erreur se produit et l'exécution
s'arrête.

\begin{figure}
\begin{center}
\def\modeK{Noyau\xspace}
\def\modeU{Utilisateur\xspace}
\begin{tabular}{cccc}
  Mode du processeur
& Privilège (code)
& Privilège (données)
& Accès possible \\
\midrule
  \modeK{} & \modeK{} & \modeK{} & \CheckedBox{} \\
  \modeK{} & \modeK{} & \modeU{} & \CheckedBox{} \\
  \modeK{} & \modeU{} & \modeK{} & \CheckedBox{} \\
  \modeK{} & \modeU{} & \modeU{} & \CheckedBox{} \\
  \modeU{} & \modeK{} & \modeK{} & \Square{} \\
  \modeU{} & \modeK{} & \modeU{} & \Square{} \\
  \modeU{} & \modeU{} & \modeK{} & \Square{} \\
  \modeU{} & \modeU{} & \modeU{} & \CheckedBox{} \\
\end{tabular}
\end{center}

\caption{ Règles d'exécution de code et d'accès à la mémoire. }
\label{fig:erreursec}

\end{figure}

En plus de cette vérification, certains types d'instructions sont explicitement
réservés au mode le plus privilégié: par exemple les lectures ou écritures sur
des ports matériels, ou celles qui permettent de définir les niveaux de
privilèges des différentes zones mémoire.

Les programmes utilisateur ne pouvant pas accéder à ces instructions de bas
niveau, ils sont très limités dans ce qu'ils peuvent faire. Puisque
l'interaction avec le matériel (comme l'écriture sur un disque dur) se fait
uniquement via des instructions privilégiées, ils sont limités à l'utilisation
du processeur et de la mémoire, permettant uniquement de réaliser des calculs.

Pour utiliser le matériel ou accéder à des abstractions de haut niveau (comme
créer un nouveau processus), ils doivent donc passer par l'intermédiaire du
noyau. La communication entre le noyau et les programmes utilisateur est
constituée par le mécanisme des \emph{appels système}.

Lors d'un appel système, une fonction du noyau est invoquée (en mode noyau) avec
des paramètres provenant de l'utilisateur. Il faut donc être particulièrement
précautionneux dans le traitement de ces données.

Par exemple, considérons un appel système de lecture depuis un disque: on passe
au noyau les arguments $(d, o, n, a)$ où $d$ est le nom du disque, $o$ (pour
\emph{offset}) l'adresse sur le disque où commencer la lecture, $n$ le nombre
d'octets à lire et $a$ l'adresse en mémoire où commencer à stocker les
résultats.

Dans le cas d'utilisation prévu, le noyau va copier la mémoire lue dans $a$. Le
processeur est en mode noyau, en train d'exécuter une instruction du noyau
manipulant des données utilisateur. D'après la figure~\ref{fig:erreursec} aucune
erreur ne se produit.

Mais même si ce cas ne produit pas d'erreur à l'exécution, il est tout de même
erroné. En effet, si on passe à l'appel système une adresse $a$ faisant partie
de l'espace noyau, que se passe-t-il?

L'exécution est presque identique: au moment de la copie on est en mode noyau,
en train d'exécuter une instruction du noyau manipulant des données noyau.
Encore une fois il n'y a pas d'erreur à l'exécution.

On peut donc écrire n'importe où en mémoire. De même, une fonction d'écriture
sur un disque (et lisant en mémoire) permettrait de lire de la mémoire du noyau.
À partir de ces primitives, on peut accéder aux autres processus exécutés, ou
détourner l'exécution vers du code arbitraire. L'isolation est totalement
brisée à cause de ces appels système.

La cause de ceci est qu'on a accédé à la mémoire en testant les privilèges du
noyau au lieu de tester les privilèges de celui qui a fait la requête
(l'utilisateur). Ce problème est connu sous le nom de \emph{confused deputy
problem}~\cite{hardy88confused}.

La bonne manière d'implanter un appel système est donc d'interdire le
déréférencement direct des pointeurs dont la valeur peut être contrôlée par
l'utilisateur. Dans le cas du passage par adresse d'un argument, il aurait fallu
vérifier à l'exécution que celui-ci a bien les mêmes privilèges que l'appelant.

Il est facile d'oublier d'ajouter cette vérification, puisque le cas ``normal''
fonctionne. Avec ce genre d'exemple on voit comment les bugs peuvent arriver si
fréquemment et pourquoi il est aussi capital de les détecter avant l'exécution.

\section{Systèmes de types}
\label{sec:types-intro}

La plupart des langages de programmation incorporent la notion de type, dont un
des buts est d'empêcher de manipuler des données incompatibles entre elles.

Au niveau du langage machine, les seules données qu'un ordinateur manipule sont
des nombres. Selon les opérations effectuées, ils seront interprétés comme des
entiers, des adresses mémoires, ou des caractères. Pourtant il est clair que
certaines opérations n'ont pas de sens: par exemple, multiplier un nombre par
une adresse, ou déréférencer le résultat d'une division sont des comportements
qu'on voudrait pouvoir empêcher.

En un mot, le but du typage est de classifier les objets et de restreindre les
opérations possibles selon la classe d'un objet: en somme, ``ne pas ajouter des
pommes et des oranges''. Le modèle qui permet cette classification est appelé
\emph{système de types} et est en général constitué d'un ensemble de
\emph{règles de typage}, comme ``un entier plus un entier égale un entier''.

\paragraph{Typage dynamique}

Dans ce cas, chaque valeur manipulée par le programme est décorée d'une
étiquette définissant comment interpréter la valeur en question. Les règles de
typage sont alors réalisées à l'exécution. Par exemple, l'opérateur ``$+$''
vérifie que ces deux opérandes ont une étiquette ``entier'', et construit alors
une valeur obtenue en faisant l'addition des deux valeurs, avec une étiquette
``entier''. Par exemple, le langage Python~\link{python} utilise cette
stratégie.

\paragraph{Typage statique}

Dans ce cas on fait les vérifications à la compilation. En quelque sorte,
l'approche dynamique est pessimiste, puisqu'elle demande de traiter à chaque
étape le cas où les types ne sont pas corrects. Intuitivement, dans le cas où
toutes les fonctions se comportent bien, faire la vérification est inutile. Pour
vérifier ceci, on donne à chaque fonction un contrat comme ``si deux entiers
sont passés, et que la fonction renvoie une valeur, alors cette valeur sera un
entier''. Cet ensemble de contrats peut être vérifié statiquement par le
compilateur, à l'aide d'un système de types statique.

Par exemple, on peut dire que l Ainsi l'opérateur ``$+$'' a pour type $(\tInt,
\tInt) → \tInt$. Cela veut dire que si on lui passe deux entiers ($\tInt$
\tInt), alors la valeur obtenue est également un entier.

\paragraph{Typage fort ou faible}

Contrairement à la distinction claire entre typage statique ou dynamique, la
séparation entre typage fort et faible est moins nette. En poussant à l'extrême,
les systèmes de types forts permettent d'éliminer totalement la nécessité de
réaliser des tests de typage. Mais souvent ce n'est pas le cas, car il peut y
avoir des constructions au sein du langage qui permettent de contourner le
système de types, comme un opérateur de transtypage.

\paragraph{Polymorphisme}

Parfois, il est trop restrictif de donner un unique type à une fonction. Si on
considère une fonction ajoutant un élément à une liste, ou une autre
triant un tableau en place, leur type doit-il faire intervenir le type des
éléments manipulés?

En première approximation, on peut imaginer fournir une version du code par type
de données à manipuler. C'est la solution retenue par les premières versions du
langage Pascal, ce qui rendait très difficile l'écriture de
bibliothèques~\cite{PascalNoEscape}. On parle alors de monomorphisme.

Une autre manière de procéder est d'autoriser plusieurs fonctions à avoir le
même nom, mais avec des types d'arguments différents. Par exemple, on peut
définir séparément l'addition entre deux entiers, entre deux flottants, ou entre
un entier et un flottant. Selon les informations connues à la compilation, la
bonne version sera choisie. C'est ainsi que fonctionnent les opérateurs en C++.
On parle de polymorphisme \emph{ad hoc}, ou de surcharge.

La dernière possibilité est le polymorphisme paramétrique, qui consiste à
utiliser le même code quelque soit le type des arguments. Dans ce cas, on
utilise une seule fonction pour traiter une liste d'entiers ou une liste de
flottants, par exemple. Au lieu d'associer à chaque fonction un type, dans
certains cas on lui associe un type paramétré, instanciable en un type concret.
Dans le cas des fonctions de traitement de liste, l'idée est que lorsqu'on ne
touche pas aux éléments, alors le traitement est valable quelque soit leur type.
Cette technique a été décrite en premier dans~\cite{Milner78}.

Pour un tour d'horizon de différents systèmes types de statiques, avec en
particulier du polymorphisme, on pourra se référer à~\cite{TAPL}.

\section{Langages d'implantation}

Puisque notre but est d'analyser du code provenant de systèmes d'exploitation,
il est nécessaire de s'intéresser aux langages de programmation dans lesquels
ils sont écrits.

\paragraph{Assembleur}

Le noyau d'un système d'exploitation nécessite d'accéder au matériel, donc il
est naturellement bas niveau. Pour accéder aux fonctionnalités spécifiques de
chaque processeur, il est donc nécessaire d'en implanter une partie dans le
langage d'assemblage natif de chaque architecture. Historiquement, les premiers
système d'exploitations étaient entièrement écrits en assembleur. Cela est de
plus en plus rare, sauf dans les cas où les ressources sont trop limitées pour
exécuter du code compilé, comme dans les systèmes embarqués ou temps-réel.

\paragraph{C}

Le système Unix, développé à partir de 1969, a tout d'abord été développé en
assembleur sur un mini-ordinateur PDP-7, puis a été porté sur d'autres
architectures matérielles. Pour aider ce portage, il a été nécessaire de créer
un ``assembleur portable'', le langage C~\cite{KandR,AnsiC}. Son but est de
fournir des abstractions au dessus du langage d'assemblage. Les structures de
contrôle (\texttt{if}, \texttt{while}, \texttt{for}) permettent d'utiliser la
programmation structurée, c'est-à-dire en limitant l'utilisation de
l'instruction \texttt{goto}. Les types de données sont également abstraits de la
machine: ainsi, \texttt{int} désigne un entier machine, indépendamment de sa
taille concrète. Son système de types est assez rudimentaire: toutes les formes
de transtypage sont acceptées, certaines conversions sont insérées
automatiquement par le compilateur, et la plupart des abstractions fournies par
le langage sont perméables. Le noyau Linux est écrit dans un dialecte du langage
C. Le noyau du système Mac OSX d'Apple est également un dérivé d'Unix, et est
donc aussi écrit dans ce langage.

\paragraph{C++}

Le manque de modularité de C a conduit à la création d'un nouveau langage plus
riche autour de celui-ci. Dans les années 1980, C++ est apparu comme un C
amélioré, avec de nombreuses fonctionnalités supplémentaires comme un modèle
objet avec des classes, du polymorphisme \emph{ad hoc} et un système de
\emph{templates} puissant permettant d'implanter une forme de polymorphisme
universel. Microsoft Windows, initialement écrit en C, s'est enrichi de plus en
plus de code C++, en faisant un des principaux systèmes d'exploitation écrits
dans ce langage.

\paragraph{Ada}

Le langage C est bas niveau, mais son système de types n'est pas suffisamment
évolué pour le rendre plus sûr que le langage d'assemblage vers lequel il est
compilé. Dans les années 1970, le département de la défense américain
(\emph{DoD}) a créé un comité dans le but de créer un nouveau langage pour ses
systèmes temps-réels embarqués. Le résultat est Ada, un langage impératif avec
un système de types statique plus contraignant mais plus sûr que celui de C. Par
exemple, deux types entiers différents (comme les entiers de 16 bits et ceux de
32 bits) ne sont pas compatibles. Il faut ajouter une conversion explicite dans
le code, qui se traduira par un test à l'exécution et éventuellement par une
erreur \texttt{Constraint\_Error}. Bien qu'il ne se soit pas beaucoup développé
dans les applications grand public, Ada garde toujours une place de choix dans
l'implantation de systèmes embarqués, notamment dans les applications spatiales
et de défense.

\paragraph{Divers} % TODO titre

Il est bien sûr possible d'implanter un système d'exploitation dans n'importe
quel langage, si celui-ci peut être compilé vers du langage machine et qu'il
permet facilement d'inclure des fragments d'assembleur pour les parties de bas
niveau. Néanmoins il s'agit en général de travaux de recherche. On peut ainsi
trouver des systèmes d'exploitation écrits en C\#~\cite{SingularityOverview},
Haskell~\cite{HaskellOS} ou Ocaml~\cite{unikernels}.

% TODO[B] TROLL

\section{Langages d'analyse}

Les langages décrits précédemments sont faits pour être facilement écrits par
des programmeurs humains. En général ils sont ambigüs ou ont des comportements
implicites. Pour analyser du code source, il est plus pratique d'avoir une
représentation intermédiaire plus simple afin d'avoir moins de traitements
dupliqués. Dans de nombreux projets, des sous-ensembles de C ont été définis
pour aller dans ce sens.

\subsection*{Middle-ends}

Les premiers candidats sont bien entendu les représentations intermédiaires
utilisées dans les compilateurs C. Elles ont l'avantage d'accepter en plus du C
standard, les diverses extensions (GNU, Microsoft, Plan9) utilisées par la
plupart des logiciels. En particulier, le noyau Linux repose fortement sur les
extensions GNU.\@

\paragraph{GCC} utilise une représentation interne nommée
GIMPLE\cite{gcc-gimple}. Il s'agit d'une structure d'arbre écrite en C, reposant
sur de nombreuses macros afin de cacher les détails d'implantation interne
pouvant varier entre deux versions. Cette représentation étant réputée difficile
à manipuler, le projet MELT\cite{gcc-melt} permet de générer un greffon de
compilateur à partir d'un dialecte de Lisp.

\paragraph{LLVM}\cite{llvm-pres} est un compilateur développé par la communauté
puis sponsorisé Apple. À la différence de GCC, sa base de code est écrite en
C++. Il utilise une représentation intermédiaire qui peut être manipulée soit
sous forme d'une structure de données C++, soit d'un fichier de code-octet
compact, soit sous forme textuelle.

\paragraph{Cmm} est une représentation interne utilisée pour la génération de
code lors de la compilation d'Objective Caml~\link{ocamlManual}, et disponible
dans les sources du compilateur sous le chemin \texttt{asmcomp/cmm.mli} (il
s'agit donc d'une structure de données OCaml). Ce langage a l'avantage d'être
très restreint, mais malheureusement il n'existe pas directement de traducteur
permettant de compiler C vers Cmm.

\paragraph{C-~-}\cite{spjcmm} \link{cmm}, dont le nom est inspiré du précédent,
est un projet qui visait à unifier les langages intermédiaires utilisés par les
compilateurs. L'idée est que si un front-end peut émettre du C-~- (sous forme de
texte), il est possible d'obtenir du code machine efficace. Le compilateur
Haskell GHC, par exemple, utilise une représentation intermédiaire très
similaire à C-~-.

\subsection*{Langages intermédiaires ad hoc}

Comme le problème de construire une représentation intermédiaire adaptée à une
analyse statique n'est pas nouveau, plusieurs projets ont déjà essayé d'y
apporter une solution. Puisque qu'ils sont développés en parallèle des
compilateurs, le support des extensions est en général moins important dans ces
langages.

\paragraph{CIL}\cite{NeculaCil} \link{kerneisCil} est une représentation en
OCaml d'un programme C, développée depuis 2002. Grâce à un mécanisme de
greffons, elle permet de prototyper rapidement des analyses statiques de
programmes.

\paragraph{Compcert} est un projet qui vise à produire un compilateur certifié
pour C. C'est-à-dire que le fait que les transformations conservent la
sémantique est prouvé. Il utilise de nombreux langages intermédiaires, dont
CIL.\@ Pour le front-end, le langage se nomme Clight\cite{cfront}. Les passes de
middle-end, quant à elles, utilisent Cminor\cite{cminorSL}.

\paragraph{Newspeak}\cite{newspeak} est un langage intermédiaire développé par
EADS Innovation Works, et qui est spécialisé dans l'analyse de valeurs par
interprétation abstraite. Il sera décrit plus en détails dans la
Section~\ref{sec:npk}.

\subsection*{La solution retenue: \langname{}}

Dans la suite de cette thèse nous nous appuierons sur \langname{}, un langage
intermédiaire minimal, pour exprimer les règles de typage de la manière la plus
simple possible.

\section{Plan de la thèse}

Cette thèse comporte trois parties.

\paragraph{La partie~\ref{part:ctx}} présente le contexte de ces travaux. Le
fonctionnement général d'un système d'exploitation y est détaillé, et les
problèmes de manipulation de pointeurs contrôlés par l'utilisateur y sont
introduits. On fait ensuite un tour d'horizon des techniques existantes
permettant de traiter ce problème. Cette partie se conclut par la proposition au
cœur de cette thèse.

\paragraph{La partie~\ref{part:lang}} décrit notre solution: \langname{}, un
langage impératif. On y décrit sa syntaxe, sa sémantique ainsi qu'un système de
types statiques. On l'étend ensuite pour capturer les problèmes d'adressage
mémoire présents dans les systèmes d'exploitation en ajoutant des pointeurs
contrôlés par l'utilisateur. Le système de types est également étendu. Pour
chacune de ces variantes, on établit la propriété de sûreté de typage reliant la
sémantique d'exécution aux types statiques.

\paragraph{La partie~\ref{part:xp}} documente la démarche expérimentale associée
à ces travaux. L'implantation du système de types est décrite, afin que la
manière de transformer automatiquement du code C en \langname{}. Un cas d'étude
est déroulé, consistant d'un bug ayant touché le noyau Linux. Il est démontré
que le système de type capture précisément ce genre d'erreur de programmation.
Enfin, les possibilités d'extension tant théoriques qu'expérimentales sont
présentées.

%  - historique des techniques
%  - analyse non intrusive
%  - méthodes de développement du noyau
%  - syntaxe, sémantique, typage

% TODO [E] + besoins EADS

% vim: spelllang=fr



\long\def\partintro{
Après avoir décrit le contexte général de ces travaux, nous décrivons leurs
enjeux.

Le chapitre~\ref{cha:os} explore plus en détail le fonctionnement d'un système
d'exploitation, y compris la séparation du code en plusieurs niveaux de
privilèges. L'architecture Intel 32 bits est prise comme support. En
particulier, le mécanisme des appels système est décrit et on montre qu'une
implantation naïve de la communication entre espaces utilisateur et noyau casse
toute isolation.

Le chapitre~\ref{cha:etatdelart} consiste en un tour d'horizon des techniques
existantes en analyses de programmes. Ces analyses se centrent autour des
problèmes liés à la vérification de code système ou embarqué, y compris le
problème de manipulation mémoire évoqué dans le chapitre~\ref{cha:os}.

On conclut en présentant notre technique: \langname, un langage permettant de
typer des programmes impératifs, plus précisément en ajoutant des types
pointeurs abstraits.
}

\part{Méthodes formelles pour la sécurité}
\label{part:ctx}


\chapter{Systèmes d'exploitation}

\label{cha:os}
\section{Rôle d'un système d'exploitation}

% Fig Archi {{{
\begin{figure}
\centering
\begin{tikzpicture}
[ every node/.style={}
, node distance=2cm
]

\node (cpu) {CPU};

\node[below of=cpu] (northb) {Northbridge};

\node[right of=northb] (mem) {Mémoire};
\node[left of=northb]  (agp) {AGP};
\node[below of=northb] (southb) {Southbridge};

\node[below left of=southb, xshift=-1cm] (sata) {SATA};
\node[below right of=southb, xshift=1cm] (usb)  {USB};

\node[below left of=sata]  (hdd)     {Disque dur};
\node[below right of=sata] (cdrom)   {CD-ROM};

\node[below left of=usb]  (memo)    {Mémorette};
\node[below right of=usb] (clavier) {Clavier};

\draw (cpu) -- (northb);
\draw (northb) -- (agp);
\draw (northb) -- (mem);
\draw (southb) -- (northb);
\draw (southb) -- (sata);
\draw (southb) -- (usb);
\draw (usb) -- (memo);
\draw (usb) -- (clavier);
\draw (sata) -- (hdd);
\draw (sata) -- (cdrom);

\end{tikzpicture}

\caption{Architecture simplifiée d'un ordinateur}
\label{fig:archi-simplifiee}
\end{figure}
% }}}

Au plus bas, un ordinateur est constitué de nombreux composants matériels :
micro-processeur, mémoire, et divers périphériques
(figure~\ref{fig:archi-simplifiee}). Et au niveau de l'utilisateur, de dizaines
de logiciels permettant d'effectuer toutes sortes de calculs et de
communication. Le système d'exploitation permet de faire l'interface entre ces
niveaux d'abstraction. Dans ce chapitre nous allons voir son rôle et une
implantation possible. Pour une description plus détaillée, on pourra se référer
à l'ouvrage \cite{tanenbaum}.

Au cours de l'histoire des systèmes informatiques, la manière de les programmer
a beaucoup évolué. Au départ, les programmeurs avaient accès au matériel dans
son intégralité : toute la mémoire pouvait être accédé, toutes les instructions
pouvaient être utilisées.

Néanmoins c'est un peu restrictif : si on est seul à utiliser un système, on est
par définition... seul à pouvoir l'utiliser. Dans la seconde moitié des années
60, sont apparus les premiers systèmes "à temps partagé", permettant à plusieurs
utilisateurs de travailler en même temps.

Permettre l'exécution de plusieurs programmes en même temps est une idée
révolutionnaire, mais elle n'est pas sans difficultés techniques : en effet les
ressources de la machine doivent être aussi partagées entre les utilisateurs et
les programmes. Par exemple, plusieurs programmes vont par exemple utiliser le
CPU les uns à la suite des autres (partage \emph{temporel}) ; et chaque
programme aura à sa disposition une partie de la mémoire principale, ou du
disque dur (partage \emph{spatial}).

Si deux programmes (ou plus) s'exécutent de manière concurrente sur le même
matériel, il faut s'assurer par exemple que les deux s'exécutent à peu près
aussi souvent, ou que l'un ne puisse pas écrire dans la mémoire de l'autre. Ce
sont des rôles du système d'exploitation.

Cela passe donc par un certain bridage des possibilités du programme : plutôt
que de le faire exécuter n'importe quel type d'instruction, il communique avec
le système d'exploitation. Bien que cela ait l'air d'une limitation, c'est aussi
bénéfique pour le programmeur puisque cela permet de définir des abstractions au
niveau du noyau.

Par exemple, si un programmeur veut copier des données depuis un CD-ROM vers la
mémoire principale, il devra interroger le bus SATA, interroger le lecteur sur
la présence d'un disque dans le lecteur, activer le moteur, calculer le numéro
de \emph{frame} des données sur le disque, demander la lecture, puis déclencher
une copie de la mémoire.

Si dans un autre cas il voulait récupérer des données depuis une mémorette USB,
il devrait interroger le bus USB, rechercher le bon numéro de périphérique, le
bon \emph{endpoint} dans celui-ci, lui appliquer une commande de lecture au bon
numéro de bloc, puis copier la mémoire.

Ces deux opérations, bien qu'elles aient le même but (copier de la mémoire
depuis un périphérique amovible), ne sont pas effectuées en pratique de la même
manière. C'est pourquoi le système d'exploitation fournit les notions de
fichier, lecteur, etc : le programmeur n'a plus qu'à utiliser des commandes de
haut niveau (``monter un lecteur'', ``ouvrir un fichier'', ``lire dans un
fichier'') et selon le type de lecteur, le système d'exploitation effectuera les
actions appropriées.

En résumé, un système d'exploitation est l'intermédiaire entre le logiciel et
le matériel, et en particulier assure les rôles suivants :

\begin{itemize}
\item
  Gestion des processus : un système d'exploitation peut permettre
  d'exécuter plusieurs programmes à la fois. Il faut alors orchestrer
  ces différents processus et les séparer en terme de temps et de
  ressources partagées.
\item
  Gestion de la mémoire : chaque processus, en plus du noyau, doit
  disposer d'un espace mémoire différent. C'est-à-dire qu'un processus
  ne doit pas pouvoir interférer avec un autre.
\item
  Gestion des périphériques : le noyau étant le seul code à s'exécuter
  en mode privilégié, c'est lui qui doit communiquer avec les
  périphériques matériels.
\item
  Abstractions : le noyau fournit aux programmes une interface unifiée,
  permettant de stocker des informations de la même manière sur un
  disque dur ou une clef USB (alors que l'accès se déroulera de manière
  très différente en pratique). C'est ici que la notion arbitraire de
  fichier sera définie, par exemple.
\end{itemize}

\section{Architecture Intel}

Dans cette section nous décrivons le fonctionnement d'un processeur utilisant
une architecture Intel 32 bits. Les morceaux de code seront écrit en syntaxe
AT\&T, celle que comprend l'assembleur GNU.

\subsection{Assembleur}

Pour faire des calculs, le processeur est composé de registres, qui sont des
petites zones de mémoire interne, et peut accéder à la mémoire principale.

Les registres principaux sont nommés \eax, \ebx, \ecx, \edx, \esi et \edi. Ils
peuvent être utilisés pour n'importe quel type d'opération, mais certains sont
spécialisés : par exemple il est plus efficace d'utiliser \eax en accumulateur,
ou \ecx en compteur. \esp est le pointeur de pile (\emph{stack pointer}), qui a
un rôle particulier dans les instructions \asminstr{push} et \asminstr{pop}.
\ebp sert de point de repère dans les appels de fonction. \eip contient
l'adresse de l'instruction courante.

Les calculs sont décrits sous forme d'une suite d'instructions. Chaque
instruction est composée d'une mnémonique et d'une liste d'opérandes. Les
mnénoniques (\asminstr{mov}, \asminstr{call}, \asminstr{sub}, etc) définissent
un type d'opération à appliquer sur les opérandes.

Ces opérandes peuvent être de plusieurs types :

\begin{itemize}
\item
  un nombre, noté \$4
\item
  le nom d'un registre, noté \texttt{\%eax}
\item
  une opérande mémoire, c'est à dire le contenu de la mémoire à une
  adresse effective. Cette adresse effective peut être exprimée de
  plusieurs manières :
  \begin{itemize}
  \item
    directement : \texttt{addr}
  \item
    indirectement : \texttt{(\%ecx)}. L'adresse effective est le contenu
    du registre.
  \item
    ``base + déplacement'' : \texttt{4(\%ecx)}. L'adresse effective est
    le contenu du registre plus le déplacement (4 ici).
  \end{itemize}
\end{itemize}

En pratique il y a des modes d'adressage plus complexes, et toutes les
combinaisons ne sont pas possibles, mais ceux-ci suffiront à décrire les
exemples suivants :

\begin{itemize}

\item \texttt{mov src, dst} copie le contenu de \texttt{src} dans \texttt{dst}.

\item \texttt{add src, dst} calcule la somme des contenus de \texttt{src} et
  \texttt{dst} et place ce résultat dans \texttt{dst}.

\item \texttt{push src} place \texttt{src} sur la pile, c'est à dire que cette
  instruction décrémente le pointeur de pile \esp de la taille de \texttt{src},
  puis place \texttt{src} à l'adresse mémoire de la nouvelle valeur \esp.

\item \texttt{pop src} réalise l'opération inverse : elle charge le contenu de
  la mémoire à l'adresse \esp dans \texttt{src} puis incrémente \esp de la
  taille correspondante.

\item \texttt{jmp addr} saute à l'adresse \texttt{addr} : c'est l'équivalent de
  \texttt{mov addr, \%eip}.

\item \texttt{call addr} sert aux appels de fonction : cela revient à
  \texttt{push \%eip} puis \texttt{jmp addr}.

\item \texttt{rest} sert à revenir d'une fonction : c'est l'équivalent de
  \texttt{pop \%eip}.

\end{itemize}

\subsection{Appel de fonction et conventions d'appel}

\begin{figure} % {{{ fig:stackframe
\begin{tikzpicture}
[ stack/.style={draw,shape=rectangle,minimum height=1cm,minimum width=5cm}
, bigbrace/.style={decorate,decoration={brace,amplitude=3mm}}
, >=latex
]

  \path (0,-8.5)   node [stack, minimum height=2cm] (gloc) {Locales de g}
  -- ++(0,1.5) node [stack] (gebp) {Cadre précédent}
  -- ++(0,1)   node [stack] (gret) {Retour de g}
  -- ++(0,1.5) node [stack, minimum height=2cm] (gpar) {Arguments de g}
  -- ++(0,1.5) node [stack] (floc) {Locales de f}
  -- ++(0,1)   node [stack] (febp) {Cadre précédent}
  -- ++(0,1)   node [stack] (fret) {Retour de f}
  -- ++(0,1)   node [stack] (fpar) {Arguments de f};

  \draw [->] (gebp) -- ++(4, 0) |- (febp.355);
  \draw [->] (febp.5) -- ++(1.5, 0) -- ++(0,4);

  \draw [dashed] (fpar.north west) -- ++ (0,2);
  \draw [dashed] (fpar.north east) -- ++ (0,2);
  \draw [dashed] (gloc.south east) -- ++ (0,-2);
  \draw [dashed] (gloc.south west) -- ++ (0,-2);

  \node at (5, 1) (ox) {};

  \node [text width=3cm] at (-0.8, 2)     {Adresses hautes (sens de \asminstr{pop})};
  \node [text width=3cm] at (-0.8, -10.5) {Adresses basses (sens de \asminstr{push})};

  \node [text width=3cm] at (-6, -3) (asmtext) {\tt
    call f\\
    ... \\
    f:  ... \\
    pushl \$2 \\
    pushl \$4 \\
    call g \\
    addl \$8, \%esp \\
    leave \\
    ret \\
    ... \\
    g: push \%ebp \\
    mov \%esp, \%ebp \\
    sub \$8, \%esp \\
    ... \\
    leave \\
    ret
    };

  \draw [->] (gret) -- ++ (-3, 0) |- ($ (asmtext.north west) + (2.5, -3) $);
  \draw [->] (fret) -- ++ (-3, 0) |- ($ (asmtext.north west) + (2.5, -0.5) $);

  \draw [<-] (gebp.west) -- ++ (-0.5,0) node [xshift=-4mm] {\ebp};
  \draw [<-] ($ (gloc.west) + (0,-0.5) $) -- ++ (-0.5,0) node [xshift=-4mm] {\esp};

  \draw [<-] ($ (asmtext.north west) + (2.5, -6.7) $) -- ++ (0.5,0) node [xshift=3mm] {\eip};

  \draw [bigbrace]
        ($ (gpar.north east) + (3, -0.1) $)
     -- node [auto,midway,xshift=3mm] {Cadre de g}
        ($ (gloc.south east) + (3, 0) $);

  \draw [bigbrace]
        ($ (fpar.north east) + (3, 0) $)
     -- node [auto,midway,xshift=3mm] {Cadre de f}
        ($ (floc.south east) + (3, 0.1) $);

\end{tikzpicture}

\caption[Cadres de pile]{ Cadres de pile. }
\label{fig:stackframe}
\end{figure} % }}}

Dans le langage d'assemblage, il n'y a pas de notion de fonction ; mais
\asminstr{call} et \asminstr{ret} permettent de sauvegarder et de restaurer une
adresse de retour, ce qui permet de faire un saut et revenir à l'adresse initiale.
Ce système permet déjà de créer des procédures, c'est-à-dire des fonctions sans
arguments ni valeur de retour.

Pour gérer ceux-ci, il faut que les deux morceaux (appelant et appelé) se
mettent d'accord sur une convention d'appel commune. La convention utilisée sous
GNU/Linux est appelée \emph{cdecl} et possède les caractéristiques suivantes :

\begin{itemize}
\item la valeur de retour d'une fonction est stockée dans \eax
\item \eax, \ecx et \edx peuvent être écrasés sans avoir à les sauvegarder
\item les arguments sont placés sur la pile (et enlevés) par l'appelant. Les
  paramètres sont empilés de droite à gauche.
\end{itemize}

Pour accéder à ses paramètres, une fonction peut donc utiliser un adressage
relatif à \esp. Cela peut fonctionner, mais cela complique les choses si elle
contient aussi des variables locales. En effet, les variables locales sont
placées sur la pile, au dessus des (c'est à dire, empilées après) paramètres,
augmentant le décalage.

Pour simplifier, la pile est organisée en cadres logiques : chaque cadre
correspond à un niveau dans la pile d'appels de fonctions. Si f appelle g, qui
appelle h, il y aura dans l'ordre sur la pile le cadre de f, celui de g puis
celui de h.

Ces cadres sont chainés à l'aide du registre \ebp : à tout moment, \ebp contient
l'adresse du cadre de l'appelant.

Prenons exemple sur la figure~\ref{fig:stackframe} : pour appeler
\texttt{g(4,2)}, \texttt{f} empile les arguments de droite à gauche.
L'instruction \asminstr{call g} empile l'adresse de l'instruction suivante sur
la pile puis saute au début de \texttt{g}.

Au début de la fonction, les trois instructions permettent de sauvegarder
l'ancienne valeur de \ebp, faire pointer \ebp à une endroit fixe dans le cadre
de pile, puis allouer 8 octets de mémoire pour les variables locales.

Dans le corps de la fonction \texttt{g}, on peut donc se référer aux variables
locales par \texttt{-4(\%ebp)}, \texttt{-8(\%ebp)}, etc, et aux arguments par
  \texttt{8(\%ebp)}, \texttt{12(\%ebp)}, etc.

À la fin de la fonction, l'instruction \asminstr{leave} est équivalente à
\texttt{mov \%ebp, \%esp} puis \texttt{pop \%ebp} et permet de défaire le cadre
de pile, laissant l'adresse de retour en haut de pile. Le \asminstr{ret} final
la dépile et y saute.

\subsection{Tâches, niveaux de privilèges}
\label{sec:taches}

Sans mécanisme particulier, le processeur exécuterait uniquement une suite
d'instruction à la fois. Pour lui permettre d'exécuter plusieurs tâches, un
système de partage du temps existe.

À des intervalles de temps réguliers, le système est programmé pour recevoir une
interruption. C'est une condition exceptionnelle (au même titre qu'une division
par zéro) qui fait sauter automatiquement le processeur dans une routine de
traitement d'interruption. Si à cet endroit le code sauvegarde les registres et
restaure un autre ensemble de registres, on peut exécuter plusieurs tâches de
manière entrelacée. Si l'alternance est assez rapide, cela peut donner
l'illusion que les programmes s'exécutent en parallèle. Comme l'interruption
peut survenir à tout moment, on parle de multitâche préemptif.

En plus de cet ordonnancement de processus, l'architecture Intel permet
d'affecter des niveaux de privilège à ces tâches, en restreignant le type
d'instructions exécutables, ou en donnant un accès limité à la mémoire aux
tâches de niveaux moins elevés.

Il y a 4 niveaux de privilèges (nommés aussi \emph{rings}) : le \ring{0} est le
plus privilégié, le \ring{3} le moins privilégié. Dans l'exemple précédent, on
pourrait isoler l'ordonnanceur de processus en le faisant s'exécuter en \ring{0}
alors que les autres tâches seraient en \ring{3}.

\begin{center}\rule{3in}{0.4pt}\end{center}


\begin{figure}
\begin{Verbatim}
+-------------------------------------+
|   Ring 3                            |
|                                     |
|       Programmes utilisateur        |
|                                     |
|                                     |
|   ^ | appels système                |
+---|-v-------------------------------+
|                                     |
|   Ring 0                            |
|                                     |
|       Noyau                         |
|                                     |
|   ^ | instructions privilégiées     |
+---|-v-------------------------------+

   Hardware
\end{Verbatim}

\begin{tikzpicture}

\begin{scope}

  %\path[clip] (0, 0.8) rectangle (5, -0.8);

  \foreach \x/\c in { 4/20, 3/30, 2/40, 1/50, 0/60 }
  \draw[fill=black!\c] ($ (-2,0) + (\x,0) $) circle (2 cm);

\end{scope}

%\draw [<-] (3.5, 0) -- ++(0,-2) node {Applications};
%\draw [<-] (0, 0) -- ++(0,-2) node {Noyau};

\end{tikzpicture}

\caption[Les différents \emph{rings}]{
  Les différents \emph{rings}. Seul le \ring{0} a accès au hardware
  via des instructions privilégiées. Pour accéder aux fonctionnalités du noyau,
  les programmes utilisateur doivent passer par des appels système.}
\label{fig:rings}
\end{figure}

La première protection nécessaire est d'isoler le système d'exploitation
lui-même des programmes qu'il va permettre d'exécuter. Pour ce faire, Le
processeur permet d'exécuter des tâches selon plusieurs niveaux de privilège,
aussi appelés \emph{rings} (figure~\ref{fig:rings}) : du \ring{3}, le
moins privilégié, jusqu'au \ring{0}, le plus privilégié. On peut configurer
le processeur de manière à ce que les instructions privilégiées (accès aux ports
d'entrée/sortie...) ne soient possibles qu'en \ring{0}. Bien que 4 niveaux
soient disponibles, Linux n'utilise que les \emph{rings} 0 et 3 : le noyau
lui-même en \ring{0} et les processus utilisateur en \ring{3}.

\subsection{Mémoire virtuelle}

À partir du moment où plusieurs processus s'exécutent de manière concurrente, un
problème d'isolation se pose : si un processus peut lire dans la mémoire d'un
autre, des informations peuvent fuiter ; et s'il peut y écrire, il peut
détourner l'exécution.

\todo{defs : intégrité, confidentialité}

Le mécanisme de mémoire virtuelle permet de donner à deux tâches une vue
différente de la mémoire : c'est à dire que vue de tâches différentes, une
adresse contiendra une valeur différente.

\begin{figure}
\begin{tikzpicture}
[ bit/.style={ draw
             , shape=rectangle
             , minimum width=4mm
             , minimum height=4mm
             , node distance=4mm
             }
, pstuff/.style={ draw
                , shape=rectangle
                , node distance=1cm
                , text width=2cm
                , text centered
                }
, >=latex
]

\node [bit] (b31) {};

% generate this ?
\node [bit, right of=b31] (b30) {}; \node [bit, right of=b30] (b29) {}; \node [bit, right of=b29] (b28) {};
\node [bit, right of=b28] (b27) {}; \node [bit, right of=b27] (b26) {}; \node [bit, right of=b26] (b25) {};
\node [bit, right of=b25] (b24) {}; \node [bit, right of=b24] (b23) {}; \node [bit, right of=b23] (b22) {};
\node [bit, right of=b22] (b21) {}; \node [bit, right of=b21] (b20) {}; \node [bit, right of=b20] (b19) {};
\node [bit, right of=b19] (b18) {}; \node [bit, right of=b18] (b17) {}; \node [bit, right of=b17] (b16) {};
\node [bit, right of=b16] (b15) {}; \node [bit, right of=b15] (b14) {}; \node [bit, right of=b14] (b13) {};
\node [bit, right of=b13] (b12) {}; \node [bit, right of=b12] (b11) {}; \node [bit, right of=b11] (b10) {};
\node [bit, right of=b10] (b9) {}; \node [bit, right of=b9] (b8) {}; \node [bit, right of=b8] (b7) {};
\node [bit, right of=b7] (b6) {}; \node [bit, right of=b6] (b5) {}; \node [bit, right of=b5] (b4) {};
\node [bit, right of=b4] (b3) {}; \node [bit, right of=b3] (b2) {}; \node [bit, right of=b2] (b1) {};
\node [bit, right of=b1] (b0) {};

\foreach \b in {0,11,12,21,22,31}
    \node [above of=b\b, node distance=5mm] {\small \b};

\newcommand{\mbrace}[3]{
    \draw [bigbrace]                      ($ (b#1.south east) + (-0.1,-0.3) $)
    -- node [midway, yshift=-2mm] (#3) {} ($ (b#2.south west) +  (0.1,-0.3) $);
} % mbrace

\mbrace{0}{11}{bpp};
\mbrace{12}{21}{bpt};
\mbrace{22}{31}{bpd};

\node[pstuff, below of=bpd] (pd) {Répertoire de pages};
\node[pstuff, below of=bpt] (pt) {Table de pages};
\node[pstuff, below of=bpp] (pp) {Page de 4 Kio};

\node[left of=pd, node distance=2cm] (cr) {\crtrois};

\draw [->] (bpp) -- (pp);
\draw [->] (bpt) -- (pt);
\draw [->] (bpd) -- (pd);

\draw [->] (cr) -- (pd);
\draw [->] (pd) -- (pt);
\draw [->] (pt) -- (pp);
\draw [->] (pp) -- ++ (3,0);

\end{tikzpicture}

\caption{Implémentation de la mémoire virtuelle}
\label{fig:pagetables}
\end{figure}

Ce mécanisme est contrôlé par valeur du registre \crtrois : les 10 premiers bits
d'une adresse sont un index dans le répertoire de pages qui commence à l'adresse contenue dans
\crtrois. À cet index, se trouve l'adresse d'une table de pages. Les 10 bits
suivants de l'adresse sont un index dans cette page, donnant l'adresse d'une
page de 4 kibi-octets (figure~\ref{fig:pagetables}).

\begin{figure} % fig:memoire-virtuelle {{{
\centering
\begin{tikzpicture}
  [every node/.style={node distance=2cm}]

  \node (u1) {};

  \node[right of=u1] (phy) {};

  \node[right of=phy] (u2) {};

  \draw (u1) rectangle ++(1,-6);
  \draw (phy) rectangle ++(1,-6);
  \draw (u2) rectangle ++(1,-6);
\end{tikzpicture}
\caption{Mécanisme de mémoire virtuelle.}
\label{fig:memoire-virtuelle}
\end{figure} % }}}

En ce qui concerne la mémoire, les différentes tâches ont une vision différente
de la mémoire physique : c'est-à-dire que deux tâches lisant à une même adresse
peuvent avoir un résultat différent. C'est le concept de mémoire virtuelle
(fig~\ref{fig:memoire-virtuelle}).

\begin{figure} % fig:memmap {{{
\centering
\fbox{
  \begin{tikzpicture}
  [scale=0.7
  ,user/.style={fill=black!20}
  ,kernel/.style={fill=black!70}
  ]

  % Memory zone
  %
  % #1 - start
  % #2 - end
  % #3 - color
  \newcommand{\mzone}[3]{
    \path[#3] (#1,0) rectangle (#2,1);
  }

  % Address label
  %
  % #1 - x position
  % #2 - text
  \newcommand{\alabel}[2]{
    \path (#1,1) -- ++(0,0.3) node [pos=1] {\small \tt #2};

  }

  % exec
  \mzone{0.5}{1}{user}

  % lib
  \mzone{2.5}{3.2}{user}

  % stack
  \mzone{3.7}{4}{user}

  % stack
  \mzone{5}{5.5}{user}

  % kernel
  \mzone{6}{8}{kernel}

  % contour
  \draw (0,0) rectangle (8,1);

  \alabel{0}{0}
  \alabel{6}{3 Go}
  \alabel{8}{4 Go}

\end{tikzpicture}

}

\caption[Espace d'adressage d'un processus]{L'espace d'adressage d'un processus.
En gris clair, les zones accessibles à tous les niveaux de privilèges : code du
programme, bibliothèques, tas, pile. En gris foncé, la mémoire du noyau,
réservée au mode privilégié.}

\label{fig:memmap}
\end{figure}
% }}}

\section{Cas de Linux}

Dans cette section, nous allons voir comment ses mécanismes sont implantés dans
le noyau Linux. Une description plus détaillée pourra être trouvée dans
\cite{UnderstandingTheLinuxKernel} par exemple.

Deux rings sont utilisés : en \ring{0}, le code noyau et en \ring{3}, le code
utilisateur.

Une notion de tâche similaire à celle décrite dans la section~\ref{sec:taches}
existe : elles s'exécutent l'une après l'autre, le changement s'effectuant sur
interruptions.

Pour faire appel aux services du noyau, le code utilisateur peut faire appel à
des appels systèmes, qui sont des fonctions exécutées par le noyau. Chaque tâche
doit donc avoir deux piles : une pile "unitisateur" qui sert pour l'application
elle-même, et une pile "noyau" qui sert aux appels système.

\subsection{Appels système}

Les programmes utilisateur s'exécutant en \ring{3}, ils ne peuvent pas
contenir d'instructions privilégiées, et donc ne peuvent pas accéder directement
au matériel (c'était le but !). Pour qu'ils puissent interagir avec le système
(afficher une sortie, écrire sur le disque...), le mécanisme des appels système
est nécessaire. Il s'agit d'une interface de haut niveau entre les \emph{rings}
3 et 0. Du point de vue du programmeur, il s'agit d'un ensemble de fonctions C
``magiques'' qui font appel au système d'exploitation pour effectuer des
opérations.

Prenons le cas de l'appel \texttt{getpid}, qui retourne le numéro de processus
courant. La bibliothèque C fournit une fonction du même nom :

\begin{Verbatim}
pid_t getpid(void);
\end{Verbatim}

A priori, rien de différent d'une fonction implantée directement en C. Par un
processus détaillé ci-après, cette fonction va invoquer la fonction, suivante,
définie dans le noyau (\texttt{kernel/timer.c}) :

\begin{Verbatim}
SYSCALL_DEFINE0(getpid)
{
        return task_tgid_vnr(current);
}
\end{Verbatim}



Le mécanisme de couplage entre ces deux fonctions est le suivant. Une
explication plus détaillée est disponible dans la documentation d'Intel
\cite{intelsys}.

Il y a bien une fonction \texttt{getpid} présente dans la bibliothèque C du
système. C'est la fonction qui est directement appelée par le programme. Cette
fonction commence par placer le numéro de l'appel système (noté
\texttt{\_\_NR\_getpid}, valant 20 ici) dans \eax, puis les arguments éventuels
dans les registres (\ebx, \ecx, \edx, \esi puis \edi). Une interruption
logicielle est ensuite déclenchée (\texttt{int 0x80}) : étant donné la
configuration du processeur, elle sera traitée en \ring{0}, à un point
d'entrée prédéfini (\texttt{arch/x86/kernel/entry\_32.S},
\texttt{ENTRY(system\_call)}).

L'exécution reprend donc en \ring{0}, avec une pile propre au processus.
Les valeurs des registres ont été préservées : elles sont mises sur la pile du
noyau. Le noyau examine le numéro d'appel système (toujours dans \eax) et
appelle la fonction correspondante dans le tableau \texttt{sys\_call\_table} :
les arguments sont en place sur la pile, comme dans le cas d'un appel de
fonction ``classique''. La fonction, nommée \texttt{sys\_getpid}, peut donc être
écrite en C.

Au retour de la fonction, la valeur de retour est placée à la place de \eax là
où les registres ont été sauvegardés sur la pile noyau. L'instruction
\texttt{iret} permet de restaurer les registres et de repasser en mode
utilisateur, juste après l'interruption. La fonction de la bibliothèque C peut
alors retourner au programme appelant.

% vim: spelllang=fr


%\chapter{Typage}

%Dans ce chapitre, nous explorons la notion de type dans les langages de
programmation. Tout d'abord pourquoi elle existe et en quoi elle aide à rendre
les programmes plus sûrs. Il y a autant de système de types que de langages de
programmation, donc nous présenterons ensuite une taxonomie de ces systèmes en
les regroupant par caractéristiques communes. Cette classification sera appuyée
par des exemples de code C\cite{AnsiC,KandR}, OCaml\cite{ocamlManual,DAOC},
Haskell\cite{haskell98,rwh}, Python\cite{pythonSite} et
Perl\cite{perlCamelBook}.

\section{Présentation et but}

Nous avons vu dans le chapitre~\ref{cha:os} qu'au niveau du langage machine, les
seules données qu'un ordinateur manipule sont des nombres. Selon les opérations
effectuées, ils seront interprétés comme des entiers, des adresses mémoires, ou
des caractères. Pourtant il est clair que certaines opérations n'ont pas de
sens: par exemple, ajouter deux adresses, ou déréférencer le résultat d'une
division sont des comportements qu'on voudrait pouvoir empêcher.

En un mot, le but du typage est de classifier les objets et de restreindre les
opérations possibles selon la classe d'un objet : "ne pas ajouter des pommes et
des oranges". Le modèle qui permet cette classification est appelé \emph{système
de types} et est en général constitué d'un ensemble de \emph{règles de typage},
comme "un entier plus un entier égale un entier".

\section{Taxonomie}

La définition d'un langage de programmation introduit la plupart du temps celle
d'un système de types. Il y a donc de nombreux systèmes de types différents,
dont nous pouvons donner une classification sommaire.

\subsection{Dynamique, statique, mixte}

Il y a deux grandes familles de systèmes de types, selon quand se fait la
vérification de types. On peut en effet l'effectuer au moment de l'exécution, ou
au contraire prévenir les erreurs à l'exécution en la faisant au moment de la
compilation (ou avant l'interprétation).

\subsubsection{Typage dynamique}

La première est le typage \emph{dynamique}. Pour différencier les différents
types de données, on ajoute une étiquette à chaque valeur. Dans tout le
programme, on ne manipulera que des valeurs étiquetées, c'est à dire des couples
(donnée, nom de type). Si on veut réaliser l'opération $(\hexa{00000001}, Int) +
(\hexa{0000f000}, Int)$, on vérifie tout d'abord qu'on peut réaliser l'opération
$+$ entre deux $Int$. Ensuite on réalise l'opération elle même, qu'on étiquette
avec le type du résultat : $(\hexa{0000f001}, Int)$. Si au contraire on tente
d'ajouter deux adresses $(\hexa{2e8d5a90}, Addr) + (\hexa{76a5e0ec}, Addr)$, la
vérification échoue et l'opération s'arrête avec une erreur.

\begin{figure}
  \insertcode{typage-dynamique.py}
  \caption{Session Python présentant le typage dynamique}
  \label{fig:typage-dynamique}
\end{figure}

La figure~\ref{fig:typage-dynamique} est une session interactive Python qui
illustre le typage dynamique. Chaque variable, en plus de sa valeur, possède une
étiquette qui permet d'identifier le type de celle-ci. On peut accéder au type
d'une valeur \texttt{x} avec la construction \texttt{type(x)}.

Au moment de réaliser une opération comme \texttt{+}, l'interpréteur Python
vérifie les types des opérandes : s'ils sont compatibles, il créé une valeur de
résultat, et sinon il lève une exception.

Comme l'implémentation elle-même des fonction a accès aux informations de type,
elle peut faire des traitements particuliers. Par exemple, pour l'addition de
\texttt{a} (de type entier) et de \texttt{c} (de type flottant), la fonction
d'addition va d'abord convertir \texttt{a} en flottant, puis réaliser l'addition
dans le domaine des flottants.

\subsubsection{Typage statique}

Le typage dynamique est simple à comprendre puisque toute les vérifications se
font dans la sémantique dynamique (ie, à l'exécution). C'est à double tranchant:
d'une part, cela permet d'être plus flexible, mais d'autre part, cela permet à
des programmes incorrects d'être exécutés.

On peut lire le code source d'un programme et essayer de "deviner" quels seront
les types des différentes expressions. Dans certains cas, cela n'est pas
possible (fig~\ref{fig:nontypable}) ; mais lorsqu'on peut conclure cela élimine
la nécessité de faire les tests de type dynamiques car on a réalisé le typage
\emph{statiquement}.

Bien sûr, deviner n'est pas suffisant : il faut formaliser cette analyse. Dans
le cas dynamique, ce sont les fonctions elles-mêmes qui réalisent les tests de
type et qui appliquent des règles comme "si les arguments ont pour type
\texttt{int} alors la valeur de retour a pour type \texttt{int}" : la fonction
qui réalise ce test sur les valeurs. Dans le cas statique, c'est le compilateur
(ou l'interpréteur) qui réalise ce test sur les expressions non évaluées. En
appliquant de proche en proche un ensemble de règles (liées uniquement au
langage de programmation), on finit par associer à chaque sous-expression du
programme un type.

Benjamin Pierce résume cette approche dans cette définition très générale :

\begin{definition}[Système de types]
Un système de types est une méthode syntaxique tractable qui vise à prouver
l'absence de certains comportements de programmes en classifiant leurs phrases
selon le genre de valeurs qu'elles produisent. \cite{TAPL}
\end{definition}

\begin{figure}
  \insertcode{non-typable.py}
  \caption{Fonction Python non typable statiquement.}
  \label{fig:nontypable}
\end{figure}

A première vue, cela semble moins puissant que le typage dynamique : en effet,
il existe des programmes qui s'exécuteront sans erreur de type mais sur lesquels
le typage statique ne peut s'appliquer. Dans la figure~\ref{fig:nontypable}, on
peut voir par une simple analyse de cas que si on fournit un booléen à
\texttt{f}, elle retourne un entier. Mais selon la valeur de \texttt{b}, la
variable \texttt{x} contiendra une valeur de type entier ou fonction.

Même si cet exemple est construit artificiellement, il illustre le problème
suivant : les types statiques demandent un certain effort et au programmeur et
au compilateur. Mais Dans le cas où le typage statique est possible, les
garanties sont importantes : les valeurs portées par une variable auront
toujours le même type. Par voie de conséquence, la vérification dynamique de
types réussira toujours, et on peut la supprimer. Il est également possible
de supprimer toutes les étiquettes de typage : on parle de \emph{type erasure}.
Une conséquence heureuse de cette suppression est que l'exécution de ce
programme se fera de manière plus rapide.

Connaître les types à la compilation permet aussi de réaliser plus
d'optimisations. Par exemple, en Python, considérons l'expression \texttt{y = x
- x}. Sans information sur le type de \texttt{x}, aucune simplification n'est
possible : l'implémentation de la différence sur ce type est une fonction
quelconque, sans propriétés particulières \emph{a priori}. Si au contraire, on
sait que \texttt{x} est un entier, on peut en déduire que \texttt{y = 0}, sans
réaliser la soustraction (si c'était la seule utilisation de \texttt{x}, le
calcul de \texttt{x} aurait alors pu être éliminé).

\subsection{Fort, faible, sound}

\begin{figure}
  \insertcode{cast.java}
  \caption{Transtypage en Java}
  \label{fig:javacast}
\end{figure}

Si un système de types statique permet d'éliminer totalement la nécessité de
réaliser des tests de typage, on dit qu'il est \emph{fort}. Mais ce n'est que
rarement le cas. En effet, il peut y avoir des constructions au sein du langage
qui permettent de contourner le système de types, comme un opérateur de
transtypage\ref{fig:javacast}. À l'exécution, une erreur de types est levée :

\begin{Verbatim}
Exception in thread "main" java.lang.ClassCastException:
    java.lang.Integer cannot be cast to java.lang.Float
        at Cast.main(Cast.java:5)
\end{Verbatim}

\clearpage

\wip{}

\subsection{Polymorphisme}

Dans le cas du typage statique, restreindre une opération à un seul type de
données peut être assez restrictif.

Par exemple, quel doit être le type d'une fonction qui trie un tableau ?

\begin{figure}
\centering
\begin{tikzpicture}
  [ ptype/.style={draw=black,shape=rectangle}
  , child/.style={ptype, xshift=2cm}
  , gen0/.style={       node distance=2cm}
  , gen1/.style={child, node distance=1cm}
  , gen2/.style={child, node distance=5mm}
  ]

  \node[ptype]                     (mono) {Mono};
  \node[ptype,below of=mono, gen0] (poly) {Poly};
  \node[above right of=poly, gen1] (univ)  {Universel};
  \node[above right of=univ, gen2] (param) {Paramétrique};
  \node[below right of=univ, gen2] (inclu) {Inclusion};
  \node[below right of=poly, gen1] (adho)  {Ad-hoc};
  \node[above right of=adho, gen2] (over)  {Surcharge};
  \node[below right of=adho, gen2] (coer)  {Coercition};

  \draw (poly) -- (univ);
  \draw (poly) -- (adho);

  \draw (univ) -- (param);
  \draw (univ) -- (inclu);

  \draw (adho) -- (over);
  \draw (adho) -- (coer);

      %node {Mono}
      %node {Poly}
          %node {Universal}
              %node {Param}
              %node {Inclusion}
          %node {Ad-Hoc}
              %node {Surcharge}
              %node {Coercition}

\end{tikzpicture}

\caption{Les différents types de polymorphisme.}
\label{fig:types-de-polymorphisme}
\end{figure}

\subsubsection{Monomorphisme}

Une première solution peut être de forcer des types concrets, c'est à dire
qu'une même fonction ne pourra s'appliquer qu'à un seul type de données.

Il est confortable pour le programmeur de n'avoir à écrire un algorithme qu'une
seule fois, indépendamment du type des éléments considérés.

Il existe deux grandes classes de systèmes de types introduisant du
polymorphisme.

\subsubsection{Polymorphisme universel}

% = {paramétrique, par inclusion}

Le polymorphisme est dit universel si toute fonction générique peut s'appliquer
à n'importe quel type.

\subsubsection{Polymorphisme ad-hoc}

% = {par surcharge, par coercition}

Le polymorphisme est \emph{ad-hoc} si les fonctions génériques ne peuvent
s'appliquer qu'à un ensemble de types prédéfini.

\subsubsection{Polymorphisme paramétrique}

\todo{Historique + citer le papier de Milner sur le polymorphisme}
\cite{Milner78}

\cite{PascalNoEscape}

\begin{figure}
  \insertcode{listappend.ml}
  \caption{Fonction de concaténation de listes en OCaml.}
  \label{fig:listappend}
\end{figure}

La fonction de la figure~\ref{fig:listappend} n'opère que sur la structure du type
liste (en utilisant ses constructeurs \texttt{{[}{]}} et \listcons ainsi que
le filtrage) : les éléments de \texttt{lx} et \texttt{ly} ne sont pas manipulés
à part pour les transférer dans le résultat.

Moralement, cette fonction est donc indépendante du type de données contenu dans
la liste : elle pourra agir sur des listes de n'importe quel type d'élément.

Plutôt qu'un type, on peut lui donner le \emph{schéma de types} suivant :

\[
  \textrm{append} : \forall a . a \textrm{list}
                             -> a \textrm{list}
                             -> a \textrm{list}
\]

C'est à dire que \texttt{append} peut être utilisé avec n'importe quel type
concret \texttt{a} en substituant les variables quantifiées (on parle d'
\emph{instanciation}).

\subsubsection{Polymorphisme par sous-typage}

\todo{héritage,sous-typage,classe,méthode,héritage multiple,late binding,Liskov}

Certains langages définissent la notion de sous-typage. C'est une relation
d'ordre partiel sur les types, qui modélise la relation "est un". Chaque
sous-classe peut redéfinir le comportement de chaque méthode de ses
superclasses.

\subsubsection{Polymorphisme par surcharge}

Considérons l'opération d'addition : \texttt{+}. On peut considérer que certains
types l'implémentent, et pas d'autres : ajouter deux flottants ou deux entiers a
du sens, mais pas ajouter deux pointeurs.

On dira que \texttt{+} est \emph{surchargé}. À chaque site d'appel, il faudra
\emph{résoudre la surcharge} pour déterminer quelle fonction appeler.

Cela rend l'inférence de types \todo{introduire l'inférence plus haut}
impossible dans le cas général, puisque certaines constructions sont ambigües.

\begin{figure}
  \insertcode{showread.hs}
  \caption{Cas d'ambigüité avec de la surcharge ad-hoc.}
  \label{fig:showread}
\end{figure}

Dans le code Haskell de la figure~\ref{fig:showread}, \texttt{show} peut
s'appliquer à toutes les valeurs de types "affichables" et renvoie une
représentation textuelle. \texttt{read} réalise le contraire avec les types
"lisibles".

Lorsqu'on compose ces deux fonctions, le type de la valeur intermédiaire est
capital puisqu'il détermine les instances de \texttt{show} et \texttt{read} à
utiliser.

\subsubsection{Polymorphisme par coercition}

\subsubsection{Polymorphisme d'ordre supérieur}

\begin{verbatim}
g f = (f true, f 2)
\end{verbatim}

\[
g : (\forall a . a -> a) -> (bool * int)
\]

Pas inférable (annotations nécessaires).

\subsection{Expressivité, garanties, types dépendants}

\section{Exemples}

\subsection{Faible dynamique : Perl}
\subsection{Faible statique : C}
\subsection{Fort dynamique : Python}
\subsection{Fort statique : OCaml}
\subsection{Fort statique à effets typés : Haskell}
\subsection{Theorem prover : Coq}

% vim: spelllang=fr


\chapter{Analyses statiques existantes}

\label{cha:etatdelart}
Dans ce chapitre, nous présentons un tour d'horizon des techniques existantes
permettant d'analyser des programmes. En particulier, on s'intéresse à la
propriété de sécurité décrite dans le chapitre~\ref{cha:os}, mais on ne se
limite pas à celle-ci: il est également intéressant de considérer des analyses
développées pour d'autres propriétés, celles-ci pouvant potentiellement
s'adapter.

L'analyse statique de programmes est un sujet de recherche actif depuis
l'apparition de la science informatique. On commence par en présenter une
classification, puis on montrera des exemples pertinents permettant d'analyser
du code système ou embarqué.

\section{Taxonomie}

\paragraph{Techniques statiques et dynamiques}

L'analyse peut être faite au moment de la compilation, ou au moment de
l'exécution. En général on peut obtenir des informations plus précises de
manière dynamique, mais cela ne prend en compte que les parties du programme qui
seront vraiment exécutées. Un autre problème des techniques dynamiques est qu'il
est souvent nécessaire d'instrumenter l'environnement d'exécution (ce qui ---
dans le cas où cela est possible --- peut se traduire par un impact en
performances). L'approche statique, en revanche, nécessite de construire à
l'arrêt une carte mentale du programme, ce qui n'est pas toujours possible dans
certains langages.

Les techniques dynamiques sont néanmoins les plus répandues, puisqu'elles sont
plus simples à mettre en œuvre et permettent de trouver des erreurs pendant le
processus de développement. De plus, on peut considérer qu'un programme avec une
forte couverture par les tests a de grandes chances d'être correct pour toutes
les entrées. C'est par exemple ce qui se passe dans l'avionique civile.

Mais pour s'assurer de la correction d'un programme, on ne peut pas s'appuyer
uniquement sur les tests --- ou de manière générale sur des analyses dynamiques
--- car il est impossible d'étudier l'ensemble complet de tous les comportements
possibles. Par exemple, si un bug se présente lors d'une interaction entre deux
composants qui n'a pas été testée, il passera inaperçu. Pour cette raison, la
plupart des analyses présentées ici sont statiques.

\paragraph{Cohérence et complétude}

Le but d'une analyse statique est de catégoriser les programmes selon s'ils
satisfont ou non un ensemble de propriétés fixées à l'avance. Malheureusement,
cela n'est que rarement possible car l'ensemble des valeurs possibles lors de
l'exécution d'un programme n'est pas un ensemble calculable (théorème de
Rice~\cite{rice}). Autrement dit, il ne peut exister une procédure de décision
prenant un programme et le déclarant correct ou incorrect.

Il n'est donc pas possible d'écrire un analyseur statique parfait, ne se
trompant jamais. Toute technique statique va donc de se retrouver dans au moins
un des cas suivants:

\begin{itemize}
\item
  un programme valide (pour une propriété donnée) est rejeté: on parle de
  \emph{faux positif}.
\item
  un programme invalide n'est pas détecté: on parle de
  \emph{faux négatif}.
\end{itemize}

En général, et dans notre cas, on préfère s'assurer que les programmes acceptés
possèdent la propriété recherchée, quitte à en rejeter certains. C'est
l'approche que nous retiendrons. Tolérer les faux négatifs n'est cependant pas
toujours une mauvaise idée. Par exemple, si le but est de trouver des
constructions dangereuses dans les programmes, on peut signaler certains cas qui
empiriquement valent d'être vérifiés manuellement.

Par ailleurs la plupart des techniques ne concernent que les programmes qui
terminent. On étudie donc la correction, ou les propriétés des termes
convergents. Prouver automatiquement que l'exécution ne boucle pas est une
propriété toute autre qui n'est pas ici considérée.

\section{Méthodes syntaxiques}

L'analyse la plus simple consiste à traiter un programme comme du texte, et à y
rechercher des motifs dangereux. Ainsi, utiliser des outils comme \texttt{grep}
permet parfois de trouver un grand nombre de vulnérabilités~\cite{SpenderGrep}.

On peut continuer cette approche en recherchant des motifs mais en étant
sensible à la syntaxe et au flot de contrôle du programme. Cette notion de
\emph{semantic grep} est présente dans l'outil Coccinelle
\cite{coccinelle09,coccinelle11}: on peut définir des
\emph{patches sémantiques} pour détecter ou modifier des constructions
particulières.

Ces techniques sont utiles parce qu'elles permettent de plonger rapidement dans
le code, en identifiant par exemple des appels à des fonctions dangereuses. En
revanche, cela n'est utile que lorsque les propriétés que l'ont recherche sont
très locales. Elles offrent également peu de garantie (ce n'est pas parce qu'une
construction dangereuse n'est pas présente que le code est correct).

\section{Interprétation abstraite}

L'interprétation abstraite est une technique d'analyse générique qui permet de
simuler statiquement tous les comportements d'un programme
\cite{Cousot77,Cousot92-1}. Un exemple d'application est de calculer les bornes
de variation des variables pour s'assurer qu'aucun débordement de tableau n'est
possible~\cite{AllamigeonHymansSSTIC07}.

L'idée est d'associer à chaque ensemble concret de valeurs, une représentation
abstraite. Sur celle-ci, on peut définir des opérations indépendantes de la
valeur exacte des données, mais préservant l'abstraction
(figure~\ref{fig:dom-sig}). Par exemple, les règles comme
``$-$'' $×$ ``$-$'' $=$ ``$+$''
définissent le domaine abstrait des signes. Les domaines ont une structure de
treillis, c'est-à-dire qu'ils possèdent les notions d'ordre partiel et d'union
de valeurs. En calculant les valeurs extremales d'une variable, on obtient le
domaine des intervalles.

\begin{figure}%{{{
\centering

\begin{minipage}{0.4\textwidth}
  \begin{tikzpicture}
  \node at (1, 2) (t) {$\top$};
  \node at (0, 1) (m) {$-$};
  \node at (1, 1) (z) {$0$};
  \node at (2, 1) (p) {$+$};
  \node at (1, 0) (b) {$\bot$};
  \draw (t) -- (p) -- (b);
  \draw (t) -- (m) -- (b);
  \draw (t) -- (z) -- (b);
  \end{tikzpicture}
\end{minipage}
\begin{minipage}{0.4\textwidth}
  \begin{align*}
  γ~(-) = & ℝ^- \\
  γ~(0) = & \{0\} \\
  γ~(+) = & ℝ^+ \\
  \end{align*}
\end{minipage}

\caption{Domaine des signes}
\label{fig:dom-sig}
\end{figure}%}}}

De tels domaines ne capturent aucune relation entre variables. Ils sont dits non
relationnels. Lorsque plusieurs variables sont analysées en même temps, utiliser
de tels domaines consiste à considérer un produit cartésien d'ensembles
abstraits (figure~\ref{fig:dom-cartesien}).

Des domaines abstraits plus précis permettent de retenir celles-ci. Pour ce
faire, il faut modéliser l'ensemble des valeurs des variables comme un tout.
Parmi les domaines relationnels courants on peut citer: le domaine des
polyèdres, permettant de retenir tous les invariants affines entre fonctions
(figure~\ref{fig:dom-poly}); le domaine des zones, permettant de représenter
des relations affines de la forme $v_i - v_j \le c$ (figure~\ref{fig:dom-zones})
; ou encore le domaine des octogones qui est un compromis entre les polyèdres et
les zones. Il permet de représenter les relations $\pm v_i \pm v_j \le c$
(figure~\ref{fig:dom-octo}).

\begin{figure}%{{{

  \centering

  %{{{
  \subbottom[Domaine non relationnel]{
\label{fig:dom-cartesien}
      \begin{tikzpicture}
      \draw[->] (0,-0.5) -- (0,3);
      \draw[->] (-0.5,0) -- (4,0);
%
      \draw[fill=red!30] (1,1) rectangle (3,2);
%
      \draw[dashed] (3,1) -- (3,0);
      \draw[dashed] (1,1) -- (1,0);
%
      \draw[dashed] (1,2) -- (0,2);
      \draw[dashed] (1,1) -- (0,1);
%
      \draw[line width=3pt] (0,1) -- (0,2);
      \draw[line width=3pt] (3,0) -- (1,0);
%
      \end{tikzpicture}
  }%}}}
  \subbottom[Domaine des polyèdres]{%{{{
\label{fig:dom-poly}
%
    \begin{tikzpicture}[scale=0.5]
    \path[use as bounding box] (0,0) rectangle (8, 6);
%
    \draw[fill=red!30] (1,2) -- (2,4) -- (4,5) -- (6,5) -- (7,3) -- (5,1) -- cycle;
%
    \end{tikzpicture}
  }%}}}

  \subbottom[Domaine des zones]{%{{{
\label{fig:dom-zones}
%
    \begin{tikzpicture}
    \path[use as bounding box] (0,0) rectangle (4,4);
%
    \draw[fill=red!30] (1,2) -- (2,3) -- (3,3) -- (3,2) -- (2,1) -- (1,1) -- cycle;
%
    \end{tikzpicture}
  }%}}}
  \subbottom[Domaine des octaèdres]{%{{{
\label{fig:dom-octo}
%
    \begin{tikzpicture}
    \path[use as bounding box] (0,0) rectangle (4,4);
%
    \draw[fill=red!30] (1,1.5) -- (1,2) -- (2,3) -- (2.5,3) -- (3,2.5) -- (3,2) -- (2,1) -- (1.5,1) -- cycle;
%
    \end{tikzpicture}
  }%}}}

  \caption{Quelques domaines abstraits}
\label{fig:dom-abstraits}
\end{figure}%}}}

En plus des domaines numériques, il est nécessaire d'employer des domaines
spécialisés dans la modélisation de la mémoire. Cela est nécessaire pour pouvoir
prendre en compte les pointeurs. Par exemple, on peut représenter un pointeur
par un ensemble de variables possiblement pointées et une valeur abstraite
représentant le décalage (\emph{offset}) du pointeur par rapport au début de la
zone mémoire. Cette valeur peut elle-même être abstraite par un domaine
numérique.

Au delà des domaines eux-mêmes, l'analyse se fait sous forme d'un calcul de
point fixe. La manière la plus simple est d'utiliser un algorithme de
\emph{liste de travail}, décrit par exemple dans~\cite{tapsoft95}. Les
raffinement en revanche sont nombreux.

Dès~\cite{Cousot77} il est remarqué que la terminaison de l'analyse n'est
assurée que si le treillis des valeurs abstraites est de hauteur finie, ou qu'un
opérateur d'élargissement (\emph{widening}) $\nabla$ est employé. L'idée est
qu'une fois qu'on a calculé quelques termes d'une suite croissante, on peut
réaliser une projection de celle ci. Par exemple, dans le domaine des
intervalles, $[0;2]~\nabla~[0;3] = [0;+\infty[$. On atteint alors un point fixe
mais qui est plus grand que celui qu'on aurait obtenu sans cette accélération:
on perd en précision. Pour en gagner, on peut redescendre sur le treillis des
points fixe avec une suite d'itérations décroissantes~\cite{granger}. Dans
l'itération de point fixe, il est possible d'obtenir les résultats de manière
plus efficace en choisissant un ordre particulier dans les calculs des
sous-itérations\cite{policy}.

En termes d'ingéniérie logicielle, implanter un analyseur statique est un défi
en soi. En plus des domaines abstraits, d'un itérateur, il faut traduire le code
source à analyser dans un langage, et traduire les résultats de l'analyse en un
ensemble d'alarmes à présenter à l'utilisateur.

Cette technique est très puissante: si un interpréteur abstrait \emph{sound}
(c'est-à-dire réalisant une surapproximation) analyse un programme et ne renvoie
pas d'erreur, alors on a prouvé que le programme est correct (par rapport aux
propriétés que vérifient les domaines abstraits). Cela a été appliqué avec
succès avec les analyseurs Astrée~\cite{Astree04,Astree05,AstreeScale} chez
Airbus ou CGS~\cite{cgs} à la NASA par exemple.

Cependant, ces analyses sont difficiles à mettre en œuvre. Avec des domaines
abstraits classiques comme ceux présentés ci-dessus, les premières analyses
peuvent remonter un nombre prohibitif de fausses alarmes. Pour ``aider''
l'analyse, il faut soit annoter le code soit développer des domaines abstraits
\emph{ad hoc} au programme à analyser.

Il existe également des analyseurs statiques par interprétation abstraite qui ne
sont pas \emph{sound}, c'est-à-dire qui peuvent manquer des comportements
erronés. Leur approche est plus d'aider le programmeur à détecter certains types
de bugs pendant le développement. On citer l'exemple de
Coverity~\cite{coverityBillion}, qui publie régulièrement des rapports de
qualité sur certains logiciels open source. Néanmoins, de part leur aspect non
\emph{sound}, les analyses réalisées ne peuvent pas être assimilées à de la
vérification formelle en tant que telle.

\section{Typage}

Le typage, introduit dans la section~\ref{sec:types-intro}, peut aussi être
utilisé pour la vérification de programmes. On peut le voir comme une manière de
catégoriser les types de données manipulés par la machine, mais également à plus
au niveau comme une manière d'articuler les différents composants d'un
programme. Mais on peut aussi programmer avec les types, c'est-à-dire utiliser
le compilateur (dans le cas statique) ou l'environnement d'exécution (dans le
cas dynamique) pour vérifier des propriétés écrites par le programmeur.

\paragraph{Systèmes ad hoc}

Les systèmes de types les plus simples expriment des contrats esssentiellement
liés à la sûreté d'exécution, pour ne pas utiliser des valeurs de types
incompatibles entre eux. Mais il est possible d'étendre le langage avec des
annotations plus riches: par exemple en vérifiant statiquement que des listes
ne sont pas vides\cite{lightweight-static-capabilities}, ou dans le domaine de
la sécurité, d'empêcher des fuites d'information~\cite{LZ06a}.

\paragraph{Qualificateurs de types}

Dans le cas particulier des vulnérabilités liées à une mauvaise utilisation de
la mémoire, les développeurs du noyau Linux ont ajouté un système d'annotations
au code source. Un pointeur peut être décoré d'une annotation
\texttt{\_\_kernel} ou \texttt{\_\_user} selon s'il est sûr ou pas. Celles-ci
sont ignorées par le compilateur, mais un outil d'analyse statique ad-hoc nommé
Sparse~\link{sparse} peut être utilisé pour détecter les cas les plus simples
d'erreurs. Il demande aussi au programmeur d'ajouter de nombreuses annotations
dans le programme.

Cette solution se rapproche de la solution décrite dans ce manuscrit. Ce système
d'annotations sur les types a été formalisé sous le nom de \emph{qualificateurs
de types}~\cite{toplas-quals}: chaque type peut être décoré d'un ensemble de
qualificateurs (à la manière de \texttt{const}), et des règles de typage
permettent d'établir des propriétés sur le programme.

Plus précisément, les jugements de typage de la forme $Γ ⊢ e : t$ sont remplacés
par des jugements de typage qualifiés $Γ ⊢ e : t~q$. Les qualificateurs $q$
permettent d'exprimer plusieurs jugements. Par exemple, pointeurs on peut
étudier le fait qu'une variable soit constante ou pas, que sa valeur soit connue
à la compilation, ou encore qu'elle puisse être nulle ou pas. La spécificité de
ce système est que les qualificateurs sont ordonnés, du plus spécifique au moins
spécifique, et que l'on forme alors un treillis à partir de ces informations.
Partant des deux caractéristiques précédentes, on forme le treillis de la
figure~\ref{fig:cqual-treillis}. Le qualificateur \texttt{const} désigne les
données dont la valeur change au cours de l'exécution; \texttt{dynamic} celles
qui peuvent être connues à la compilation; et \texttt{nonzero} celles qui ne
peuvent jamais être nulles. Le cube sur lequel se trouvent correspond à une
relation d'ordre, du plus spécifique (en bas) au plus général (en haut).

\begin{figure}
\centering
\begin{tikzpicture}

\node                                (A) {\texttt{dynamic}};
\node[below of=A]                    (C) {\texttt{const dynamic}};
\node[left of=C, node distance=4cm]  (B) {ø};
\node[right of=C, node distance=4cm] (D) {\texttt{dynamic nonzero}};
\node[below of=B]                    (E) {\texttt{const}};
\node[below of=C]                    (F) {\texttt{nonzero}};
\node[below of=D]                    (G) {\texttt{const dynamic nonzero}};
\node[below of=F]                    (H) {\texttt{const nonzero}};

\foreach \s/\d in {H/E,H/F,H/G,E/B,E/C,F/B,F/D,G/C,G/D,B/A,C/A,D/A}
  \draw[->] (\s) -- (\d);

\end{tikzpicture}

\caption{Treillis de qualificateurs}
\label{fig:cqual-treillis}
\end{figure}

Cette relation d'ordre $\preceq$ entre qualificateurs, induit une relation de
sous-typage $\sqsubseteq$ entre les types qualifiés: si $q \preceq q'$, alors
$t~q \sqsubseteq t~q'$.

Ces analyses ont été implantée dans l'outil CQual. Ce système peut servir à
inférer les annotations \texttt{const}~\cite{pldi99}, à l'analyse de souillure
pour les chaînes de format~\cite{usenix01} (pouvant poser des problèmes de
sécurité~\cite{format-string-attacks}) et des propriétés dépendantes du flot de
contrôle, comme des invariants sur les verrous~\cite{pldi02}, à rapprocher du
concept de \emph{typestates}~\cite{tse12-typestate}. Il a également été appliqué
à la classe de vulnérabilités sur les pointeurs utilisateurs dont il est ici
l'objet~\cite{cquk-usenix04}.

Cette approche étant assez proche de la nôtre, la conclusion de la
partie~\ref{part:lang} (page~\pageref{cha:ccl2}) sera dédiée à un comparatif
entre les deux.

%TODO[E]: dis tout de suite un mot sur pourquoi c'est mieux chez toi

Ces techniques de typage sont séduisantes parce qu'elles sont en général simples
à mettre en place: à l'aide d'un ensemble de règles, on attribue un type à
chaque expression. Si le typage se termine sans erreur, alors on est assuré de
la correction du programme (par rapport aux propriétés capturées par le système
de types).

Le typage statique peut également être implanté de manière efficace. La plupart
des systèmes de types peuvent être vérifiés en pratique dans un temps linéaire
en la taille du programme considéré. L'inférence, en revanche, peut parfois
``exploser'' en temps quadratique. % TODO check le coup du quadratique

\section{Analyse de code système}

Les logiciels système demandent des garanties de sécurité et de fiabilité. Pour
les assurer, de nombreuses analyses \emph{ad-hoc} ciblent les noyaux de systèmes
d'exploitation.

Ajouter un système de types forts à C est l'idée centrale de
CCured~\cite{ccured-toplas}. Dans les cas où il n'est pas possible de conclure,
des vérifications à l'exécution sont ajoutées. Cependant, cela nécessite une
instrumentation dynamique qui se paye en performances.

Saturn~\cite{paste07} est un système pour analyser du code système écrit en C.
Il traite le problème des pointeurs utilisateur en utilisant une analyse de
forme ``pointe-sur''~\cite{oakland08}. Comme l'interprétation abstraite, son but
est d'être très précis, au détriment d'un temps de calcul important dans
certains cas.

%TODO[E]: à préciser
% TODO ccl

\section{Langages sûrs}

Une autre approche est de concevoir un langage à la fois bas niveau et sûr,
permettant d'exprimer des programmes proches de la machine tout en interdisant
les constructions dangereuses.

Le langage Cyclone~\cite{cyclone-safety} est conçu comme un C ``sûr''. Afin
d'apporter plus de sûreté au modèle mémoire de C, des tests dynamiques sont
ajoutés, par exemple aux endroits ou des conversions implicites peuvent poser
problème. Le langage se distingue par le fait qu'il possède plusieurs types de
pointeurs: des pointeurs classiques (\texttt{int *}), des pointeurs ``jamais
nuls'' (\texttt{int @}; un test à l'exécution est alors inséré), et des
``pointeurs lourds'' (\verb!int ?!; qui contiennent des informations sur la
zone mémoire pointée). L'arithmétique des pointeurs n'est autorisée que sur ces
derniers, rendant impossibles les débordements de tableaux (ceux-ci étant
détectés au pire à l'exécution). Le problème des \emph{dangling pointers} est
résolu en utilisant un système de régions~\cite{cyclone-regions}, inspiré des
travaux de Jouvelot, Talpin et Tofte~\cite{jfp92,ToTa1993,popl94}. Cela permet
d'interdire statiquement les constructions où l'on déréférence un pointeur
faisant référence à une région de mémoire qui n'est plus allouée.

Le langage Rust~\link{rust} développé par Mozilla prend un approche similaire en
distinguant plusieurs types de pointeurs pour gérer la mémoire de manière plus
fine. Les \emph{managed pointers} (notés \texttt{@int}) utilisent un
ramasse-miette pour libérer la mémoire allouée lorsqu'ils ne sont plus
accessibles. Les \emph{owning pointers} (notés \texttt{\textasciitilde{}int})
décrivent une relation 1 à 1 entre deux objets, comme les
\texttt{std::unique\_ptr} de C++: la mémoire est libérée lorsque le pointeur
l'est. Les \emph{borrowed pointers} (notés \texttt{\&int}) correspondent aux
pointeurs obtenus en prenant l'adresse d'un objet, ou d'un champ d'un objet. Une
analyse statique faite lors de la compilation s'assure que la durée de vie de
ces pointeurs est plus longue que l'objet pointé, afin d'éviter les
\emph{dangling pointers} (pointeurs fous). Cette analyse est également basée sur
les régions. Une fonction qui retourne l'adresse d'une variable locale sera donc
rejetée par le compilateur. Enfin, le dernier type est celui des \emph{raw
pointers} (notés \texttt{*int}), pour lesquels le langage n'apporte aucune
garantie (il faut d'ailleurs encapsuler chaque utilisation dans un bloc marqué
explicitement \texttt{unsafe}). Ils sont équivalents aux pointeurs de C.

Ces techniques sont utiles pour créer des nouveaux programmes sûrs, mais on ne
peut pas les appliquer pour étudier la correction de logiciels existants. Leurs
systèmes de types apportent dans le langage différents types de pointeurs,
permettant de manipuler finement la mémoire, à la manière des \emph{smart
pointers} de C++. De cette approche on retient surtout l'analyse de régions de
Rust qui permet de manipuler de manière sûre les adresses des variables locales,
et les pointeurs lourds de Cyclone qui apportent une sûreté à l'arithmétique de
pointeurs, au prix d'un test dynamique.

\section{Logique de Hoare}

Une technique pour vérifier statiquement des propriétés sur la sémantique d'un
programme a été formalisée par Robert Floyd~\cite{FloydMeaning} et Tony
Hoare~\cite{hoare}.

Elle consiste à écrire les invariants qui sont maintenus à un point donné du
programme. Ces propositions sont écrites dans une logique $\mathcal{L}$. Chaque
instruction $i$ est annotée d'une pré-condition $P$ et d'une post-condition $Q$,
ce que l'on note $\hoare{P}{i}{Q}$. Cela signifie que si $P$ est vérifiée et que
l'exécution de $i$ se termine, alors $Q$ sera vérifiée.

En plus des règles de $\mathcal{L}$, des règles d'inférence traduisent la
sémantique du programme; par exemple la règle de composition est:

\begin{mathpar}
  \irule{Hoare-Seq}
    { \hoare{P}{i_1}{Q} \\
      \hoare{Q}{i_2}{R}
    }{
      \hoare{P}{i_1;i_2}{R}
    }
\end{mathpar}

Les pré-conditions peuvent être renforcées et les post-conditions relâchées:

\begin{mathpar}
  \irule{Hoare-Consequence}
    { ⊢_{\mathcal{L}} P  ⇒ P' \\
      \hoare{P}{i}{Q} \\
      ⊢_{\mathcal{L}} Q' ⇒ Q
    }
    { \hoare{P'}{i}{Q'} }
\end{mathpar}

Il est alors possible d'annoter le programme avec ses invariants formalisés de
manière explicite dans $\mathcal{L}$. Ceux-ci seront vérifiés à la compilation
lorsque c'est possible, sinon à l'exécution.

%TODO[B] Le rapport avec la prog sys? un cite ? (ça a forcément été fait)

La règle de conséquence permet de découpler les propriétés du programme lui-même
: plusieurs niveaux d'annotations sont possibles, du moins précis au plus
précis. En fait, il est même possible d'annoter chaque point de contrôle par
l'ensemble d'annotations vide: \hoare{T}{i}{T} est toujours vrai.

Augmenter graduellement les pré- et post-conditions est néanmoins assez
difficile, puisqu'il peut être nécessaire de modifier l'ensemble des conditions
à la fois. Cette difficulté est mentionnée dans~\cite{cssv}, où un système de
programmation par contrats est utilisé pour vérifier la correction de routines
de manipulation de chaînes en C.

Ce type d'annotations a été implanté par exemple pour le langage Java dans le
système JML~\cite{jmlkluwer} ou pour le langage C\# dans Spec\#~\cite{krml136}.

% TODO [E] ce sont juste des langages d'annotations: qu'en fait-on après ?

\section{Assistants de preuve}

%TODO[E] globalement trop rapide

Avec un système de types classique, le fait qu'un terme soit bien typé amène
quelques propriétés sur son exécution. Par exemple, le fait que seulement un
ensemble réduit d'erreurs puisse arriver (comme la division par zéro).

En enrichissant le langage des types, on peut augmenter l'expressivité du
typage. Par exemple, on peut former des types ``entier pair'', ``vecteur de $n$
entiers'', ou encore ``liste triée d'entiers''.

Habituellement, les termes peuvent dépendre d'autres termes (par composition) ou
de types (par des annotations). Les types peuvent également dépendre d'autres
types (par composition de types: par exemple, un couple de $a$ et de $b$ a pour
type $a * b$). Enrichir l'expressivité du typage revient essentiellement à
introduire des termes dans les types, comme $n$ dans l'exemple précédent du
vecteurs de $n$ entiers. C'est pourquoi on parle de types dépendants. Parmi les
langages proposant ces types on peut citer Coq~\cite{coqmanual},
Agda~\cite{agdatutorial} ou Isabelle~\cite{isabelletutorial}.

Dans un langage classique, la plupart des types sont habitables, c'est à dire
qu'il existe des termes ayant ces types. En revanche, avec les types dépendants
ce n'est pas toujours vrai: par exemple ``vecteur de $-1$ entiers'' n'a pas
d'habitants. Ainsi, pouvoir construire un terme d'un type donné est une
information en soi.

On peut voir ce phénomène sous un autre angle: les termes sont à leur type ce
que les preuves sont à leur théorème. C'est la correspondance de Curry-Howard.
Prenons un exemple: si on dispose d'un terme de type $a$, et d'une fonction de
type $a → b$, alors on peut en déduire (par simple application) un terme de type
$b$. Vu sous forme logique: si on a une preuve de $A$, et une preuve de $A ⇒
B$, alors on peut en tirer une preuve de $B$. À l'aide de cette correspondance,
on peut donc voir les systèmes de types dépendants comme des assistants de
preuves généralistes où l'on peut faire vérifier par un compilateur des
théorèmes mathématiques comme le théorème des 4 couleurs~\cite{4colproof}.

Cette technique fait se rejoindre les mondes du typage et des annotations par
triplets de Hoare: on enrichit le langage des types suffisamment pour encoder
une logique, et on vérifie suffisamment les triplets statiquement pour éliminer
tous les tests dynamiques. Cependant cette technique est très complexe à mettre
en œuvre, puisqu'il faut encoder toutes les propriétés voulues dans un
formalisme de très bas niveau (similaire à la théorie des ensembles).

\section{Conclusion}

Il existe de nombreuses techniques pour vérifier du code système ou embarqué. Il
y a divers choix à faire entre l'expressivité, l'intégration de tests
dynamiques, ou la facilité de mise en œuvre.

%TODO[E]: il faut conclure par rapport à ce que tu désires rechercher comme
%garanties

Le typage statique est une solution assez pragmatique, puisqu'elle peut
s'appliquer à des programmes existants, et est efficace à utiliser. Son
expressivité limitée nous empêche de reposer entièrement sur elle pour garantir
l'absence d'erreur dans les programmes système. C'est pourquoi nous approchons
la sûreté de la manière suivante dans notre solution, \langname.

%TODO[E]: phrase

Tout d'abord, on utilise le typage pour manipuler les données de manière
compatible: les types des opérations et fonctions sont vérifiés à la
compilation.

Ensuite, les accès aux tableaux et aux pointeurs sont vérifiés dynamiquement.
Dans le cas où une erreur est déclenchée, l'exécution s'arrête plutôt que de
corrompre la mémoire. La pile est également nettoyée à chaque retour de fonction
afin d'éviter les pointeurs fous (\emph{dangling pointers}).

Enfin, les pointeurs provenant de l'espace utilisateur sont repérés statiquement
afin que leur déréférencement se fasse au sein de fonctions sûres. Cela permet
de préserver l'isolation entre le noyau et l'espace utilisateur.

% vim: spelllang=fr


\partccl{part:ctx}

Nous avons montré que l'écriture de noyaux de systèmes d'exploitation nécessite
de manipuler des données provenant d'une zone non sûre, l'espace utilisateur.
Parmi ces données, il arrive de récupérer des pointeurs qui servent à passer des
données par référence à l'appelant, dans certains appels système. Si on
déréférence ces pointeurs sans vérifier qu'ils pointent bien vers une zone
mémoire également contrôlée par l'appelant, on risque de lire ou d'écrire dans
des zones mémoires réservées au noyau seul.

Nous proposons une technique de typage pour détecter ces cas dangereux. Pour ce
faire, il faudra tout d'abord définir un langage impératif bien typé que nous
appellerons \langname.
Celui-ci s'appuie sur le langage \newspeak, qui est un langage intermédiaire
développé par EADS dans le but de vérifier la sûreté de programmes C embarqués.
À ce titre, il existe un compilateur qui est capable de traduire du code C vers
\newspeak.

Une fois \langname défini, nous lui ajoutons des constructions permettant
d'écrire du code noyau, et en particulier on lui ajoute des pointeurs
utilisateur: il s'agit de pointeurs dont la valeur est contrôlée par un
utilisateur interagissant avec le programme via un appel système. Ces pointeurs
ont un type distinct des pointeurs habituels.

Cela revient à introduire une séparation entre le type des données et sa
représentation, c'est-à-dire définir un type abstrait. Les pointeurs
utilisateur sont en fait des pointeurs classiques déguisés, pour lesquels on
interdit l'opérateur de déréférencement.

Cette technique est en fait générique : on peut également l'appliquer à certains
types d'entiers. En C, il est commun d'utiliser des \texttt{int} pour tout et
n'importe quoi : pour des entiers bien sûr (au sens de $ℤ$), mais aussi comme
identificateurs pour lesquels les opérations usuelles comme l'addition n'ont pas
de sens. Par exemple, sous Linux, l'opération d'ouverture de fichier renvoie un
entier, dit \emph{descripteur de fichier}, qui identifie ce fichier pour ce
processus. Il est passé à toutes les fonctions de lecture et d'écriture de
fichiers.
Le langage autorise donc par exemple de multiplier entre
eux deux descripteurs de fichiers, mais le résultat n'a pas de raison \emph{a
priori} d'être un descripteur de fichier valide.


%Un
%autre idiome est d'utiliser des entiers comme des ensembles de bits (on parle de
%\emph{bitmask}).
%Un prototype d'analyse statique sur ces \emph{bitmasks} est
%présenté en annexe~\ref{anx:bitmasks}.

En n'offrant pas cette distinction, le langage C permet d'écrire du code qui
peut s'exécuter mais dont la sémantique n'est, quelque part, pas bien fondée.
Par exemple, en additionnant deux descripteurs de fichiers, ou en déréférençant
un pointeur contrôlé par l'utilisateur. En effet, le système de types de C est
trop primitif pour pouvoir garantir une véritable isolation entre deux types de
même représentation : il n'y a pas de types abstraits. Certes, \texttt{typedef}
permet d'introduire un nouveau nom pour un type, mais ce n'est qu'un raccourci
syntaxique. Le compilateur ne peut en effet pas considérer un programme sans
avoir la définition quasi-complète des types qui y apparaissent. La seule
exception concerne les pointeurs sur structures : si on ne fait que les
affecter, il n'est pas nécessaire de connaître la taille ni la disposition de la
structure, donc la définition peut ne pas être visible. Cette technique, connue
sous le nom de \emph{pointeurs opaques}, n'est pas applicable aux autres types.

En ajoutant une couche de typage, on interdit ces opérations à la compilation.
Cela permet deux choses : pour le code déjà écrit, de détecter et corriger les
manipulations dangereuses ; et, pour le nouveau code, de s'assurer qu'il est
correct. Par exemple, si on écrit un éditeur de texte, on peut éviter de
nombreuses erreurs de programmation en définissant un type \enquote{indice de
ligne} et un type \enquote{indice de colonne} incompatibles entre eux.

Le but de cette thèse est donc de définir un langage intermédiaire proche de C,
mais bien typable, et de lui adjoindre un système de types tel que les
programmes bien typés manipulent les pointeurs utilisateur sans causer de
problèmes de sécurité.

% TODO veuve

% vim: spelllang=fr


\long\def\partintro{
Dans cette partie, nous allons présenter un langage impératif modélisant une
sous-classe \enquote{bien typable} du langage C. Le chapitre~\ref{cha:lang}
décrit sa syntaxe, ainsi que sa sémantique d'exécution. À ce point, de nombreux
programmes acceptés peuvent provoquer des erreurs à l'exécution.

Afin de rejeter ces programmes incorrects, on définit ensuite dans le
chapitre~\ref{cha:typbase} une sémantique statique s'appuyant sur un système de
types simples. Des propriétés de sûreté de typage sont ensuite établies,
permettant de catégoriser l'ensemble des erreurs à l'exécution possibles.

Le chapitre~\ref{cha:qualifs} commence par étendre notre langage avec une
nouvelle classe d'erreurs à l'exécution, modélisant les accès à la mémoire
utilisateur catégorisés comme dangereux dans le chapitre~\ref{cha:os}. Une
extension au système de types du chapitre~\ref{cha:typbase} est ensuite établie,
et on prouve que les programmes ainsi typés ne peuvent pas atteindre ces cas
d'erreur.

Trois types d'erreurs à l'exécution sont possibles:

\begin{itemize}
\item
les erreurs liées aux valeurs: lorsqu'on tente d'appliquer à une opération des
valeurs incompatibles (additionner un entier et une fonction par exemple).
L'accès à des variables qui n'existent pas rentre aussi dans cette catégorie.
\item
les erreurs mémoire, qui résultent d'un débordement de tableau, du
déréférencement d'un pointeur invalide ou d'arithmétique de pointeur invalide.
\item
les erreurs de sécurité, qui consistent en le déréférencement d'un pointeur dont
la valeur est contrôlée par l'espace utilisateur. Celles-ci sont uniquement
possibles en contexte noyau.
\end{itemize}

L'introduction des types simples enlève la possibilité de rencontrer le premier
cas. Il reste en revanche toujours possible de rencontrer des erreurs mémoire
ainsi que des divisions par zéro. Éliminer ces erreurs dépasse le cadre de ce
travail.

En présence d'extensions permettant de manipuler des pointeurs utilisateurs, une
extension naïve du système de types ne suffit pas à empêcher la présence
d'erreurs de sécurité. Celles-ci sont évitées par l'ajout de règles de typage
supplémentaires.
}

\part{Un langage pour l'analyse de code système : \langname}
\label{part:lang}


\chapter{Syntaxe et sémantique d'évaluation}

\label{cha:lang}
Dans ce chapitre, on décrit le cœur de notre travail: le langage \langname, sur
lequel s'appuieront les analyses de typage des chapitres~\ref{cha:typbase}
et~\ref{cha:qualifs}.

Le langage C~\cite{KandR} est un langage impératif, conçu pour être un
«assembleur portable». Ses types de données et les opérations associées sont
donc naturellement très bas niveau.

Ses types de données sont établis pour représenter les mots mémoire manipulables
par les processeurs : essentiellement des entiers et flottants de plusieurs
tailles. Les types composés correespondent à des zones de mémoire contigües,
homogènes (dans le cas des tableaux) ou hétérogènes (dans le cas des
structures).

Une des spécificités de C est qu'il expose au programmeur la notion de pointeur,
c'est-à-dire des variables qui représentent directement une adresse en mémoire.
Les pointeurs peuvent être typés (on garde une indication sur le type de l'objet
stocké à cette adresse) ou non typés.

Son système de types rudimentaire ne permet pas d'avoir beaucoup de garanties
sur la sûreté du programme. En effet, aucune vérification n'est effectuée en
dehors de celles faites par le programmeur.

Le but ici est de définir \langname, un langage plus simple mais qui permettra
de raisonner sur une certaine classe de programmes C.

Tout d'abord, on commence par présenter les notations qui accompagneront le
reste des chapitres. Cela inclut la définition de la notion de lentille, qui est
utilisée dans certaines définitions.

Ensuite, on présente \langname en soi, c'est-à-dire sa syntaxe, ainsi qu'une
description de ses caractéristiques principales. En particulier, le modèle
mémoire est expliqué, ainsi que les valeurs manipulées par le langage.

Enfin, on décrit une sémantique opérationnelle pour ce langage. Cela permet de
définir précisément l'exécution d'un programme \langname au niveau de la
mémoire.

L'implantation de ces analyses est faite dans le chapitre~\ref{cha:implem}. En
particulier, on y verra comment traduire un sous-ensemble de C vers Newspeak,
sur lequel notre implantation est faite.

\section{Notations}

\subsection*{Inférence}

La sémantique opérationnelle consiste en la définition d'une relation de
transition $\cdot\rightarrow\cdot$ entre états de l'interpréteur\footnote{Dans le
chapitre~\ref{cha:typbase}, la relation de typage $\cdot ⊢ \cdot : \cdot$ sera
définie par la même technique.}.

Cette relation est définie inductivement sur la syntaxe du programme. Plutôt que
de présenter l'induction explicitement, elle est représentée par des jugements
logiques et des règles d'inférences, de la forme:

\[
\irule{Nom}{P_1 \\ … \\ P_n}{C}
\]

Les $P_i$ sont les prémisses, et $C$ la conclusion. Cette règle s'interprète de
la manière suivante: si les $P_i$ sont prouvées, alors $C$ est prouvée.

Certaines règles n'ont pas de prémisse, ce sont des axiomes:

\[
\iaxiom{Ax}{A}
\]

Compte-tenu de la structure des règles, la preuve d'un jugement pourra donc être
vue sous la forme d'un arbre où les axiomes sont les feuilles, en haut, et la
conclusion est la racine, en bas.

\[
  \irule{r1}{
    \irule{r2}
          {
            \iaxiom{r3}{A_1}
              \\
            \iaxiom{r4}{A_2}
          }
          {B_1}
    \\
    \irule{r5}
      {
        \iaxiom{r6} {A_3}
      }{B_2}
      }{C}
\]

\subsection*{Listes}

$X^*$ est l'ensemble des suites finies de $X$, indexées à partir de 1. Si $u ∈
X^*$, on note $|u|$ le nombre d'éléments de $u$ (le cardinal de son domaine de
définition). Pour $i ∈ [1 ; |u|]$, on note $u_i = u(i)$ le i-ème élément de la
suite.

On peut aussi voir les suites comme des listes: on note $[]$ la suite vide,
telle que $|[]| = 0$. On définit en outre la construction de suite de la manière
suivante: si $x ∈ X$ et $u ∈ X^*$, la liste $x::u ∈ X^*$ est la liste $v$ telle
que:

\begin{align*}
                       v_1 & = x \\
  ∀ i ∈ [1; |u|] , v_{i+1} & = u_i
\end{align*}

Cela signifie que la tête de liste ($x$ dans la liste $x :: u$) est toujours
accessible à l'indice $1$.

\subsection*{Lentilles}

Dans la définition de la sémantique de \langname, on utilise des \emph{lentilles
bidirectionnelles}. Cette notion n'est pas propre à la sémantique des
programmes. Il s'agit d'une technique permettant de relier la modification d'un
objet à la modification d'un de ses sous-composants. On peut en trouver des
descriptions plus détaillées décrite dans~\cite{PierceLenses}
et~\cite{LaarhovenLenses}.

\begin{definition}[Lentille]

Étant donnés deux ensembles $R$ et $A$, une \emph{lentille} $ℒ ∈ \setLens{R}{A}$
(ou \emph{accesseur}) est un moyen d'accéder en lecture ou en écriture à
sous-valeur de type $A$ au sein d'une valeur de type $R$ (pour \emph{record}).
Elle est consistuée des opérations suivantes:

\begin{itemize}
\item
  une fonction de lecture $\mathrm{get}_ℒ : R → A$
\item
  une fonction de mise à jour $\mathrm{put}_ℒ : (A × R) → R$
\end{itemize}

telles que pour tous $a∈A, a'∈A, r∈R$:

\begin{align*}
\tag{GetPut}
\mathrm{put}_ℒ(\mathrm{get}_ℒ(r), r) &= r \\
\tag{PutGet}
\mathrm{get}_ℒ(\mathrm{put}_ℒ(a, r)) &= a \\
\tag{PutPut}
\mathrm{put}_ℒ(a', \mathrm{put}_ℒ(a, r)) &= \mathrm{put}_ℒ(a', r) \\
\end{align*}

On note $ℒ =
\mkLens{\mathrm{get}_{ℒ}}{\mathrm{put}_{ℒ}}$.

\textsc{GetPut} signifie que si on lit une valeur puis qu'on la réécrit, l'objet
n'est pas modifié; \textsc{PutGet} décrit l'opération inverse: si on écrit
une valeur dans le champ, c'est la valeur qui sera lue; enfin, \textsc{PutPut}
évoque le fait que chaque écriture est totale: quand deux écritures se suivent,
seule la seconde compte.

\end{definition}

Une illustration se trouve dans la figure~\ref{fig:lens-howto}.

\begin{figure}[h]

\tikzstyle{lenspatA}=[fill=white,pattern=crosshatch dots]
\tikzstyle{lenspatB}=[fill=black!50]
\tikzstyle{lenspatC}=[fill=black!20]

  \begin{align*}
  \lensGet{ℒ}{\lensNodeBig{lenspatA}{lenspatB}} &= \lensInner{lenspatB} \\
  \lensPut{ℒ}{\lensInner{lenspatC}}{\lensNodeBig{lenspatA}{lenspatB}}
        &= \lensNodeBig{lenspatA}{lenspatC} \\
  \end{align*}

\caption{Fonctionnement d'une lentille}
\label{fig:lens-howto}
\end{figure}

\begin{example}[Lentilles de tête et de queue de liste]

Soit $E$ un ensemble. On considère $L(E)$, l'ensemble des listes d'éléments de
$E$.

On définit les fonctions suivantes. Notons qu'elles ne sont pas définies sur la
liste vide $[]$, qui pourra être traité comme un cas d'erreur.

\begin{align*}
  \mathrm{get}_T     (t::q) &= t   & \mathrm{put}_T (t', t::q) &= t'::q  \\
  \mathrm{get}_Q     (t::q) &= q   & \mathrm{put}_Q (q', t::q) &= t::q'
\end{align*}

Alors
$T = \mkLens{\mathrm{get}_T}{\mathrm{put}_T} ∈ \setLens{L(E)}{E}$
et
$Q = \mkLens{\mathrm{get}_Q}{\mathrm{put}_Q} ∈ \setLens{L(E)}{L(E)}$.

On a par exemple:

\begin{align*}
\mathrm{get}_T (1::6::1::8::[]) &= 1   &   \mathrm{put}_Q (4::2::[], 3::6::1::5::[]) = 3::4::2::[]
\end{align*}

\end{example}

\begin{definition}[Lentille indexée]

Les objets de certains ensembles $R$ sont composés de plusieurs sous-objets
accessibles à travers un indice $i ∈ I$. Une lentille indexée est une fonction
$Δ$ qui associe à un indice $i$ une lentille entre $R$ et un de ses champs
$A_i$:

\[
  ∀ i ∈ I, ∃ A_i, Δ(i) ∈ \setLens{R}{A_i}
\]

On note alors:

\begin{align*}
    r {[ i ]}_Δ     & \eqdef \mathrm{get}_{Δ(i)}(r) \\
    r {[ i ← a ]}_Δ & \eqdef \mathrm{put}_{Δ(i)}(a, r) \\
\end{align*}

\end{definition}

Un exemple est illustré dans la figure~\ref{fig:lens-idx-ex}.

\begin{figure}[h]

\tikzstyle{lenspatA}=[fill=white,pattern=crosshatch dots]
\tikzstyle{lenspatB}=[fill=black!50]
\tikzstyle{lenspatC}=[fill=black!20]

  \begin{align*}
    \lensGet{Δ(b)}{\lensNodeBigIdx{lenspatA}{lenspatB}{lenspatB}{lenspatB}{lenspatB}} &= \lensInnerStar{lenspatB} \\
    \lensGet{Δ(c)}{\lensNodeBigIdx{lenspatA}{lenspatB}{lenspatB}{lenspatB}{lenspatB}} &= \lensInner{lenspatB} \\
    \lensPut{Δ(b)}{\lensInnerStar{lenspatC}}{\lensNodeBigIdx{lenspatA}{lenspatB}{lenspatB}{lenspatB}{lenspatB}} &=
      \lensNodeBigIdx{lenspatA}{lenspatB}{lenspatC}{lenspatB}{lenspatB} \\
    \lensPut{Δ(c)}{\lensInner{lenspatC}}{\lensNodeBigIdx{lenspatA}{lenspatB}{lenspatB}{lenspatB}{lenspatB}} &=
      \lensNodeBigIdx{lenspatA}{lenspatB}{lenspatB}{lenspatC}{lenspatB}
  \end{align*}

\caption{Fonctionnement d'une lentille indexée}
\label{fig:lens-idx-ex}
\end{figure}

\begin{example}[Lentille «n\ieme élément d'un tuple»]

Soient $n ∈ ℕ$, et $n$ ensembles $E_1, …, E_n$.

Pour tout $i ∈ [1; n]$, on définit:

\begin{align*}
   g_i((x_1, …, x_n)) &= x_i \\
p_i(y, (x_1, …, x_n)) &= (x_1, …, x_{i-1}, y, x_{i+1}, …, x_n)\\
\end{align*}

Définissons $T(i) = \mkLens{g_i}{p_i}$. Alors $T(i) ∈ \setLens{(E_1×…×E_n)}{E_i}$.

Donc $T$ est une lentille indexée, et on a par exemple:

\begin{align*}
(3,1,4,1,5) {[2]}_T &= \mathrm{get}_{T(2)} ((3, 1, 4, 1, 5)) \\
                    &= 1 \\
\\
(9,2,6,5,3) {[3 ← 1]}_T &= \mathrm{put}_{T(3)} (1, (9,2,6,5,3)) \\
                        &= (9,2,1,5,3)
\end{align*}
\end{example}

La notation $3 ← 1$ peut surprendre, mais il s'interprète de la manière
suivante: en remplaçant l'élément d'indice $3$ par $1$.

\begin{definition}[Composition de lentilles]
\label{def:lens-comp}

  Soient $ℒ_1 ∈ \setLens{A}{B}$ et $ℒ_2 ∈ \setLens{B}{C}$.

  La composition de $ℒ_1$ et $ℒ_2$ est la lentille
  $ℒ ∈ \setLens{A}{C}$ définie de la manière suivante:

  \begin{align*}
    \mathrm{get}_{ℒ} (r) =&
        \mathrm{get}_{ℒ_2}
        (\mathrm{get}_{ℒ_1} r) \\
    \mathrm{put}_{ℒ} (a, r) =&
        \mathrm{put}_{ℒ_1} (\mathrm{put}_{ℒ_2}
        (a, \mathrm{get}_{ℒ_1} r), r) \\
  \end{align*}

On notera alors $ℒ = ℒ_1 \ggg ℒ_2$ («$ℒ_1$ flèche $ℒ_2$»).

\end{definition}

\begin{figure}[h]

\centering

\tikzstyle{lensInnerBefore}=[fill=black!85]
\tikzstyle{lensInnerAfter}=[fill=black!10]
\tikzstyle{lensMid}=[fill=black!50]
\tikzstyle{lensExt}=[fill=white,pattern=crosshatch dots]

\shorthandoff{!}
\begin{tikzpicture}
[node distance=2cm
,bignode/.style={draw,shape=rectangle,minimum size=1cm,lensExt}
,smallnode/.style={draw,shape=circle,minimum size=8mm,lensMid}
,trinode/.style={draw,shape=regular polygon,regular polygon sides=3}
]
\node[bignode] (A) {};
\node[right of=A,smallnode] (B) {};
\node[right of=B,trinode,lensInnerAfter] (C) {};
\node[below of=B,smallnode] (D) {};
\node[below of=A,node distance=4cm,bignode] (E) {};

\draw[->] (A) to node[auto] {$\mathrm{get}_{ℒ_1}$} (B);
\draw[->] (B) to (D);
\draw[->] (A) to (E);
\draw (C) |- node[auto,near end,swap] {$\mathrm{put}_{ℒ_2}$} ($(B)!0.5!(D)$);
\draw (D) |- node[auto,near end,swap] {$\mathrm{put}_{ℒ_1}$} ($(A)!0.75!(E)$);

\node[smallnode] at (A) {};
\node[trinode,lensInnerBefore] at (A) {};

\node[trinode,lensInnerBefore] at (B) {};

\node[trinode,lensInnerAfter] at (D) {};

\node[smallnode] at (E) {};
\node[trinode,lensInnerAfter] at (E) {};

\node[left of=A, node distance=4cm,bignode] (GA) {};
\node[below of=GA,smallnode] (GB) {};
\node[below of=GB,trinode,lensInnerBefore] (GC) {};

\node[smallnode] at (GA) {};

\node[trinode, lensInnerBefore] at (GA) {};
\node[trinode, lensInnerBefore] at (GB) {};

\draw [->] (GA) to node[auto] {$\mathrm{get}_{ℒ_1}$} (GB);
\draw [->] (GB) to node[auto] {$\mathrm{get}_{ℒ_2}$} (GC);

\end{tikzpicture}
\shorthandoff{!}

\caption{Composition de lentilles}
\label{fig:compo-lens}
\end{figure}

Cette définition est illustrée dans la figure~\ref{fig:compo-lens}. Une preuve
que la composition est une lentille est donnée en annexe~\ref{proof:compo-lens}.

\section{Syntaxe}

Les figures~\ref{fig:stx-data} et~\ref{fig:stx} présentent
notre langage intermédiaire. Il contient la plupart des fonctionnalités
présentes dans les langages impératifs comme C.

\begin{figure}%{{{

  \figstxdata{}

  \caption{Syntaxe des expressions}
\label{fig:stx-data}
\end{figure}%}}}

\begin{figure}%{{{

  \figstxinstr{}

  \caption{Syntaxe des instructions}
\label{fig:stx}
\end{figure}%}}}

Un programme est organisé en fonctions, qui contiennent des instructions, qui
elles-mêmes manipulent des expressions.

Le flot de contrôle est simplifié par rapport à C:\@ il ne contient que
l'alternative («if») et la boucle «while». Les autres formes de boucle
(«do/while» et «for») peuvent être émulées par une boucle «while».

Les fonctionnalités manquantes, et comment les émuler, seront examinées dans le
chapitre~\ref{cha:conclusion}.

Les opérateurs sont donnés dans la figure~\ref{fig:stx-ops}.

\begin{figure}[h]%{{{

  \figstxops{}

  \caption{Syntaxe des opérateurs}
\label{fig:stx-ops}
\end{figure}%}}}

\section{Mémoire et valeurs}

L'interprète que nous nous apprêtons à définir manipule des valeurs qui sont
associées aux variables du programme.

La mémoire est constituée de variables (toutes mutables), qui contiennent des
valeurs. Ces variables sont organisées, d'une part en un ensemble de variables
globales, et d'autre part en une pile de contextes d'appel (qu'on appellera donc
aussi cadres de pile, ou \emph{stack frames} en anglais). Cette structure
empilée permet de représenter les différents contextes à chaque appel de
fonction: par exemple, si une fonction s'appelle récursivement, plusieurs
instances de ses variables locales sont présentes dans le programme.

La structure de pile des locales permet de les organiser en niveaux
indépendants: à \linebreak chaque appel de fonction, un nouveau cadre de pile
est créé, comprenant ses paramètres et ses variables locales. Au contraire, pour
les globales il n'y a pas de système d'empilement, puisque ces variables sont
accessibles depuis tout point du programme.

Pour identifier de manière non ambigüe une variable, on note simplement $x$ la
variable globale nommée $x$, et $(n, x)$ la variable locale nommée $x$ dans le
$n$\ieme cadre de pile\footnote{Les paramètres de fonction sont traités comme
des variables locales et se retrouvent dans le cadre correspondant.}.

Les affectations peuvent avoir la forme $x ← e$ où $x$ est une variable et $e$
est une expression, mais pas seulement. En effet, à gauche de $←$ on trouve en
général non pas une variable mais une left-value. Pour représenter quelle partie
de la mémoire doit être accédée par cette left-value, on introduit la notion de
chemin $φ$. Un chemin est une left-value évaluée: les cas sont similaires, sauf
que tous les indices sont évalués. Par exemple, $φ = (5, x).p$ représente le
champ «$p$» de la variable $x$ dans le 5\ieme cadre de pile. C'est à ce moment
qu'on évalue les déréférencements qui peuvent apparaître dans une left-value.

Les valeurs, quant à elles, peuvent avoir les formes suivantes (résumé sur la
figure~\ref{fig:interp-val}):

\begin{itemize}
\item

$\widehat{c}$: une constante. La notation circonflexe permet de distinguer
les constructions syntaxique des constructions sémantiques. Par exemple, à la
syntaxe $3$ correspond la valeur $\widehat{3}$.

Les valeurs entières sont les entiers signés sur 32 bits, c'est-à-dire entre
$-2^{31}$ à $2^{31}-1$. Mais ce choix est arbitraire: on
aurait pu choisir des nombres à 64 bits, par exemple.
Les flottants sont les flottants IEEE 754 de 32 bits~\cite{ieee754}.

Il n'y a pas de distinction entre procédures et fonction, toutes les fonction
doivent renvoyer une valeur. Celles qui ne retournent pas de valeur
«intéressante» renvoient alors une valeur d'un type à un seul élément noté
\eUnit\footnote{Cette notation évoque un $n$-uplet à 0 composante.}, et donc le
type sera noté \tUnit.

\item

$\widehat{\&}~φ$: une référence mémoire. Ce chemin correspond à un pointeur sur
une left-value. Par exemple, l'expression $\&x$ s'évalue en $\widehat{\&}~φ =
\widehat{\&}~(5, x)$ si $x$ désigne lexicalement une variable dans le 5\ieme
cadre de pile.

\item

$\widehat{ \eArray{v_1 ;…; v_n} }$: un tableau. C'est une valeur composée qui
contiennt un certain nombre de valeurs d'un même type. Par exemple, 100 entiers.
On accède à ces valeurs par un indice entier, compilation. C'est une erreur
($\serr{array}$) d'accéder à un tableau en dehors de ses bornes, c'est-à-dire en
dehors de $[0;n-1]$ pour un tableau à $n$ éléments. Pareillement,
$\widehat{\eArray{\cdot}}$ permet de désigner les valeurs tableau. Par exemple,
si $x$ vaut $2$ et $y$ vaut $3$, l'expression $\eArray{ x, y }$ s'évaluera en
valeur $\widehat{\eArray{ 2, 3 }}$

\item

$\widehat{ \eStruct{ l_1 : v_1 ; … ; l_n : v_n } }$: une structure. C'est une
valeur composée mais hétérogène. Les différents éléments (appelés \emph{champs})
sont désignés par noms $l$ (pour \emph{label}). Dans le programme, le nom de
champ $l$ est décoré de la définition complète de la structure $S$. Celle-ci
n'est pas utilisée dans l'évaluation et sera donc décrite au
chapitre~\ref{cha:typbase}. Comme précédemment, on note
$\widehat{\eStruct{\cdot}}$ pour dénoter les valeurs.

\item

$\widehat{f}$: une fonction. On garde en mémoire l'intégralité de la définition
de la fonction (liste de paramètres, de variables locales et corps). Même si les
fonctions locales sont possibles, il n'est pas possible d'accéder aux variables
de la portée entourante depuis la fonction intérieure (il n'y a pas de
fermetures). Contrairement à C, les fonctions ne sont pas des cas spéciaux. Par
exemple, les fonctions globales sont simplement des variables globales de type
fonctionnel, et «les pointeurs sur fonction» de C sont remplacés par des
variables de type fonction.

\item $Ω$: une erreur. Par exemple le résultat de $5 / 0$ est $\serr{div}$.

\end{itemize}

\begin{figure}%{{{

  \begin{align*}
  \gramdef{Valeurs}{v}
      { \widehat{c}     }{Constante}
      { \widehat{\&}~φ  }{Référence mémoire}
      { \widehat{
         \eStruct{ l_1 : v_1 ;
       … ; l_n : v_n }
       }                }{Structure}
      { \widehat{
        \eArray{v_1 ;…; v_n}
        }               }{Tableau}
      { \widehat{f}     }{Fonction}
      { Ω               }{Erreur}
      {END}
  \\
  \\
  \gramdef{Chemins}{φ}
     { a    }{Adresse}
     { φ\widehat{.}l  }{Accès à un champ}
     { φ\widehat{[n]} }{Accès à un élément}
     {END}
  \\
  \\
  \gramdef{Adresses}{a}
     { (n, x) }{Variable locale}
     { (x)    }{Variable globale}
     {END}
  \\
  \\
  \gramdef{Erreur}{Ω}
    { \serr{array} }{Débordement de tableau}
    { \serr{ptr}   }{Erreur de pointeur}
    { \serr{div}   }{Division par zéro}
    { \serr{field} }{Erreur de champ}
    { \serr{var}   }{Variable inconnue}
    { \serr{typ}   }{Données incompatibles}
    {END}
  \end{align*}

  \caption{Valeurs}
\label{fig:interp-val}
\end{figure}%}}}

Les erreurs peuvent être classifiées en deux grand groupes: d'une part,
$\serr{field}$, $\serr{var}$ et $\serr{typ}$ sont des erreur de typage
dynamique, qui arrivent lorsqu'on accède dynamiquement a des données qui
n'existent pas où qu'on manipule des types de données incompatibles. D'autre
part, $\serr{div}$, $\serr{array}$ et $\serr{ptr}$ correspondent à des valeurs
mal utilisées. Le but du système de types du chapitre~\ref{cha:typbase} sera
d'éliminer complètement les erreurs du premier groupe.

\section{Interprète}

La figure~\ref{fig:interp-stack} résume comment ces valeurs sont organisées. Une
pile est une liste de cadre de piles, et un cadre de pile est une liste de
couples (nom, valeur). Un état mémoire $m$ est un couple $(s, g)$ où $s$ est une
pile et $g$ un cadre de pile (qui représente les variables globales).
On note $|m| = |s|$ la hauteur de la pile (en nombre de cadres).

\begin{figure}%{{{

  \begin{align*}
  \gramdef{Pile}{s}
    { [] }{Pile vide}
    { \{ x_1 ↦v_1; … ; x_n ↦v_n\} :: s }{Ajout d'un cadre}
    {END}
  \\
  \\
  \gramdef{État mémoire}{m}
    { (s, \{ x_1 ↦v_1; … ; x_n ↦v_n\})}{Pile, globales}
    {END}
  \\
  \\
  \gramdef[4cm]{État d'interprète}{Ξ}
    { \msi{m}{e} }{Expression, mémoire}
    { \msi{m}{i} }{Instruction, mémoire}
    { Ω          }{Erreur}
    {END}
  \end{align*}

  \caption{Composantes d'un état mémoire}
\label{fig:interp-stack}
\end{figure}%}}}


Enfin, l'interprétation est définie comme une relation $\cdot → \cdot$ entre
états $Ξ$; ces états sont d'une des formes suivantes:

\begin{itemize}
\item
  un couple $\msi{m}{e}$ où $e$ est une expression et $m$ un état
  mémoire. $m$ est l'état mémoire sous lequel l'évaluation sera
  réalisée. Par exemple $\mm{([], [x↦3])}{3}{([], [x↦3])}{\widehat{3}}$
  L'évaluation des expressions est détaillée dans la
  section~\ref{sec:eval-exp}.
\item
  un couple $\msi{m}{i}$ où $i$ est une instruction et $m$ un état
  mémoire. La réduction des instructions est traitée dans la
  section~\ref{sec:eval-instr}.
  Par exemple, $ \msi{m}{(x ← 3 ; y ← x)} → \msi{m [x ↦ \widehat{3}]}{y ← x}
                           → \msi{m [x ↦ \widehat{3}][y ↦
    \widehat{3}]}{\iPass}$.
% TODO continuation, citer Blazy

\item
  un couple $\msi{m}{lv}$ où $lv$ est une left-value et $m$ un état mémoire.
  L'évaluation des left-values est décrite en section~\ref{sec:eval-lv}.
\item
  une erreur $Ω$. La propagation des erreurs est détaillée dans la
  section~\ref{sec:eval-errors}.
\end{itemize}

L'évaluation des expressions, left-values et instruction se fait à petits pas.
C'est-à-dire qu'on simplifie d'étape en étape leur forme, jusqu'à arriver à un
cas de base:

\begin{itemize}
\item pour les expressions, une valeur $v$.
\item pour les instructions, l'instruction $\iPass$ ou $\iReturn{v}$ où $v$ est
    une valeur.
\item pour les left-values, un chemin $φ$.
\end{itemize}

On considère en fait la clôture transitive de cette relation. Cela revient à
ajouter une règle :

\begin{mathpar}
    \irule{Trans}
        { Ξ_1 → Ξ_2
       \\ Ξ_2 → Ξ_3
        }
        { Ξ_1 → Ξ_3}
\end{mathpar}

\section{Opérations sur les valeurs}
\label{sec:sem-ops}

Un certain nombre d'opérations est possible sur les valeurs
(figure~\ref{fig:stx-ops}):

\begin{itemize}
\item
  les opérations arithmétiques $+$, $-$, $×$, $/$ et $\%$ sur les entiers.
  L'opérateur $\%$ correspond au modulo (reste de la division euclidienne).
  En cas de division par zéro, l'erreur $\serr{div}$ est levée.
\item
  les versions «pointées»
  $\floatop{+}$,
  $\floatop{-}$,
  $\floatop{×}$
  et
  $\floatop{/}$
  sur les flottants
\item
  les opérations d'arithmétique de pointeur $+_p$ et $-_p$ qui à un chemin
  mémoire et un entier associent un chemin mémoire.
\item
  les opérations d'égalité $=$ et $≠$. L'égalité entre entiers ou entre
  flottants est immédiate. Deux valeurs composées (tableaux ou structures) sont
  égales si elles ont la même forme (même taille pour les tableaux, ou même
  champs pour les structures) et que toutes leurs sous-valeurs sont égales deux
  à deux. Deux références mémoire sont égales \linebreak lorsque les chemins
  qu'elles décrivent sont syntaxiquement égaux.
\item
  les opérations de comparaison $≤,≥,<,>$ sont définies avec leur sémantique
  habituelle sur les entiers et les flottants. Sur les références mémoires,
  elles sont définies dans le cas où les deux opérandes sont de la forme
  $φ[\cdot]$ par: $φ[n]~\opbin~φ[m] \eqdef n~\opbin~m$. Dans les autres cas,
  l'erreur $\serr{ptr}$ est renvoyée.
\item
  les opérateurs bit à bit sont définis sur les entiers. $\&$, $|$ et $\opxor$
  représentent respectivement la conjonction, la disjonction et la disjonction
  exclusive (XOR).
\item
  des versions logiques de la conjonction ($\&\&$) et de la disjonction ($||$)
  sont également présentes. Leur sémantique est donnée par le tableau suivant:

  \begin{center}
    \begin{tabular}{rr@{\hskip 1cm}cc}
      \toprule
       $n$ &  $m$ & $n~\&\&~m$ & $n~||~m$ \\
      \midrule
       $0$ &  $0$ & $0$        & $0$      \\
       $0$ & $≠0$ & $0$        & $1$      \\
      $≠0$ &  $0$ & $0$        & $1$      \\
      $≠0$ & $≠0$ & $1$        & $1$      \\
      \bottomrule
    \end{tabular}
  \end{center}

\item
  des opérateurs de décalage à gauche ($\ll$) et à droite ($\gg$) sont présents.
  Eux aussi ne s'appliquent qu'aux entiers.

\item
  les opérateurs arithmétiques unaires $+$, $-$, $\floatop{+}$ et $\floatop{-}$
  sont équivalents par définition à l'opération binaire correspondante.
  Par exemple $\widehat{-}4 \eqdef 0\widehat{-}4$.

\item
  $\sim$ inverse tous les bits de son opérande.
  $!$ est une version logique, % chktex 40
  c'est-à-dire que $!0 = 1$ et si $n ≠ 0$, $!n = 0$.
\end{itemize}

Si ces opérateurs sémantiques reçoivent des données incompatibles (par exemple
si on tente d'ajouter une fonction et un entier), l'erreur spéciale $\serr{typ}$
est renvoyée.

\section{Opérations sur les états mémoire}
\label{sec:memops}

\begin{definition}[Recherche de variable]

  La recherche de variable permet d'associer à une variable $x$ une adresse $a$.

  Chaque fonction peut accéder aux variables locales de la fonction en cours,
  ainsi qu'aux variables globales.

  Remarque le cadre de variables locales le plus récent a toujours l'indice $1$.

\[
  \cLookup{(s, g)}{x} =
        \begin{cases}
            (|s|, x) & \mbox{ si } |s| > 0
                       \mbox{ et } (x ↦ v) ∈ s_1 \\
            x        & \mbox{ si } (x ↦ v) ∈ g \\
            \serr{var} & \mbox{ sinon }
        \end{cases}
\]

\end{definition}

En entrant dans une fonction, on rajoutera un cadre de pile qui contient les
paramètres de la fonction ainsi que ses variables locales. En retournant à
l'appelant, il faudra supprimer ce cadre de pile.

\begin{definition}[Manipulations de pile]

  On définit l'empilement d'un cadre de pile $c = ((x_1↦ v_1), …, (x_n↦ v_n))$
  sur un état mémoire $m = (s, g)$ (figure~\ref{fig:stackops-push}):

  \[
      \cPush{(s, g)}{c} = (c::s, g)
  \]

  On définit aussi l'extension du dernier cadre de pile, qui sert aux
  déclarations de variables locales (figure~\ref{fig:stackops-extend}):

  \[
    \cExtend{(c::s, g)}{x}{v} = (((x↦v :: c) :: s), g)
  \]

  L'opération inverse de $\phxxx\cExtend$ sera simplement notée «$-$» : $m - x$,
  par exemple.

  De même on définit le dépilement (figure~\ref{fig:stackops-pop}):

  \[
      \cPop{(c::s, g)} = (s, g)
  \]

\end{definition}

\begin{definition}[Hauteur d'une valeur]
\label{def:hauteur-val}
Une valeur peut être, ou contenir, une référence vers une variable de la pile.
La hauteur d'une valeur est l'indice du plus haut cadre qu'elle référence, ou
$-1$ sinon.

\begin{minipage}{0.3\textwidth}
\begin{align*}
    \cH{\widehat{c}} &= -1 \\
    \cH{\widehat{f}} &= -1 \\
    \cH{\widehat{\&}~φ} &= \cHphi{φ} \\
    \cH{\widehat{\eStruct{l_1:v_1;…;l_n:v_n}}} &= \max_{i∈[1;n]} \cH{v_i} \\
    \cH{\widehat{\eArray{v_1,…,v_n}}} &= \max_{i∈[1;n]} \cH{v_i}
\end{align*}
\end{minipage}
où:
\begin{minipage}{0.3\textwidth}
\begin{align*}
    \cHphi{(x)} &= -1 \\
    \cHphi{(n,x)} &= n \\
    \cHphi{φ\widehat{.}l} &= \cHphi{φ} \\
    \cHphi{φ\widehat{[n]}} &= \cHphi{φ}
\end{align*}
\end{minipage}

\end{definition}


Les opérations $\mathrm{Extend}$ et $\mathrm{Pop}$ ne sont définies que pour une
pile non vides. Néanmoins cela ne pose pas de problème, puisque lors de
l'exécution, la pile n'est vide que lors de l'évaluation d'expressions dans les
phrases de programme. À cet endroit, seules des expressions peuvent apparaître,
et leur évaluation ne manipule jamais la pile avec ces opérations.

\begin{figure}%{{{

  \centering

  \subbottom[Empilement]{
\label{fig:stackops-push}
  \begin{tikzpicture}
    [ stack/.style={draw,shape=rectangle,minimum height=5mm,minimum width=2cm}
    , every node/.style={font=\footnotesize}
    , faded/.style={fill=black!20}
    ]
    \path (8,0)   node [stack, faded] {}
    -- ++ (0,0.5) node [stack, faded] {}
    -- ++ (0,0.5) node [stack, faded] {}
    ;
    \path (12,0)  node [stack, faded] {}
    -- ++ (0,0.5) node [stack, faded] {}
    -- ++ (0,0.5) node [stack, faded] {}
    -- ++ (0,0.5) node [stack] {$x↦0$}
    ;
    \draw [->] (9.2, 1) -- node[auto] {$\cPush{\cdot}{(x ↦ 0)}$} ++(1.6, 0);
  \end{tikzpicture}
  }

  \subbottom[Extension de cadre]{
\label{fig:stackops-extend}
  \begin{tikzpicture}
    [ stack/.style={draw,shape=rectangle,minimum height=5mm,minimum width=2cm}
    , every node/.style={font=\footnotesize}
    , faded/.style={fill=black!20}
    ]
    \path (4,-3)  node [stack, faded] {}
    -- ++ (0,0.5) node [stack, faded] {}
    -- ++ (0,0.5) node [stack, faded] {}
    -- ++ (0,0.5) node [stack] {$x↦0$}
    ;
    \path (8,-3)  node [stack, faded] {}
    -- ++ (0,0.5) node [stack, faded] {}
    -- ++ (0,0.5) node [stack, faded] {}
    -- ++ (0,0.5) node [stack] {$x↦0, y↦3$}
    ;
    \draw [->] (5.2, -2) -- node[auto] {$\cExtend{\cdot}{y}{3}$} ++(1.6, 0);
  \end{tikzpicture}
  }

  \subbottom[Dépilement]{
\label{fig:stackops-pop}
  \begin{tikzpicture}
    [ stack/.style={draw,shape=rectangle,minimum height=5mm,minimum width=2cm}
    , every node/.style={font=\footnotesize}
    , faded/.style={fill=black!20}
    ]
    \path (0,0)   node [stack, faded] {}
    -- ++ (0,0.5) node [stack, faded] {}
    -- ++ (0,0.5) node [stack, faded] {}
    -- ++ (0,0.5) node [stack] {$…$}
    ;
    \path (4,0)   node [stack, faded] {}
    -- ++ (0,0.5) node [stack, faded] {}
    -- ++ (0,0.5) node [stack, faded] {}
    ;
    \draw [->] (1.2, 1) -- node[auto] {Pop} ++(1.6, 0);
  \end{tikzpicture}
  }

  \caption{Opérations de pile}
\label{fig:stackops}
\end{figure}%}}}

On définit aussi une opération de nettoyage de pile, qui sera utile pour les
retours de fonction.

En effet, si une référence au dernier cadre est toujours présente après
le retour, elle pourra se résoudre en un objet différent plus tard dans
l'exécution du programme.

On commence par définir une fonction générique supprimant les valeurs satisfiant
un prédicat sur les chemins que la mémoire contient:

\begin{align*}
                \cP{p}{(s, g)} &= (s', g') \\
                 \mbox{où } s' &= (\cLP{p}{s_1}, …, \cLP{p}{s_{|s|}}) \\
                            g' &=  \cLP{p}{g} \\
\cLP{p}{(x_1 ↦ v_1, …, x_n ↦ v_n)} &= (x_1 ↦ \cVP{p}{v_1}, …, x_n ↦ \cVP{p}{v_n}) \\
\cVP{p}{\widehat{c}} &= \widehat{c} \\
\cVP{p}{\widehat{f}} &= \widehat{f} \\
\cVP{p}{\wphi{φ}} &= \begin{cases}
                       \eNull   & \mbox{ si }    p(φ) \\
                       \wphi{φ} & \mbox{ sinon}
                     \end{cases} \\
\cVP{p}{\widehat{\eStruct{l_1:v_1, …, l_n:v_n}}} &=
    \widehat{\eStruct{l_1:\cVP{p}{v_1}, …, l_n:\cVP{p}{v_n}}} \\
\cVP{p}{\widehat{\eArray{v_1, …, v_n}}} &=
    \widehat{\eStruct{\cVP{p}{v_1}, …, \cVP{p}{v_n}}}
\end{align*}

La fonction $\mathrm{Cleanup}$, et une version pour les valeurs, sont données
par:

\begin{align*}
\cCleanup{m} &=  \cP{λφ: \cHphi{φ} > |m|)}{m} \\
\cV{p}{v}    &= \cVP{λφ: \cHphi{φ} > |m|)}{v}
\end{align*}

On définit aussi une opération pour supprimer une variable de la mémoire:

\begin{align*}
    \cCleanVar{m}{a}  &= \cP{λφ: φ = a}{m} \\
    \cVV{v}{a}        &= \cVP{λφ: φ = a}{v}
\end{align*}

Sans ces opérations, examinons le programme suivant de la
figure~\ref{fig:prog-cleanup}.

\begin{SaveVerbatim}[]{LeBonProgramme}
f = Fun () {
    Decl x = 0 in
    return (&x);
}

g = Fun (p) {
    Decl x=0.0 in
    *p = 1;
}

h = Fun () {
    Decl p=f() in
    g(p);
}
\end{SaveVerbatim}

\begin{figure}[h]
\hspace{1cm}\BUseVerbatim{LeBonProgramme}
\caption{Cassage du typage par un pointeur fou}
\label{fig:prog-cleanup}
\end{figure}

L'exécution de $h()$ donne à $p$ la valeur $(1, x)$. Donc en arrivant dans $g$, le
déréférencement de $p$ va modifier $x$.

\section{Accesseurs}
\label{sec:mem-access}

Le but de cette section est de définir rigoureusement les accès à la mémoire. À
partir d'un état mémoire $m$ et d'une left-value $φ$, on veut pouvoir obtenir:

\begin{itemize}
    \item la valeur accessible au chemin $φ$: $m[φ]_Φ$
    \item l'état mémoire obtenu en remplaçant celle-ci par une nouvelle valeur $v'$ :
            $m[φ ← v']_Φ$
\end{itemize}

Pour définir cette lentille indexée $Φ$, on commence par définir des lentilles
élémentaires, et on les compose pour pouvoir définir des lentilles entre
valeurs.

\subsection*{Accès à une liste par indice : $I$}

On définit une lentille indexée $I : ℕ → \setLens{\textrm{List}(α)}{α}$
permettant d'accéder aux éléments d'une liste par leur indice.

\begin{align*}
    l{[n]}_I     &= l_n \mbox{ si } n ∈ [1; |l|] \\
    l{[n ← x]}_I &= l' \\
           \mbox{ où } & l'_n = x \\
                       & ∀ i ≠ n, l'_i = l_i
\end{align*}

I n'est définie que pour $n ∈ [1 ;|l|]$.

\subsection*{Accès à une liste d'associations : $L$}

  Une liste d'association est une liste de paires (clef, valeur) avec
  l'invariant supplémentaire que les clefs sont uniques. Il est donc possible de
  trouver au plus une valeur associée à une clef donnée. L'écriture est
  également possible, en remplaçant un couple par un couple avec une valeur
  différente.

\begin{align*}
l{[x]}_L  &=
    \begin{cases}
        v          & \mbox{ si } ∃! n ∈ [1;|l|]. ∃v. l_n = (x ↦ v) \\
        \serr{var} & \mbox{ sinon }
    \end{cases} \\
l{[x ← v]}_L  &=
    \begin{cases}
        l{[n ← (x ↦ v')]}_I & \mbox{ si } ∃! n ∈ [1;|l|]. ∃v. l_n = (x ↦ v) \\
        \serr{var} & \mbox{ sinon }
    \end{cases} \\
\end{align*}

\subsection*{Accès par adresse : $A$}

Les états mémoire sont constitués des listes d'association (nom, valeur).

L'accesseur par adresse ${[\cdot]}_A$ permet de généraliser l'accès à ces
valeurs en utilisant comme clef non pas un nom mais une adresse.

Selon cette adresse, on accède soit à la liste des globales, soit à une
des listes de la pile des locales.

On pose $m = (s, g)$.

Les accès aux globales se font de la manière suivante. Si la variable n'existe
pas, notons que $L$ retourne $\serr{var}$.

\begin{align*}
m {[x]}_A     &= g{[x]}_L \\
m {[x ← v]}_A &= (s, g{[x←v]}_L) \\
\end{align*}

Les accès aux locales reviennent à accéder à la bonne variable du bon cadre de
pile. Cela revient naturellement à composer les lentilles $L$ et $I$. On définit
donc une lentille $ℒ_{n,x} = I(|s|-n+1) \ggg L(x)$ qui accède à la variable $x$
du $n$\ieme cadre de pile.

% TODO vérifier les off by one
% TODO lentilles fst et snd?

\begin{align*}
m {[(n, x)]}_A     & = \begin{cases}
                           \lensGet{ℒ_{n,x}}{s} & \mbox{ si $n ∈ [1 ; |s|]$} \\
                           \serr{var}          & \mbox{ sinon}
                       \end{cases} \\
m {[(n, x) ← v]}_A & = \begin{cases}
                            (\lensPut{ℒ_{n,x}}{v}{s}, g) & \mbox{ si $n ∈ [1;|s|]$} \\
                            \serr{var}          & \mbox{ sinon}
                       \end{cases} \\
\end{align*}

Les numéros de cadre qui permettent d'identifier les locales (le $n$ dans $(n,
x)$) croissent avec la pile. D'autre part, l'empilement se fait en tête de liste
(près de l'indice 1). Donc pour accéder aux plus vieilles locales (numérotées
1), il faut accéder au dernier élément de la liste. Ceci explique pourquoi un
indice $|s|-n+1$ apparaît dans la définition précédente.

\subsection*{Accès par champ : $F$}

  Les valeurs qui sont des structures possèdent des sous-valeurs, associées à
  des noms de champ.

  L'accesseur ${[ \cdot ]}_F$ permet de lire et de modifier un champ de ces
  valeurs.

  L'erreur $\serr{field}$ est levée si on accède à un champ non existant.

  \begin{align*}
    \eStruct{ l_1 : v_1; … ; l_n : v_n }{[l]}_F &= v_i \mbox{ si } ∃i∈[1;n], l=l_i\\
    \eStruct{ l_1 : v_1; … ; l_n : v_n }{[l ← v]}_F &=
        \eStruct{ l_1 : v_1 \\
           &\quad ; … \\
           &\quad ; l_{p-1} : v_{p-1} \\
           &\quad ; l_p : v \\
           &\quad ; l_{p+1} : v_{p+1} \\
           &\quad ; … \\
           &\quad ; l_n : v_n } \mbox{ si } ∃p ∈ [1;n], l = l_p \\
    \eStruct{ l_1 : v_1; … ; l_n : v_n }{[l]}_F     &= \serr{field} \mbox{ sinon} \\
    \eStruct{ l_1 : v_1; … ; l_n : v_n }{[l ← v]}_F &= \serr{field} \mbox{ sinon}
  \end{align*}

\subsection*{Accès par indice de tableau : $T$}

  On définit de même un accesseur ${[\cdot]}_T$ pour les accès par indice à des
  valeurs tableaux. Néanmoins le paramètre indice est toujours un entier et pas
  une expression arbitraire.

  \begin{align*}
    \eArray{ v_1 ; … ; v_n } {[i]}_T   &= v_{i+1} \textrm{ si } i ∈ [0;n-1] \\
    \eArray{ v_1 ; … ; v_n } {[i]}_T   &= \serr{array} \textrm{ sinon} \\
    \eArray{ v_1 ; … ; v_n } {[i←v]}_T &= \eArray{ v'_1 ; … ; v'_n } \textrm{ si } i ∈ [0;n-1] \\
                      \mbox{ où } & \begin{cases}
                                      v'_i = v \\
                                      ∀j≠i, v'_j = v_j \\
                                    \end{cases} \\
    \eArray{ v_1 ; … ; v_n } {[i←v]}_T &= \serr{array} \textrm{ sinon} \\
  \end{align*}

\subsection*{Accès par chemin : $Φ$}
\label{subsec:acces-phi}

  L'accès par chemin $Φ$ permet de lire et de modifier la mémoire en profondeur.

  On peut accéder directement à une variable:

  \begin{align*}
    Φ(a) &= A(a)
  \end{align*}

  Les accès à des sous-valeurs se font en composant les accesseurs
  (définition~\ref{def:lens-comp}, page~\pageref{def:lens-comp}):

  \begin{align*}
    Φ(φ.l)  &= Φ(φ) \ggg F(l) \\
    Φ(φ[i]) &= Φ(φ) \ggg T(i)
  \end{align*}

  Enfin, l'accès à la mémoire par le pointeur nul provoque une erreur:

  \begin{align*}
    m{[\eNull    ]}_Φ &= \serr{ptr} \\
    m{[\eNull ← v]}_Φ &= \serr{ptr}
  \end{align*}

\section{Contextes d'évaluation}

L'évaluation des expressions repose sur la notion de contextes d'évaluation.
L'idée est que si on peut évaluer une expression, alors on peut évaluer une
expression qui contient celle-ci.

Par exemple, supposons que $\mm{m}{f(3)}{m}{2}$. Alors on peut ajouter la
constante $1$ à gauche de chaque expression sans changer le résultat:
$\mm{m}{1+f(3)}{m}{1+2}$. On a utilisé le même contexte
$C = 1+\ctxEmpty$.

Pour pouvoir raisonner en termes de contextes, 3 points sont nécessaires:

\begin{itemize}
\item comment découper une expression selon un contexte
\item comment appliquer une règle d'évaluation sous un contexte
\item comment regrouper une expression et un contexte
\end{itemize}

\begin{figure}
\figctx{}

\caption{Contextes d'évaluation}
\label{fig:eval-ctx}
\end{figure}

Le premier point consiste à définir les contextes eux-mêmes
(figure~\ref{fig:eval-ctx}).

Le deuxième est résolu par la règle d'inférence suivante.

\begin{mathpar}
  \semrule{Ctx}
\end{mathpar}

Enfin, le troisième revient à définit l'opérateur de substitution
$\phxx{\ctxSub}$ présent dans la règle précédente. Afin de pouvoir appliquer des
substitution au niveau des left-values et des instructions, on définit aussi
respectivement $\phxx{\ctxSubL}$ et $\phxx{\ctxSubI}$.

Notons que puisque $i \gramisa e$ et $e \gramisa lv$, on peut aussi l'appliquer
aux expressions et aux left-values: l'opération $\phxx\ctxSub$ est purement
syntaxique.

Dans la définition de l'ensemble des contextes, chaque cas hormis le cas de base
fait apparaître exactement un «$C$». Chaque contexte est donc constitué
d'exactement une occurrence de $\ctxEmpty$ (une dérivation de $C$ est toujours
linéaire). L'opération de substitution consiste à remplacer ce trou:
$\ctxSub{C}{X}$ est l'objet syntaxique (instruction, expression ou left-value)
obtenu en remplaçant l'unique $\ctxEmpty$ dans $C$ par $X$. Par exemple,
$\ctxSub{\iDecl{x}{2+\ctxEmpty}}{5} = \iDecl{x}{2+5}$

À titre d'illustration, décomposons l'évaluation de $e_1~\opbin~e_2$ en $v =
v_1~\widehat{\opbin}~v_2$ depuis un état mémoire $m$:

\begin{enumerate}
\item
  on commence par évaluer, l'expression $e_1$ en une valeur $v_1$. Le nouvel
  état mémoire est noté $m'$. Soit donc $\mm{m}{e_1}{m'}{v_1}$.
\item
  En appliquant la règle \textsc{Ctx} avec $C = \ctxOp{\ctxEmpty}{e_2}$ (qui est
  une des formes possibles pour un contexte d'évaluation), on déduit de 1.\ que
  $\mm{m}{e_1~\opbin~e_2}{m'}{v_1~\opbin~e_2}$
\item
  D'autre part, on évalue $e_2$ depuis $m'$. En supposant encore que
  l'évaluation converge, notons $v_2$ la valeur calculée et $m''$ l'état mémoire
  résultant: $\mm{m'}{e_2}{m''}{v_2}$.
\item
  Appliquons la règle~\textsc{Ctx} à 3.\ avec $C = \ctxOp{v_1}{\ctxEmpty}$. On
  obtient $\mm{m}{v_1~\opbin~e_2}{m'}{v_1~\opbin~v_2}$.
\item
  En combinant les résultats de 2.\ et 4.\ on en déduit que
  $\mm{m}{e_1~\opbin~e_2}{m''}{v_1\opbin~v_2}$.
\item D'après la règle~\textsc{Exp-Binop},
  $ \mm{m''}{v_1\opbin~v_2}{m''}{v_1~\widehat{\opbin}~v_2}$
\item D'après 5.\ et 6., on a par combinaison
  $\mm{m}{e_1~\opbin~e_2}{m''}{v}$
  en posant
  $v = v_1~\widehat{\opbin}~v_2$.
\end{enumerate}

\section{Left-values}
\label{sec:eval-lv}


Obtenir un chemin à partir d'un nom de variable revient à résoudre le nom de
cette variable: est-elle accessible? Le nom désigne-t-il une variable locale
ou une variable globale?

\begin{mathpar}
  \semrule{Phi-Var}
\end{mathpar}

Les règles portant sur le déréférencement et l'accès à un champ de structure
sont similaires: on commence par évaluer la left-value sur laquelle porte ce
modificateur, et on place le même modificateur sur le chemin résultant. Dans le
cas des champs de structure, l'annotation de structure $S$ n'est pas prise
en compte.

\begin{mathpar}
  \semrule{Phi-Struct}
\end{mathpar}

Enfin, pour évaluer un chemin dans un tableau, on commence par procéder comme
précédemment, c'est-à-dire en évaluant la left-value sur laquelle porte
l'opération d'indexation. Puis on évalue l'expression d'indice en une valeur qui
permet de construire le chemin résultant.

\begin{mathpar}
  \semrule{Phi-Array}
\end{mathpar}

Notons qu'en procédant ainsi, on évalue les left-values de gauche à droite:
dans l'expression $x[e_1][e_2][e_3]$, $e_1$ est évalué en premier, puis
$e_2$, puis $e_3$.

La règle portant sur le déréférencement est particulière. On peut penser que la
bonne définition de $φ$ consiste à se calquer sur la définition de $lv$, en
remplaçant les noms de variable par leur adresse résolue et en évaluant les
indices de tableau, et à ajouter une règle qui transforme $*φ$ en
$\widehat{*}φ$.

Cela ne fonctionne pas, car alors les déréférencements sont évalués trop tard:
au moment de l'affectation dans la left-value plutôt qu'à sa définition. La
figure~\ref{fig:lazy-deref} illustre ce problème.

%{{{
\begin{SaveVerbatim}[]{LeLazy}
Decl s0 = { .f = 0 } in
Decl s1 = { .f = 1 } in
Decl x  = & s0 in
Decl p = & ((*x).f) in
/* (a) */
x <- & s1
/* (a) */
\end{SaveVerbatim}

\begin{figure}[h]

\hspace{1cm}\BUseVerbatim{LeLazy}

\caption{Évaluation stricte ou paresseuse des left-values}
\label{fig:lazy-deref}
\end{figure}%}}}

On s'intéresse à l'évaluation
de l'expression \texttt{*p} aux points \texttt{(a)} et \texttt{(b)}. Avec une
sémantique paresseuse (en ajoutant un $\widehat{*}φ$), la valeur de
\texttt{p}
est $\widehat{\&}~((*(0,x)).f)$, donc
\texttt{*p} est évalué à $0$ en
\texttt{(a)}
et $1$ en
\texttt{(b)}.
Au contraire, avec une sémantique stricte (correcte),
\texttt{p} vaut
$\widehat{\&}~(((0,s0).f)$ et donc
\texttt{*p} est évalué à $0$ en
\texttt{(a)}
et en
\texttt{(b)}.
Les règles d'évaluation sont donc les suivantes. Dans le cas où la valeur
référencée n'a pas la forme $m[φ']_Φ$, aucune règle ne peut s'appliquer (comme
lorsqu'on cherche à réduire l'addition d'une fonction et d'un entier, par
exemple). Cela est préférable à renvoyer $\serr{ptr}$ car on montrera que la
première est toujours évitée dans les programmes typés
(théorème~\ref{thm:progres}).

\begin{mathpar}
  \semrule{Phi-Deref}
\end{mathpar}

Un exemple d'évaluation est donné dans la figure~\ref{fig:eval-lv}.

\begin{figure}[h]%{{{

  \centering

  \begin{tikzpicture}
    \node              (a1) {$x$};
    \node[right of=a1] (a2) {$.l_S$};
    \node[right of=a2] (a3) {$[2*n]$};
    \node[right of=a3] (a4) {$.g$};

    \node[below of=a1] (b1) {$(2, x)$};
    \node[below of=a2] (b2) {$.l$};
    \node[below of=a3] (b3) {$[4]$};
    \node[below of=a4] (b4) {$.g$};

    \draw[->] (a1) -- (b1);
    \draw[->] (a2) -- (b2);
    \draw[->] (a3) -- (b3);
    \draw[->] (a4) -- (b4);
  \end{tikzpicture}

  \caption{Évaluation des left-values.}
\label{fig:eval-lv}
\end{figure}%}}}

\section{Expressions}
\label{sec:eval-exp}

Évaluer une constante est le cas le plus simple, puisqu'en quelque sorte
celle-ci est déjà évaluée. À chaque constante syntaxique $c$, on peut associer
une valeur sémantique $\widehat{c}$. Par exemple, au chiffre (symbole) $3$, on
associe le nombre (entier) $\widehat{3}$.

\begin{mathpar}
  \semrule{Exp-Cst}
\end{mathpar}

De même, une fonction est déjà évaluée:

\begin{mathpar}
  \semrule{Exp-Fun}
\end{mathpar}

Pour lire le contenu d'un emplacement mémoire (left-value), il faut tout d'abord
l'évaluer en un chemin.

\begin{mathpar}
  \semrule{Exp-Lv}
\end{mathpar}

Pour évaluer une expression constituée d'un opérateur, on évalue une expression,
puis l'autre (l'ordre d'évaluation, est encore imposé: de gauche à droite). À
chaque opérateur $\opbin$, correspond un opérateur sémantique $\widehat{\opbin}$
qui agit sur les valeurs. Par exemple, l'opérateur $\widehat{+}$ est l'addition
classique entre entiers.
Comme précisé dans la section~\ref{sec:sem-ops}, la division par zéro via $/$,
$\%$ ou $\floatop{/}$ provoque l'erreur $\serr{div}$.

\begin{mathpar}
  \semrule{Exp-UnOp}

  \semrule{Exp-BinOp}
\end{mathpar}

Il est nécessaire de dire un mot sur les opérations $\widehat{+_p}$ et
$\widehat{-_p}$ définissant l'arithmétique des pointeurs. Celles-ci sont
uniquement définies pour les références mémoire à un tableau, c'est-à-dire
celles qui ont la forme $\widehat{\&}~φ[n]$ (par souci de clarté on ne note pas
tous les~$\widehat{\cdot}$ ci-dessous). On a alors:

\begin{align*}
  \widehat{\&}~φ[n]~+_p~i & = \widehat{\&}~φ[n+i] \\
  \widehat{\&}~φ[n]~-_p~i & = \widehat{\&}~φ[n-i] \\
\end{align*}

Cela implique qu'on ne peut pas faire faire d'arithmétique de pointeurs au sein
d'une même structure. Autrement c'est une erreur de manipulation de pointeurs
\label{page:def-arith-ptr-error}
$\serr{ptr}$.
\footnote{
Cela est cohérent avec la norme C99:
    «If the pointer operand points to an element of an array object, and the
    array is large enough, […] ; otherwise, the behavior is undefined.
    »~\cite[6.5.6~§8]{AnsiC}}
et l'opérateur $\widehat{\opbin}$ renvoie

Pour prendre l'adresse d'une variable, il suffit de résoudre celle-ci dans
l'état mémoire courant.

\begin{mathpar}
  \semrule{Exp-AddrOf}
\end{mathpar}

L'affectation se déroule 3 étapes: d'abord, l'expression est évaluée en une
valeur $v$. Ensuite, la left-value est évaluée en un chemin $φ$. Enfin, un
nouvel état mémoire est construit, où la valeur accessible par $φ$ est remplacée
par $v$. Comme dans le langage C, l'expression d'affectation produit une valeur,
qui est celle qui a été affectée.

\begin{mathpar}
  \semrule{Exp-Set}
\end{mathpar}

\subsection*{Expressions composées}

Les littéraux de structures sont évalués en leurs constructions syntaxiques
respectives. Puisque les contextes d'évaluation sont de la forme $ [ v_1 ; … ; C
; … ; e_n ] $, l'évaluation se fait toujours de gauche à droite.

\begin{mathpar}
  \semrule{Exp-Struct}

  \semrule{Exp-Array}
\end{mathpar}

%\shorthandoff{!}
%\begin{figure}%{{{

  %\centering

  %\begin{tikzpicture}
    %[node distance=2cm]

    %\node (m0) {$m_0$};
    %\node[right of=m0, node distance=3cm] (m1) {$m_1$};
    %\node[below of=m1]                    (m2) {$m_2$};
    %\node at ($(m0)!(m2)!(0,1)$)          (m3) {$m_3$};
    %\node[below of=m3]                    (m4) {$m_4$};

    %\draw [->] (m0) -- node[auto] {$\mathrm{Push}(\vec{a}↦\vec{v})$}    (m1);
    %\draw [->] (m1) -- node[auto] {$i → \iReturn{v_r}$}                 (m2);
    %\draw [->] (m2) -- node[auto] {$\mathrm{Pop}$}                      (m3);
    %\draw [->] (m3) -- node[auto] {$\mathrm{Cleanup}$}                  (m4);

    %\draw [dashed, ->] ($ (m0) + (-5mm, -3mm) $)
                    %-- node[auto, swap] {$f(\vec{e}) → v$}
                       %($ (m4) + (-5mm, 3mm) $);

  %\end{tikzpicture}

  %\caption[Appel d'une fonction]
  %{ Appel d'une fonction.
    %La taille de la pile croît de gauche à droite,
    %et les réductions se font de haut en bas.
  %}
%\label{fig:fcall-details}

%\end{figure}%}}}
%\shorthandon{!}

L'appel de fonction est traité de la manière suivante. On ne peut pas facilement
relier un pas d'évaluation de $i$ à un pas d'évaluation de $\eFun{a}{i}
(v_1,…,v_n)$, donc un contexte $C \gramisa \eFun{a}{\ctxEmpty} (v_1,…,v_n)$
n'est pas à considérer. En effet, l'empilement suivi du dépilement modifie la
mémoire.

On emploie donc une règle \textsc{Exp-Call-Ctx} qui relie un pas
interne $\mm{m_1}{i}{m_2}{i'}$ à un pas externe. Si une fois l'instruction
interne réduite d'un pas, on évalue les arguments en des valeurs $v'_i$. Ils
correspondent aux nouvelles valeurs à passer à la fonction.

Les autres fonctions permettent de transférer le flot de contrôle: en retournant
la même instruction pour une instruction terminale, ou en propageant une erreur.
Dans le cas où on retourne de la fonction par $i = \iReturn{v}$, il faut alors
supprimer les références aux variables qui ont disparu grâce aux opérateurs
$\phx\cCleanup$ et $\phxx\cV$.

\label{page:return-fonction}
On suppose deux choses sur chaque fonction: d'une part, les noms de ses
arguments sont deux à deux différents, et d'autre part, son corps se termine par
une instruction $\iReturn{\cdot}$. Cela veut dire que la dernière instruction
doit être soit de cette forme, soit par exemple une alternative dans laquelle
les deux branches se terminent par un $\iReturn{\cdot}$. C'est une propriété qui
peut être détectée statiquement avant l'exécution. Néanmoins, dans la syntaxe
concrète, on peut supposer qu'un $\iReturn{\eUnit}$ est inséré automatiquement
en fin de fonction lorsqu'aucun $\iReturn{\cdot}$ n'est présent dans son corps.

\begin{mathpar}
    \semrule{Exp-Call-Ctx}

    \semrule{Exp-Call-Err}

    \semrule{Exp-Call-Return}
\end{mathpar}

\section{Instructions}
\label{sec:eval-instr}

Les cas de la séquence et de l'évaluation d'une expression sont sans surprise.

\begin{mathpar}
  \semrule{Seq}

  \semrule{Exp}
\end{mathpar}

L'évaluation de $\iDecl{x}{v}{i}$ sous $m$ se fait de la manière suivante,
similaire à l'appel de fonction.
La règle principale est \textsc{Decl-Ctx} qui relie un pas d'évaluation sous une
déclaration à un pas d'évaluation externe : pour ce faire, on étend l'état
mémoire en ajoutant $x$, on effectue le pas, puis on enlève $x$. L'instruction
résultante est la déclaration de $x$ avec la nouvelle valeur $v'$ de $x$ après
le pas d'exécution\footnote{
    On peut remarquer qu'il est impossible de définir un contexte d'évaluation
    $C \gramisa \iDecl{x}{v}{C}$. En effet, puisque celui-ci nécessiterait
    d'ajouter une variable, il ne préserve pas la mémoire.
}.

\label{page:decl-masquage}
On suppose qu'il n'y a pas de masquage au sein d'une fonction, c'est-à-dire que
le nom d'une variable déclarée n'est pas visible avant cette déclaration.

Si $i$ est terminale ($\iPass$ ou $\iReturn{v}$), alors on peut s'évaluer en $i$
en nettoyant l'espace mémoire des références à $x$ qui peuvent subsister.

Enfin, si une erreur se produit elle est propagée.

\begin{mathpar}
 \semrule{Decl-Pass}

 \semrule{Decl-Return}

 \semrule{Decl-Ctx}

 \semrule{Decl-Err}
\end{mathpar}

Pour traiter l'alternative, on a besoin de 2 règles. Elles commencent de la même
manière, en évaluant la condition. Si le résultat est 0 (et seulement dans ce
cas), c'est la règle \textsc{If-False} qui est appliquée et l'instruction
revient à évaluer la branche «\emph{else}». Dans les autres cas, c'est la règle
\textsc{If-True} qui s'applique et la branche «\emph{then}» qui est prise.

\begin{mathpar}
  \semrule{If-False}

  \semrule{If-True}
\end{mathpar}

On exprime la sémantique de la boucle comme une simple règle de réécriture:

\begin{mathpar}
  \semrule{While}
\end{mathpar}

Enfin, si un $\phx\iReturn$ apparaît dans une séquence, on peut supprimer la suite:

\begin{mathpar}
  \semrule{Return}
\end{mathpar}

\section{Erreurs}
\label{sec:eval-errors}

Les erreurs se propagent des données vers l'interprète; c'est-à-dire que si
une expression ou instruction est réduite en une valeur d'erreur $Ω$, alors une
transition est faite vers cet état d'erreur.

Cela est aussi vrai d'une sous-expression ou sous-instruction: si l'évaluation
de $e_1$ provoque une erreur, l'évaluation de $e_1 + e_2$ également. La notion
de sous-expression ou sous instruction est définie en fonction des contextes
$C$. Notons que dans \textsc{Eval-Err}, $\ctxSub{C}{e}$ peut être une expression
ou une instruction.

\begin{mathpar}
  \semrule{Exp-Err}

  \semrule{Eval-Err}
\end{mathpar}

\section{Phrases et exécution d'un programme}

Un programme est constitué d'une suite de phrases: déclarations de fonctions,
de variables et de types, et évaluation d'expressions.

Donc l'évaluation d'une phrase $p$ fait passer d'un état mémoire $m$
à un autre $m'$, ce que l'on note $\ph{m}{p}{m'}$.

L'évaluation d'une expression est uniquement faite pour ses effets de bord. Par
exemple, après avoir défini les fonctions du programme, on pourra appeller
$\textrm{main}()$. La déclaration d'une variable globale (avec un initialiseur),
quant à elle consiste à évaluer cet initialiseur et à étendre l'état mémoire
avec ce couple (variable, valeur). On suppose que les variables globales ont
toutes des noms différents. Notons que ces évaluations se font à grands pas.

Enfin, l'exécution d'un programme est sans surprise l'exécution de ses phrases,
les unes à la suite des autres.

\begin{mathpar}
  \semrule{ET-Exp}

  \semrule{ET-Var}

  \semrule{T*-Nil}

  \semrule{T*-Cons}

  \semrule{Prog}
\end{mathpar}

\section*{Conclusion}

Nous venons de décrire un langage impératif, \langname. D'une part, sa syntaxe
et une description nous permettent de voir ce qui le rapproche de C. Sa
sémantique est explicitée entièrement, ce qui nous permet de prouver des
propriétés dynamiques sur les analyses de typage décrites dans les
chapitres~\ref{cha:typbase} et~\ref{cha:qualifs}.

% vim: spelllang=fr


\chapter{Typage}

\label{cha:typbase}
Dans ce chapitre, nous enrichissons le langage défini dans le
chapitre~\ref{cha:lang} d'un système de types. Celui-ci permet de séparer les
programmes bien formés, comme celui de la figure~\ref{fig:progwf:good} des
programmes mal formés comme celui de la figure~\ref{fig:progwf:bad}.

\begin{SaveVerbatim}[]{VerbProgGood}
f = Fun() {
  Decl x = 0 in
  x <- 1
  return x
}
\end{SaveVerbatim}

\begin{SaveVerbatim}[]{VerbProgBad}
f = Fun() {
  Decl x = 0 in
  x <- 1
  return (*x)
}
\end{SaveVerbatim}

\begin{figure}[h]

  \centering

  \subbottom[Programme bien formé]{
\label{fig:progwf:good}
    \begin{minipage}{5cm}
    \BUseVerbatim{VerbProgGood}
    \end{minipage}
  }
  \hspace{2cm}
  \subbottom[Programme mal formé]{
\label{fig:progwf:bad}
    \begin{minipage}{5cm}
    \BUseVerbatim{VerbProgBad}
    \end{minipage}
  }

  \caption{Programmes bien et mal formés}
\label{fig:progwf}

\end{figure}

Le but d'un tel système de types est de rejeter les programmes pour lesquels ont
peut facilement déterminer qu'ils sont faux, c'est-à-dire dont on peut prouver
qu'il provoqueraient des erreurs à l'exécution dues à une incompatibilité entre
valeurs. En ajoutant cette étape, on restreint la classe d'erreurs qui
pourraient bloquer la sémantique.

On emploie un système de types monomorphes: à chaque expression, on associe un
type, et non un schéma un de types. Il y a deux exceptions : tout d'abord, le
pointeur $\eNull$ a pour type $t~*$. Mais cela ne pose pas de complications, car
il n'y a pas d'étape de généralisation de $t$. Cela indique seulement que
$\eNull$ est compatible avec n'importe quel type pointeur. L'autre cas concerne
les tests d'égalité, qu'on veut pouvoir faire sur les valeurs de types scalaires
et tous les types qui sont composés de ces types (on veut par exemple éviter de
tester l'égalité de deux fonctions). Là encore, cela ne pose pas de problème,
puisqu'on peut statiquement savoir si un type a cette propriété ou pas.

En plus des types de base $\tInt$, $\tFloat$ et $\tUnit$, on peut construire des
types composés: pointeurs, tableaux, structures et fonctions.

Le principe du typage est d'associer à chaque construction syntaxique une
étiquette représentant le genre de valeurs qu'elle produira. Dans le programme
de la figure~\ref{fig:progwf:good}, la variable $x$ est initialisée avec la
valeur $0$, c'est donc un entier. Cela signifie que dans tout le programme,
toutes les instances de cette variable
\footnote{Deux variables peuvent avoir le même nom dans deux fonctions
  différentes, par exemple. Dans ce cas il n'y a aucune contrainte particulière
  entre ces deux variables. L'analyse de typage se fait toujours dans un
  contexte précis.
}
porteront ce type. La première instruction est l'affectation de la constante $1$
(entière) à $x$ dont on sait qu'elle porte des valeurs entières, ce qui est donc
correct. Le fait de rencontrer $\iReturn{x}$ permet de conclure que le type de
la fonction est $(~) → \tInt$.

Dans la seconde fonction, au contraire, l'opérateur $*$ est appliqué à $x$ (le
début de l'analyse est identique et permet de conclure que $x$ porte des valeurs
entières). Or cet opérateur prend un argument d'un type pointeur de la forme
$t*$ et renvoie alors une valeur de type $t$. Ceci est valable pour tout $t$
($\tInt$, $\tFloat$ ou même $t'*$: le déréférencement d'un pointeur sur pointeur
donne un pointeur), mais le type de $x$, $\tInt$, n'est pas de cette forme. Ce
programme est donc mal typé.

Dans ce chapitre, on commence par poser les notations qui vont servir à définir
la relation de typage. Ensuite, on explique les différentes règles de typage sur
les composantes de \langname: expressions, instructions et phrases. Enfin, dans
le reste du chapitre on établit des propriétés qui sont respectées par les
programmes bien typées. On conclut par les théorèmes de progrès et de
préservation qui établissent la sûreté du typage.

\section{Environnements et notations}

Les types associés aux expressions sont décrits dans la
figure~\ref{fig:les-types}. Tous sont des types concrets: il n'y a pas de
polymorphisme.

\begin{figure}[h]

  \begin{align*}
  \gramdef{Type}{t}
      { \tInt                       }{Entier}
      { \tFloat                     }{Flottant}
      { \tUnit                      }{Unité}
      { t*                          }{Pointeur}
      { t[]                         }{Tableau}
      { S                           }{Structure}
      { (t_1, …, t_n) \rightarrow t }{Fonction}
      {END} \\
  \\
  \gramdef{Structure}{S}
      { \tStruct{ l_1 : t_1; … ; l_n : t_n } }{Structure simple}
      {END} \\
  \end{align*}

  \caption{Types et environnements de typage}

\label{fig:les-types}

\end{figure}

Pour maintenir les contextes de typage, un environnement $Γ$ associe un type à
un ensemble de variables.

\label{page:gamma-split}
Plus précisément, un environnement $Γ$ est composé de deux listes de couples (variable,
type): une pour les variables locales, et une pour les variables
globales
\footnote{
    Cette distinction est nécessaire pour les définition de fonction: on
    remplace la liste des variables locales, mais on conserve le type des
    variables globales.
}

Si $Γ = (Γ_G, Γ_L) = ({(γ_i)}_{i∈[1;n]}, {(η_i)}_{i∈[1;m]})$,
avec $γ_i = (g_i, t_i)$ et $η_i = (l_i, u_i)$,
on utilise les notations suivantes:

\begin{align*}
x:t ∈ Γ &\eqdef ∃i ∈ [1;n], γ_i = (x, t) ∨ ∃i ∈ [1;m], η_i = (x, t) \\
\cDom{Γ_G} &\eqdef \{ g_i / i ∈ [1;n] \} \\
\cDom{Γ_L} &\eqdef \{ l_i / i ∈ [1;m] \} \\
\cDom{Γ} &\eqdef \cDom{Γ_G} ∪ \cDom{Γ_L} \\
\extendGlobal{Γ}{x}{t} &\eqdef ({(γ'_i)}_{i∈[1;n+1]}, Γ_L) ~\textrm{tel que}
                        \begin{cases}
                          ∀ i ∈ [1;n], γ'_i    = γ_i \\
                                       γ_{n+1} = (x, t) \\
                        \end{cases} \\
\extendLocal{Γ}{x}{t} &\eqdef (Γ_G, {(η'_i)}_{i∈[1;m+1]}) ~\textrm{tel que}
                        \begin{cases}
                          ∀ i ∈ [1;m], η'_i    = η_i \\
                                       η_{n+1} = (x, t) \\
                        \end{cases}
\end{align*}

Le type des fonctions semble faire apparaître un n-uplet $(t_1, …, t_n)$ mais ce
n'est qu'une notation: il n'y a pas de n-uplets de première classe, ils sont
toujours présents dans un type fonctionnel.

Le typage correspond à la définition des trois jugements suivants. Les deux
premiers sont mutuellement récursifs car une instruction peut consister en
l'évaluation d'une expression, et la définition d'une fonction repose sur le
typage de son corps.

\paragraph{Typage d'une expression:} on note de la manière suivante le fait
qu'une expression $e$ (telle que définie dans la figure~\ref{fig:stx-data}) ait
pour type $t$ dans le contexte $Γ$.

  \[
    \ty{Γ}{e}{t}
  \]

\paragraph{Typage d'une instruction:} les instructions n'ont en revanche pas de
type. Mais il est tout de même nécessaire de vérifier que toutes les
sous-expressions apparaissant dans une instruction sont cohérentes ensemble.

On note de la manière suivante le fait que sous l'environnement $Γ$
l'instruction $i$ est bien typée:

  \[
    \tyi{Γ}{i}
  \]

\paragraph{Typage d'une phrase:} De par leur nature séquentielle, les phrases
qui composent un programme altèrent l'environnement de typage. Par exemple, la
déclaration d'une variable globale ajoute une valeur dans l'environnement.

On note

  \[
    \typh{Γ}{p}{Γ'}
  \]

si le typage de la phrase $p$ transforme l'environnement $Γ$ en $Γ'$.

On étend cette notation aux suites de phrases, ce qui définit le typage d'un
programme, ce que l'on note $⊢ P$.

\section{Expressions}

\subsection*{Littéraux}

Le typage des littéraux numériques ne dépend pas de l'environnement de typage:
ce sont toujours des entiers ou des flottants.

\begin{mathpar}

  \disprule{Cst-Int}

  \disprule{Cst-Float}

\end{mathpar}

Le pointeur nul, quant à lui, est compatible avec tous les types pointeur.

\begin{mathpar}
  \disprule{Cst-Null}
\end{mathpar}

Enfin, le littéral unité a le type $\tUnit$.

\begin{mathpar}
  \disprule{Cst-Unit}
\end{mathpar}

\subsection*{Left-values}

Rappelons que l'environnement de typage $Γ$ contient le type des variables
accessibles du programme. Le cas où la left-value à typer est une variable est
donc direct: il suffit de retrouver son type dans l'environnement.

\begin{mathpar}
  \disprule{Lv-Var}
\end{mathpar}

Dans le cas d'un déréférencement, on commence par typer la left-value
déréférencée. Si elle a un type pointeur, la valeur déréférencée est du type
pointé.

\begin{mathpar}
  \disprule{Lv-Deref}
\end{mathpar}

Pour une left-value indexée (l'accès à tableau), on s'assure que l'indice soit
entier, et que la left-value a un type tableau: le type de l'élement est encore
une fois le type de base du type tableau ($t$ pour $t[]$).

\begin{mathpar}
  \disprule{Lv-Index}
\end{mathpar}

Le typage de l'accès à un champ est facilité par le fait que dans le programme,
le type complet de la structure est accessible sur chaque accès.

Dans la définition de cette règle on utilise la notation:

\[
(l, t) ∈ \eStruct{l_1 : t_1 ; … ; l_n : t_n }
\eqdef
∃ i ∈ [1;n],
l = l_i ∧ t = t_i
\]

\begin{mathpar}
  \disprule{Lv-Field}
\end{mathpar}

\subsection*{Opérateurs}

Un certain nombre d'opérations est possible sur le type \tInt.

\begin{mathpar}
  \disprule{Op-Int}
\end{mathpar}

De même sur \tFloat.

\begin{mathpar}
  \disprule{Op-Float}
\end{mathpar}

Les opérateurs de comparaison peuvent s'appliquer à deux opérandes qui sont d'un
type qui supporte l'égalité. Ceci est représenté par un jugement
$\textsc{Eq}(t)$ qui est vrai pour les types $\tInt$, $\tFloat$ et pointeurs,
ainsi que les types composés si les types de leurs composantes
(figure~\ref{fig:jugement-eq}). Les opérateurs $=$ et $≠$ renvoient alors un
\tInt:

\begin{mathpar}
  \disprule{Op-Eq}
\end{mathpar}

\begin{figure}[h]

  \ruleheader{$\textsc{Eq}(t)$}

  \begin{mathpar}
    \disprule{Eq-Num}

    \disprule{Eq-Ptr}

    \disprule{Eq-Array}

    \disprule{Eq-Struct}
  \end{mathpar}

\caption{Jugements d'égalité sur les types}
\label{fig:jugement-eq}
\end{figure}

Les opérateurs unaires «$+$» et «$-$» appliquent aux entiers, et leurs
équivalents «$\floatop{+}$» et «$\floatop{-}$» aux flottants.

\begin{mathpar}
\disprule{Unop-Plus-Int}
\and
\disprule{Unop-Plus-Float}

\disprule{Unop-Minus-Int}
\and
\disprule{Unop-Minus-Float}
\end{mathpar}

Les opérateurs de négation unaires, en revanche, ne s'appliquent qu'aux
entiers.

\begin{mathpar}
  \disprule{Unop-Not}
\end{mathpar}

L'arithmétique de pointeurs préserve le type des pointeurs.

\begin{mathpar}
  \disprule{Ptr-Arith}
\end{mathpar}

\subsection*{Autres expressions}

Prendre l'adresse d'une left-value rend un type pointeur sur le type de
celle-ci.

\begin{mathpar}
  \disprule{Addr}
\end{mathpar}

Pour typer une affectation, on vérifie que la left-value (à gauche) et
l'expression (à droite) ont le même type. C'est alors le type résultat de
l'expression d'affectation.

\begin{mathpar}
  \disprule{Set}
\end{mathpar}

Un littéral tableau a pour type $t[]$ où $t$ est le type de chacun de ses
éléments.

\begin{mathpar}
  \disprule{Array}
\end{mathpar}

Un littéral de structure est bien typé si ses champs sont bien typés.

\begin{mathpar}
  \disprule{Struct}
\end{mathpar}

Pour typer un appel de fonction, on s'assure que la fonction a bien un type
fonctionnel. On type alors chacun des arguments avec le type attendu. Le
résultat est du type de retour de la fonction.

\begin{mathpar}
  \disprule{Call}
\end{mathpar}

\section{Instructions}

La séquence est simple à traiter: l'instruction vide est toujours bien typée,
et la suite de deux instructions est bien typée si celles-ci le sont également.

\begin{mathpar}
  \disprule{Pass}

  \disprule{Seq}
\end{mathpar}

Une instruction constituée d'une expression est bien typée si celle-ci peut être
typée dans ce même contexte.

\begin{mathpar}
  \disprule{Exp}
\end{mathpar}

Une déclaration de variable est bien typée si son bloc interne est bien typé
quand on ajoute à l'environnement la variable avec le type de son initialiseur.

\begin{mathpar}
  \disprule{Decl}
\end{mathpar}

Les constructions de contrôle sont bien typées si leurs sous-instructions sont
bien typées, et si la condition est d'un type entier.

\begin{mathpar}
  \disprule{If}

  \disprule{While}
\end{mathpar}

\section{Fonctions}

Le typage des fonctions fait intervenir une variable virtuelle $\vRet$. Cela
revient à typer l'instruction $\iReturn{e}$ comme $\vRet ← e$. Cela rappelle le
langage Pascal, où pour retourner une valeur on l'affecte à une variable nommée
comme la fonction courante
\footnote{
    Si on n'avait pas introduit la restriction que chaque fonction doit terminer
    par un $\iReturn{\cdot}$ (page~\pageref{page:return-fonction}), alors le type de
    $\vRet$ pourrait rester inconnu. En pratique cela veut dire que la valeur de
    retour d'une telle fonction serait compatible avec n'importe quel type $t$, ce
    qui briserait la sûreté du typage.
}.

\begin{mathpar}
  \disprule{Return}
\end{mathpar}

Pour typer une définition de fonction, on commence par créer un nouvel
environnement de typage $Γ'$ obtenu par la suite d'opérations suivantes:

\begin{itemize}
\item
  on enlève l'ensemble des locales. Cela inclut le couple $\vRet : t_f$
  correspondant à la valeur de retour de la fonction appelante.
\item
  on ajoute les types des arguments $a_i : t_i$
\item
  on ajoute le type de la valeur de retour de la fonction appelée,
  $\vRet : t$
\end{itemize}

Si le corps de la fonction est bien typé sous $Γ'$, alors la fonction est
typable en $(t_1, …, t_n) → t$ sous $Γ$.

\begin{mathpar}
  \disprule{Fun}
\end{mathpar}

\section{Phrases}

Le typage des phrases est détaillé dans la figure~\ref{fig:typ-ph}. Le typage
d'une expression est le cas le plus simple. En effet, il y a juste à vérifier
que celle-ci est bien typable (avec ce type) dans l'environnement de départ:
l'environnement n'est pas modifié. En revanche, la déclaration d'une variable
globale commence de la même manière, mais on enrichit l'environnement de typage
des globales de cette nouvelle association.

\begin{figure}[h]

  \ruleheader{$\typh{Γ}{p}{Γ'}$}

  \begin{mathpar}
    \disprule{T-Exp}

    \disprule{T-Var}
  \end{mathpar}

  \ruleheader{$\typhstar{Γ}{ps}{Γ'}$}

  \begin{mathpar}
    \irule{T*-Nil}
      { }
      { \typhstar{Γ}{[~]}{Γ} }

    \irule{T*-Cons}
      { \typh{Γ_1}{p}{Γ_2}
     \\ \typhstar{Γ_2}{ps}{Γ_3}
      }
      { \typhstar{Γ_1}{p::ps}{Γ_3} }
  \end{mathpar}

  \ruleheader{$Γ ⊢ P$}

  \begin{mathpar}
    \irule{Prog}
      { \typhstar{[~]}{P}{Γ} }
      { ⊢ P }
  \end{mathpar}

  \caption{Typage des phrases et programmes}
\label{fig:typ-ph}

\end{figure}

\section{Sûreté du typage}

Comme nous l'évoquions au début de ce chapitre, le but du typage est de rejeter
certains programmes afin de ne garder que ceux qui ne provoquent pas un certain
type d'erreurs à l'exécution.

Dans la suite, nous donnons des propriétés que respectent tous les programmes
bien typés. Il est traditionnel de rappeler l'adage de Robin Milner:

\begin{quote}
  Well-typed programs don't go wrong.
\end{quote}

\emph{To go wrong} reste bien sûr à définir! Cette sûreté du typage repose sur
les deux théorèmes:

\begin{itemize}
\item progrès:
  si un terme est bien typé, il y a toujours une règle
  d'évaluation qui s'applique.
\item
  préservation (ou \emph{subject reduction}):
  l'évaluation transforme un terme bien typé en un terme du même type.
\end{itemize}

\section{Typage des valeurs}

Puisque nous allons manipuler les propriétés statiques et dynamiques des
programmes, nous allons avoir à traiter des environnements de typage $Γ$ et des
états mémoires $m$. La première chose à faire est donc d'établir une
correspondance entre ces deux mondes.

Étant donné un état mémoire $m$, on associe un type de valeur $τ$ aux valeurs
$v$. Cela est fait sous la forme d'un jugement $\semtyp{m}{v}{τ}$.

Ces types de valeurs ne sont pas exactement les mêmes que les types statiques.
Pour les calculer, on n'a pas accès au code du programme, seulement à ses
données. Il est par exemple possible de reconnaître le type des constantes, mais
pas celui des fonctions. Celles-ci sont en fait le seul cas qu'il est impossible
de déterminer statiquement. On le remplace donc par un cas plus simple où seul
l'arité est conservée.

Le cas des références (règle \textsc{S-Ptr}) utilise le typage des left-values,
codéfini par:

\[
    \semtypphi{m}{φ}{τ}
    \eqdef
    \semtyp{m}{m[φ]_Φ}{τ}
\]

Les règles de définition du typage des valeurs sont données dans la
figure~\ref{fig:types-semantiques}.

\begin{figure}[h]%{{{
\begin{align*}
\gramdef[2.5cm]{Type\\de valeur}{τ}
    { \tInt                       }{Entier}
    { \tFloat                     }{Flottant}
    { \tUnit                      }{Unité}
    { τ~*                         }{Pointeur}
    { τ[~]                        }{Tableau}
    { \tStruct{l_1:τ_1;…;l_n:τ_n} }{Structure}
    { \stFun{n}                   }{Fonction}
    {END}
\end{align*}

\caption{Types de valeurs}
\label{fig:types-semantiques}
\end{figure}%}}}

Les règles sont détaillées dans la figure~\ref{fig:regles-typ-sem}: les types
des constantes sont simples à retrouver car il y a assez d'information en
mémoire. Pour les références, ce qui peut être déréférencé en une valeur de type
$τ$ est un $τ~*$. Le typage des valeurs composées se fait en profondeur. Enfin,
la seule information restant à l'exécution sur les fonctions est son arité.

\begin{figure}[h]%{{{
  \ruleheader{$\semtyp{m}{v}{τ}$}

  \begin{mathpar}
    \disprule{S-Int}

    \disprule{S-Float}

    \disprule{S-Unit}

    \disprule{S-Null}

    \disprule{S-Ptr}

    \disprule{S-Array}

    \disprule{S-Struct}

    \disprule{S-Fun}
  \end{mathpar}

  \caption{Règles de typage des valeurs}
\label{fig:regles-typ-sem}

\end{figure}%}}}

La prochaine étape est de définir une relation de compatibilité entre les types
de valeurs $τ$ et statiques $t$. Nous noterons ceci sous la forme d'un jugement
$\tComp{τ}{t}$. Les règles sont décrites dans la
figure~\ref{fig:regles-comp-typ}, la règle importante étant \textsc{Comp-Fun}.
Notons qu'on garde le même nom pour les types de base, et que par exemple
$\tInt$ peut être vu soit comme un type statique, soit comme un type de valeur.
Il y a donc un abus de notation dans la règle \textsc{Comp-Ground}: quand on
note $\tComp{\tInt}{\tInt}$, le premier désigne le type des valeurs à
l'exécution, et le second le type statique.

\begin{figure}[h]%{{{
  \ruleheader{$\tComp{τ}{t}$}

  \begin{mathpar}
    \irule{Comp-Ground}
      { t ∈ \{ \tInt{}, \tFloat{}, \tUnit{} \} }
      { \tComp{t}{t} }

    \irule{Comp-Ptr}
      { \tComp{τ}{t} }
      { \tComp{τ~*}{t~*} }

    \irule{Comp-Array}
      { \tComp{τ}{t} }
      { \tComp{τ[~]}{t[~]} }

    \irule{Comp-Struct}
      { ∀ i ∈ [1;n]. \tComp{τ_i}{t_i} }
      { \tComp{\tStruct{l_1:τ_1;…;l_n:τ_n}}
              {\tStruct{l_1:t_1;…;l_n:t_n}}
      }

    \irule{Comp-Fun}
      { }
      { \tComp{\stFun{n}}
              {(t_1, …, t_n) → t}
      }
  \end{mathpar}

  \caption{Compatibilité entre types de valeurs et statiques}
\label{fig:regles-comp-typ}
\end{figure}%}}}

%Grâce à ce jugement, on peut donner la définition suivante.

%\begin{definition}[État mémoire bien typé]
  %On dit qu'un état mémoire $m$ est bien typé sous un environnement $Γ$, ce que
  %l'on note $\mcomp{Γ}{m}$, si les domaines de $Γ$ et $m$ coïncident sur la
  %partie visible, et que les types des valeurs associées aux variables visibles coïncident
  %avec leurs types statiques.

  %\[
    %\mcomp{Γ}{m} \eqdef
        %\begin{cases}
            %\cDom{Γ} = \{ \cVarname{a} / a ∈ \cVisible{m} \} \\
            %∀ a ∈ \cVisible{m}.
            %∃ τ,t.
            %\begin{cases}
                %Γ ⊢ \cVarname{a} : t
                %\\
                %\semtyp{m}{m[a]_A}{τ}
                %\\
                %\tComp{τ}{t}
            %\end{cases}
        %\end{cases}
  %\]

    %$\cVisible{m}$ désigne l'ensemble des adresses du cadre de pile le plus
    %récents ainsi que celles des variables globales:

%\[
%\cVisible{(s, g)} \eqdef \{ (x_i) / (x_i ↦ v_i) ∈ g \}
                       %∪ \{ (|s|, x_i) / (x_i ↦ v_i) ∈ s_0 \}
%\]

    %$\cVarname{a}$ désigne le nom simple d'une variable :

%\begin{align*}
%\cVarname{(x)} &\eqdef x \\
%\cVarname{(n,x)} &\eqdef x
%\end{align*}

%\end{definition}

On définit enfin la notion d'état mémoire bien typé. Cela se fait par induction
sur la forme de $Γ$ et $m$. Fonctionnellement, cela veut implique que les accès
à la mémoire retournent des valeurs en accord avec le type statique
(lemme~\ref{lemma:mem-typ}). Les règles définissant cette relation sont données
dans la figure~\ref{fig:comp-mem}.

\begin{figure}[h]%{{{

\centering

    \begin{mathpar}
        \irule{M-Empty}
            { }
            { \mcomp{[~]}{([~],[~])} }

        \irule{M-Global}
            { \mcomp{Γ}{(s, g)}
           \\ m ⊧ v : τ
           \\ \tComp{τ}{t}
            }
            { \mcomp{\extendGlobal{Γ}{x}{t}}{(s, ((x↦v)::g))} }

        \irule{M-Push}
            { \mcomp{Γ}{m}
           \\ Γ  = (Γ_G, Γ_L)
           \\ Γ' = (Γ_G, [x_1:t_1, …, x_n:t_n, \vRet:t] )
           \\ m ⊧ v_1 : τ_1 \\ \tComp{τ_1}{t_1}
           \\ …
           \\ m ⊧ v_n : τ_n \\ \tComp{τ_n}{t_n}
            }
            { \mcomp{Γ'}{\cPush{m}{(x_1↦v_1,…,x_n↦v_n)}}
            }

        \irule{M-Pop}
            { Γ = (Γ_G, Γ_L)
           \\ \mcomp{Γ}{m}
           \\ m' = \cPush{m}{x_1↦v_1, …, x_n↦v_n}
           \\ \mcomp{(Γ_G, [x_1:t_1, …, x_n:t_n, \vRet:t])}{m'}
           %\\ m' ⊧ v_1 : τ_1 \\ \tComp{τ_1}{t_1}
           %\\ …
           %\\ m' ⊧ v_n : τ_n \\ \tComp{τ_n}{t_n}
            }
            { \mcomp{Γ}{\cCleanup{\cPop{m'}} } }

        \irule{M-Decl}
            { \mcomp{Γ}{m}
           \\ m ⊧ v : τ
           \\ \tComp{τ}{t}
           \\ x ∉ Γ
            }
            { \mcomp{\extendLocal{Γ}{x}{t}}{\cExtend{m}{x}{v}} }

        \irule{M-DeclClean}
            { \mcomp{\extendLocal{Γ}{x}{t}}{m}
            }
            { \mcomp{Γ}{\cCleanVar{m - x}{x}} }

        \irule{M-Write}
            { \mcomp{Γ}{m}
           \\ Γ ⊢ φ : t
           \\ m ⊧ v : τ
           \\ \tComp{τ}{t}
           \\ m' = m[φ ← v]_Φ
            }
            {
              \mcomp{Γ}{m'}
            }
    \end{mathpar}

    \caption{Compatibilité entre états mémoire et environnements de typage}
\label{fig:comp-mem}

\end{figure}%}}}

\section{Propriétés du typage}

On commence par énoncer quelques lemmes utiles dans la démonstration de ces
théorèmes.

\begin{lemma}[Inversion]
\label{lemma:inversion}

  À partir d'un jugement de typage, on peut en déduire des informations sur les
  types de ses sous-expressions.

\begin{itemize}
\item
  Constantes
  \begin{itemize}
    \item si $Γ ⊢ n : t$, alors $t = \tInt$
    \item si $Γ ⊢ d : t$, alors $t = \tFloat$
    \item si $Γ ⊢ \eNull : t$, alors $∃ t', t = t'*$
    \item si $Γ ⊢ \eUnit : t$, alors $t = \tUnit$
  \end{itemize}

\item Références mémoire:
  \begin{itemize}
    \item
      si $Γ ⊢ x : t$, alors $x : t ∈ Γ$
    \item
      si $Γ ⊢ *e : t$, alors $Γ ⊢ e : t*$
    \item
      si $Γ ⊢ lv[e] : t$, alors $Γ ⊢ lv : t[]$ et $Γ ⊢ e : \tInt$
    \item
      si $Γ ⊢ lv.l_S : t$, alors $Γ ⊢ lv : S$

  \end{itemize}

\item Opérations:
  \begin{itemize}
    \item si $ Γ ⊢ \opun~e : t $, alors on est dans un des cas suivants:
      \begin{itemize}
        \item
          $\opun ∈ \{+,-,\sim, ! \}$,
          $t = \tInt$,
          $Γ ⊢ e : \tInt$
        \item
          $\opun ∈ \{+.,-.\}$,
          $t = \tFloat$,
          $Γ ⊢ e : \tFloat$
      \end{itemize}
    \item si $ Γ ⊢ e_1~\opbin~e_2 : t $, un des cas suivants se présente:
      \begin{itemize}
        \item
          $\opbin ∈ \opbinintset{}$,
          $Γ ⊢ e_1 : \tInt$,
          $Γ ⊢ e_2 : \tInt$,
          $t = \tInt$
        \item
          $\opbin ∈ \opbinfloatset{}$,
          $Γ ⊢ e_1 : \tFloat$,
          $Γ ⊢ e_2 : \tFloat$,
          $t = \tFloat$
        \item
          $\opbin ∈ \{=,≠\}$,
          $Γ ⊢ e_1 : t'$,
          $Γ ⊢ e_2 : t'$,
          $\textsc{Eq}(t')$,
          $t = \tInt$
        \item
          $\opbin ∈ \{≤,≥,<,>\}$,
          $t = \tInt$,
          $Γ ⊢ e_1 : t'$,
          $Γ ⊢ e_2 : t'$,
          $t' ∈ \{ \tInt, \tFloat \}$
        \item
          $\opbin ∈ \{+_p, -_p\}$,
          $∃t', t = t'*$,
          $Γ ⊢ e_1 : t'*$,
          $Γ ⊢ e_2 : \tInt$
      \end{itemize}
  \end{itemize}
\item
  Appel de fonction: si $Γ ⊢ e (e_1, …, e_n) : t$, il existe $(t_1, …, t_n)$
  tels que:

  \[
    \begin{cases}
      Γ ⊢ e : (t_1, …, t_n) → t \\
      ∀ i ∈ [1;n], Γ ⊢ e_i : t_i
    \end{cases}
  \]

\item Fonction: si $(Γ_G, Γ_L) ⊢ \eFun{a_1, …, a_n}{i} : t$, alors il existe
$(t_1, …, t_n)$ et $t'$ tels que:

  \[ \begin{cases}
    t' = (t_1, …, t_n) → t \\
    (Γ_G, [a_1:t_1, …, a_n:t_n, \vRet:t]) ⊢ i
  \end{cases} \]

\item Si $Γ ⊢ lv ← e : t$, alors $Γ ⊢ lv : t$ et $Γ ⊢ e : t$.
\item Si $Γ ⊢ \&~lv : t$, alors il existe $t'$ tel que
$Γ ⊢ lv : t'$ et $t = t'~*$.

\item Instructions:
\begin{itemize}
\item Si $Γ ⊢ i_1;i_2$, alors $Γ ⊢ i_1$ et $Γ ⊢ i_2$.
\item Si $Γ ⊢ e$, alors il existe $t$ tel que $Γ ⊢ e : t$.
\item Si $Γ ⊢ \iDecl{x}{e}{i}$, alors il existe $t$ tel que
      $Γ ⊢ e : t$ et $\extendLocal{Γ}{x}{t} ⊢ i$.
\item Si $Γ ⊢ \iIf{e}{i_t}{i_f}$, alors $Γ ⊢ e : \tInt$,
      $Γ ⊢ i_t$ et $Γ ⊢ i_f$.
\item Si $Γ ⊢ \iWhile{e}{i}$, alors $Γ ⊢ e : \tInt$ et $Γ ⊢ i$.
\item Si $Γ ⊢ \iReturn{e}$, alors il existe $t$ tel que $Γ ⊢ e : t$ et
$Γ ⊢ \vRet : t$.
\end{itemize}

\end{itemize}

\end{lemma}

\begin{proof}[Démonstration (esquisse)]
Pour chaque forme de jugement de typage, on liste les règles qui peuvent amener
à cette conclusion.
\end{proof}

Afin de pouvoir inverser la relation de typage pour les expressions de la forme
$\ctxSub{C}{e}$, on commence par donner un type aux contextes $C$.

\begin{definition}[Typage des contextes d'évaluation]

On définit une relation de typage pour les contextes d'évaluation $C$. Elle
permet de relier le type d'un sous-élément au type complet. Plusieurs cas se
distinguent, selon la nature (instruction d'un côté, expression ou left-value de
l'autre) du «trou» et de l'élement complet.

\begin{align*}
& \tyctx{Γ}{C}{t_i}{t_e} \\
& \tyctx{Γ}{C}{}{t_e} \\
& \tyctx{Γ}{C}{t_i}{} \\
& \tyctx{Γ}{C}{}{}
\end{align*}

Ces cas sont définis par les règles suivantes:

\begin{mathpar}
\irule{T-Ctx-EE}
    {\extendLocal{Γ}{x}{t_i} ⊢ \ctxSub{C}{x} : t_2}
    {\tyctx{Γ}{C}{t_i}{t_e}}

\irule{T-Ctx-IE}
    { ∀i. Γ ⊢ i ⇒ Γ ⊢ \ctxSub{C}{i} : t_e }
    { \tyctx{Γ}{C}{}{t_e} }
\\
\irule{T-Ctx-EI}
    { \extendLocal{Γ}{x}{t_i} ⊢ \ctxSub{C}{x} }
    { \tyctx{Γ}{C}{t_i}{} }

\irule{T-Ctx-II}
    { ∀i. Γ ⊢ i ⇒ Γ ⊢ \ctxSub{C}{i} }
    { \tyctx{Γ}{C}{}{} }
\end{mathpar}

\end{definition}

Par exemple, on peut avoir:

\begin{itemize}
\item $\tyctx{Γ}{1+\ctxEmpty}{\tInt}{\tInt}$
\item $\tyctx{Γ}{\iReturn{\ctxEmpty}}{\tInt}{}$
\item $\tyctx{Γ}{\eFun{x}{\ctxEmpty}}{}{(\tInt) → \tInt}$
\item $\tyctx{Γ}{\ctxEmpty;x ← 2}{}{}$
\end{itemize}

\begin{lemma}[Typage sous un contexte]
\label{lemma:typ-ctx}

Partant d'un jugement de typage de $C$ et de son «trou», on peut en déduire un
jugement de typage sur la totalité.

\begin{center}
\begin{tabular}{ll@{\hspace{1.5cm}}l}
\toprule
Sous-élément & Contexte & Conclusion \\
\midrule
$Γ ⊢ e : t_i$ &
$\tyctx{Γ}{C}{t_i}{t_e}$ &
$Γ ⊢ \ctxSub{C}{e} : t_e$ \\
$Γ ⊢ i$ &
$\tyctx{Γ}{C}{}{t_e}$ &
$Γ ⊢ \ctxSub{C}{e} : t_e$ \\
$Γ ⊢ e : t_i$ &
$\tyctx{Γ}{C}{t_i}{}$ &
$Γ ⊢ \ctxSub{C}{e}$ \\
$Γ ⊢ i$ &
$\tyctx{Γ}{C}{}{}$ &
$Γ ⊢ \ctxSub{C}{i}$ \\
\bottomrule
\end{tabular}
\end{center}

\end{lemma}

\begin{proof}[Démonstration (esquisse)]
Pour chaque cas, on procède en deux parties. D'une part, le jugement de typage
du contexte ne peut venir que d'une règle, ce qui permet d'inverser celui-ci.
La suite est différente selon les cas.

\paragraph{\textsc{T-Ctx-IE} et \textsc{T-Ctx-II}:}

Il suffit alors d'instancier le $i$ correspondant dans la prémisse de la règle
($∀i, Γ ⊢ i ⇒ …$).

\paragraph{\textsc{T-Ctx-EE} et \textsc{T-Ctx-EI}:}

Ce cas repose sur un lemme de substitution admis ici. Si une construction est
bien typée sous $\extendLocal{Γ}{x}{t}$, alors reste bien typée sous $Γ$ si l'on
y remplace $x$ par une expression de type $t$.
\end{proof}

\begin{lemma}[Inversion du typage sous un contexte]
\label{lemma:inv-ctx}

Partant d'une relation de typage d'une entité syntaxique (expression,
instruction ou left-value) de la forme $\ctxSub{C}{x}$, on peut en déduire une
information de typage sur $C$ et $x$.

\begin{center}
\begin{tabular}{l@{\hspace{1.5cm}}ll}
\toprule
Hypothèse & \multicolumn{2}{l}{Conclusion: $∃ t',$ …} \\
\midrule
$Γ ⊢ \ctxSub{C}{e} : t$ &
$Γ ⊢ e : t'$ &
$\tyctx{Γ}{C}{t'}{t}$ \\
$Γ ⊢ \ctxSub{C}{i} : t$ &
$Γ ⊢ i$ &
$\tyctx{Γ}{C}{}{t}$ \\
$Γ ⊢ \ctxSub{C}{e}$ &
$Γ ⊢ e : t'$ &
$\tyctx{Γ}{C}{t'}{}$ \\
$Γ ⊢ \ctxSub{C}{i}$ &
$Γ ⊢ i$ &
$\tyctx{Γ}{C}{}{}$ \\
\bottomrule
\end{tabular}
\end{center}

\end{lemma}

\begin{proof}[Démonstration (esquisse)]
On procède par induction sur la forme de $C$.
\end{proof}

Il est aussi possible de réaliser l'opération inverse: à partir du type d'une
valeur, on peut déterminer sa forme syntaxique. C'est bien sûr uniquement
possible pour les valeurs, pas pour n'importe quelle expression (par exemple
l'expression $x$ (variable) peut avoir n'importe quel type $t$ dans le contexte
$Γ = x:t$).

\begin{lemma}[Formes canoniques]
\label{lemma:canon}

  Il est possible de déterminer la forme syntaxique d'une valeur étant donné son
  type, comme décrit dans le tableau suivant. Par exemple, d'après la première
  ligne, si $Γ ⊢ v : \tInt$, alors $∃ n, v = d$.

  \begin{center}
  \begin{tabular}{cc}
  \toprule
  Type de $v$                     & Forme de $v$ \\
  \midrule
  $\tInt$                         & $n$ \\
  $\tFloat$                       & $d$ \\
  $\tUnit$                        & $\eUnit$ \\
  $t*$                            & $\wphi{φ}$ ou $\eNull$ \\
  $t[~]$                          & $\eArray{v_1; …; v_n}$ \\
  $\tStruct{l_1 : t_1;…;l_n:t_n}$ & $\eStruct{l_1 : v_1;…;l_n:v_n}$ \\
  $(t_1, …, t_n) → t$             & $\eFun{a_1, …, a_n}{i}$ \\
  \bottomrule
  \end{tabular}
  \end{center}

\end{lemma}

\begin{proof}[Démonstration (esquisse)]
On fait l'inventaire, pour chaque forme syntaxique, des règles pouvant arriver à
cette dérivation.
\end{proof}

\begin{lemma}[Hauteur des chemins typés]
\label{lemma:hauteur-chem}

    Une valeur typée ne peut jamais pointer au dessus du niveau courant de pile.
($\phx\cH$ provient de la définition~\ref{def:hauteur-val},
page~\pageref{def:hauteur-val}).

    Si $m ⊧ v : τ$, alors $\cH{v} ≤ |m|$.

\end{lemma}

\begin{proof}%{{{
On procède par induction sur la forme de $v$.

\paragraph{$\widehat{c}$:} Alors $\cH{v} = -1$. Comme $|m| ≥ 0$, ce cas est
établi.

\paragraph{$\widehat{f}$:} Idem.

\paragraph{$\widehat{\&}~(x)$:} Idem.

\paragraph{$\widehat{\&}~(n, x)$:}
D'après la forme de $v$, la dernière règle appliquée dans la dérivation de
$m ⊧ v : τ$ est \textsc{S-Ptr}, donc $m[(n, x)]_Φ$ est une valeur.
D'après la définition de $φ$, $n ≤ |m|$.

\paragraph{$\widehat{\&}~φ.l$:} %
On procède par induction sur $v' = \widehat{\&}~φ$.
Comme $\cH{\widehat{\&}~φ} ≤ |m|$ et
$\cH{\widehat{\&}~φ.l} = \cH{\widehat{\&}~φ}$,
on en déduit que
$\cH{\widehat{\&}~φ.l} ≤ |m|$

\paragraph{$\widehat{\&}~φ[n]$:}
Idem

\paragraph{$\widehat{\eStruct{l_1:v_1;…;l_n:v_n}}$:}
Par induction, $∀i ∈ [1;n], \cH{v_i} ≤ |m|$. Donc il en est de même pour leur
maximum, et $\cH{v} ≤ |m|$.

\paragraph{$\widehat{\eArray{v_1,…,v_n}}$:}
Idem.

\end{proof}%}}}


\begin{lemma}[Accès à une mémoire bien typée]
\label{lemma:mem-typ}

Si $\mcomp{Γ}{m}$, et $Γ ⊢ φ : t$, alors
$m[φ]_Φ$ est une valeur $v$ et
$m ⊧ v : τ$
où
$\tComp{τ}{t}$.

\end{lemma}

\begin{proof}

\def\cVar#1{\mathrm{Var}(#1)}
\def\cVars#1{\mathrm{Vars}(#1)} % TODO utile?

Commençons par donner un résultat préliminaire. On utilise les notations
suivantes:

\begin{itemize}

\item $\cVar{φ}$ la variable qui apparaît dans $φ$:

\begin{align*}
    \cVar{a}    &= a \\
    \cVar{φ.l}  &= \cVar{φ} \\
    \cVar{φ[n]} &= \cVar{φ}
\end{align*}

\item $\cVars{m}$ est l'ensemble des variables de $m$:

\[
\cVars{([s_1, …, s_n], g)} =
\{ (x) / x↦v ∈ g \}
\cup
\bigcup_{i ∈ [1;n]} \{ (n - i + 1, x) / x↦v ∈ s_i \}
\]

\end{itemize}

Si $\mcomp{Γ}{m}$ et $Γ ⊢ φ : t$, alors $\cVar{φ}$
est soit une variable globale $(x)$ avec $x ∈ \cDom{Γ_G}$, soit une variable locale
$(|m|, x)$ du plus haut cadre de pile avec $x ∈ \cDom{Γ_L}$. Cela se prouve par
récurrence sur la dérivation de $\mcomp{Γ}{m}$.

% TODO présenter informellement l'autre

À l'aide de ce résultat, on prouve le lemme~\ref{lemma:mem-typ} par induction
sur une dérivation de $\mcomp{Γ}{m}$.

\paragraph{\textsc{M-Empty}:} % {{{
$Γ_G = Γ_L = [~]$, la prémisse $Γ ⊢ φ : t$ est donc impossible à satisfaire.
% }}}
\paragraph{\textsc{M-Global}:} % {{{
Soient $φ$ tel que $\extendGlobal{Γ}{x}{t'} ⊢ φ : t$ et $m' = (s, ((x↦v)::g))$.
Alors la variable référencée par $φ$ est soit $(x)$, soit $(y)$ avec $y ∈
\cDom{Γ_G}$, soit $(|m|, y)$ avec $y ∈ \cDom{Γ_L}$.

Dans le premier cas, $m'[φ]_Φ = v$, ce qui permet de conclure.

Dans les autres cas, $\cVar{φ} ∈ \cVars{m}$. Donc $m'[φ]_Φ = m[φ]_Φ$, ce qui
nous permet de conclure grâce à l'hypothèse d'induction.
% }}}
\paragraph{\textsc{M-Decl}:} % {{{
On part de $Γ ⊢ φ : t'$. Alors $\cVar{φ}$ est soit $x$, soit une autre variable
locale, soit une globale. Dans le premier cas, $m'[φ]_Φ = v$ et les prémisses
nous permettent de conclure. Dans tous les autres cas, $m'[φ]_Φ = m[φ]$ et on
applique l'hypothèse d'induction.
% }}}
\paragraph{\textsc{M-DeclClean}:} % {{{

On suppose que $Γ ⊢ φ : t'$. Alors $\extendLocal{Γ}{x}{t} ⊢ φ : t'$ par
affaiblissement. On peut donc appliquer l'hypothèse d'induction: $m[φ]_Φ = v$ où
$m ⊧ v : τ'$ avec $\tComp{τ'}{t'}$. On distingue alors selon la forme de $v$.
Si $v = \wphi{φ'}$ où $\cVar{φ'} = (|m|, x)$, alors $m'[φ]_Φ = \eNull$ par
l'opération $\phxx\cCleanVar$. Le type $τ'$ étant un type pointeur par le
lemme~\ref{lemma:canon}, on peut conclure.
Dans les autres cas $m[φ]_Φ = m'[φ]_Φ$ ce qui termine ce cas.

% }}}
\paragraph{\textsc{M-Push}:} % {{{

On procède d'une manière similaire. $φ$ peut faire référence soit à un des
$x_i$, auquel cas la valeur $v_i$ convient ; soit à une variable globale, auquel
cas on applique l'hypothèse de récurrence.

% }}}
\paragraph{\textsc{M-Pop}:} % {{{

On part de $\mcomp{Γ}{m''}$ où $m'' = \cCleanup{\cPop{m'}}$. Deux cas se
produisent selon la forme de $\cVar{φ}$.

Si $\cVar{φ} = (x)$ avec $x ∈ \cDom{Γ_G}$, alors $Γ' ⊢ φ : t$ où $Γ' = (Γ_G,
[x_1:t_1, …, x_n:t_n, \vRet:t])$. On applique alors l'hypothèse de récurrence
sur $Γ ⊧ m'$: il vient que $m'[φ]_Φ = v$ où $m ⊧ v : τ'$ avec $\tComp{τ'}{t'}$.
Comme $\cH{v} ≤ |m'|$ (lemme~\ref{lemma:hauteur-chem}), deux cas peuvent se
produire. Si $\cH{v} = |m'|$, alors $m''[φ]_Φ = \eNull$ et on a bien la
compatibilité mémoire (l'argument est similaire au cas \textsc{Decl-Clean}).
Sinon, $m''[φ]_Φ = m'[φ]_Φ$ et on conclut directement.

Si $\cVar{φ} = (|m''|, x)$, on procède de la mème manière sauf qu'on invoque alors
le cas d'induction sur $\mcomp{Γ}{m}$.

% }}}
\paragraph{\textsc{M-Write}:} % {{{
On partant de $\mcomp{Γ}{m'}$ où $m' = m[φ ← v]_Φ$, et on suppose que
$Γ ⊢ φ' : t'$.

Si $φ = φ'$: alors il suffit d'appliquer \textsc{GetPut} à la lentille $Φ$:
$m[φ ← m[φ]_Φ]_Φ = m$, ce qui donne directement la conclusion.

Si $φ ≠ φ'$:
$Γ ⊢ φ' : t'$ donc $\cVar{φ'}$ est soit une locale soit une globale de $m'$. En
particulier $\cVar{φ'} ∈ \cVars{m'}$. Comme $φ ≠ φ'$, $m'[φ']_Φ = m[φ']_Φ$ et on
conclut grâce à l'hypothèse d'induction.
% }}}
\end{proof}


\section{Progrès et préservation}

Ces lemmes étant établis, on énonce maintenant le théorème de progrès.
Contrairement aux langages où tout est expression, il faut traiter séparément
les trois constructions principales de \langname : les expressions, les
left-values et les instructions. Celles-ci sont mutuellement dépendantes de par
le graphe suivant.

\begin{center}
\begin{tikzpicture}
[node distance=5cm]

\node[draw] (I) {Instructions};

\node[draw, right of=I] (E) {Expressions};

\node[draw, right of=E] (L) {Left-values};

\draw[->] (I) edge[bend left] node[auto] {Évaluation} (E);
\draw[->] (E) edge[bend left] node[auto] {Abstraction} (I);
\draw[->] (E) edge[bend left] node[auto] {Accès mémoire} (L);
\draw[->] (L) edge[bend left] node[auto] {Indices} (E);

\end{tikzpicture}
\end{center}

\begin{theorem}[Progrès]
\label{thm:progres}

Supposons que $Γ ⊢ i$. Soit $m$ un état mémoire tel que $\mcomp{Γ}{m}$.

Alors l'un des cas suivants est vrai:
\begin{itemize}
\item $i = \iPass$
\item $∃v, i = \iReturn{v}$
\item $∃ (i', m'), \mm{m}{i}{m'}{i'}$
\item $∃ Ω ∈ \{\serr{div},\serr{array},\serr{ptr}\}, \msi{m}{i} → Ω$
\end{itemize}

\jolibreak

  Supposons que $Γ ⊢ e : t$. Soit $m$ un état mémoire tel que $\mcomp{Γ}{m}$.
  Alors l'un des cas suivant est vrai:

\begin{itemize}
  \item $∃ v ≠ Ω, e = v$
  \item $∃ (e', m'), \mm{m}{e}{m'}{e'}$
  \item $∃ Ω ∈ \{\serr{div},\serr{array},\serr{ptr}\}, \msi{m}{e} → Ω$
\end{itemize}

\jolibreak

Supposons que $Γ ⊢ lv : t$. Soit $m$ un état mémoire tel que $\mcomp{Γ}{m}$.

Alors l'un des cas suivants est vrai:
\begin{itemize}
\item $∃φ, lv = φ$
\item $∃ (lv', m'), \mm{m}{lv}{m'}{lv'}$
\item $∃ Ω ∈ \{\serr{div},\serr{array},\serr{ptr}\}, \msi{m}{lv} → Ω$
\end{itemize}

\end{theorem}

C'est-à-dire, soit:

\begin{itemize}
  \item l'entité (instruction, expression ou left-value) est complètement
évaluée.
  \item un pas d'évaluation est possible.
  \item une erreur de division, tableau ou pointeur se produit.
\end{itemize}

La preuve du théorème~\ref{thm:progres} se trouve en annexe~\ref{proof:progres}.

\begin{lemma}[Représentabilité]
\label{lemma:repr}

On définit un opérateur de représentation d'un type statique à l'exécution:

    \begin{align*}
        \cRepr{ \tInt   } &= \tInt   \\
        \cRepr{ \tFloat } &= \tFloat \\
        \cRepr{ \tUnit  } &= \tUnit  \\
        \cRepr{ t'*      } &= \cRepr{t'}* \\
        \cRepr{ t'[]     } &= \cRepr{t'}[] \\
        \cRepr{ \tStruct{ l_1 : t_1; … ; l_n : t_n } }
        &= \tStruct{ l_1 : \cRepr{t_1}; … ; l_n : \cRepr{t_n} }\\
        \cRepr{ (t_1, …, t_n) \rightarrow t } &= \stFun{n}
    \end{align*}

Supposons que $Γ ⊢ v : t$ et $\mcomp{Γ}{m}$.
On pose $τ = \cRepr{t}$.
Alors $m ⊧ v : τ$ et $\tComp{τ}{t}$.

\end{lemma}

\begin{proof}%{{{

On procède par induction sur la forme de $t$.

\begin{itemize}
\item $\tInt$:
    D'après le lemme des formes canoniques, $v = n$.
    On conclut avec \textsc{S-Int} et \textsc{Comp-Ground}.

\item $\tFloat$:
    Idem avec $v = d$ et \textsc{S-Float}.

\item $\tUnit$:
    Idem avec $v = \eUnit$ et \textsc{S-Unit}.

\item $t = t'*$:
    Soient $τ' = \cRepr{t'}$ et $τ = τ'~*$.
    D'après le lemme des formes canoniques, deux cas sont possibles:

    \begin{itemize}
        \item $v = \widehat{\&}~φ$:

            Par inversion, $Γ ⊢ φ : t'$

            Puisque $\mcomp{Γ}{m}$ et $Φ ⊢ φ : t'$, on obtient par le
            lemme~\ref{lemma:mem-typ} que
            $m[φ]_Φ$ est une valeur telle que
            $m ⊧ m[φ]_Φ : τ'$ où $\tComp{τ'}{t'}$.
            D'après \textsc{S-Ptr}, $m ⊧ \widehat{\&}~φ : τ$.

            De plus par \textsc{Comp-Ptr}
            $\tComp{τ}{t}$.

        \item $v = \eNull$:

            Par induction, $\tComp{τ'}{t'}$. Alors par \textsc{Comp-Ptr},
            $\tComp{τ~*}{t~*}$.

            En outre, grâce à \textsc{S-Null}, on obtient $m ⊧ \eNull: τ~*$.

    \end{itemize}

\item $t = t'[]$:
    Par le lemme des formes canoniques,
    $v = \widehat{\eArray{v_1,…,v_n}}$.
    Par inversion on obtient que $∀i, Γ ⊢ v_i : t'$.

    Soient $τ' = \cRepr{t'}$ et $τ = τ'~[]$.

    Alors par induction $∀i, m ⊧ v_i : τ'$ et $\tComp{τ}{t}$.
    De la première propriété il vient (via \textsc{S-Array})
    $m ⊧ v : τ$, et de la seconde (via \textsc{Comp-Array})
    $\tComp{τ}{t}$.

\item $\tStruct{ l_1 : t_1; … ; l_n : t_n }$:
    Par le lemme des formes canoniques,
    $v = \widehat{\eStruct{l_1:v_1;…;l_n:v_n}}$.
    Et par inversion, $∀i, Γ ⊢ v_i : t_i$.

    Soient $τ_i = \cRepr{t_i}$ et
    $τ = \tStruct{ l_1 : τ_1; … ; l_n : τ_n }$.

    Alors par induction, $∀i, m ⊧ v_i : τ_i$ et $\tComp{τ_i}{t_i}$.

    On déduit de \textsc{S-Struct} que $m ⊧ v : τ$, et de
    \textsc{Comp-Struct} que $\tComp{τ}{t}$.

\item $t = (t_1, …, t_n) \rightarrow t'$:
    Par formes canoniques, on a
    $v = \widehat{\eFun{a_1,…,a_n}{i}}$.

    Soit $τ = \stFun{n}$: par \textsc{S-Fun} on obtient que
    $m ⊧ v : τ$. On conclut d'autre part que $\tComp{τ}{t}$
    grâce à \textsc{Comp-Fun}.

\end{itemize}
\end{proof}%}}}

\begin{theorem}[Préservation]
\label{thm:preservation}

Soient $Γ$ un environnement de typage, et $m$ un état mémoire tels que
$\mcomp{Γ}{m}$.

Alors:

\begin{itemize}
\item
    Si $Γ ⊢ lv : t$ et $\mm{m}{lv}{m'}{φ}$,
    alors $\mcomp{Γ}{\cCleanup{m'}}$ et $\semtypphi{m'}{φ}{τ}$ où $\tComp{τ}{t}$.

\item
    Si $Γ ⊢ lv : t$ et $\mm{m}{lv}{m'}{lv'}$,
    alors $\mcomp{Γ}{\cCleanup{m'}}$ et $Γ ⊢ lv' : t$.

\item
    Si $Γ ⊢ e : t$ et $\mm{m}{e}{m'}{v}$,
    alors $\mcomp{Γ}{\cCleanup{m'}}$ et $m' ⊧ v : τ$ où $\tComp{τ}{t}$.

\item
    Si $Γ ⊢ e : t$ et $\mm{m}{e}{m'}{e'}$,
    alors $\mcomp{Γ}{\cCleanup{m'}}$ et $Γ ⊢ e' : t$.

\item
    Si $Γ ⊢ i$ et $\mm{m}{i}{m'}{i'}$,
    alors $\mcomp{Γ}{\cCleanup{m'}}$ et $Γ ⊢ i'$.
\end{itemize}

  Autrement dit, si une construction est typable, alors un pas d'évaluation ne
  modifie pas son type et préserve le typage de la mémoire.

\end{theorem}

\paragraph{Remarque}

Dans la formulation classique de ce théorème, on indique que $\mcomp{Γ}{m}$
implique $\mcomp{Γ}{m'}$. Ici, la conclusion est moins forte en indiquant
seulement que $\mcomp{Γ}{\cCleanup{m'}}$. Cela indique que la compatibilité
mémoire est établie mais peut localement introduire des pointeurs fous. En fait,
comme une étape de $\phx\cCleanup$ est faite après chaque appel de fonction et
chaque déclaration, la propriété classique est vraie mais uniquement sur un plus
grand pas d'exécution.

La preuve de ce théorème se trouve en annexe~\ref{proof:preservation}.

Cela prouve qu'aucun terme ne reste «bloqué» parce qu'aucune règle ne
s'applique, et que la sémantique respecte le typage. En quelque sorte, les types
sont un contrat entre les expressions et les fonctions: si leur évaluation
converge, alors une valeur du type inféré sera produite.

Enfin, on donne une version de ces propriétés pour les phrases de programme.

\begin{theorem}[Progrès pour les phrases]
\label{thm:prog-phr}

Soient $Γ$ un environnement de typage, $m$ un état mémoire et $p$ une phrase de
programme. Supposons que $\typh{Γ}{p}{Γ'}$ et $\mcomp{Γ}{m}$.

On suppose en outre que l'évaluation de $p$ termine.

Alors $∃ m'. \ph{m}{p}{m'}$.

\end{theorem}

\begin{proof}

Ici il n'y a pas de difficulté puisque la contrainte (forte) de terminaison se
lit $\mm{m}{e}{m'}{v'}$ où $e$ est l'expression apparaissant dans $p$.

Selon la forme de $p$, il suffit alors d'appliquer la règle
\textsc{ET-Exp} ou \textsc{ET-Var}.

\end{proof}

\begin{theorem}[Préservation pour les phrases]
\label{thm:presa-phr}

On suppose que les trois propriétés suivantes sont vérifiées:

\[
\begin{cases}
    \mcomp{Γ}{m} \\
    \typh{Γ}{p}{Γ'} \\
    \ph{m}{p}{m'}
\end{cases}
\]

Alors $\mcomp{Γ'}{m'}$.

\end{theorem}

\begin{proof}

On distingue selon la dernière règle appliquée dans la dérivation de
$\ph{m}{p}{m'}$.

\paragraph{\textsc{ET-Exp}:}%{{{
La dernière règle appliquée pour dériver $\typh{Γ}{p}{Γ'}$ est donc
\textsc{T-Exp}. D'après les prémisses de ces deux règles, on a donc
$Γ ⊢ e : t$ et $\mm{m}{e}{m'}{v}$. Alors, d'après le théorème de préservation,
$\mcomp{Γ}{m'}$.
%}}}
\paragraph{\textsc{ET-Var}:}%{{{

Ici, la dérivation de $\typh{Γ}{p}{Γ'}$ termine par \textsc{T-Var}.
D'après leurs prémisses, on a donc :
$Γ ⊢ e : t$, $Γ' = \extendGlobal{Γ}{x}{t}$, et
$m'' = (s, (x ↦ v)::g)$ où $(s, g) = m'$
(on cherche à prouver que $\mcomp{Γ'}{m''}$).

En appliquant le théorème de préservation, on obtient que $Γ ⊢ v : t$ et
$\mcomp{Γ}{m'}$. D'après le lemme~\ref{lemma:repr}, il existe $τ$ tel que
$m' ⊧ v : τ$ où $\tComp{τ}{t}$. On peut alors appliquer \textsc{M-Global} qui
nous donne que $\mcomp{Γ'}{m''}$.
% }}}

\end{proof}

\section*{Conclusion}

Le langage \langname, ainsi que décrit dans le chapitre~\ref{cha:lang}, ne
possède pas de garanties de sûreté. C'est-à-dire qu'on ne peut rien prédire sur
l'exécution d'un programme. Par exemple, $\eFun{x}{\iPass} + 1$ est une
expression correcte mais pour laquelle il n'y a pas de règle d'évaluation qui
s'applique. En ajoutant un système de type, on s'assure que les termes peuvent
être évalués, et que les valeurs produites sont en accord avec les types données
aux différentes parties du programme. Cela permet surtout de s'assurer que les
programmes ne peuvent provoquer une erreur d'exécution que dans certains cas
particuliers, comme les division ou les accès aux tableaux.

% TODO trop court

% vim: spelllang=fr


\chapter{Extensions de typage}

\label{cha:qualifs}
% TODO

Nous venons de définir un système de types sûrs dans le
chapitre~\ref{cha:typbase}. Cela permet de mettre en relation les types des
expressions avec les valeurs qui leur seront associées. Cela permet une forme
d'analyse de flot: si peu de constructions permettent de créer des valeurs d'un
type $t$, alors toutes les valeurs de type $t$ proviennent de ces
\enquote{sources}.

On se propose ici d'enrichir le système de types de plusieurs extensions
permettant d'explorer cette idée, en ajoutant de la \enquote{signification} dans
les types de données des programmes. Ces extensions permettront de détecter des
erreurs de programmation communes, appuyées sur des exemples réels.

On commence par un exemple permettant distinguer plusieurs utilisations de types
entiers, selon s'ils sont utilisés comme entiers arithmétiques ou ensemble de
bits. Cela permet de détecter une erreur courante qui consiste à mélanger les
opérateurs logiques et bit à bit.

Ensuite, on étend de manière indépendante le système de types, cette fois au
niveau des pointeurs. Plus précisément, dans le contexte des systèmes
d'exploitation, introduit une différence entre les pointeurs dont la valeur est
contrôlée par l'utilisateur et les autres.

\section{Exemple préliminaire: les entiers utilisés comme bitmasks}

Dans le langage C, les types de données décrivent uniquement l'agencement en
mémoire des valeurs. Ils n'ont pas de signification plus sémantique permettant
d'exprimer ce que les données représentent. Par exemple, dans un programme
manipulant des dates, on sera amené à manipuler des numéros de mois et d'années,
représentés par des types entiers. Le langage C permet de définir des nouveaux
types:

\begin{verbatim}
typedef int month_t;
typedef int year_t;
\end{verbatim}

Cependant, rien ne distingue le nouveau type de l'ancien. Il ne s'agit que d'une
aide à la documentation. Dans cet exemple, \texttt{month\_t} et \texttt{year\_t}
sont tous les deux des nouveaux noms pour le type \texttt{int}, donc ils sont en
fait compatibles. Le compilateur ne peut donc pas détecter qu'on utilise un
numéro de mois là où un numéro d'année était attendu (ou \emph{vice versa}).

Cet idiome est commun en C. On manipule notamment certaines données abstraites
par des clés entières, et un \texttt{typedef} particulier permet de désigner
celles-ci. Par exemple sous Linux, les numéros de processus sont des indices
dans la table de processus interne au noyau, et on y accède par une valeur de
type \texttt{pid\_t}. De même, les utilisateurs sont représentés par un nombre
entier du type \texttt{uid\_t}.

Un autre idiome est répandu: l'utilisation d'entiers comme représentation d'un
ensemble de booléens. En effet, un nombre $a = \sum_{i=0}^{N-1} a_i 2^i$ peut
s'interpréter comme l'ensemble de ses bits égaux à 1: $\{ i ∈ [0;N-1] / a_i = 1
\}$. Un entier de 32 bits peut donc représenter une combinaison de 32 options
indépendantes.

\begin{figure}
    \insertcode{unix-open.c}
    \caption{Interface permettant d'ouvrir un fichier sous Unix}
    \label{fig:unix-open}
\end{figure}

C'est de cette manière que fonctionne l'interface qui permet d'ouvrir un fichier
sous Unix (figure~\ref{fig:unix-open}). Le paramètre \texttt{flags} est un
entier qui encode les options liées à l'ouverture du fichier. On précise son
mode (lecture, écriture ou les deux) par les bits 1 et 2, s'il faut créer le
fichier ou non s'il n'existe pas par le bit 7, si dans ce cas il doit être
effacé par le bit 8, etc. On obtient le paramètre complet en réalisant un
\enquote{ou} bit à bit entre des constantes. Le paramètre \texttt{mode} encode
de la même manière les permissions que doit avoir le fichier créé, le cas
échéant (\texttt{mode\_t} désigne en fait \texttt{unsigned int}).

Ces fonctions retournent un entier, qui est un \emph{descripteur de fichier}. Il
correspond à un indice numérique dans une table interne au processus. Par
exemple, $0$ désigne son entrée standard, $1$ sa sortie standard, et $2$ son
flux d'erreur standard.

On identifie donc au moins trois utilisations du type \texttt{int}:

\begin{itemize}
    \item entier: c'est l'utilisation classique pour représenter des valeurs
        numériques. Toutes les opérations sont possibles.
    \item bitmask: on utilise un entier comme ensemble de bits. Seules les
        opérations bit à bit ont du sens.
    \item entier opaque: on utilise un entier de manière purement abstraite.
        C'est l'exemple des descripteurs de fichier.
\end{itemize}

Ces utilisations du type n'ont rien à voir ; il faudrait donc empêcher
d'utiliser un descripteur de fichier comme un mode, et vice-versa. De même,
aucun opérateur n'a de sens sur les descripteurs de fichier, mais l'opérateur
\texttt{\textbar{}} du \enquote{ou} bit à bit doit rester possible pour les
modes.

On décrit ici une technique de typage pour détecter et interdire ces mauvaises
utilisations en proposant une version \enquote{bien typée} de la fonction
\texttt{open}. Plus précisément, on donne à ses deuxième et troisième arguments
(respectivement \texttt{flags} et \texttt{mode}) le nouveau type \tBits qui
correspond aux entiers utilisés comme bitmasks. Le type de retour n'est pas
modifié (il reste \tInt), mais on décrit comment il est possible de rendre ce
type opaque.

\subsection{Modifications}

On commence par ajouter deux types: d'une part \tBits bien sûr, mais également
\tChar qui apparaît dans les chaînes de caractères. On ne spécifie pas plus ce
dernier mais on suppose qu'il existe des littéraux de chaînes qui retournent un
pointeur vers le premier élément d'un tableaux de caractères. Pour rester
compatible avec C, on suppose qu'un caractère nul \texttt{'\textbackslash{}0'}
est inséré à la fin de la chaîne. On ajoute ces chaînes uniquement dans le but
de pouvoir représenter les noms de fichiers.

Au niveau des valeurs, les entiers utilisés comme bitmasks sont représentés par
des valeurs entières classiques $\widehat{n}$. En particulier, on n'ajoute pas
de nouveau type sémantique, mais on ajoute une règle de compatibilité entre le
type de valeurs \tInt et le type statique \tBits (cela signifie qu'une valeur de
type \tBits est représentée par un \tInt en mémoire,
figure~\ref{fig:ajouts-bitmasks}). Par ailleurs, on change le type des
\enquote{constructeurs} (\texttt{O\_*}) et du \enquote{consommateur}
\texttt{open} (figure~\ref{fig:bitmasks-types}).

\begin{figure}[h]

\begin{align*}
\gramdef{Type}{t}
  { … }{}
  { \tChar }{Caractère}
  { \tBits }{Entier utilisé comme \emph{bitmask}}
  {END}
\end{align*}

\[
  \irule{Comp-Bits}
    { }
    { \tComp{\tInt}{\tBits} }
\]

\caption{Ajouts liés aux entiers utilisés comme bitmasks}

\label{fig:ajouts-bitmasks}

\end{figure}

\begin{figure}[h]

\begin{align*}
    [~] ⊢ \texttt{O\_RDONLY} &: \tBits \\
    [~] ⊢ \texttt{O\_RDWR}   &: \tBits \\
    [~] ⊢ \texttt{O\_APPEND} &: \tBits \\
    \\
    [~] ⊢ \texttt{open} &: (\tChar~*, \tBits) → \tInt
\end{align*}

\caption{Nouvelles valeurs liées aux bitmasks}
\label{fig:bitmasks-types}

\end{figure}

Pour que les opérations bit à bit puissent s'appliquer aux bitmasks, on remplace
aux règles s'appliquant à $\tInt$ les règles suivantes.

\begin{mathpar}
    \irule{Op-Bits}
        { \opbin ∈ \{~|~, ~\&~, ~\opxor~ \}
       \\ Γ ⊢ e_1 : \tBits
       \\ Γ ⊢ e_2 : \tBits
       }{ Γ ⊢ e_1~\opbin~e_2 : \tBits
       }

    \irule{Not-Bits}
        { Γ ⊢ e : \tBits }
        { Γ ⊢ \sim e : \tBits }
\end{mathpar}

Il reste à permettre d'utiliser les bitmasks dans les contextes où on attend un
entier. Par exemple, pour écrire $\iIf{x~\&~\hexa{80}}{…}{…}$ (test du bit
numéro 7). On veut donc exprimer que \enquote{un \tBits est un \tInt}. Cette
relation entre différents types d'entier correspond à un cas particulier de sous-typage.

On ajoute la règle suivante. Elle permet d'utiliser une expression de type
$\tBits$ là où une expression de type $\tInt$ est attendue
\footnote{
    Cela modifie légèrement l'implantation de l'inférence de types. Le type
    d'une expression utilisée comme opérande de l'opérateur $+$ n'est donc pas
    \emph{a priori} de type \tInt, mais \tBits ou \tInt.
}.
% TODO[E] cela veut dire qu'on peut additionner un \tBits et un \tInt pour obtenir
% un \tInt ?

\begin{mathpar}
    \irule{Sub-BitsInt}
        { Γ ⊢ e : \tBits }
        { Γ ⊢ e : \tInt }
\end{mathpar}

Ainsi, si $Γ ⊢ e : \tBits$, on a par exemple $Γ ⊢ ~!~e : \tInt$. On rappelle que
la règle permettant de typer $!$ est inchangée et reste :

\[
    \disprule{Unop-Not}
\]

% Alternative
%
% ----------------------------------------------
% ! : s'applique à bits ou int mais retourne int
% ----------------------------------------------

\subsection{Exemple: $!~x~\&~y$}

% - !E1 & E2
% + !(E1 & E2)

Les nombreux opérateurs de C (repris en \langname) posent plusieurs problèmes:

\begin{itemize}
\item il sont nombreux et il est facile de confondre $\&\&$ avec $\&$, ou
      $!$ avec $\sim$;
\item il y a un opérateur \enquote{ou exclusif} bit à bit ($\opxor$) mais pas
      d'équivalent logique;
\item la priorité des opérateurs semble parfois arbitraire. Par exemple, les
      opérateurs de décalage sont plus prioritaires que les additions, donc
      $x \ll 2 + 1$ est interprété comme $(x \ll 2) + 1$.
\end{itemize}

Le premier et le dernier point permettent d'expliquer une erreur courante: celle
qui consiste à écrire $!~x~\&~y$ au lieu de $!~(x~\&~y)$.

En effet, la première expression est équivalente à $(!~x)~\&~y$. Comme $!~x$
vaut $0$ ou $1$, l'expression résultante vaut $y~\&~1$ si $x = 0$, ou $0$ sinon.
Il s'agit probablement d'une erreur de programmation. L'alternative $!~(x~\&~y)$
a plus de sens: elle vaut $0$ si $x$ et $y$ ont un bit en commun, $1$ sinon.

On vérifie enfin que la première n'est pas bien typée alors que la seconde
l'est. Dans les deux cas suivants on se place dans un environnement $Γ$
comportant deux variables globales $x$ et $y$ de type $\tBits$. Alors $(!~x) \&
y$ n'est pas bien typée. En effet, $Γ ⊢~!~x : \tInt$ et la seule règle qui
s'applique à l'opérateur $\&$ ne peut pas s'appliquer. En revanche la seconde
est bien typée (figure~\ref{fig:notxy-ok}).

\begin{figure}[h]
\begin{mathpar}
    Γ = ([x ↦ \tBits, y ↦ \tBits], [~])

    \irule{Unop-Not}
        {
            \irule{Sub-BitsInt}
                {
                    \irule{Op-Bits}
                        { Γ ⊢ x : \tBits
                       \\ Γ ⊢ y : \tBits
                        }
                        {Γ ⊢ x~\&~y : \tBits}
                }
                {Γ ⊢ x~\&~y : \tInt}
        }
        {Γ ⊢~!~(x~\&~y) : \tInt}
\end{mathpar}

\caption{Dérivation montrant que $!~(x~\&~y)$ est bien typée}
\label{fig:notxy-ok}
\end{figure}

Cela conclut cet exemple préliminaire, qui permet de voir en quoi \langname est
adapté à des analyses de typage légères. Puisque le typage est sûr, on en déduit
que les valeurs d'un certain type ne peuvent être créées que par un certain
nombre de constructeurs. Par exemple ici les bitmasks ne proviennent que de
combinaisons de constantes. C'est précisément cette idée de détection de source
qui est au cœur de l'analyse suivante.

\section{Analyse de provenance de pointeurs}

Jusqu'ici \langname{} est un langage impératif généraliste, ne prenant pas en
compte les spécificités de l'adressage utilisé dans les systèmes d'exploitation.

Dans cette section, on commence par l'étendre en ajoutant des constructions
modélisant les variables présentes dans l'espace utilisateur (cf.
chapitre~\ref{cha:os}). Pour accéder à celles-ci, on ajoute un opérateur de
déréférencement sûr qui vérifie à l'exécution que l'invariant suivant est
respecté:

\begin{quote}
Les pointeurs dont la valeur est contrôlée par l'utilisateur pointent vers
l'espace utilisateur.
\end{quote}

La terminologie mérite d'être détaillée:

Un pointeur contrôlé par l'utilisateur, ou \emph{pointeur utilisateur}, est une
référence mémoire dont la valeur est modifiable par le code utilisateur (opposé
au code noyau, que nous analysons ici). Ceci correspond à des données provenant
de l'extérieur du système vérifié. C'est une propriété statique, qui peut être
déterminée à la compilation à partir de considérations syntaxiques. Par exemple,
l'adresse d'une variable locale au sein de code noyau est toujours considérée
comme étant contrôlée par le noyau.

Un pointeur pointant vers l'espace utilisateur fait référence à une variable
allouée en espace utilisateur. Cela veut dire qu'y accéder ne risque pas de
mettre en péril l'isolation du noyau en faisant fuiter des informations
confidentielles ou en déjouant son intégrité. Cette propriété est dynamique: un
pointeur utilisateur peut \emph{a priori} pointer vers l'espace utilisateur, ou
non.

Pour prouver que l'invariant précédent est bien respecté, on procède en
plusieurs étapes.

Tout d'abord, on définit une nouvelle erreur $\serr{sec}$ (pour
\enquote{sécurité}), déclenchée lorsqu'un pointeur contrôlé par l'utilisateur et
pointant vers le noyau est déréférencé (le cas que l'on cherche à éviter). Il
est important de noter que ce cas d'erreur est \enquote{virtuel}, c'est-à-dire
qu'on l'ajoute à la sémantique pour pouvoir le détecter facilement comme un cas
d'erreur, mais dans une sémantique de plus bas niveau, comme en C, l'erreur ne
serait pas déclenchée. D'un point de vue opérationnel, cela équivaut à ajouter
un test dynamique à chaque déréférencement, ce qui est sûr mais se paye en
performances. Ajouter ce cas d'erreur virtuel dans la sémantique d'évaluation
permet de transformer un problème de sécurité (empêcher les fuites
d'information) en problème de sûreté (empêcher les erreurs à l'exécution).

Ensuite, on montre qu'avec cet ajout, si on étend naïvement le système de types
en donnant le même type aux pointeurs contrôlés par l'utilisateur et le noyau,
le théorème de progrès (\ref{thm:progres}) n'est plus valable. Cela signifie que
le système de types classique présenté dans le chapitre~\ref{cha:typbase} ne
suffit pas à capturer les propriétés de sécurité que nous voulons interdire.

L'étape suivante est d'étendre, à son tour, le système de types de \langname{}
en distinguant les types des pointeurs contrôlés par l'utilisateur des pointeurs
contrôlés par le noyau. Puisqu'on veut interdire le déréférencement des premiers
par l'opérateur \texttt{*}, on introduit les constructions $\cfuF$ et $\ctuF$
qui réaliseront le déréférencement sûr de ces pointeurs.

Enfin, une fois ces modifications faites, on prouve que les propriétés de
progrès et de préservation sont rétablies.

\subsection{Extensions noyau pour \langname}
\label{sec:extensions-noyau}

On ajoute à \langname{} la notion de valeur provenant de l'espace utilisateur.
Pour marquer la séparation entre les deux espaces d'adressage, on ajoute une
construction $φ \gramisa \vTainted{φ'}$. Le chemin interne $φ'$ désigne une
variable classique (un pointeur noyau) et l'opérateur $\phx\vTainted$ permet de
l'interpréter comme un pointeur vers l'espace utilisateur. En quelque sorte, on
ne classifie pas les valeurs selon la variable pointée mais selon la
construction du pointeur.

Remarquons qu'on n'introduit pas de sous-typage: les pointeurs noyau ne peuvent
être utilisés qu'en tant que pointeurs noyau, et les pointeurs utilisateurs
qu'en tant que pointeurs utilisateurs.

En plus du déréférencement par \texttt{*} (qui devra donc renvoyer
$\serr{sec}$ pour les valeurs de la forme $\vTainted{φ'}$), il faut aussi
ajouter des constructions de lecture et d'écriture à travers les pointeurs
utilisateurs. Ceci sera fait sous forme de deux fonctions, $\cfuF$ et \linebreak
$\ctuF$. Celles-ci prennent deux pointeurs en paramètre et renvoient un booléen
indiquant si la copie a pu être faite (si le paramètre contrôlé par
l'utilisateur pointe en espace noyau, les fonctions ne font pas la copie et
signalent l'erreur).

Illustrons ceci par un exemple. Imaginons un appel système fictif qui renvoie la
version du noyau, en remplissant par pointeur une structure contenant les champs
entiers \texttt{major}, \texttt{minor} et \texttt{patch} (un équivalent dans
Linux est l'appel système \verb!uname()!). Celui-ci peut être alors écrit
comme dans la figure~\ref{fig:ex-sys-getver}. Une fois la structure noyau $v$
remplie, il faut la copier vers l'espace utilisateur. La fonction $\ctuF$ va
réaliser cette copie (de la même manière qu'avec un \verb!memcpy()!), mais
après avoir testé dynamiquement que $p$ pointe en espace utilisateur (dans le
cas contraire, la copie n'est pas faite).

\begin{figure}[h]
\setlength{\parindent}{1cm}
\begin{align*}
  & \mathrm{sys\_getver} = \mathrm{fun}(p) \{ \\
  & \indent \textsc{Decl}~v= \{~\texttt{major} : 3
                              ; \texttt{minor} : 14
                              ; \texttt{patch} : 15~\}~\textsc{in} \\
  & \indent \uPut{p}{\& v} \\
  & \}
\end{align*}

\caption{Implantation d'un appel système qui remplit une structure par pointeur}
\label{fig:ex-sys-getver}
\end{figure}

On peut remarquer que, contrairement aux fonctions présentes dans le noyau
Linux, les fonctions $\cfuF$ et $\ctuF$ n'ont pas de paramètre indiquant la
taille à copier. Cela est dû au fait que le modèle mémoire de \langname{} est de
plus haut niveau. L'information de taille est déjà présente dans chaque valeur.

Une autre remarque à faire est qu'il n'y a pas de manière de copier des données
de l'espace utilisateur vers l'espace utilisateur. Il est nécessaire de passer
par l'espace noyau. La raison est que, puisqu'il faut réaliser deux tests
dynamiques, les erreurs peuvent arriver à ces deux endroits. Plutôt que de
proposer un opérateur qui réalise cette copie, on laisse le programmeur faire
les deux copies manuellement.

On commence donc par ajouter aux instructions des constructions $\phxx{\uPut}$
et $\phxx{\uGet}$ de copie sûre. Afin de leur donner une sémantique, il faut
étendre l'ensemble des valeurs pointeur $φ$ aux constructions de la forme
$\vTainted{φ'}$. Pour créer des termes s'évaluant en de telles valeurs, il faut
une construction syntaxique $\eTaint{e}$ telle que, si $e$ s'évalue en
$\wphi{φ}$, $\eTaint{e}$ s'évalue en $\wphi{\vTainted{φ}$}. Cela demande 2
autres ajouts: un nouveau contexte d'évaluation $\eTaint{C}$ et une règle
d'évaluation. Enfin, on ajoute une nouvelle erreur $\serr{sec}$ à déclencher
lorsqu'on déréférence directement un pointeur utilisateur. Ces étapes sont
résumées dans la figure~\ref{fig:qualif-changes}.

\begin{figure}%{{{

\begin{align*}
\gramdef{Expressions}{e}
  { … }{}
  { \eTaint{lv} }{Adresse utilisateur}
  { \uGet{e_d}{e_s} }{Lecture depuis l'espace utilisateur}
  { \uPut{e_d}{e_s} }{Écriture vers l'espace utilisateur}
  {END}
\\
\\
\gramdef{Contextes}{C}
  { … }{}
  { \eTaint{C}}{}
  { \uGet{C}{e} }{}
  { \uGet{v}{C} }{}
  { \uPut{C}{e} }{}
  { \uPut{v}{C} }{}
  {END}
\\
\\
\gramdef{Chemins}{φ}
  { … }{}
  { \vTainted{φ} }{Pointeur utilisateur}
  {END}
\\
\\
\gramdef{Erreurs}{Ω}
  { … }{}
  { \serr{sec} }{Erreur de sécurité}
  {END}
\end{align*}

\caption{Ajouts liés aux pointeurs utilisateurs (par rapport à l'interprète du
chapitre~\ref{cha:lang})}

\label{fig:qualif-changes}

\end{figure}%}}}

En résumé, on a deux constructions pour créer des pointeurs à partir d'une
valeur gauche: $\&~\cdot$ crée un pointeur noyau, et $\eTaint{\cdot}$ crée un
pointeur utilisateur. Seule la première est faite pour être utilisée dans le
code à analyser. La seconde sert uniquement à modéliser les points d'entrée du
noyau. Par exemple, la fonction $\mathrm{sys\_getver}$ de la
figure~\ref{fig:ex-sys-getver} peut être appelée par un utilisateur de la
manière suivante:

{
\setlength{\parindent}{1cm}
\begin{align*}
  & \indent \textsc{Decl}~v= \{~\texttt{major} : 0
                              ; \texttt{minor} : 0
                              ; \texttt{patch} : 0~\}~\textsc{in} \\
  & \indent \mathrm{sys\_getver} (\eTaint{v})
\end{align*}
}

\subsection{Extensions sémantiques}

En ce qui concerne l'évaluation des expressions $\phx\eTaint$, on ajoute la
règle suivante:

\[
  \semrule{Phi-User}
\]

Dans $\wphi{(\vTainted{φ})}$, l'opérateur $\phx\wphi$ indique que la valeur
créée est une référence mémoire. Cette référence mémoire, $\vTainted{φ}$, est
contrôlée par l'utilisateur. C'est ce qu'indique le constructeur $\phx\vTainted$

Ensuite, il est nécessaire d'adapter les règles d'accès à la mémoire pour
déclencher une erreur $\serr{sec}$ en cas de déréférencement d'un pointeur
utilisateur. Les accès mémoire en lecture proviennent de la règle
\textsc{Exp-Lv} et ceux en lecture, de la règle \textsc{Exp-Set}, rappellées
ici:

\begin{mathpar}
  \semrule{Exp-Lv}

  \semrule{Exp-Set}
\end{mathpar}

Les accès à la mémoire sont en effet faits par le biais de la lentille $Φ$. Il
suffit donc d'adapter sa définition (page~\pageref{subsec:acces-phi}) de
celle-ci en rajoutant les cas suivant:

\begin{align*}
    \lensGet{Φ}{\vTainted{φ}}    &= \serr{sec} \\
    \lensPut{Φ}{\vTainted{φ}}{v} &= \serr{sec}
\end{align*}

Enfin, il est nécessaire de donner une sémantique aux fonctions $\cfuF$ et
\linebreak $\ctuF$. L'idée est que celles-ci testent dynamiquement la
\emph{valeur} du paramètre \linebreak contrôlé par l'utilisateur afin de
vérifier que celui-ci pointe vers l'espace utilisateur (c'est-à-dire, qu'il est
de la forme $\vTainted{φ}$).

Deux cas peuvent se produire. Soit la partie à vérifier a la forme
$\vTainted{φ'}$, soit non (et dans ce cas $∄φ', φ = \vTainted{φ'}$). Dans le
premier cas (règles \textsc{User-*-OK}), alors la copie est faite et l'opération
de copie retourne la valeur entière $0$. Dans le second (règles
\textsc{User-*-Err}), aucune copie n'est faite et la valeur $-1$ est retournée.
Ce comportement est calé celui des fonctions \texttt{copy\_\{from,to\}\_user} du
noyau Linux: en cas de succès elles renvoient $0$, et en cas d'erreur
\texttt{-EFAULT} ($= -14$).

\begin{mathpar}
  \semrule{User-Get-OK}

  \semrule{User-Get-Err}

  \semrule{User-Put-OK}

  \semrule{User-Put-Err}
\end{mathpar}

Ces règles sont à appliquer en priorité de la règle d'appel de fonction
classique, puisqu'il s'agit d'élements de syntaxe différents. En effet ces
\enquote{fonctions} ne sont pas implantables directement en \langname, puisqu'il
n'y a pas par exemple d'opérateur permettant d'extraire $φ$ depuis une valeur
$\vTainted{φ}$. L'opération en \enquote{boîte noire} de ces deux fonctions
permet d'assurer que l'accès à l'espace utilisateur est toujours couplé à un
test dynamique.

\subsection{Insuffisance des types simples}

Étant donné \langname{} augmenté de cette extension sémantique, on peut étendre
naïvement le système de types avec la règle suivante:

\begin{mathpar}
  \irule{Addr-User-Ignore}
    { Γ ⊢ lv : t }
    { Γ ⊢ \eTaint{lv} : t~* }
\end{mathpar}

Cette règle est compatible avec l'extension, sauf qu'elle introduit des termes
qui sont bien typables mais dont l'évaluation provoque une erreur autre que
$\serr{div}$, $\serr{array}$ ou $\serr{ptr}$, violant ainsi le
théorème~\ref{thm:progres}. Posons:

\[
\begin{cases}
  e = * \eTaint{x} \\
  Γ = x : \tInt \\
  m = ([[x↦0]], [~]) \\
\end{cases}
\]

Les hypothèses du théorème de progrès sont bien vérifiées, mais cependant la
conclusion n'est pas vraie:

\begin{itemize}
  \item
    On a bien $\mcomp{Γ}{m}$. En effet:

    \begin{mathpar}
        \irule{M-Push}
            {
                \irule{M-Empty}
                    { }
                    { \mcomp{[~]}{([~], [~])} }
            \\
            [~] ⊧ 0 : \tInt
            \\
            \tComp{\tInt}{\tInt}
            }
            { \mcomp{x:\tInt}{([[x↦0]], [~])} }
    \end{mathpar}

  \item
    $e$ est bien typée sous $Γ$:

      \begin{mathpar}
        \irule{Lv-Deref}
          {
            \irule{Addr-User-Ignore}
              {
                \irule{Lv-Deref}
                  {
                    \irule{Lv-Var}
                      { x:\tInt ∈ Γ }
                      { Γ ⊢ x : \tInt }
                  }
                  { Γ ⊢ \& x : \tInt*}
              }
              { Γ ⊢ \eTaint{x} : \tInt*}
          }
          { Γ ⊢ * \eTaint {x} : \tInt}
      \end{mathpar}

  \item
    L'évaluation de $e$ sous $m$ provoque une erreur différente de
    $\serr{div}$, $\serr{array}$, ou $\serr{ptr}$:

      \begin{mathpar}
      \inferrule*
        {
          \irule{Exp-Lv}
            {
              \inferrule*
                { }
                {m[\vTainted{x}] = \serr{sec}}
            }
            { \mm{m}{* \eTaint {x}}{m}{\serr{sec}} }
          \\
          \irule{Eval-Err}
            { }
            {
              \msi{m}{\serr{sec}} → \serr{sec}
            }
        }
        {\msi{m}{e} → \serr{sec}}
      \end{mathpar}

\end{itemize}

Cela montre que le typage n'apporte plus de garantie de sûreté sur l'exécution:
le système de types naïvement étendu par une règle comme
\textsc{Addr-User-Ignore} n'est pas en adéquation avec les extensions présentées
dans la section~\ref{sec:extensions-noyau}. Il faut donc raffiner les règles de
typage pour interdire ce cas.

\subsection{Extensions du système de types}
\label{sec:extension-types}

On présente ici un système de types plus expressif permettant de capturer les
extensions de sémantique. \emph{In fine}, cela permettra de prouver le
théorème~\ref{thm:progres-qual} qui est l'équivalent du
théorème~\ref{thm:progres} mais pour le nouveau jugement de typage.

Définir un nouveau système de types revient à étendre le jugement de typage
$\cdot ⊢ \cdot : \cdot$, en modifiant certaines règles et en en ajoutant
d'autres. Naturellement, la plupart des différences porteront sur le traitement
des pointeurs.

\subsubsection*{Pointeurs utilisateurs}

Le changement clef est l'ajout de \emph{pointeurs utilisateur}. En plus des
types pointeurs habituels $\ptrK{t}$, on ajoute des types pointeurs utilisateur
$\ptrU{t}$. La différence entre les deux représente \emph{qui} contrôle leur
valeur (section~\ref{sec:secu-syscalls}).

Les différences sont les suivantes (figure~\ref{fig:qualif-changes-typ}):

\begin{figure}[b]%{{{

\begin{align*}
\gramdef{Type}{t}
  { … }{}
  { \ptrU{t} }{Pointeur utilisateur}
  {END}
\\
\\
\gramdef[2.5cm]{Type\\de valeur}{τ}
  { … }{}
  { \ptrU{τ} }{Pointeur utilisateur}
  {END}
\end{align*}

\caption{Ajouts liés aux pointeurs utilisateur (par rapport aux
figures~\ref{fig:les-types} et~\ref{fig:types-semantiques})}

\label{fig:qualif-changes-typ}
\end{figure}%}}}

\begin{itemize}

\item Les types \enquote{$\ptrK{t}$} s'appliquent aux pointeurs contrôlés par le
noyau. Par exemple, prendre l'adresse d'un objet de la pile noyau donne un
pointeur noyau.

\item Les types \enquote{$\ptrU{t}$}, quant à eux, s'appliquent aux pointeurs
qui proviennent de l'espace utilisateur. Ces pointeurs proviennent toujours
d'interfaces particulières, comme les appels système ou les paramètres
passés aux implantations de la fonction \texttt{ioctl}.

\end{itemize}

L'ensemble des notations est résumé dans le tableau suivant:

\begin{center}
\begin{tabular}{ccc}
\toprule
        & Noyau        & Utilisateur \\
\midrule
Syntaxe & $\&~x$       & $\eTaint{x}$ \\
Valeur  & $\wphi{(x)}$ & $\wphi{\vTainted{(x)}}$ \\
Type    & $\ptrK{t}$   & $\ptrU{t}$ \\
Accès   & $*~x$        & copy\_*\_user \\
\bottomrule
\end{tabular}
\end{center}

Puisqu'on s'intéresse à la provenance des pointeurs, détaillons les règles qui
créent, manipulent et utilisent des pointeurs.

\subsubsection*{Sources de pointeurs}

La source principale de pointeurs est l'opérateur $\&$ qui prend l'adresse d'une
variable. Celle-ci est bien entendue contrôlée par le noyau (dans le sens où son
déréférencement est toujours sûr). Cette construction crée donc des pointeurs
noyau, et on maintient la règle suivante:

\begin{mathpar}
  \disprule{Addr}
\end{mathpar}

\subsubsection*{Manipulations de pointeurs}

L'avantage du typage est que celui-ci suit le flot de données: si à un endroit
une valeur de type $t$ est affectée à une variable, que le contenu de cette
variable est placé puis retiré d'une structure de données, il conserve ce type
$t$. En particulier un pointeur utilisateur reste un pointeur utilisateur.

Une seule règle consomme un pointeur et en retourne un. Elle concerne
l'arithmétique des pointeurs. On ne l'étend pas aux pointeurs utilisateur, car
pour effectuer de l'arithmérique, il faut observer la forme du pointeur
sous-jacent. Si on veut laisser $\phx\vTainted$ opaque, il faut donc interdire
l'arithmétique sur les pointeurs utilisateur.

\subsubsection*{Utilisations de pointeurs}

La principale restriction est que seuls les pointeurs noyau peuvent être
déréférencés de manière sûre. La règle capitale est donc la suivante (déjà
introduite dans le chapitre~\ref{cha:typbase}):

\begin{mathpar}
  \disprule{Lv-Deref}
\end{mathpar}

Ainsi, on interdit le déréférencement des expressions de type $\ptrU{t}$ à la
compilation.

L'opérateur $\eTaint{\cdot}$ transforme un pointeur selon la règle suivante:

\[
    \disprule{Addr-User}
\]

Les \enquote{fonctions} $\cfuF$ et $\ctuF$ sont typées de la manière suivante.
Il est à remarquer que ce ne sont pas vraiment des fonctions et qu'elles n'ont
pas un type en $(t_1, t_2) → t$, car il faudrait un type polymorphe pour pouvoir
les appliquer à n'importe quel type de pointeurs. Leur typage est donc plus
proche de celui d'un opérateur.

\begin{mathpar}
    \disprule{GetU}

    \disprule{PutU}
\end{mathpar}

\subsection{Sûreté du typage}

\subsubsection*{Typage sémantique}

La définition du typage sémantique doit aussi être étendue au cas $φ =
\vTainted{φ'}$. En essence, \textsc{S-UserPtr} énonce que traverser un
constructeur $\phx\vTainted$ transforme un pointeur en pointeur utilisateur.

\begin{mathpar}
  \irule{S-UserPtr}
      { \semtyp{m}{\wphi{          φ }}{\ptrK{τ}} }
      { \semtyp{m}{\wphi{\vTainted{φ}}}{\ptrU{τ}} }

  \irule{Comp-Ptr}
    { \tComp{τ}{t} }
    { \tComp{\ptrU{τ}}{\ptrU{t}} }
\end{mathpar}

\subsubsection*{Propriétés du typage}

\begin{lemma}[Inversion du typage]
\label{lemma:inv-qualif}

En plus des cas présentés dans le lemme~\ref{lemma:inversion}, les cas suivants
permettent de remonter un jugement de typage.

\begin{itemize}
\item
    Si $Γ ⊢ \eTaint{e} : t$, alors il existe $t'$ tel que
    $t = t'~@$ et $Γ ⊢ e : t'$.
\item
    Si $Γ ⊢ \uGet{e_d}{e_s} : t$, alors
    $t = \tInt$ et il existe $t'$ tel que $Γ ⊢ e_d : t~*$ et
    $Γ ⊢ e_s : t~@$.
\item
    Si $Γ ⊢ \uPut{e_d}{e_s} : t$, alors
    $t = \tInt$ et il existe $t'$ tel que $Γ ⊢ e_d : t~@$ et
    $Γ ⊢ e_s : t~*$.
\end{itemize}
\end{lemma}

\begin{proof}[Démonstration]
Pour chaque forme syntaxique, on liste les règles qui ont comme conclusion
un jugement de typage portant sur celle-ci. Comme aucune autre règle ne
convient, on peut en déduire que c'est l'une de celles-ci qui a été
appliquée, et donc qu'une des prémisses est vraie.
\end{proof}

\subsubsection*{Progrès et préservation}

La propriété que nous cherchons à prouver est que le déréférencement d'un
pointeur dont la valeur est contrôlée par l'utilisateur ne peut se faire qu'à
travers une fonction qui vérifie la sûreté de celui-ci.

En fait il s'agit des théorèmes de sûreté du chapitre précédent.

\begin{theorem}[Progrès pour les extensions noyau]
\label{thm:progres-qual}

Le théorème~\ref{thm:progres} reste valable avec les extensions de ce chapitre.

\end{theorem}

La preuve de ce théorème est en annexe~\ref{proof:progres-qualif}.


\begin{theorem}[Préservation pour les extensions noyau]

\label{thm:preservation-qualif}

Le théorème~\ref{thm:preservation} reste valable avec les extensions de ce
chapitre.

\end{theorem}

La preuve de ce théorème est en annexe~\ref{proof:preservation-qualif}.

Ces extensions ne modifient pas les théorème de progrès et préservation sur les
phrases (théorèmes~\ref{thm:prog-phr} et~\ref{thm:presa-phr}).

La sûreté du typage étant à nouveau établie, on a montré que l'ajout de types
pointeurs utilisateur suffit pour avoir une adéquation entre les extensions de
sémantique de la section~\ref{sec:extensions-noyau} et les extensions du système
de type de la section~\ref{sec:extension-types}.

\subsection*{Conclusion}

En partant de \langname tel que décrit dans les chapitres~\ref{cha:lang}
et~\ref{cha:typbase}, on décrit une extension de sa syntaxe et de sa sémantique.
Cela permet d'exprimer les pointeurs vers
l'espace utilisateur, qui sont utilisés pour l'implantation d'appels système
(chapitre~\ref{cha:os}).

Une première idée pour le typage de ces nouveaux pointeurs est de leur donner le
même type que les pointeurs classiques. On a montré ensuite que ce typage naïf
ne suffit pas: il permet en effet de faire fuiter de l'information, ce qu'on
note par un cas d'erreur $\serr{sec}$. En termes de systèmes de types, cela
signifie que le théorème de progrès (théorème~\ref{thm:progres},
page~\pageref{thm:progres}) n'est plus vérifié.

Le langage des types est donc enrichi pour séparer les pointeurs utilisateur des
pointeurs noyau: les premiers sont explicitement construits par un ensemble de
sources bien déterminé, et les autres sont créés par exemple en prenant
l'adresse d'une variable. La règle de typage \textsc{Lv-Deref} assure que seuls
les pointeurs noyau peuvent être déréférencés. Pour accéder aux pointeurs
utilisateur, il faut appeler les constructions $\ctuF$ et $\cfuF$, qui sont
typées adéquatement et vérifient dynamiquement que les pointeurs dont la valeur
est contrôlée par l'utilisateur pointent vers l'espace utilisateur.

% vim: spelllang=fr


\partccl{part:lang}
\label{cha:ccl2}
L'idée derrière CQual se rapprochant de la nôtre, et cet outil ayant été
appliqué aux types de vulnérabilités qui nous intéressent~\cite{cquk-usenix04},
il est utile de détailler pourquoi notre approche est différente de ces travaux.

Tout d'abord, CQual modifie fondamentalement l'ensemble du système de types (via
le jugement de typage) et insère des qualificateurs à tous les niveaux de types.
Au contraire, nous ne modifions le système de types que là où cela est
nécessaire, c'est-à-dire sur les types pointeurs.

D'autre part, le système de types de CQual fait intervenir de manière
fondamentale une relation de sous-typage. Le problème de déréférencement des
pointeurs utilisateurs peut être traité dans ce cadre en posant $\qKernel
\preceq \qUser$ pour restreindre certaines opérations aux pointeurs $\qKernel$.

Notre approche, au contraire, n'utilise pas de sous-typage, mais consiste à
définir un type abstrait $t~\qKernel~*$ partageant certaines propriétés avec
$t~\qUser~*$ (comme la taille et la représentation) mais incompatible avec
certaines opérations. C'est à rapprocher du système de masquage par interface
des langages ML, ou encore de la définition de types abstraits par le mot-clef
\texttt{newtype} en Haskell.

En C, il est notamment commun d'utiliser des \texttt{int} pour tout et n'importe
quoi : pour des entiers bien sûr (au sens de $ℤ$), mais aussi comme
identificateurs pour lesquels les opérations usuelles comme l'addition n'ont pas
de sens. Par exemple, sous Linux, l'opération d'ouverture de fichier renvoie un
entier, dit \emph{descripteur de fichier}, qui identifie ce fichier pour ce
processus. Il est passé à toutes les fonctions de lecture et d'écriture. Un
autre idiome est d'utiliser des entiers comme des ensembles de bits (on parle de
\emph{bitmask}). Par exemple, en ouvrant un fichier on précise le mode du
fichier (lecture, écriture ou les deux) par les bits 1 et 2, s'il faut créer le
fichier ou non s'il n'existe pas par le bit 7, s'il dans ce cas il doit être
effacé par le bit 8), etc. On obtient un mode en réalisant un ``ou'' bit à bit
entre des constantes. Ces deux utilisations du type \texttt{int} n'ont rien à
voir ; il faudrait donc empêcher d'utiliser un descripteur de fichier comme un
mode, et vice-versa. De même, aucun opérateur n'a de sens sur les descripteurs
de fichier, mais l'opérateur \texttt{\textbar{}} du ``ou'' bit à bit doit rester
possible pour les modes.

Dans le cas des pointeurs, même si le noyau Linux (et la plupart des systèmes
d'exploitation) ne comportent que deux espaces d'adressage, il est commun dans
les systèmes embarqués de manipuler des pointeurs provenant d'espaces mémoire
indépendants : par exemple, de la de la mémoire flash, de la RAM, ou une EEPROM
de configuration. Ces différentes mémoires possèdent des adresses, et un
pointeur est interprété comme faisant référence à une ou l'autre selon le code
dont il est tiré. Dans ces cas, la notion de sous-typage n'est pas adaptée.


\long\def\partintro{
Après avoir décrit notre solution dans la partie~\ref{part:lang}, on présente
ici son implantation.

Le chapitre~\ref{cha:implem} décrit l'implantation en elle-même: un prototype
d'analyseur de types, distribué avec le langage \newspeak sur
\url{http://penjili.org/}. Il s'agit d'un logiciel libre, distribué sous la
license LGPL. La compilation depuis C est réalisée par l'utilitaire
\ctonewspeak. Celui-ci, comme le langage \newspeak, est antérieur à ce projet,
mais le support de plusieurs extensions GNU C a été développé spécialement pour
pouvoir analyser le code du noyau Linux.

L'analyse en elle-même est implantée de la manière classique avec une variation
de l'algorithme W de Damas et Milner. Pour des raisons de simplicité et
d'efficacité, l'unification est faite en utilisant le partage de références
plutôt que des substitutions explicites. L'algorithme d'inférence ne pose pas de
problèmes de performance.

Ensuite, dans le chapitre~\ref{cha:etudedecas}, on applique cette implantation
au noyau Linux. On commence par décrire l'implantation des appels système sous
ce noyau, et les problèmes de sécurité qui peuvent arriver (ceci a déjà été
abordé dans le chapitre~\ref{cha:os}). Dans une deuxième partie, on décrit le
cas d'un \emph{bug} de pilote graphique dans le noyau Linux. On montre que les
analyses précédentes permettent de distinguer statiquement entre le cas
incorrect et le cas corrigé.
}

\part{Expérimentation}
\label{part:xp}

\chapter{Implantation}

\label{cha:implem}
Dans ce chapitre, nous décrivons la mise en œuvre des analyses statiques
précédentes. Nous commençons par un tour d'horizon des représentations
intermédiaires possibles, avant de décrire celle retenue : Newspeak. La chaîne
de compilation est explicitée, partant de C pour aller au langage impératif
décrit dans le chapitre~\ref{cha:lang}. Enfin, nous donnons les détails d'un
algorithme d'inférence de types à la Hindley-Milner, reposant sur l'unification
et le partage de références.

%\begin{tabular}{| l || l | l | l | l | l | l |}
\hline
Nom           & Langage hôte    & Description       & But                 & Langage source                   &  Sémantique & Types \\ \hline \hline
C             & Texte           & \cite{AnsiC}      & Compilation         & C                                &  Non        & Oui   \\ \hline
GIMPLE        & C               & \cite{gcc-gimple} & Compilation         & C + extensions, C++, Ada, \ldots &  Non        & Non   \\ \hline
LLVM          & C++             & \cite{llvm-pres}  & Compilation         & C + extensions, C++, Ada, \ldots &             & Oui   \\ \hline
Cminusminus   & Texte           & \cite{spjcmm}     & Compilation         &                                  &             & Oui   \\ \hline
Cmm (Haskell) & Texte / Haskell &                   & Compilation         &                                  &             &       \\ \hline
Cmm (OCaml)   & OCaml           &                   & Compilation         &                                  &             &       \\ \hline
CIL           & OCaml           & \cite{NeculaCil}  & Analyse             &                                  &             &       \\ \hline
Clight        & Coq             & \cite{cfront}     & Compilation/Analyse &                                  &  Oui        & Oui   \\ \hline
Cminor        & Coq             & \cite{cminorSL}   & Compilation/Analyse &                                  &  Oui        & Oui   \\ \hline
Newspeak      & Ocaml           & \cite{newspeak}   & Analyse             &                                  &  Oui        & Oui   \\ \hline
\end{tabular}

% vim: textwidth=0


\section{Newspeak}
\label{sec:npk}

Newspeak~\cite{newspeak} est un langage conçu pour être à la fois :

\begin{itemize}
  \item \textbf{Précis} : sa sémantique est définie formellement
  \cite{newspeak}
  \item \textbf{Expressif} : la plupart des primitives présentes dans les
    langages de bas niveau sont compilables en Newspeak \footnote{Penjili étant
    spécialisé dans l'aéronautique, les constructions complexes comme
    \texttt{longjmp} / \texttt{setjmp} ou les exceptions ne sont pas nécessaires.}
    % TODO adapter la footnote
  \item \textbf{Simple} : peu de primitives sont présentes.
  \item \textbf{Minimal} : aucun élément syntaxique ne peut être exprimé
    comme combinaison d'autres.
  \item \textbf{Explicite} : les constructions sont toutes indépendantes du
    contexte.
  \item \textbf{Orienté analyse} : les primitives sont décorées d'informations
    reflétant leur validité (tests de bornes, etc)
  \item \textbf{Indépendant de l'architecture} : toutes les caractéristiques comme
    la taille des types ou l'alignement des structures sont rendues explicites.
\end{itemize}

\section{Chaîne de compilation}

\begin{figure}
  \centering
  \begin{tikzpicture}\shorthandoff{!}
\tikzstyle{file}=[draw, shape=rectangle, node distance=2.2cm, minimum
height=1cm, shade, top color=white,
    bottom color=blue!50!black!20, draw=blue!40!black!60, very thick];

\node [file] (c1) {\textcolor{black}{.c}};
\node [file, below of=c1] (c2) {\textcolor{black}{.c}};
\node [node distance=2.2cm, below of=c2] (c3) {};

\node [file, minimum height=0, node distance=5mm, above of=c3,draw] (c3b) {.adb};
\node [file, minimum height=0, node distance=5mm, below of=c3,draw] (c3s) {.ads};

\path (c3b.north west) ++(-3mm,3mm) [draw,dotted] rectangle ($(c3s.south east)+(3mm,-3mm)$);

\node [below of=c3, node distance=2.2cm] (c4){};

\node [file, right of=c1] (cc1) {\textcolor{black}{.c}};
\node [file, right of=c2] (cc2) {\textcolor{black}{.c}};
\node [node distance=2.2cm, right of=c3] (cc3) {};

\node [below of=cc3, node distance=2.2cm](cc4){};

\draw[->] (c1) -- node[above] {{\tiny \ttfamily cpp}} (cc1);
\draw[->] (c2) -- (cc2);

\node [file, right of=cc1] (no1) {\textcolor{black}{.no}};
\node [file, right of=cc2] (no2) {\textcolor{black}{.no}};
\node [file, right of=cc3] (no3) {\textcolor{black}{.no}};

\node [below of=no3, node distance=2.2cm](no4){};

\draw[->] (cc1) -- node[above] {{\tiny \ttfamily c2newspeak -c}} (no1);
\draw[->] (cc2) --  (no2);

\draw[->] ($ (c3b.north east)!0.5!(c3s.south east) + (3mm,0) $) -- node[above] {\tiny \ttfamily ada2newspeak -c} (no3);


\node [file, right of=no2] (npk) {\textcolor{black}{.npk}};

\node [right of=no3, node distance=2.2cm](npk2){};
\node [right of=no4, node distance=2.2cm](npk3){};

\node[draw, ellipse, right of=no2, minimum height=3cm]{};

\draw[->] (no1) -- node {} (npk);
\draw[->] (no2) -- node[above] {{\tiny \ttfamily c2newspeak}} (npk);
\draw[->] (no3) -- node {} (npk);


\node [file, right of=npk, node distance=3cm, yshift=1cm](warn)
    {\textcolor{black}{\parbox{2cm}{\centering Programme \linebreak typé}}};
\node [right of=npk, node distance=3cm, yshift=-1cm](warnn)
    {\textcolor{black}{Erreurs}};

\node [right of=npk3, node distance=3cm](analyser2){};

\node [right of=npk, node distance=3cm] {\small ou};

\draw[->] (npk) -- node[above, yshift=2mm] {{\tiny \ttfamily ptrtype}} (warn);

\draw[->] (npk) -- (warnn);

\draw[->] (c4)   -- node[above, text depth=3pt] {\footnotesize Pr\'etraitement} (cc4);
\draw[->] (cc4)  -- node[above, text depth=3pt] {\footnotesize Compilation} (no4);
\draw[->] (no4)  -- node[above, text depth=3pt] {\footnotesize \'Edition de liens} (npk3);
\draw[->] (npk3) -- node[above, text depth=3pt] {\footnotesize Analyse} (analyser2);
\end{tikzpicture}

  \caption{Compilation depuis Newspeak}
  \label{fig:compil-npk}
\end{figure}

%TODO Mettre à jour la figure

La compilation vers C est faite en trois étapes (figure~\ref{fig:compil-npk}) :
prétraitement du code source, compilation de C prétraité vers \newspeak{}, puis
compilation de \newspeak{} vers ce langage.

\section*{Trad}

La première étape consiste à prétraiter les fichiers C source avec le logiciel
\texttt{cpp}, comme pour une compilation normale. Cette étape interprète les
directives comme \texttt{\#include}, \texttt{\#ifdef}, À cet étape, les
commentaires sont aussi supprimés.

% TODO comment interpréter les /*!npk xxx */ ?

Une fois cette passe effectuée, nous avons des fichiers C prétraités,
c'est à dire des unités de compilation autocontenues. % TODO bof

Puisque la directive \texttt{\#include} est textuelle, ces fichiers sont très
grands et donnent lieu à beaucoup de duplication dans les passes suivantes.

À ce niveau, les fichiers sont passés à l'outil \texttt{c2newspeak} qui les
traduit vers Newspeak. Dans cette étape, les types et les noms sont résolus, et
le programme est annoté de manière à rendre les prochaines étapes indépendante
du contexte. Par exemple, chaque déclaration de variable est adjointe d'une
description complète du type.

Lors de cette étape, le flôt de contrôle est également simplifié. C en effet
propose de nombreuses constructions ambigües ou redondantes.

Au contraire, Newspeak propose un nombre réduit de constructions. Rappelons que
le but de ce langage est de faciliter l'analyse statique : des constructions
orthogonales permettent donc d'éviter la duplication de règles sémantique, ou de
code lors de l'implémentation d'un analyseur.

Par exemple, plutôt que de fournir une boucle \emph{while}, une boucle
\emph{do/while} et une boucle for, Newspeak fournit une unique boucle
\npkWhile{}. La sortie de boucle est compilée vers un \npkGoto{}\cite{goto}, qui est
toujours un saut vers l'avant (similaire à un "break" généralisé).

La sémantique de Newspeak et la traduction de C vers Newspeak sont décrites dans
\cite{newspeak}.

Newspeak est conçu pour l'analyse statique par interprétation abstraite. Il a
donc une vue de bas niveau sur les programmes. Par exemple, aucune distinction
n'est faite entre l'accès à un champ et l'accès un l'élément d'un tableau (tous
deux sont traduits par un décalage numérique depuis le début de la zone
mémoire). Pour supprimer cette ambiguïté, il faut s'interfacer dans les
structures internes de \texttt{c2newspeak}, où les informations nécessaires sont
encore présentes.

Ensuite, les différents fichiers sont liés ensemble. Cet étape consiste
principalement à s'assurer que les hypothèses faites par les différentes unités
de compilation sont cohérentes entre elles. Les objets marqués \texttt{static},
invisibles à l'extérieur de leur unité de compilation, sont renommés afin qu'ils
aient un nom unique.

Enfin, l'implantation d'un algorithme d'inférence pour les systèmes de types
décrits dans les chapitres~\ref{cha:typbase} et~\ref{cha:qualifs} est assez
simple. Puisqu'ils sont suffisament proches du lambda calcul simplement typé, on
peut utiliser une variante de l'algorithme W de Damas et
Milner~\cite{DamasMilner}. On utilise l'optimisation classique qui consiste à se
reposer sur le partage de références pour réaliser l'unification, plutôt que de
faire des substitutions explicites. Puisque ces systèmes de types sont
monomorphes, on présente une erreur si des variable de type libres sont
présentes.

À la fin de cette étape, on obtient soit un programme complètement annoté, soit
une erreur de type.

\subsection{Prétraitement}

\ctonewspeak{} travaillant uniquement sur du code prétraité (dans directives de
préprocesseur), la première étape consiste donc à faire passer le code par \cpp:
les macros sont développées, les constantes remplacées par leurs valeurs, les
commentaires supprimés, les fichiers d'en-tête inclus, etc.

\subsection{Compilation (levée des ambigüités)}

Cette passe est réalisée par l'utilitaire \ctonewspeak{}. L'essentiel de la
compilation consiste à mettre à plat les définition de types, et \subsection{Annotations}

\newspeak{} a de nombreux avantages, mais pour une analyse par typage il est
trop bas niveau. Par exemple, dans le code suivant

\begin{Verbatim}[commandchars=\\\{\}]
\PY{k}{struct} \PY{n}{s} \PY{p}{\PYZob{}}
    \PY{k+kt}{int} \PY{n}{a}\PY{p}{;}
    \PY{k+kt}{int} \PY{n}{b}\PY{p}{;}
\PY{p}{\PYZcb{}}\PY{p}{;}

\PY{k+kt}{int} \PY{n+nf}{main}\PY{p}{(}\PY{k+kt}{void}\PY{p}{)}
\PY{p}{\PYZob{}}
    \PY{k}{struct} \PY{n}{s} \PY{n}{x}\PY{p}{;}
    \PY{k+kt}{int} \PY{n}{y}\PY{p}{[}\PY{l+m+mi}{10}\PY{p}{]}\PY{p}{;}
    \PY{n}{x}\PY{p}{.}\PY{n}{b} \PY{o}{=} \PY{l+m+mi}{1}\PY{p}{;}
    \PY{n}{y}\PY{p}{[}\PY{l+m+mi}{1}\PY{p}{]} \PY{o}{=} \PY{l+m+mi}{1}\PY{p}{;}
    \PY{k}{return} \PY{l+m+mi}{0}\PY{p}{;}
\PY{p}{\PYZcb{}}
\end{Verbatim}


\wip{}

\subsection{Implantation de l'algorithme de typage}

On prend l'exemple d'un lambda-calcul simplement typé avec entiers, flottants et
couples (figure~\ref{fig:stlc}).

\begin{figure}

\gramlr{Termes}{
\begin{align*}
  t \gramisa & x              & \textrm{Variable} \\
    \gramor  & λ x . t        & \textrm{Abstraction} \\
    \gramor  & t              & \textrm{Application} \\
    \gramor  & n              & \textrm{Entier} \\
    \gramor  & f              & \textrm{Flottant} \\
    \gramor  & (t, t)         & \textrm{Couple} \\
    \gramor  & \textrm{fst}~t & \textrm{Projection gauche} \\
    \gramor  & \textrm{snd}~t & \textrm{Projection droite}
\end{align*}
}

\gramlr{Types}{
\begin{align*}
  τ \gramisa &  \tInt          & \textrm{Entier}\\
    \gramor  &  \tFloat        & \textrm{Flottant}\\
    \gramor  & τ \rightarrow τ & \textrm{Fonction} \\
    \gramor  & τ \times τ      & \textrm{Produit}
\end{align*}
}

\gramlr{Contextes}{
\begin{align*}
  Γ \gramisa & ε     & \textrm{Contexte vide}\\
    \gramor  & Γ,x:τ & \textrm{Extension}
\end{align*}
}

\gramlr{Règles}{
\begin{mathpar}
\irule{Int}{ }{Γ⊢n:\tInt}
\and
\irule{Float}{ }{Γ⊢f:\tFloat}
\and
\irule{Var}{x:τ∈Γ}{Γ⊢x:τ}
\and
\irule{App}{Γ⊢f:τ_1 \rightarrow τ_2 \\ Γ ⊢ x : τ_1}{Γ ⊢ f x : τ_2}
\and
\irule{Abs}{Γ,x:τ_1⊢ y : τ_2}{Γ⊢ λx.y:τ_1 \rightarrow τ_2}
\and
\irule{Proj-g}{Γ⊢x:τ_1\times τ_2}{Γ⊢\textrm{fst}~x:τ_1}
\and
\irule{Proj-d}{Γ⊢x:τ_1\times τ_2}{Γ⊢\textrm{snd}~x:τ_2}
\and
\irule{Tup}{Γ⊢x:τ_1 \\ Γ⊢y:τ_2}{Γ⊢(x,y):τ_1 \times τ_2}
\end{mathpar}
}


\caption{Lambda calcul simplement typé avec entiers, flottants et couples}
\label{fig:stlc}

\end{figure}

Prenons l'exemple de la fonction suivante\footnote{ On suppose que \texttt{plus}
est une fonction de l'environnement global qui a pour type $\tInt \rightarrow
\tInt \rightarrow \tInt$.} :

\[
f = λx.λy. \textrm{plus} (\textrm{plus} (\textrm{fst} x) (\textrm{snd} x)) y
\]

Soit en syntaxe ML :

\begin{Verbatim}
let f x y =
  plus
    (plus
      (fst x)
      (snd x)
    )
    y
\end{Verbatim}

Puisque \texttt{fst} et \texttt{snd} sont appliqués à \texttt{x}, ce doit être
un tuple. En outre on additionne ces deux composantes ensemble, donc elles
doivent être de type \tInt (et le résultat aussi). Par le même argument,
\texttt{y} doit aussi être de type \tInt. En conclusion, \texttt{x} est de type
$\tInt × \tInt$ et \texttt{y} de type $\tInt$, donc f est de type $\tInt × \tInt
→ \tInt → \tInt$.

Pour implanter cette analyse, on peut remarquer qu'étant donné la forme d'un
terme, on peut savoir quelle règle de typage a été utilisée en dernier. Il est
ainsi possible de ``remonter'' l'arbre d'inférence afin de savoir quelles règles
ont été employées (figure~\ref{fig:inftree-rules})\footnote{Par souci de clarté,
les prémisses des applications de \textsc{(Var)} ne sont pas notées.}. Pour le
moment, on ne connait pas les types.

% TODO quel prop cela utilise t-il?

\begin{figure} % {{{ Fig règles
\def\disptypeR#1#2{:#1}

\def\gammaXY{\ensuremath{Γ_2}}
\def\gammaX{\ensuremath{Γ_1}}
\def\gammaOneDef{\ensuremath{Γ^0, x\disptypeL{\tInt \times \tInt}, y\disptypeL{\tInt}}}
\def\gammaTwoDef{\ensuremath{Γ^0, x\disptypeL{\tInt \times \tInt}}}
\def\termFst#1{\ensuremath{\textrm{fst}~x}}
\def\termSnd#1{\ensuremath{\textrm{snd}~x}}
\def\fstx{\termFst{x}}
\def\sndx{\termSnd{x}}
\begin{mathpar}
\irule
  {Abs}
  {
    \irule
      {Abs}
      {
        \inferrule*[right=(App)]
          {
            \inferrule*[right=(Var), rightskip=4em]
              { }
              {\gammaXY ⊢ y
              \disptypeR{\tInt}
              }
            \\
            \inferrule*[right=(App),vdots=4em,leftskip=5em,rightskip=5em]
              {
                \inferrule*[right=(Var), leftskip=3em, rightskip=7em]
                  { }
                  {\gammaXY ⊢ \textrm{plus}
                  \disptypeR{\tInt \rightarrow \tInt \rightarrow \tInt}
                  }
                \\
                \inferrule*[right=(App),vdots=3em,leftskip=7em,rightskip=3em]
                  {
                    \inferrule*[right=(App), vdots=4em, rightskip=5em]
                      {
                        \inferrule*[right=(Var), rightskip=5em]
                          { }
                          {\gammaXY ⊢ \textrm{plus}
                          \disptypeR{\tInt \rightarrow \tInt \rightarrow \tInt}
                          }
                        \\
                        \inferrule*[right=(Proj-G), vdots=3em, leftskip=5em]
                          {
                            \irule
                              {Var}
                              { }
                              {\gammaXY ⊢ x
                              \disptypeR{\tInt \times \tInt}
                              }
                          }
                          {\gammaXY ⊢ \fstx
                          \disptypeR{\tInt}
                          }
                      }
                      {\gammaXY ⊢ \textrm{plus} (\fstx)
                      \disptypeR{\tInt \rightarrow \tInt}
                      }
                    \\
                    \inferrule*[right=(Proj-D), leftskip=5em, rightskip=4em]
                      {
                        \irule
                          {Var}
                          { }
                          {\gammaXY ⊢ x
                          \disptypeR{\tInt \times \tInt}
                          }
                      }
                      {\gammaXY ⊢ \sndx
                      \disptypeR{\tInt}
                      }
                  }
                  {\gammaXY ⊢ \textrm{plus} (\fstx) (\sndx)
                  \disptypeR{\tInt}
                  }
              }
              {\gammaXY ⊢ \textrm{plus} (\textrm{plus} (\fstx) (\sndx))
              \disptypeR{\tInt \rightarrow \tInt}
              }
          }
          {\gammaXY ⊢ \textrm{plus} (\textrm{plus} (\fstx) (\sndx)) y
          \disptypeR{\tInt}
          }
      }
      {\gammaX ⊢ λy. \textrm{plus} (\textrm{plus} (\fstx) (\sndx)) y
      \disptypeR{\tInt \rightarrow \tInt}
      }
  }
  {Γ_0 ⊢ λx.λy. \textrm{plus} (\textrm{plus} (\fstx) (\sndx)) y
  \disptypeR{\tInt \times \tInt \rightarrow \tInt \rightarrow \tInt}
  }
\end{mathpar}

\[\gammaX = \gammaTwoDef \qquad \gammaXY = \gammaOneDef\]


\caption{Arbre d'inférence : règles à utiliser}
\label{fig:inftree-rules}
\end{figure} % }}}

Une fois à cette étape, on peut donner un nom à chaque type inconnu : $t_1, t_2,
\ldots$. L'utilisation qui en est faite permet de générer un ensemble de
contraintes d'unification (figure~\ref{fig:stlc-const}).

Le symbole $?$ qui apparait dans les règles correspond à ``n'importe quel
type'', dans le sens où l'unification avec ce symbole est toujours possible.

\begin{figure} % {{{
  \def\stlcConst#1#2{
  \fbox{
    \begin{tabular}{m{3.5cm}||c}
      \raggedleft \fbox{\raggedleft \ensuremath{#1}} & \fbox{\ensuremath{#2}}
    \end{tabular}
  }
  }
  \centering
  \begin{tabular}{m{4cm}||lm{4cm}||l}
    \raggedleft
    $
    \irule{Var}
      { (x, t_x) ∈ Γ }
      { Γ ⊢ x : t_1 }
    $
    &
    $
      t_1 = t_x
    $
    &
    \raggedleft
    $
    \irule{Abs}
      { Γ, (x:t_x) ⊢     e : t_2 }
      { Γ          ⊢ λx. e : t_1 }
    $
    &
    $
      t_1 = t_x → t_2
    $
    \\
    \multicolumn{4}{c}{} \\
    \raggedleft
    $
    \irule{App}
      { Γ ⊢ e_1 : t_1
     \\ Γ ⊢ e_2 : t_2
      }
      { Γ ⊢ e_1~e_2 : t_3 }
    $
    &
    $
      t_1 = t_2 → t_3
    $
    &
    \raggedleft
    $
    \irule{Proj-G}
      { Γ ⊢ e : t_2 }
      { Γ ⊢ \termFst{e} : t_1 }
    $
    &
    $
      t_2 = t_1 × ?
    $
    \\
    \multicolumn{4}{c}{} \\
    \raggedleft
    $
    \irule{Proj-D}
      { Γ ⊢ e : t_2 }
      { Γ ⊢ \termSnd{e} : t_1 }
    $
    &
    $
      t_2 = ? × t_1
    $
    &
    \raggedleft
    $
    \irule{Tup}
       { Γ⊢e_1:t_1
      \\ Γ⊢e_2:t_2
      }{
        Γ⊢(e_1,e_2):t_3
      }
    $
    &
    $
    t_3 = t_1 \times t_2
    $
  \end{tabular}
\caption{Contraintes crées par les applications de règles}
\label{fig:stlc-const}
\end{figure} %}}}

\begin{figure}
\centering

\subbottom[Contraintes]{
  \label{fig:inf-constr-ex}
  \begin{minipage}{0.5\linewidth}
  \begin{align}
  t_1    &= t_x → t_2             \tag{Abs1} \\
  t_2    &= t_y → t_3             \tag{Abs2} \\
  t_5    &= t_4 → t_3             \tag{App1} \\
  t_4    &= t_y                   \tag{Var1} \\
  t_6    &= t_7 → t_5             \tag{App2} \\
  t_6    &= \tInt → \tInt → \tInt \tag{Var2} \\
  t_{10} &= t_8 → t_7             \tag{App3} \\
  t_9    &= ? × t_8               \tag{Proj-D1} \\
  t_9    &= t_x                   \tag{Var3} \\
  t_{11} &= t_{12} → t_{10}       \tag{App4} \\
  t_{11} &= \tInt → \tInt → \tInt \tag{Var4} \\
  t_{13} &= t_{12} × ?            \tag{Proj-G1} \\
  t_{13} &= t_x                   \tag{Var5}
  \end{align}
  \end{minipage}
}
\subbottom[Solution]{
  \label{fig:inf-sol-ex}
  \begin{minipage}{0.3\linewidth}
  \begin{align*}
  t_{1}  &= \tInt × \tInt → \tInt \\
  t_{2}  &= \tInt → \tInt         \\
  t_{3}  &= \tInt                 \\
  t_{4}  &= \tInt                 \\
  t_{5}  &= \tInt → \tInt         \\
  t_{6}  &= \tInt → \tInt → \tInt \\
  t_{7}  &= \tInt                 \\
  t_{8}  &= \tInt                 \\
  t_{9}  &= \tInt × \tInt         \\
  t_{10} &= \tInt → \tInt         \\
  t_{11} &= \tInt → \tInt → \tInt \\
  t_{12} &= \tInt                 \\
  t_{13} &= \tInt × \tInt         \\
  t_{x}  &= \tInt × \tInt         \\
  t_{y}  &= \tInt
  \end{align*}
  \end{minipage}
}

\caption{Contraintes d'égalité et solution obtenues à partir de la
figure~\ref{fig:inftree-rules}}

\end{figure}

% TODO mal numéroté

On en déduit donc un ensemble de contraintes d'égalité
(figure~\ref{fig:inf-constr-ex}). La prochaine étape est de résoudre ce système.
D'après (App1) et (App2), $t_6 = t_7 → t_4 → t_3$ d'où par (Var2), $t_7 → t_4 →
t_3 = \tInt → \tInt → \tInt$. Le constructeur $→$ étant injectif, on a donc $t_3
= t_4 = t_7 = \tInt$. En procédant de même avec (App3) et (App4), on obtient
$t_8 = t_{12} = \tInt$. D'après (Proj-G1), (Var5), (Var3), et (Proj-D1) on a
$t_{12} × ? = t_{13} = t_{x} = t_9 = ? × t_{8}$, donc $t_x = t_{12} × t_8 =
\tInt × \tInt$. En remplaçant, on obtient tous les autres types
(figures~\ref{fig:inf-sol-ex}) et donc un arbre de typage complet
(figure~\ref{fig:inftree-full}).

\begin{figure} % {{{ Fig full tree
\def\disptypeR#1#2{:#2}

\def\gammaXY{\ensuremath{Γ_2}}
\def\gammaX{\ensuremath{Γ_1}}
\def\gammaOneDef{\ensuremath{Γ^0, x\disptypeL{\tInt \times \tInt}, y\disptypeL{\tInt}}}
\def\gammaTwoDef{\ensuremath{Γ^0, x\disptypeL{\tInt \times \tInt}}}
\def\termFst#1{\ensuremath{\textrm{fst}~x}}
\def\termSnd#1{\ensuremath{\textrm{snd}~x}}
\def\fstx{\termFst{x}}
\def\sndx{\termSnd{x}}
\begin{mathpar}
\irule
  {Abs}
  {
    \irule
      {Abs}
      {
        \inferrule*[right=(App)]
          {
            \inferrule*[right=(Var), rightskip=4em]
              { }
              {\gammaXY ⊢ y
              \disptypeR{\tInt}
              }
            \\
            \inferrule*[right=(App),vdots=4em,leftskip=5em,rightskip=5em]
              {
                \inferrule*[right=(Var), leftskip=3em, rightskip=7em]
                  { }
                  {\gammaXY ⊢ \textrm{plus}
                  \disptypeR{\tInt \rightarrow \tInt \rightarrow \tInt}
                  }
                \\
                \inferrule*[right=(App),vdots=3em,leftskip=7em,rightskip=3em]
                  {
                    \inferrule*[right=(App), vdots=4em, rightskip=5em]
                      {
                        \inferrule*[right=(Var), rightskip=5em]
                          { }
                          {\gammaXY ⊢ \textrm{plus}
                          \disptypeR{\tInt \rightarrow \tInt \rightarrow \tInt}
                          }
                        \\
                        \inferrule*[right=(Proj-G), vdots=3em, leftskip=5em]
                          {
                            \irule
                              {Var}
                              { }
                              {\gammaXY ⊢ x
                              \disptypeR{\tInt \times \tInt}
                              }
                          }
                          {\gammaXY ⊢ \fstx
                          \disptypeR{\tInt}
                          }
                      }
                      {\gammaXY ⊢ \textrm{plus} (\fstx)
                      \disptypeR{\tInt \rightarrow \tInt}
                      }
                    \\
                    \inferrule*[right=(Proj-D), leftskip=5em, rightskip=4em]
                      {
                        \irule
                          {Var}
                          { }
                          {\gammaXY ⊢ x
                          \disptypeR{\tInt \times \tInt}
                          }
                      }
                      {\gammaXY ⊢ \sndx
                      \disptypeR{\tInt}
                      }
                  }
                  {\gammaXY ⊢ \textrm{plus} (\fstx) (\sndx)
                  \disptypeR{\tInt}
                  }
              }
              {\gammaXY ⊢ \textrm{plus} (\textrm{plus} (\fstx) (\sndx))
              \disptypeR{\tInt \rightarrow \tInt}
              }
          }
          {\gammaXY ⊢ \textrm{plus} (\textrm{plus} (\fstx) (\sndx)) y
          \disptypeR{\tInt}
          }
      }
      {\gammaX ⊢ λy. \textrm{plus} (\textrm{plus} (\fstx) (\sndx)) y
      \disptypeR{\tInt \rightarrow \tInt}
      }
  }
  {Γ_0 ⊢ λx.λy. \textrm{plus} (\textrm{plus} (\fstx) (\sndx)) y
  \disptypeR{\tInt \times \tInt \rightarrow \tInt \rightarrow \tInt}
  }
\end{mathpar}

\[\gammaX = \gammaTwoDef \qquad \gammaXY = \gammaOneDef\]


\caption{Arbre d'inférence complet}
\label{fig:inftree-full}
\end{figure} % }}}


Pour effectuer une résolution de manière automatique, il faut une technique
d'unification. La manière la plus classique de procéder est due à
Robinson~\cite{robinson} et se déroule par substitution.


\clearpage

Pour résoudre ces contraintes, on commence par les simplifier : si $t_a
\rightarrow t_b = t_c \rightarrow t_d$, alors $t_a = t_c$ et $t_b = t_d$. De
même si $t_a \times t_b = t_c \times t_d$. Au contraire, si $t_a \rightarrow t_b
= t_c \times t_d$, il est impossible d'unifier les types et il faut abandonner
l'inférence de types. D'autre cas sont impossibles, par exemple $\tInt = t_1
\rightarrow t_2$ ou $\tInt = \tFloat$.

Une fois ces simplifications réalisées, les contraintes restantes sont d'une des
formes suivantes :

\begin{itemize}
\item
  $t_i = t_i$. Il n'y a rien à faire, cette contrainte peut être supprimée.
\item
  $t_i = t_j$ avec $i \ne j$ : toutes les occurrences de $t_j$ dans les autres
  contraintes peuvent être remplacées par $t_i$.
\item
  $t_i = x$ (ou $x = t_i$) où $x$ est un type concret : idem.
\end{itemize}

%TODO C'est faux

\clearpage

\begin{figure} % fig:unifpartage {{{

  \subbottom{
  \label{fig:unifpartage:a}
  \begin{tikzpicture}
  \node               (var) {Var};
  \node[right of=var] (ref) {ref};
  \node[right of=ref, node distance=1.7cm] (u0) {Unknown 0};
  \node[below of=var] (ptr) {Ptr};
  \node[left of=var]  (x) {x :};
  \node[left of=ptr] (p) {p :};
  \draw[->] (x) -- (var);
  \draw[->] (p) -- (ptr);
  \draw[->] (ptr) -- (var);
  \draw[->] (var) -- (ref);
  \draw[->] (ref) -- (u0);
  \end{tikzpicture}
  }
  \subbottom{
  \label{fig:unifpartage:b}
  \begin{tikzpicture}
  \node               (var) {Var};
  \node[right of=var] (ref) {ref};
  \node[right of=ref, node distance=1.7cm] (u0) {Instanciated};
  \node[right of=u0, node distance=1.7cm] (ii) {Int};
  \node[below of=var] (ptr) {Ptr};
  \node[left of=var]  (x) {x :};
  \node[left of=ptr] (p) {p :};
  \draw[->] (x) -- (var);
  \draw[->] (p) -- (ptr);
  \draw[->] (ptr) -- (var);
  \draw[->] (var) -- (ref);
  \draw[->] (ref) -- (u0);
  \draw[->] (u0) -- (ii);
  \end{tikzpicture}
  }
  \subbottom{
  \label{fig:unifpartage:c}
  \begin{tikzpicture}
  \node               (var) {Int};
  \node[below of=var] (ptr) {Ptr};
  \node[left of=var]  (x) {x :};
  \node[left of=ptr] (p) {p :};
  \node[right of=ptr] (pi) {Int};
  \draw[->] (x) -- (var);
  \draw[->] (p) -- (ptr);
  \draw[->] (ptr) -- (pi);
  \end{tikzpicture}
  }

  \caption{Unification par partage}
  \label{fig:unifpartage}

\end{figure} % }}}

Plutôt que de modifier toutes les occurrences d'un type $t_i$, on va affecter à
$t_i$ la valeur du nouveau type.

L'implémentation de cet algorithme utilise le partage et les références
(figure~\ref{fig:unifpartage}).

D'abord \ref{fig:unifpartage:a}, ensuite \ref{fig:unifpartage:b}, et enfin
\ref{fig:unifpartage:c}.

\begin{figure} % fig:exunif:c {{{

  \insertcode{ex-unif-c.c}

  \caption{Compilation d'un programme C -- avant}
  \label{fig:exunif:c}
\end{figure} % }}}

Prenons l'exemple de la figure~\ref{fig:exunif:c} et typons-le "à la main". On
commence par oublier toutes les étiquettes de type présentes dans le programme.
Celui-ci devient alors :

\insertcode{exunif-code.c}

La premiere ligne introduit deux variables. On peut noter leurs types respectifs
(inconnus pour le moment) $t_1$ et $t_2$. La première affectation \texttt{p =
\&x} implique que les deux côtés du signe "=" ont le même type. À gauche, le
type est $t_2$, et à droite $Ptr(t_1)$. On applique le même raisonnement à la
seconde affectation : à gauche, le type est $t_1$ et à droite Int. On en déduit
que le type de x est Int et celui de p est Ptr(Int).


\insertcode{lambda-types.ml}

Pour implanter cet algorithme, on représente les types de données du programmes
à typer par une valeur de type \texttt{ml\_type}. En plus des constantes de
types comme int ou float, et des constructeurs de type comme pair et fun, le
constructeur Var permet d'exprimer les variables de types (inconnues ou non).

Celles-ci sont numérotées par un int, on suppose avoir à disposition deux
fonctions manipulant un compteur global d'inconnues.

\insertcode{exunif-counter.ml}

De plus, on a un module gérant les environnements de typage. Il pourra être
implanté avec des listes d'association ou des tables de hachage, par exemple. Sa
signature est :

\insertcode{exunif-env.ml}

Reprenons l'exemple précédent. Partant d'un environnement vide (Env.empty), on
commence par l'étendre de deux variables. Comme on n'a aucune information, il
fait allouer des nouveaux noms d'inconnues (qui correspondent à $t_1$ et $t_2$):

\insertcode{exunif-1.ml}

La première instruction indique que les deux côtés de l'affectation doivent
avoir le même type.

\insertcode{exunif-2.ml}

\begin{figure}
  \centering
  \begin{tikzpicture}
    [scale=2.3]

    \node at (0,  0) (n0a) {t1:};
    \node at (0, -1) (n0b) {t2:};

    \node at (1,  0) (n1a) {Var\_type};
    \node at (2,  0) (n2a) {ref};
    \node at (3,  0) (n3a) {Unknown 0};

    \node at (1, -1) (n1b) {Var\_type};
    \node at (2, -1) (n2b) {ref};
    \node at (3, -1) (n3b) {Instanciated};
    \node at (4, -1) (n4b) {Ptr};

    \draw [->] (n1a) -- (n2a);
    \draw [->] (n2a) -- (n3a);

    \draw [->] (n1b) -- (n2b);
    \draw [->] (n2b) -- (n3b);
    \draw [->] (n3b) -- (n4b);

    \draw [->] (n4b) -- ++ (0, 0.5) -| (n1a);

  \end{tikzpicture}
  \caption{Unification : partage}
  \label{fig:unifsharing}
\end{figure}

Ici il se passe plusieurs choses intéréssantes. D'une part nous faisont appel à
une fonction externe typeof qui retourne le type d'une expression sous un
environnement (dans une implantation complète il s'agirait d'un appel récursif).
Dans ce cas, \texttt{typeof lhs1 env} est identique à \texttt{Env.get lhs1 env}
et \texttt{typeof rhs1 env} à \texttt{Ptr\_type t1}. L'autre aspect intéressant
est la dernière ligne : la fonction \texttt{unify} va modifier en place les
représentations des types afin de les rendre égales. L'implantation de
\texttt{unify} sera décrite plus tard. Dans ce cas précis, elle va faire pointer
la référence dans t2 vers t1 (figure~\ref{fig:unifsharing}).

Enfin, la seconde affectation se déroule à peu près de la même manière.

\insertcode{exunif-3.ml}

Ici \texttt{typeof lhs2 env} est identique à \texttt{Ptr\_type (Env.get "p"
env)} et \texttt{typeof lhs2 env} à \texttt{Const\_type Int\_type}. Et dans cas,
l'unification doit se faire entre t1 et \texttt{Const\_type Int\_type} : cela
mute la référence derrière t1 (figure~\ref{fig:typeunifref}).

\begin{figure}
  \begin{tikzpicture}
    [scale=2.3]

    \node at (0,  0) (n0a) {t1:};
    \node at (0, -1) (n0b) {t2:};

    \node at (1,  0) (n1a) {Var\_type};
    \node at (2,  0) (n2a) {ref};
    \node at (3,  0) (n3a) {Instanciated};
    \node at (4,  0) (n4a) {Const\_type};
    \node at (5,  0) (n5a) {Int\_type};

    \node at (1, -1) (n1b) {Var\_type};
    \node at (2, -1) (n2b) {ref};
    \node at (3, -1) (n3b) {Instanciated};
    \node at (4, -1) (n4b) {Ptr};

    \draw [->] (n1a) -- (n2a);
    \draw [->] (n2a) -- (n3a);
    \draw [->] (n3a) -- (n4a);
    \draw [->] (n4a) -- (n5a);

    \draw [->] (n1b) -- (n2b);
    \draw [->] (n2b) -- (n3b);
    \draw [->] (n3b) -- (n4b);

    \draw [->] (n4b) -- ++ (0, 0.5) -| (n1a);

  \end{tikzpicture}
  \caption{Unification par mutation de références}
  \label{fig:typeunifref}
\end{figure}

L'essence de l'algorithme d'inférence se situe donc dans 2 fonctions. D'une
part, \texttt{unify} qui réalise l'unification des types grâce à au partage des
références. D'autre part, la \texttt{typeof} qui encode les règles de typage
elles-mêmes et les applique à l'aide de \texttt{unify}.

\subsection{Algorithme d'unification}

Voici une implantation de la fonction \texttt{unify}.

Celle-ci prend en entrée deux types $t_1$ et $t_2$. À l'issue de l'exécution de
\texttt{unify}, ces deux types doivent pouvoir être considérés comme égaux. Si
ce n'est pas possible, une erreur sera levée.

La première étape est de réduire ces deux types, c'est à dire à transformer les
constructions \texttt{Var (ref (Instanciated t))} en \texttt{t}.

Ensuite, cela dépend des formes qu'ont les types réduits :

\begin{figure}
  \centering
  \begin{tikzpicture}
    [every node/.style={node distance=2cm}
    ,>=latex
    ]

    \node (n1) {Ptr};
    \node[right of=n1] (n2) {Var};
    \node[right of=n2] (n3) {ref};
    \node[right of=n3] (n4) {Instanciated};

    \draw [->] (n1) -- (n2);
    \draw [->] (n2) -- (n3);
    \draw [->] (n3) -- (n4);

    \draw [->] (n4) -- ++ (0, -1cm) -| (n1);

  \end{tikzpicture}
  \caption{Cycle dans le graphe de types}
  \label{fig:typecycle}
\end{figure}

\begin{itemize}

\item si les deux types sont inconnus (de la forme \texttt{Var (ref
(Instanciated t))}), on fait pointer l'une des deux références vers le premier
type. Notons que cela crée un type de la forme \texttt{Var (ref (Instanciated
(Var (ref (Unknown n)))))} qui sera réduit lors d'une prochaine étape
d'unification.

\item si un type est inconnu et pas l'autre, il faut de la même manière affecter la
référence. Mais en faisant ça inconditionnellement, cela peut poser problème :
par exemple en tentant d'unifier \texttt{a} avec \texttt{Ptr(a)} on pourrait
créer un cycle dans le graphe (figure~\ref{fig:typecycle}).
Pour éviter cette situation, il suffit de s'assurer que le type inconnu n'est
pas présent dans le type à affecter.

\item si les deux types sont des types de base (comme \tInt ou \tFloat) égaux,
on ne fait rien.

\item si les deux types sont des constructeurs de type, il faut que les
constructeurs soient égaux. On unifie en outre leurs arguments deux à deux.

\item dans les autres cas, l'algorithme échoue.

\item TODO sous typage pour les structures

\end{itemize}

TODO :

\begin{itemize}
\item implem du polymorphisme
\item implem du sous-typage
\item généralisation depuis le toy language
\end{itemize}

\clearpage

\begin{figure} % fig:exunif:tpk {{{

  \insertcode{ex-unif-tpk.ml}

  \caption{Compilation d'un programme C -- après}
  \label{fig:exunif:tpk}
\end{figure} % }}}

Le programme C (figure~\ref{fig:exunif:c}) est compilé ainsi en Tyspeak
(figure~\ref{fig:exunif:tpk}).

\clearpage

\section{Architecture de \ptrtype}
\label{sec:ptrtype-archi}

L'outil \ptrtype lit un programme Newspeak (ou un fichier C), et réalise
l'inférence de qualificateurs. En sortie, il affiche soit le programme
complètement annoté, soit une erreur. Le cœur de l'outil est dans la

Si l'argument à \ptrtype est un fichier C, il est tout d'abord compilé vers
Newspeak grâce à l'utilitaire \texttt{c2newspeak}. Ensuite, les autres passes
travaillent sur une représentation intermédiaire proche de Newspeak, mais où des
étiquettes de type supplémentaires sont ajoutées. Ce type de représentation
intermédiaire (polymorphe en le type des étiquettes) est \texttt{'a Tyspeak.t}.

\begin{figure}
\insertcode{implem-process.ml}
\caption{Implantation -- fonction principale de \ptrtype}
\label{fig:implem-process}
\end{figure}

Le reste de l'outil est résumé dans la fonction
\texttt{process\_npk} (figure~\ref{fig:implem-process}) :

\begin{itemize}

\item Grâce à la fonction \texttt{convert_unit : Newspeak.t -> unit Tyspeak.t},
  on ajoute des étiquettes ``vides'' (toutes égales à \texttt{() : unit}).

\item L'ensemble des fonctions du programme est trié topologiquement selon la
  relation~$\preceq$ définie par $f \preceq g \eqdef \textrm{``} g
  \textrm{ apparaît dans la définition de } f \textrm{''}$. Cela est fait en
  construisant une représentation de $\preceq$ sous forme de graphe, puis en
  faisant un parcours en largeur de celui-ci.

\item Les annotations extérieures sont alors lues (variable \texttt{exttbl}), ce
  qui permet de créer un environnement initial.

\item Les types de chaque fonction sont inférés, par le biais de la fonction
  suivante :

\begin{Verbatim}
val infer : Newspeak.fid list
         -> Types.simple Env.t
         -> 'a Tyspeak.t
         -> Types.simple Tyspeak.t
\end{Verbatim}

\item Ces types sont vérifiés. Cette étape est normalement redondante et est
  donc optionnelle (l'option \texttt{-{}-no-check} permet de la désactiver).

\item Le programme obtenu, de type \texttt{Types.simple Tyspeak.t}, est affiché
  sur le terminal.

\end{itemize}

\section{Inférence de types}

L'inférence de types consiste à remplacer les étiquettes de type \texttt{unit}
par des étiquettes de type \texttt{simple} (autrement dit de vraies
représentations de types).

Cette étape se fait de manière impérative : on peut créer de nouveaux types avec
\texttt{new\_unknown} et unifier deux types
avec \texttt{unify}. Leurs types sont :

\begin{verbatim}
val new_unknown : unit -> Types.simple
val unify : Types.simple -> Types.simple -> unit
\end{verbatim}

La fonction \texttt{infer} s'appuie sur un ensemble de fonctions récursivement
définies portant sur chaque type de fragment : \texttt{infer\_fdec} pour les
déclarations de fonction, \texttt{infer\_exp} pour les expressions,
\texttt{infer\_stmtkind} pour les instructions, etc.

\begin{figure}

% TODO expliquer pourquoi il y a plusieurs left-values, et toutes les
% différences avec safespeak ainsi que décrit dans les chapitres précédents

\insertcode{infer-unif-stmt.ml}

\caption{Implantation -- inférence des déclarations de variable et appels de
         fonction}

% TODO ces "Implantation -- " sont ils vraiment indispensables ?
\label{fig:implem-unif-stmt}
\end{figure}

% TODO s/lval_type/infer_lval/ ?

Les règles de typage sont implantées par \texttt{new\_unknown} et
\texttt{unify}. Par exemple, pour typer une déclaration
(figure~\ref{fig:implem-unif-stmt}), on crée un nouveau
type \texttt{t0}. On étend l'environnement courant avec cette nouvelle
association et sous ce nouvel environnement, on type le bloc de portée de la
déclaration.

De même, pour typer un appel de fonction, on infère le type de ses arguments et
left-values de retour. On obtient également le type de la fonction (à partir du
type de la fonction présent dans l'environnement, ou du type du pointeur de
fonction qui est déréférencé), et on unifie ces deux informations.

On peut donner quelques exemples, comme :

\begin{itemize}
  \item addition de deux flottant (dans \texttt{infer\_binop}) :

\begin{verbatim}
let infer_binop op (_, a) (_, b) =
  match op with
    (* [...] *)
    | N.PlusF _ ->
        unify a Float;
        unify b Float;
        Float
\end{verbatim}

  \item adresse d'une left-value (dans \texttt{lval\_type}) :

\begin{verbatim}
| T.AddrOf lv ->
    let lv' = infer_lv env lv in
    let ty = lval_type env lv in
    (T.AddrOf lv', Ptr (Kernel, ty))
\end{verbatim}

  \item déréférencement d'une left-value (dans \texttt{lval\_type}) :

\begin{verbatim}
| T.Deref(e, _sz) ->
    let (_, te) = infer_exp env e in
    let t = new_unknown () in
    unify (Ptr (Kernel, t)) te;
    t
\end{verbatim}

\end{itemize}

\section{Vérification de types}

La passe de vérification de types, optionnelle et en théorie redondante, est
présente pour s'assurer que les types inférés ne sont pas aberrants.

Là encore, on travaille par effet de bord : si tout ne se passe pas bien, on
lève une erreur. La fonction principale a le type suivant :

\begin{verbatim}
val check : Types.simple Env.t -> Types.simple Tyspeak.t -> unit
\end{verbatim}

% TODO

\section{Unification}

\begin{figure}
\insertcode{implem-types.ml}
\caption{Implantation -- types}
\label{fig:implem-typ}
\end{figure}

Les types de données utilisés sont donnés dans la figure~\ref{fig:implem-typ}.
Les types \langname sont représentés soit par un constructeur de type ``résolu''
immutable comme \texttt{Int}, soit par une référence dans le cas d'une variable
inconnue (placée alors derrière le constructeur \texttt{Var}).

Ces références contiennent une valeur du type
\texttt{simple var\_type}\footnote{
  Ce type est polymorphe pour pouvoir être repris dans l'unification des
  qualificateurs (figure~\ref{fig:implem-qualifs}).
}, c'est à
dire :

\begin{itemize}
\item soit un numéro d'inconnue (constructeur \texttt{Unknown})\footnote{On
  place le numéro dans un enregistrement pour abstraire l'entier sous-jacent, en
  empêchant par exemple de faire de l'arithmétique sur celui-ci.}.

\item soit un type résolu (constructeur \texttt{Instanciated}) si cette
  inconnue a été unifiée avec un type concret.
\end{itemize}

Ce système peut créer des représentations de types arbitrairement longues, comme
par exemple :

\texttt{Var (ref (Instanciated (Var (ref (Instanciated Int)))))}

Cela est dû au fait que \texttt{fun x -> Var (ref (Instanciated x))} est typée
\texttt{simple -> simple} et peut donc être appliquée à loisir. Pour éviter cet
effet d'allongement, on définit une fonction \texttt{shorten} qui supprime ces
chaînes (figure~\ref{fig:implem-shorten}).

% TODO décrire la première branche du pattern matching

\begin{figure}

  \insertcode{implem-shorten.ml}

\caption{Implantation -- fonction de racourcissement des représentations de types}
\label{fig:implem-shorten}
\end{figure}

La fonction d'unification, quant à elle, commence par raccourcir ses deux
arguments puis faire une analyse de cas par filtrage.

Les cas principaux pour unifier deux types réduits \texttt{sta} et \texttt{stb}
sont :

\begin{itemize}

\item \texttt{sta} et \texttt{stb} sont deux inconnues. Alors on modifie l'un
  pour pointer sur l'autre. Les références étant uniques et partagées, cela
  revient à substituer l'un par l'autre dans toutes les représentations de
  types.

\item Un type est concret, l'autre une inconnue. Dans ce cas on modifie le
  second comme étant un \texttt{Instanciated} du premier. Il faut vérifier que
  le type inconnu remplacé n'apparaît pas dans le type concret, sinon on créé un
  cycle. Ce cas a lieu par exemple quand on cherche à unifier $a$ avec
  $a~\qKernel~*$. Il faut alors signaler une erreur, ce que fait la fonction
  \texttt{occurs\_check\_failed}.

\item Les deux types sont des types concrets. Alors ils sont de la forme
  respective
  $C (t_1, …, t_n)$
  et
  $D (u_1, …, u_m)$ où $C$ et $D$ sont des constructeurs de type avec
  respectivement $n$ et $m$ arguments
  \footnote{
    Le cas $n = 0$ ou $m = 0$ correspond aux types concrets comme $\tInt$ ou
    $\tFloat$.
  }. Si $C = D$ et
  $n = m$, alors on unifie récursivementchaque $t_i$ avec $u_i$. Sinon on lève une erreur
  (fonction \texttt{type\_clash}).

  % TODO structures

\end{itemize}

Cette analyse de cas est implantée dans la fonction \texttt{unify_now}
(figure~\ref{fig:implem-unify-now}).

\begin{figure}
  \insertcode{implem-unify.ml}
\caption{Implantation -- fonction d'unification}
\label{fig:implem-unify-now}
\end{figure}

cas structure (figure~\ref{fig:implem-unify-structs}).

\begin{figure}
  \insertcode{implem-structs.ml}
\caption{Implantation -- structures}
\label{fig:implem-unify-structs}
\end{figure}

Pour les pointeurs, il est nécessaire de définir des fonctions similaires sur
les qualificateurs (figure~\ref{fig:implem-qualifs}).

\begin{figure}

  \insertcode{implem-qualifs.ml}

  \caption{Implantation -- qualificateurs}
\label{fig:implem-qualifs}

\end{figure}

\begin{figure}

\insertcode{unify-main.ml}

\caption{Implantation -- unification directe ou retardée}
\label{fig:implem-lazy}
\end{figure}

La fonction appellée directement par le reste du code, appellée \texttt{unify},
peut retarder l'unification (figure~\ref{fig:implem-lazy}). Dans ce cas, la
paire de types à unifier est mise dans une liste d'attente qui sera unifiée
après le parcours du programme. Le but est d'instrumenter l'inférence de types
afin de pouvoir en faire une exécution ``pas à pas''.

\section{Exemple}

Lançons l'analyse sur un petit exemple :

\begin{verbatim}
int f(int *x) { return (*x + 1); }
\end{verbatim}

L'exécution de notre analyseur affiche un programme complètement annoté :

\begin{verbatim}
 % ptrtype example.c
f : (Ptr (Int)) -> (Int)
Int (example.c:1#4)^f(Ptr (Int) x) {
  (.c:3#4)^!return =(int32)
    (coerce[-2147483648,2147483647]
      ( ( ([(x_Ptr (Int) : Ptr (Int))]32_Int
            : Int
          )
          + (1 : Int)
        ) : Int
      ) : Int
    );
}
\end{verbatim}

L'opérateur \texttt{coerce[a,b]} est un artefact de Newspeak, destiné à détecter
les débordements d'entiers lors d'une analyse de valeurs par interprétation
abstraite. Dans le cas de notre analyse, les valeurs ne sont pas pertinentes et
cet opérateur peut être vu que comme un ``plus'' unaire typé $(\tInt) → \tInt$.

\emph{A contrario}, lorsqu'il n'est pas possible d'inférer des types
compatibles, l'analyseur s'arrête avec une erreur.

\begin{verbatim}
void f(int *p) {
    /*!npk userptr p */
    *p = 3;
}
\end{verbatim}

Le commentaire \texttt{/*!npk userptr p */} est interpréter par l'analyseur et
le fait unifier le type de \texttt{p} avec $t~\qUser~*$, c'est à dire qu'il
force son qualificateur à être \qUser.

\begin{verbatim}
04-addrof.c:4#4 - Cannot unify qualifiers:
  Kernel
  User
\end{verbatim}

On signale l'emplacement où la dernière unification a échoué.

% vim: spelllang=fr


\chapter{Étude de cas : le noyau Linux}

\label{cha:etudedecas}
\changed{le titre était "étude de cas : un pilote de carte graphique"}

Le noyau Linux, abordé dans le chapitre~\ref{cha:os}, est un noyau de système
d'exploitation développé depuis le début des années 90 et \enquote{figure de
proue} du mouvement \emph{open-source}. Au départ écrit par Linus Torvalds sur
son ordinateur personnel, il a été porté au fil des années sur de nombreuses
architectures et s'est enrichi de nombreux pilotes de périphériques. Dans la
version 3.13.1 (2014), son code source comporte 12 millions de lignes de code
(en grande majorité du C) dont 58\% de pilotes.

\changed{màj des statistiques}

Même si le noyau est monolithique (la majeure partie des traitements s'effectue
au sein d'un même fichier objet), les sous-systèmes sont indépendants. C'est ce
qui permet d'écrire des pilotes de périphériques et des modules.

Ces pilotes manipulent des données provenant de l'utilisateur, notamment par
pointeur. Comme on l'a vu, cela peut poser des problèmes de sécurité si on
déréférence ces pointeurs sans vérification.

Dans ce chapitre, on met en œuvre sur le noyau Linux le système de types décrit
dans le chapitre~\ref{cha:qualifs}, ou plus précisément l'outil \ptrtype du
chapitre~\ref{cha:implem}.

Pour montrer que le système de types capture cette propriété et que
l'implantation est utilisable, on étudie les cas de deux bugs qui ont touché le
noyau Linux. À chaque fois, dans une routine correspondant à un appel système,
un pointeur utilisateur est déréférencé directement, pouvant provoquer une fuite
d'informations confidentielles dans le noyau.

On commence par décrire les difficultés rencontrées pour analyser le code du
noyau Linux. On décrit ensuite l'implantation du mécanisme d'appels système dans
ce noyau, et en quoi cela peut poser des problèmes. On détaille enfin les bugs
étudié, et comment les adapter pour traiter le code en question.

\section{GNU C}
\label{sec:gnuc}

Linux est écrit dans le langage C, mais pas dans la version qui correspond à la
norme. Il utilise le dialecte GNU C qui est celui que supporte GCC. Une première
difficulté pour traiter le code du noyau est donc de le compiler.

Pour traduire ce dialecte, il a été nécessaire d'adapter \ctonewspeak. La
principale particularité est la notation \verb!__attribute__((...))! qui peut
décorer les déclarations de fonctions, de variables ou de types.

Par exemple, il est possible de manipuler des étiquettes de première classe: si
\enquote{\texttt{lbl:}} est présent avant une instruction, on peut capturer
l'adresse de celle-ci avec \texttt{void *p = \&\&lbl} et y sauter indirectement
avec \texttt{goto *p}.

Une autre fonctionnalité est le concept d'instruction-expression:
\texttt{(\{bloc\})} est une expression, dont la valeur est celle de la dernière
expression évaluée lors de \texttt{bloc}.

Les attributs, quant à eux, rentrent dans trois catégories:

\begin{itemize}
  \item les annotations de compilation; par exemple, \texttt{used} désactive
  l'avertissement \enquote{cette variable n'est pas utilisée}.

  \item les optimisations; par exemple, les objets marqués \texttt{hot} sont
    groupés de telle manière qu'ils se retrouvent en cache ensemble.

  \item les annotations de bas niveau; par exemple, \verb!aligned(n)!
    spécifie qu'un objet doit être aligné sur au moins \texttt{n} bits.
\end{itemize}

Dans notre cas, toutes ces annotations peuvent être ignorées, mais il faut tout
de même adapter l'analyse syntaxique pour les ignorer. En particulier, pour le
traitement du noyau Linux, il a fallu traiter certaines formes de la
construction \texttt{typeof} qui n'étaient pas supportées.

De plus, pour que le code noyau soit compilable, il est nécessaire de définir
certaines macros. En particulier, le système de configuration de Linux utilise
des macros nommées \texttt{CONFIG\_*} pour inclure ou non certaines
fonctionnalités. Il a donc fallu faire un choix; nous avons choisi la
configuration par défaut. Pour analyser des morceaux plus importants du noyau,
il faudrait définir un fichier de configuration plus important.

\section{Appels système sous Linux}
\label{sec:linux-sys}

Dans cette section, nous allons voir comment ces mécanismes sont implantés dans
le noyau Linux. Une description plus détaillée pourra être trouvée
dans~\cite{UnderstandingTheLinuxKernel}, ou, pour le cas de la mémoire
virtuelle, dans~\cite{LinuxVMM}.

Deux rings sont utilisés: en \ring{0}, le code noyau et, en \ring{3}, le code
utilisateur.

Une notion de tâche similaire à celle décrite dans la section~\ref{sec:taches}
existe: elles s'exécutent l'une après l'autre, le changement s'effectuant sur
interruptions.

Pour faire appel aux services du noyau, le code utilisateur doit faire appel à
des appels système, qui sont des fonctions exécutées par le noyau. Chaque tâche
doit donc avoir deux piles: une pile \enquote{utilisateur}, qui sert pour
l'application elle-même, et une pile \enquote{noyau}, qui sert aux appels
système.

Grâce à la mémoire virtuelle, chaque processus possède sa propre vue de la
mémoire dans son espace d'adressage (figure~\ref{fig:memmap}), et donc chacun
gère un ensemble de tables de pages et une valeur de \crtrois{} associée (ce
mécanisme a été abordé page~\pageref{page:mem-virt}). Au moment de changer le
processus en cours, l'ordonnanceur charge donc le \crtrois{} du nouveau
processus.

\begin{figure}[b] % fig:memmap {{{
\centering
\begin{tikzpicture}
  [scale=0.7
  ,user/.style={fill=black!20}
  ,kernel/.style={fill=black!70}
  ]

  % Memory zone
  %
  % #1 - start
  % #2 - end
  % #3 - color
  \newcommand{\mzone}[3]{
    \path[#3] (#1,0) rectangle (#2,1);
  }

  % Address label
  %
  % #1 - x position
  % #2 - text
  \newcommand{\alabel}[2]{
    \path (#1,1) -- ++(0,0.3) node [pos=1] {\small \tt #2};

  }

  % exec
  \mzone{0.5}{1}{user}

  % lib
  \mzone{2.5}{3.2}{user}

  % stack
  \mzone{3.7}{4}{user}

  % stack
  \mzone{5}{5.5}{user}

  % kernel
  \mzone{6}{8}{kernel}

  % contour
  \draw (0,0) rectangle (8,1);

  \alabel{0}{0}
  \alabel{6}{3 Go}
  \alabel{8}{4 Go}

\end{tikzpicture}


\caption[Espace d'adressage d'un processus]{L'espace d'adressage d'un processus.
En gris clair, les zones accessibles à tous les niveaux de privilèges: code du
programme, bibliothèques, tas, pile. En gris foncé, la mémoire du noyau,
réservée au mode privilégié.}

\label{fig:memmap}
\end{figure}
% }}}

Les adresses basses (inférieures à \texttt{PAGE\_OFFSET} = 3 Gio =
\hexa{c0000000}) sont réservées à l'utilisateur. On y trouvera par exemple:
le code du programme,
les données du programme (variables globales),
la pile utilisateur,
le tas (mémoire allouée par \texttt{malloc} et fonctions similaires), ou encore
les bibliothèques partagées.

\changed{paragraphe plutôt que liste}

Au dessus de \texttt{PAGE\_OFFSET}, se trouve la mémoire réservée au noyau.
Cette zone contient le code du noyau, les piles noyau des processus, etc.

\label{sec:impl-syscall}

Les programmes utilisateur s'exécutant en \ring{3}, ils ne peuvent pas contenir
d'instructions privilégiées, et donc ne peuvent pas accéder directement au
matériel. Pour que ces programmes puissent interagir avec le système (afficher
une sortie, écrire sur le disque…) le mécanisme des appels système est
nécessaire. Il s'agit d'une interface de haut niveau entre les \emph{rings} 3 et
0. Du point de vue du programmeur, il s'agit d'un ensemble de fonctions C
\enquote{magiques} qui font appel au système d'exploitation pour effectuer des
opérations.

Par exemple, le programmeur peut appeller la fonction \texttt{getpid} pour
connaître le numéro du processus courant. Cela passe par une fonction
\texttt{getpid} dans la bibliothèque C, en espace utilisateur. Celle-ci va
invoquer (via un mécanisme non pertinent ici) la fonction \texttt{sys\_getpid}
du noyau (figure~\ref{fig:syscall-def}).

Comme les piles sont différentes entre les espaces, la convention d'appel est
différente: les arguments sont copiés directement par les registres.

\begin{figure}[h]
    \insertcode{syscall-definition.c}
    \caption{Fonction de définition d'un appel système}
    \label{fig:syscall-def}
\end{figure}

La macro \texttt{SYSCALL\_DEFINE0} permet de nommer la fonction
\texttt{sys\_getpid}, et définit entre autres des points d'entrée pour les
fonctionnalités de débogage du noyau. Le corps de la fonction fait directement
référence aux structures de données internes du noyau pour retourner le résultat
voulu.

\section{Risques}
\label{sec:risques-dos}

Ainsi que décrit dans la section~\ref{sec:secu-syscalls}, cela peut poser un
problème de manipuler des pointeurs contrôlés par l'utilisateur au sein d'une
routine de traitement d'appel système.

Si le déréférencement est fait sans vérification, un utilisateur mal intentionné
peut forger un pointeur vers le noyau (en déterminant des adresses valides dans
l'espace noyau entre \hexa{c0000000} et \hexa{ffffffff}). En provoquant
une lecture sur ce pointeur, des informations confidentielles peuvent fuiter ;
et, en forçant une écriture, il est possible d'augmenter ses privilèges, par
exemple en devenant super-utilisateur (\emph{root}). En pratique, il n'est pas
toujours possible d'accéder à la mémoire. La mémoire utilisateur peut par
exemple avoir été placée en zone d'échange sur le disque, ou \emph{swap}. À ce
moment là, l'erreur provoquera tout de même un déni de service. Plus de détails
sur ce mécanisme, et le fonctionnement de la mémoire virtuelle dans Linux,
peuvent être trouvés dans~\cite{userspaceaccess}.

\section{Premier exemple de bug : pilote Radeon KMS}

On décrit le cas d'un pilote vidéo qui contenait un bug de pointeur utilisateur.
Il est répertorié sur \url{http://freedesktop.org} en tant que bug \#29340.

Pour changer de mode graphique, les pilotes de GPU peuvent supporter le
\emph{Kernel Mode Setting} (KMS).

Pour configurer un périphérique, l'utilisateur communique avec le pilote noyau
avec le mécanisme d'\emph{ioctls} (pour \emph{Input/Output Control}).
Ils sont similaires à des appels système, mais spécifiques à un périphérique
particulier. Le transfert de contrôle est similaire à ce qui a été décrit dans
la section précédente: les applications utilisateurs appellent la fonction
\texttt{ioctl()} de la bibliothèque standard, qui provoque une interruption.
Celle-ci est traitée par la fonction \texttt{sys\_ioctl()} qui appelle la
routine de traitement dans le bon pilote de périphérique.

Les fonctions du noyau implantant un \emph{ioctl} sont donc vulnérables à la
même classe d'attaques que les appels système, et donc doivent être écrits avec
une attention particulière.

Le code de la figure~\ref{fig:kms-buggy} est présent dans le pilote KMS pour les
GPU AMD Radeon.

\begin{figure}[h]
\insertcode{kms-buggy.c}
\caption{Code de la fonction \texttt{radeon\_info\_ioctl}}
\label{fig:kms-buggy}
\end{figure}

On peut voir que l'argument \texttt{data} est converti en un \texttt{struct
drm\_radeon\_info *}. Un pointeur \texttt{value\_ptr} est extrait de son champ
\texttt{value}, et finalement ce pointeur est déréferencé.

Cependant, l'argument \texttt{data} est un pointeur vers une structure (allouée
en espace noyau) du type donné dans la figure~\ref{fig:kms-info}, dont les
champs proviennent d'un appel utilisateur de \verb!ioctl()!.

\begin{figure}[h]
\insertcode{kms-info.c}
\caption{Définition de \texttt{struct drm\_radeon\_info}}
\label{fig:kms-info}
\end{figure}

Pour mettre ce problème en évidence, nous avons annoté la fonction
\texttt{radeon\_info\_ioctl} de telle manière que son second paramètre soit un
pointeur noyau vers une structure contenant un champ contrôlé par l'utilisateur,
\texttt{value}.

L'intégralité de ce code peut être trouvée en annexe~\ref{cha:code-noyau}.

\begin{figure}
  \insertcode{radeon-git-mini.diff}

  \caption[Patch résolvant le problème de pointeur utilisateur.]
      {Patch résolvant le problème de pointeur utilisateur.
       La ligne précédée par un signe \texttt{-} est supprimée et remplacée
       par les lignes précédées par un signe \texttt{+}.
      }
\label{fig:linux-patch}
\end{figure}

La bonne manière de faire a été publiée avec le numéro de \emph{commit}
\texttt{d8ab3557} (figure~\ref{fig:linux-patch}) (\texttt{DRM\_COPY\_FROM\_USER}
étant une simple macro pour \texttt{copy\_from\_user}). Dans ce cas, on
n'obtient pas d'erreur de typage.

\added{Toute la section}
\section{Second exemple : \texttt{ptrace} sur architecture Blackfin}

Le noyau Linux peut s'exécuter sur l'architecture Blackfin, qui est spécialisée
dans le traitement du signal. Le problème de manipulation des pointeurs
utilisateur auquel nous nous intéressons peut également s'y produire.

En particulier nous nous intéressons à l'appel système \texttt{ptrace}. Il
permet à un processus d'accéder à la mémoire et de contrôler l'exécution d'un
autre processus, par exemple à des fins de débogage. Ainsi,
\texttt{ptrace(PTRACE\_PEEKDATA, p, addr)} renvoie la valeur du mot mémoire à
l'adresse \texttt{addr} dans l'espace d'adressage du processus \texttt{p}.

Comme pour la plupart des appels système, la fonction \texttt{ptrace} est
dépendante de l'architecture. Le deuxième exemple que nous présentons concerne
l'implantation de celle-ci pour les processeurs Blackfin,
figure~\ref{fig:ptrace-blackfin}.

Dans d'anciennes versions de Linux\footnote{Jusqu'à la version 2.6.28 --- ce bug
a été corrigé dans le \emph{commit} \texttt{7786ce82} en remplaçant l'appel à
\texttt{memcpy} par un appel à \texttt{copy\_from\_user\_page}}, cette fonction
appelle \texttt{memcpy} au lieu de \nolinkurl{copy\_from\_user} pour lire dans
la mémoire du processus. La ligne problématique est préfixée par
\texttt{/*=\textgreater{}*/}. En théorie, si un utilisateur passe un pointeur
vers une adresse du noyau à la fonction \texttt{ptrace}, il pourra lire des
données du noyau. L'appel \texttt{ptrace (PTRACE\_PEEKDATA, p, addr)} permet
ainsi non seulement de lire les variables du processus \texttt{p} si
\texttt{addr} est une adresse dans l'espace utilisateur (ce qui est le
comportement attendu), mais aussi de lire dans l'espace noyau si \texttt{addr} y
pointe (ce qui est un bug de sécurité).

\begin{figure}
    \insertcode{blackfin-ptrace.c}
    \caption{Implantation de \texttt{ptrace} sur architecture Blackfin}
    \label{fig:ptrace-blackfin}
\end{figure}

% TODO[E] même question : est ce possible de faire tourner le proto dessus

On peut repérer ce bug par simple relecture pour commencer. On commence par
remarquer que l'argument \texttt{addr}, malgré son type \texttt{long}, est en
réalité un \texttt{void *} provenant directement de l'espace utilisateur. C'est
en effet le même argument \texttt{addr} de l'appel système \texttt{ptrace}.
Cet argument correspond à l'adresse à lire dans l'espace mémoire du processus.
Comme il est passé à \texttt{memcpy}, aucune vérification n'est faite avant la
copie. La valeur pointée par \texttt{addr} sera copiée, même si elle est en
espace noyau.

En annotant correctement les types, on peut donc détecter ce bug: le type
correct de \texttt{addr} est \ptrU{\tInt}, et celui de \texttt{memcpy} est
$(\ptrK{\tInt}, \ptrK{\tInt}, \tInt) → \ptrK{\tInt}$. Il est donc impossible de
lui passer cet argument. Remarquons que le type de \texttt{memcpy} en C utilise
des pointeurs de type \texttt{void *}. Pour les traiter correctement on pourrait
utiliser du polymorphisme, mais dans ce cas précis utiliser le type \ptrK{\tInt}
est suffisant.

\paragraph{Remarque}

En pratique, le problème de sécurité n'est pas si important. En effet,
la copie se fait sous un test forçant \texttt{addr} à être entre
\nolinkurl{FIXED\_CODE\_START} et \texttt{FIXED\_CODE\_END}. Cette zone est
incluse en espace utilisateur, cela empêche donc le problème de fuite de
données.

Mais cela reste un problème de sécurité: contrairement à
\texttt{copy\_from\_user}, la fonction \texttt{memcpy} ne vérifie pas que
l'espace utilisateur est chargé en mémoire. Si ce n'est pas le cas, une faute
mémoire sera provoquée dans le noyau. Il s'agit alors d'un déni de service
(section~\ref{sec:risques-dos}), qui est tout de même un comportement à
empêcher.

\section{Procédure expérimentale}
\label{sec:demo-unif}

Pour utiliser notre système de types, plusieurs étapes sont nécessaires en plus
de traduire le noyau Linux en \langname{}.

Afin de réaliser l'analyse, il faut annoter les sources pour créer un
environnement initial (via la variable \texttt{exttbl} décrite en
section~\ref{sec:ptrtype-archi}). Plus précisément, pour chaque source de
pointeurs utilisateur, on ajoute un commentaire spécial \texttt{!npk
userptr\_fieldp x f}, qui indique que \texttt{x} est un pointeur vers une
structure contenant un pointeur utilisateur dans le champ \texttt{f}. En fait,
il unifie le type de \texttt{x} avec $\{ f: \ptrU{t} ; … \}~*$ où $t$ est une
inconnue de type. Cette annotation est nécessaire car c'est le moyen d'indiquer
que la structure contient un pointeur utilisateur.

Par rapport au code complet présent dans l'annexe~\ref{cha:code-noyau},
l'expression calculant \texttt{value_ptr} est également simplifiée. Dans le code
d'origine, \texttt{info->value} est transtypé en \texttt{unsigned long} puis en
\texttt{uint32\_t *}. En \newspeak, cela correspond à des opérateurs
\texttt{PtrToInt} \linebreak et \texttt{IntToPtr} mais, si on les autorise, on
casse le typage puisqu'il est alors possible de transformer n'importe quel type
en un autre. De plus, on modifie la définition du type \texttt{struct
drm\_radeon\_info} pour que son champ \texttt{value} ait pour type
\texttt{uint32\_t *} plutôt que \texttt{uint64\_t}. En effet, dans ce cas
d'étude, cet entier est uniquement utilisé en tant que pointeur au cours de
toute l'exécution.

\begin{figure}[p]
\insertcode{ex-drm.c}

\caption{Cas d'étude minimisé et annoté}
\label{fig:ex-drm}
\end{figure}

\begin{figure}[p]

    \insertcode{ex-drm-ok.c}

    \caption{Cas d'étude minimisé et annoté -- version correcte}
    \label{fig:ex-drm-ok}

\end{figure}

En ce qui concerne les fonctions de manipulation de pointeurs fournies
par le noyau (\texttt{get\_user}, \texttt{put\_user},
\texttt{copy\_from\_user}, \texttt{copy\_to\_user}, etc.), on ajoute à
l'environnement global leur type correct.

Enfin, on peut lancer l'inférence de type. Ainsi, sur l'exemple de la
figure~\ref{fig:ex-drm}, on obtient la sortie suivante:

\begin{Verbatim}
05-drm.c:18#8 - Type clash between :
  KPtr (_a15)
  UPtr (_a8)
\end{Verbatim}

Cela indique qu'on a essayé d'unifier un type de la forme $\ptrK{t}$ avec un
type de la forme $\ptrU{t}$, en précisant l'emplacement où la dernière
unification a échoué (les \texttt{\_aN} correspondent à des inconnues de type).

% TODO idem sur #2
% TODO[E] ce serait bien de mettre Blackfin aussi

La version correcte minimisée correspond à la figure~\ref{fig:ex-drm-ok}. Pour
celle-ci, l'inférence se fait sans erreur.

% TODO le faire

\section*{Conclusion}

Après voir décrit l'implantation de notre solution, on a montré comment celle-ci
peut s'appliquer à détecter un bug dans le noyau Linux. La première difficulté
est de traduire en \newspeak le code source écrit dans le dialecte GNU C.

En partant d'un bug existant, on montre que la version originale du code
(incluant une erreur de programmation) ne peut pas être typée, alors que sur la
version corrigée on peut inférer des types compatibles.

Le prototype décrit dans le chapitre~\ref{cha:implem} peut donc s'adapter à
détecter des bugs dans le noyau Linux. Pour le moment, il nécessite du code
annoté, mais des travaux sont en cours pour permettre de passer automatiquement
des portions plus importantes du noyau Linux. Le principal obstacle est de
devoir réécrire certaines parties du code pour supprimer les constructions non
typables.

% TODO[E] à étendre, +++

% vim: spelllang=fr


\def\partintro{}
\partccl{part:xp}
Après avoir décrit notre solution théorique dans la partie~\ref{part:lang}, nous
avons présenté ici notre démarche expérimentale. Dans le
chapitre~\ref{cha:implem}, nous avons détaillé l'implantation de notre
prototype. Pour ce faire, nous avons ajouté des étiquettes de type au langage
\newspeak et adapté l'algorithme W de Damas et Milner pour réaliser l'inférence
de types. Ce prototype est distribué sur~\link{penjili} sous le nom de
\ptrtype. Il est constitué d'environ 1600 lignes d'OCaml.

Ensuite, le but du chapitre~\ref{cha:etudedecas} est d'appliquer notre analyse
(à l'aide de ce prototype) sur le noyau Linux. Après avoir décrit le
fonctionnement des appels système sur ce noyau, on présente deux bugs qui ont
touché respectivement un pilote de carte graphique et l'implantation d'un appel
système. Ils sont la manifestation d'un problème de pointeur utilisateur mal
déréférencé dans le noyau, ainsi que décrit dans le chapitre~\ref{cha:os}. En
lançant notre analyse sur le code présentant un problème, l'erreur est détectée.
Au contraire, en la lançant sur le code après application du correctif, aucune
erreur n'est trouvée.

En s'appuyant sur le langage \newspeak, on gagne beaucoup par rapport à d'autres
représentations intermédiaires. Le fait d'avoir un langage avec peu de
constructions permet de ne pas avoir à exprimer plusieurs fois la même règle
(par exemple, une fois sur la boucle \emph{for} et une autre sur la boucle
\emph{while}).

\label{page:ccl-npk-spk}

Un des inconvénients est que le modèle mémoire utilisé par \newspeak est assez
différent de celui de \langname (ainsi que décrit dans le
chapitre~\ref{cha:lang}). \newspeak est en effet prévu pour implanter des
analyses précises de valeur reposant sur l'interprétation abstraite, et
nécessite donc un modèle mémoire de plus bas niveau (où on peut créer des
valeurs à partir d'une suite d'octets, par exemple).

Le prototype d'implantation peut évoluer dans deux directions: d'une part, en
continuant à s'appuyer sur \newspeak, on peut réaliser des pré-analyses de
typage qui permettent de guider une analyse de valeurs plus précise, par exemple
en choisissant un domaine abstrait différent en fonction des types de données
rencontrés. D'autre part, il est possible de faire une implantation plus fidèle
à \langname, qui permette d'ajouter de nouvelles fonctionnalités plus éloignées
de C. Par exemple, un système de régions comme~\cite{jfp92} permettrait de
simplifier l'environnement d'exécution en enlevant l'opération de nettoyage
mémoire $\phx\cCleanup$. Le système de types peut également être enrichi, pour
ajouter par exemple du polymorphisme. Cela rapprocherait le langage source de
Rust. Le chapitre~\ref{cha:conclusion} présente quelques unes de ces extensions
possibles.

L'expérimentation, quant à elle, est pour le moment limitée, mais on peut
l'étendre à des domaines de plus en plus importants dans le noyau Linux. Tout
d'abord, le module graphique définit d'autres fonctions implantant des
\emph{ioctls}. Celles-ci reçoivent donc également des pointeurs utilisateur et
sont susceptibles d'être vulnérables à ce genre d'erreurs de programmation.
Ensuite, d'autres modules exposent une interface similaire, à commencer par les
autres pilotes de cartes graphiques. Ceux-ci sont également un terrain sur
lequel appliquer cette analyse.

De manière générale, toutes les interfaces du noyau manipulant des pointeurs
utilisateur gagnent à être analysées. Outre les implantations des \emph{ioctls}
dans chaque pilote et les appels système, les systèmes de fichiers manipulent
aussi de tels pointeurs via leurs opérations de lecture et d'écriture.

% vim: spelllang=fr


\part{Conclusion}
% TODO pas besoin d'une partie pour ça. Faire comme l'intro

\chapter{Conclusion}

\label{cha:conclusion}

% TODO [E] pour la conclusion il y a deux plans possibles :
% 1/ limitation et travaux futurs
% 2/ parler de chaque point que l'on rencontre actuellement : transtypage...

% TODO[E] il faut dans la conclusion indiquer ce que tu voulais traiter et ce
% que tu as traité dans ta thèse

\section{Limitations}

Un seul type d'entier est supporté.

% TODO

% TODO [E] tu parles de safespeak et de la traduction c-safespeak?

\paragraph{Transtypage}

Puisque l'approche retenue est basée sur le typage statique, il est impossible
de capturer de nombreuses constructions qui sont possibles, où même idiomatiques
en C : les unions, les conversions de types (explicites ou implicites) et le
\emph{type punning}. Les deux premières sont équivalentes. Bien qu'on puisse
remplacer chaque conversion explicite d'un type $t_1$ vers un type $t_2$ par
l'appel à une fonction $\mathrm{cast}_{t_1,t_2}$, on ajoute alors un ``trou''
dans le système de types. Cette fonction devrait en effet être typée $(t_1) →
t_2$, autrement dit le type ``maudit'' $α → β$ de \texttt{Obj.magic} en OCaml ou
\texttt{unsafeCoerce} en Haskell.

% TODO [E] à définir

Le \emph{type punning} consiste à modifier directement la suite de bits de
certaines données pour la manipuler d'une manière efficace. Par exemple, il est
commun de définir un ensemble de macros pour accéder à la mantisse et à
l'exposant de flottants IEEE754. Ceci peut être fait avec des unions ou des
masques de bits.

Dans de tels cas, le typage statique est bien sûr impossible. Pour traiter ces
cas, il faudrait encapsuler la manipulation dans une fonction et y ajouter une
information de type explicite, comme $\texttt{float\_exponent} : (\tFloat) →
\tInt$.

\paragraph{Environnement d'exécution}

La sémantique opérationnelle utilise un environnement d'exécution pour certains
cas. Contrairement à C, les débordements de tampon et les déréférencements de
pointeurs sont vérifiés dynamiquement. Mais ce n'est pas une caractéristique
cruciale de cette approche : en effet, si on suppose que les programmes que l'on
analyse sont corrects de ce point de vue, on peut désactiver ces vérifications
et le reste des propriétés est toujours valable.

Un autre endroit, plus problématique, où des tests dynamiques sont faits est
lorsqu'on recherche en mémoire des pointeurs référençant un cadre d'appel qui
n'est plus valide (à l'aide de l'opérateur $\mathrm{Cleanup}(\cdot)$). Supprimer
ce test rend l'analyse incorrecte, car il est alors possible de faire référence
à une variable avec un type différent.

De même, si on peut avoir une garantie statique que les adresses des variables
locales ne seront plus accessibles au retour d'une fonction, alors on peut
supprimer le nettoyage en posant $\mathrm{Cleanup}(m) = m$. Cette garantie peut
être obtenue avec une analyse statique préalable~\cite{ifm10}.

\paragraph{Allocation dynamique}

La plupart des programmes, et le noyau Linux en particulier, utilisent la notion
d'allocation dynamique de mémoire. C'est une manière de créer dynamiquement une
zone de mémoire qui restera accessible après l'exécution de la fonction
courante. Cette mémoire pourra être libérée à l'aide d'une fonction dédiée. Dans
l'espace utilisateur, les programmes peuvent utiliser les fonctions
\texttt{malloc()}, \texttt{calloc()}, et \texttt{realloc()} pour allouer des
zones de mémoire et \texttt{free()} pour les libérer. Dans le noyau Linux, ces
fonctions existent sous la forme de \texttt{kmalloc()}, \texttt{kfree()}, etc.

% TODO citer the linux memory manager ?

Ces fonctions manipulent les données en tant que zones mémoires opaques, à
en renvoyant un pointeur vers une zone mémoire d'un nombre d'octets donnés.
Cela présuppose un modèle mémoire plus bas niveau. Pour se rapprocher de la
sémantique de \langname, une manière de faire est de définir un opérateur de
plus haut niveau prenant une expression et retournant l'adresse d'une cellule
mémoire contenant cette valeur. Le typage est alors direct :

% TODO [E] sous entendu il y a un sizeof(t)

\begin{mathpar}
  \irule{Malloc}
    { Γ ⊢ e : t }
    { Γ ⊢ \textsc{Alloc}(e) : t~\qKernel~* }

  \irule{Free}
    { Γ ⊢ e : t~\qKernel~* }
    { Γ ⊢ \textsc{Free}(e) }
\end{mathpar}

En ce qui concerne l'exécution, on peut ajouter une troisième composante aux
états mémoire : $m = (s, g, h)$ où $h$ associe à une étiquette numérique une
valeur (comme les globales mais indexées par des entiers plutôt que par leur
nom). La libération de mémoire est problématique parce qu'il faut faire
confiance au programmeur pour ne pas accéder aux zones mémoires libérées. Encore
une fois, on peut utiliser une analyse préalable comme~ \cite{ifm10}. Il est sûr
(mais pas possible en pratique) d'ignorer les commandes de libérations de
mémoire. 

% TODO [E] pointeurs fantômes

\paragraph{Structures récursives}

Une autre limitation est que seules les structures simples sont supportées. Les
structures récursives, contenant un pointeur vers un objet du même type
structure, ne sont pas typables. La raison est syntaxique : l'accès à un champ
est noté $lv.l_S$ afin d'avoir à disposition le type complet de la structure.

% TODO [E] bigre!!!

Par exemple, une liste chaînée d'entiers peut être décrite de la manière
suivante en C :

\begin{verbatim}
struct int_list {
    int value;
    struct int_list *next;
};
\end{verbatim}

% TODO et les listes chainées de Linux ? safety + typage

La compilation est possible car la définition complète de \texttt{struct
int\_list} n'est pas nécessaire pour connaître la taille d'un pointeur vers un
objet de ce type.

Pour compiler \texttt{l.next} en \langname, il faudrait écrire le type en
extension :

\[
  l.next_{ value:\tInt ; next : \{ value : \tInt ; next : … \}~*}
\]

On peut fermer cette récursion infinie en ajoutant directement un opérateur de
point fixe qui rend accessible le type d'un pointeur vers le type structure en
train d'être défini.

\begin{align*}
    S \gramisa & … \gramorx \mathrm{fix}(p → S)
\end{align*}

On écrit alors :

\[
  l.next_{ \mathrm{fix}(p → \{ value : \tInt ; next : p \}) }
\]

Les structures mutuellement récursives peuvent être également exprimée de cette
manière.

\begin{verbatim}
struct a {
    int value;
    struct b *next;
};

struct b {
    float value;
    struct a *next;
};
\end{verbatim}

Ces structures sont respectivement traduites en :

\begin{align*}
  A &= \mathrm{fix} (p → \{ value : \tInt   ; next : \{ value : \tFloat ; next : p \}~* \}) \\
  B &= \mathrm{fix} (p → \{ value : \tFloat ; next : \{ value : \tInt   ; next : p \}~* \})
\end{align*}

Ce schéma suffit à compiler toutes les types de structures possibles en C. En
effet, les seuls cas posant problème arrivent lorsqu'on fait référence à la
structure en cours de définition ; et il n'est alors possible que d'y accéder
par pointeur.

% TODO est-ce vrai ? Y'a-t-il des cas où les forward decl sont nécessaires?

\section{Travaux futurs}

\subsection{Transtypage}

% TODO [E] pourquoi?

Les casts entre types entiers ne posent pas de problèmes. Ceux entre types
pointeurs ne posent pas de problème non plus puisqu'au moment de la traduction
en \langname on efface les étiquettes de types sur les pointeurs.

Mais les casts entre entiers et pointeurs posent problème. Si on autorise les
casts, la seule manière sensible de les typer est la suivante.

\begin{mathpar}
\irule{PtrInt-Bad}
  { Γ ⊢ e : t~q~*}
  { Γ ⊢ (\tInt)~e : \tInt }

\irule{IntPtr-Bad}
  { Γ ⊢ e : \tInt }
  { Γ ⊢ (\textsc{Ptr})~e : t~q~*}
\end{mathpar}

Cela pose problème car il est alors possible de créer une fonction qui convertit
n'importe quel type pointeur en n'importe quel autre type pointeur :

\[
  ⊢ \eFun{p}{ \iReturn{(\textsc{Ptr})~(\tInt)~p} }
  : (t_a~q_a~*) → t_b~q_b~*
\]

% TODO et à l'exec?
% TODO on peut faire un dynamic check ou un type UInt ou un qualif sur tous les
% types

% TODO le modèle mémoire de Newspeak est bcp plus bas niveau que celui de
% Safespeak, ex types de deref, tailles dans Lv etc

L'opération problématique est en fait l'opérateur $(\tInt)$ de conversion d'un
pointeur vers un entier, car l'entier résultant ne pose pas de problème de
sûreté ni de sécurité. Pour lui donner une sémantique, on peut supposer que
l'environnement d'exécution est paramétré par une fonction de placement en
mémoire $\addrInt : Φ → \tInt$ qui à un chemin associe un entier.
On étend alors la sémantique par la règle suivante.

% TODO [E] ??

\[
  \irule{CastInt}
    { }
    { \mms{(\tInt)~\widehat{\&}φ}{\addrInt(φ)} }
\]

% TODO nommer les ensembles
% TODO mettre la syntaxe en annexe

Quant au transtypage dans le sens inverse, il est plus compliqué à traiter.
À l'exécution il est en effet nécessaire d'avoir à disposition une fonction
$\interpPtr{t~q~*} : \tInt → t~q~*$ qui permette d'interpréter un entier comme
pointeur qualifié $q$ vers un objet de type $t$ de manière qu'on ajoute la règle
suivante.

\[
  \irule{CastPtr}
    { }
    { \mms{(\textsc{Ptr})~\widehat{n}}{\interpPtr{t~q~*}(\widehat{n})} }
\]

% TODO les widehat...
% TODO [E] ??

L'exécution de l'analogue de $\addrInt$ demanderait de créer un pointeur
depuis un entier.

% TODO il n'y a pas de t q en mémoire, doncil faut juste un interp : Int -> Φ
% TODO problème car ça peut correspondre à plusieurs pointeurs
% TODO problème de layout mémoire : ça marche qd tout est boxé mais pas si c'est
% unboxé

La solution la plus sûre est donc d'interdire totalement ces conversions.
% TODO [E] violent mais plus sûr

\subsection{Analyse du noyau Linux}

Ici nous avons présenté l'analyse expérimentale d'une fonction problématique
d'une interface programmation particulière. Mais le principe de l'analyse est
applicable à toutes les fonctions de traitement d'\texttt{ioctl}s des pilotes
KMS, et même à toutes les fonctions faisant partie des différentes interfaces
recevant un pointeur de l'espace utilisateur.

\subsection{Autres types abstraits}

Notre approche est basée sur le fait de rendre abstrait un type de C et
d'y interdire certaines opérations : ici, on marque un certain type de pointeur
comme ``utilisateur'' et on interdit l'opérateur \texttt{*} dessus.

Le langage C n'ayant pas de types abstraits
\footnote{
  La seule abstraction possible est lorsqu'on manipule une structure par
  pointeur. Il n'est alors pas nécessaire de connaître sa définition totale.
  L'idiome est alors de placer uniquement une déclaration en avance
  (\texttt{struct s;}) dans l'en-tête (\texttt{.h}) et de renseigner
  la définition complète dans l'implantation (\texttt{.c}).
}
, on ne peut pas séparer la réprésentation d'un type (par exemple : entier signé
de 32 bits) des opérations qui y sont attachées.

Dans de nombreuses interfaces, on emploie des types entiers qui servent
d'étiquettes. Par exemple, les descripteurs de fichiers renvoyés par la fonction
\texttt{open()} et passés au fonctions \texttt{read()} et \texttt{write()} ont
le type \texttt{int}. Le langage autorise donc par exemple de multiplier entre
eux deux descripteurs de fichiers, ce qui ne correspond pas à une opération
concrète.

On peut aussi distinguer plusieurs types abstraits entre eux. Par exemple, si un
encodeur vidéo manipule des numéros de \emph{frame} et des identifiants de codec
tous les deux entiers, on peut interdire d'utiliser avec un identifiant de codec
une fonction prenant en paramètre un numéro de \emph{frame}.

% TODO check : les footnotes commençent par une majuscule


\def\partintro{}
\appendix
\renewcommand{\appendixpagename}{Annexes}
\renewcommand{\appendixtocname}{\appendixpagename}
\appendixpage*

\chapter{Module Radeon KMS}

\label{cha:code-noyau}
On inclut ici le code analysé dans le chapitre~\ref{cha:etudedecas}. On inclut à
la suite le contexte nécessaire pour comprendre ce code.

{ \small
\insertcode{radeon-buggy.c}
}


\chapter{Syntaxe et règles d'évaluation}
\label{anx:eval}

% TODO chapeaux anx

On rappelle ici pour référence la syntaxe de \langname, ainsi que sa sémantique
d'évaluation.

\section*{Syntaxe des expressions}
\vspace{-1em}

\setlength{\FrameSep}{3pt}

\begin{framed}
\figstxdata{}
\end{framed}

\section*{Syntaxe des instructions}

\begin{framed}
\figstxinstr{}
\end{framed}

\section*{Syntaxe des opérateurs}

\begin{framed}
\figstxops{}
\end{framed}

\section*{Contextes d'évaluation}

\begin{framed}
    \figctx{}
\end{framed}

\section*{Règles d'évaluation des erreurs}

\begin{framed}
\begin{mathpar}
    \ruleheader{$Ξ → Ω$}

    \semrule{Exp-Err}

    \semrule{Eval-Err}

    \semrule{Decl-Err}

    \semrule{Exp-Call-Err}
\end{mathpar}

\end{framed}

\section*{Règles d'évaluation des left-values et expressions }

\begin{framed}

\begin{mathpar}
    \ruleheader{$\mm{m}{e}{m'}{e'}$}

    \semrule{Ctx}

    \semrule{Exp-Call-Ctx}

    \ruleheader{$\mm{m}{lv}{m}{φ}$}

    \semrule{Phi-Var}

    \semrule{Phi-Deref}

    \semrule{Phi-Struct}

    \semrule{Phi-Array}

    \ruleheader{$\mm{m}{e}{m}{v}$}

    \semrule{Exp-Cst}

    \semrule{Exp-Fun}

    \semrule{Exp-Lv}

    \semrule{Exp-UnOp}

    \semrule{Exp-BinOp}

    \semrule{Exp-AddrOf}

    \semrule{Exp-Set}

    \semrule{Exp-Struct}

    \semrule{Exp-Array}

    \semrule{Exp-Call-Return}
\end{mathpar}

\end{framed}

\newpage
\section*{Règles d'évaluation des instructions, phrases et programmes}

\begin{framed}

\begin{mathpar}
    \ruleheader{$\mm{m}{i}{m}{i'}$}

    \semrule{Seq}

    \semrule{Exp}

    \semrule{Decl-Pass}

    \semrule{Decl-Return}

    \semrule{Decl-Ctx}

    \semrule{If-False}

    \semrule{If-True}

    \semrule{While}

    \semrule{Return}

    \ruleheader{$\ph{m}{p}{m'}$}

    \semrule{ET-Exp}

    \semrule{ET-Var}

    \ruleheader{$\phstar{m}{p}{m'}$}

    \semrule{T*-Nil}

    \semrule{T*-Cons}

    \ruleheader{$\phprog{P}{m}$}

    \semrule{Prog}
\end{mathpar}

\end{framed}

\newpage
\section*{Règles d'évaluation des extensions noyau}

\begin{framed}

\begin{mathpar}
  \semrule{Expr-Tainted}

  \semrule{User-Get-OK}

  \semrule{User-Get-Err}

  \semrule{User-Put-OK}

  \semrule{User-Put-Err}
\end{mathpar}

\end{framed}


\chapter{Règles de typage}
\label{anx:typage}

On rappelle ici l'ensemble des règles de typage décrites dans les
chapitres~\ref{cha:typbase} et~\ref{cha:qualifs}.

\section*{Règles de typage des constantes et left-values}

\begin{framed}

\begin{mathpar}
    \ruleheader{$Γ ⊢ c : t$}

    \disprule{Cst-Int}

    \disprule{Cst-Float}

    \disprule{Cst-Null}

    \disprule{Cst-Unit}

    \ruleheader{$Γ ⊢ lv : t$}

    \disprule{Lv-Var}

    \disprule{Lv-Deref}

    \disprule{Lv-Index}

    \disprule{Lv-Field}
\end{mathpar}
\end{framed}

\newpage
\section*{Règles de typage des opérateurs}
\begin{framed}
\begin{mathpar}
    \ruleheader{$Γ ⊢ \opun~e : t$}

    \disprule{Unop-Plus-Int}

    \disprule{Unop-Plus-Float}

    \disprule{Unop-Minus-Int}

    \disprule{Unop-Minus-Float}

    \disprule{Unop-Not}

    \ruleheader{$Γ ⊢ e_1~\opbin~e_2 : t$}

    \disprule{Op-Int}

    \disprule{Op-Float}

    \disprule{Op-Eq}

    \disprule{Ptr-Arith}

    \ruleheader{$\textsc{Eq}(t)$}

    \disprule{Eq-Num}

    \disprule{Eq-Ptr}

    \disprule{Eq-Array}

    \disprule{Eq-Struct}
\end{mathpar}
\end{framed}

\newpage
\section*{Règles de typage des expressions et instructions}
\begin{framed}
\begin{mathpar}
    \ruleheader{$Γ ⊢ e : t$}

    \disprule{Addr}

    \disprule{Struct}

    \disprule{Call}

    \disprule{Set}

    \disprule{Array}

    \disprule{Fun}

    \ruleheader{$Γ ⊢ i$}

    \disprule{Pass}

    \disprule{Seq}

    \disprule{Exp}

    \disprule{Decl}

    \disprule{If}

    \disprule{While}

    \disprule{Return}

    \ruleheader{$\typh{Γ}{p}{Γ'}$}

    \disprule{T-Exp}

    \disprule{T-Var}
\end{mathpar}

\end{framed}

\newpage
\section*{Règles de typage des valeurs}
\begin{framed}
  \ruleheader{$\semtyp{m}{v}{τ}$}

  \begin{mathpar}
    \disprule{S-Int}

    \disprule{S-Float}

    \disprule{S-Unit}

    \disprule{S-Null}

    \disprule{S-Ptr}

    \disprule{S-Array}

    \disprule{S-Struct}

    \disprule{S-Fun}
  \end{mathpar}

\end{framed}

\section*{Règles de typage des extensions noyau}
\begin{framed}

\begin{mathpar}
  \disprule{Taint}

  \disprule{GetU}

  \disprule{PutU}

  \disprule{AddrOf-Kernel}

  \disprule{Ptr-Arith-User}
\end{mathpar}

\end{framed}


\chapter{Preuves}

\wip{}

\section{Composition de lentilles}

\label{proof:compo-lens}

\begin{proof}
Pour prouver que $ℒ_1 \ggg ℒ_2 ∈ \setLens{A}{C}$, il suffit de prouver les trois
propriétés caractéristiques.

\paragraph{GetPut}%{{{

{

\def\xa{
  \lensPut{ℒ}{
    \lensGet{ℒ}{r}
    }{
      r
    }
  }


\def\xb{
  \lensPut{ℒ}
    {
      \lensGet{ℒ_2}{
        \lensGet{ℒ_1}{r}
      }
    }
    { r }
}

\def\xab{ définition de \lensGetX{ℒ} }

\def\xc{
  \lensPut{ℒ_1}
          {
            \lensPut{ℒ_2}
                    {
                      \lensGet{ℒ_2}
                        {
                          \lensGet{ℒ_1}{r}
                        }
                    }
                    { \lensGet{ℒ_1}
                              {r}
                    }
          }
          {r}
}

\def\xbc{ définition de \lensPutX{ℒ} }

\def\xd{
  \lensPut{ℒ_1}
          {
            \lensGet{ℒ_1}{r}
          }
          {r}
}

\def\xcd{ \textsc{GetPut} sur $ℒ_2$ }

\def\xe{r}

\def\xde{\textsc{GetPut} sur $ℒ_1$}

\begin{conteq}
    \xa{} \\
  = \xb{} & \xab{} \\
  = \xc{} & \xbc{} \\
  = \xd{} & \xcd{} \\
  = \xe{} & \xde{}
\end{conteq}

}

%}}}

\paragraph{PutGet}%{{{

\begin{conteq}
  \lensGet{ℒ}
    {
      \lensPut{ℒ}
        {a}
        {r}
    } \\
= \lensGet{ℒ_2}
    {
      \lensGet{ℒ_1}
        {
          \lensPut{ℒ}
            {a}
            {r}
        }
    }
& définition de \lensGetX{ℒ} \\
= \lensGet{ℒ_2}
    {
      \lensGet{ℒ_1}
        {
          \lensPut{ℒ_1}
            {
              \lensPut{ℒ_2}
                {a}
                { \lensGet{ℒ_1}{r} }
            }
            { r }
        }
    }
& définition de \lensPutX{ℒ} \\
= \lensGet{ℒ_2}
    {
      \lensPut{ℒ_2}
        {a}
        { \lensGet{ℒ_1}{r} }
    }
& \textsc{PutGet} sur $ℒ_1$ \\
= a
& \textsc{PutGet} sur $ℒ_2$ \\
\end{conteq}
%}}}

\paragraph{PutPut}%{{{

\begin{conteq}[onecolumn]
  \lensPut{ℒ}{a'}{
    \lensPut{ℒ}{a}{r}
  } \\
= \lensPut{ℒ}{a'}{
    \lensPut{ℒ_1}
      {
        \lensPut{ℒ_2}
          {a}
          { \lensGet{ℒ_1}{r} }
      }
      { r }
  }
& définition de \lensPutX{ℒ} \\
= \lensPut{ℒ_1}
    {
      \lensPut{ℒ_2}
        {a'}
        { \lensGet{ℒ_1}{
            \lensPut{ℒ_1}
              {
                \lensPut{ℒ_2}
                  {a}
                  { \lensGet{ℒ_1}{r} }
              }
              { r }
          }
        }
    }
    {
        \lensPut{ℒ_1}
          {
            \lensPut{ℒ_2}
              {a}
              { \lensGet{ℒ_1}{r} }
          }
          { r }
    }
& définition de \lensPutX{ℒ} \\
= \lensPut{ℒ_1}
    {
      \lensPut{ℒ_2}
        {a'}
        {
          \lensPut{ℒ_2}
            {a}
            { \lensGet{ℒ_1}{r} }
        }
    }
    {
        \lensPut{ℒ_1}
          {
            \lensPut{ℒ_2}
              {a}
              { \lensGet{ℒ_1}{r} }
          }
          { r }
    }
& \textsc{GetPut} sur $ℒ_1$ \\
= \lensPut{ℒ_1}
    {
      \lensPut{ℒ_2}
        {a'}
        { \lensGet{ℒ_1}{r} }
    }
    {
        \lensPut{ℒ_1}
          {
            \lensPut{ℒ_2}
              {a}
              { \lensGet{ℒ_1}{r} }
          }
          { r }
    }
& \textsc{PutPut} sur $ℒ_2$ \\
= \lensPut{ℒ_1}
    {
      \lensPut{ℒ_2}
        {a'}
        { \lensGet{ℒ_1}{r} }
    }
    { r
    }
& \textsc{PutPut} sur $ℒ_1$ \\
= \lensPut{ℒ}{a'}{r}
& définition de $\ggg$ \\
\end{conteq}
%}}}

\end{proof}

\section{Progrès}
\label{proof:progres}

On prouve de manière récursive mutuelle les
Théorèmes~\ref{thm:progres}
et~\ref{thm:progres-lv}.

% TODO mouais, à préciser mieux.

\begin{proof}

  On procède par induction sur la dérivation de $Γ ⊢ e : t$, et plus précisément
  sur la dernière règle utilisée.

  \paragraph{\textsc{Cst-Int}:} % {{{
$e$ est alors de la forme $n$, qui est une valeur.

% TODO en fait n → \widehat{n} non? idem pour float null unit
%}}}
  \paragraph{\textsc{Cst-Float}:} % {{{
$e$ est alors de la forme $d$, qui est une valeur.
%}}}
  \paragraph{\textsc{Cst-Null}:} % {{{
$e$ est alors égale à $\eNull$, qui est une valeur.
%}}}
  \paragraph{\textsc{Cst-Unit}:}%{{{
$e$ est alors égale à \eUnit, qui est une valeur.
%}}}
\paragraph{\textsc{Lv-Var}:}%{{{


Puisque $(x, t) ∈ Γ$ et $\mcomp{Γ}{m}$,
on peut appliquer le Lemme~\ref{lemma:invsem}. La construction $a =
\cLookup{x}{m}$ correspond bien à une adresse valide et on peut appliquer
\textsc{Phi-Var}.

% TODO mouais
%}}}
\paragraph{\textsc{Lv-Deref}:}%{{{

  Appliquons l'hypothèse de récurrence à $lv$ (vue en tant qu'expression).

\begin{itemize}
\item
  $lv = v$. Puisque $Γ ⊢ v : t*$, on déduit du
  Lemme~\ref{lemma:canon} que $v = φ$ ou $v = \eNull$.

  Dans le premier cas, la règle \textsc{Phi-Deref} s'applique:
  $\mms{e}{\widehat{*}φ}$.
  Dans le second, puisque $\msi{m}{*\eNull} → \serr{ptr}$, on a
  $\msi{m}{e} → \serr{ptr}$.

\item
  $\mm{m}{lv}{m'}{e'}$.
  De \textsc{Ctx} avec $C = *\ctxEmpty$, on obtient
  $\mm{m}{e}{m'}{*e'}$.

\item
  $\msi{m}{lv} → Ω$.
  En appliquant \textsc{Eval-Err} avec $C = *\ctxEmpty$, on obtient
  $\msi{m}{e} → Ω$.

\end{itemize}

% }}}
\paragraph{\textsc{Lv-Index}:} % %{{{

De même, on applique l'hypothèse de récurrence à $lv$.

\begin{itemize}
\item $lv = v$.

Comme $Γ ⊢ v : t[]$, on déduit du Lemme~\ref{lemma:canon} que
$v = \eArray{v_1; …; v_p}$.
Appliquons l'hypothèse de récurrence à $e$.

\begin{itemize}
\item $e = v'$. Puisque $Γ ⊢ e : \tInt$, on réapplique le
Lemme~\ref{lemma:canon} et $v' = n$.
D'après \textsc{Phi-Array}, $ \mms{lv[e]}{\eArray{v_1; …; v_p}\widehat{[n]}} $.
Deux cas sont à distinguer:
si $n ∈ [0;p-1]$, la partie droite vaut $v_{n+1}$ et donc
$\mm{m}{lv[e]}{m}{v_{n+1}}$.
Sinon elle vaut $\serr{array}$ et $\msi{m}{lv[e]} → \serr{array}$ par \textsc{Exp-Err}.

% TODO attention à l'off by one

\item $\mm{m}{e}{m'}{e'}$.
En appliquant \textsc{Ctx} avec $C = v[\ctxEmpty]$, on en déduit
\item $\mm{m}{lv[e]}{m'}{lv[e']}$.

\item $\msi{m}{e} → Ω$.
Avec \textsc{Eval-Err} sous ce même contexte,
$\msi{m}{lv[e]} → Ω$
\end{itemize}

\item $\mm{m}{lv}{m'}{e'}$.
On applique alors \textsc{Ctx} avec $C = \ctxEmpty[e]$, et
$\mm{m}{lv[e]}{m'}{e'[e]}$.

\item $\msi{m}{lv} → Ω$.
Toujours avec $C = \ctxEmpty[e]$, de \textsc{Eval-Err} il vient
$\msi{m}{lv[e]} → Ω$.

\end{itemize}
%}}}
\paragraph{\textsc{Lv-Field}:}%{{{

On applique l'hypothèse de récurrence du Théorème~\ref{thm:progres-lv} à $lv$.

\begin{itemize}

\item $lv = φ$
Alors \textsc{Phi-Struct} s'applique. Puisque $(l, t) ∈ S$, l'accès au champ $l$
ne provoque pas d'erreur $\serr{field}$. Donc $\mm{m}{e}{m}{φ[l]}$.

\item $\mm{m}{lv}{m'}{lv'}$
En appliquant \textsc{Ctx} avec $C = \ctxEmpty.l_S$, il vient
$\mm{m}{lv}{m'}{lv'}$.
% TODO Ctx-Lv en fait

\item $\msi{m}{lv} → Ω$
En appliquant \textsc{Eval-Err} avec $C = \ctxEmpty.l_S$, on a
$\msi{m}{lv} → Ω$.
% TODO Pas une sorte de Err-Lv ?

\end{itemize}

%}}}
  \paragraph{\textsc{Op-Int}:} % {{{

  Cela implique que $e = e_1~\opbin~e_2$. Par le Lemme~\ref{lemma:inversion}, on
  en déduit que $Γ ⊢ e_1 : \tInt$ et $Γ ⊢ e_2 : \tInt$.

  Appliquons l'hypothèse de récurrence sur $e_1$. Trois cas peuvent se produire.

\begin{itemize}

  \item $e_1 = v_1$. On a alors $\mm{m}{e_1}{m'}{v_1}$ avec $m' = m$.

    On applique l'hypothèse de récurrence à $e_2$.

      \begin{itemize}

        \item $e_2 = v_2$: alors $\mm{m'}{e_2}{m''}{v_2}$ avec $m'' = m$. On
          peut alors appliquer \textsc{Exp-BinOp}, sauf dans le cas d'une
          division par zéro ($ \opbin ∈ \{ / ; \% ; /. \} $ et
          $ v_2 = 0 $) où alors $v_1~\widehat{\opbin}~v_2 = \serr{div}$. Dans ce cas, on a
          alors par \textsc{Exp-Err} $\msi{m}{e} → \serr{div}$.

        \item $∃ (e'_2, m''), \mm{m'}{e_2}{m''}{e'_2}$.

          En appliquant \textsc{Ctx} avec $C = \ctxOp{v_1}{\ctxEmpty}$, on
          en déduit $\mm{m'}{v_1~\opbin~e_2}{m''}{v_1~\opbin~e'_2}$ soit
          $\mm{m}{e}{m''}{v_1~\opbin~e'_2}$.

        \item $\msi{m'}{e_2} → Ω$.
          De \textsc{Eval-Err} avec $C = \ctxOp{v_1}{\ctxEmpty}$
          vient alors $\msi{m}{e} → Ω$.

      \end{itemize}

  \item $∃(e_1', m'), \mm{m}{e_1}{m'}{e'_1}$.
    En appliquant \textsc{Ctx} avec $C = \ctxOp{\ctxEmpty}{e_2}$, on obtient
    $\mm{m}{e_1~\opbin~e_2}{m'}{e'_1~\opbin~e_2}$, ou
    $\mm{m}{e}{m'}{e'_1~\opbin~e_2}$.

  \item $\msi{m}{e_1} → Ω$.
    D'après \textsc{Eval-Err} avec $C = \ctxOp{\ctxEmpty}{e_2}$, on a
    $\msi{m}{e} → Ω$.

\end{itemize}

% }}}
\paragraph{\textsc{Op-Float}:} % {{{
Ce cas est similaire à \textsc{Op-Int}.
% TODO quid du lemme d'inversion (premiere ligne de Op-Int)?
%}}}
% TODO Eq-Num
% TODO Eq-Ptr
% TODO Eq-Array
% TODO Eq-Struct
\paragraph{\textsc{Op-Eq}:} %{{{
Ce cas est similaire à \textsc{Op-Int}.
% TODO quid du lemme d'inversion (premiere ligne de Op-Int)?
%}}}
\paragraph{\textsc{Unop-Plus-Int}:} % {{{

Alors $e = +~e_1$. En appliquant l'hypothèse d'induction sur $e_1$:

\begin{itemize}
\item
  soit $e_1 = v_1$. Alors en appliquant \textsc{Exp-UnOp},
  $\mm{m}{+~v_1}{m}{\widehat{+}~v_1}$, c'est-à-dire en posant $v =
  \widehat{+}~v_1$, $\mm{m}{e}{m}{v}$.
\item
  soit $∃ e'_1, m', \mm{m}{e_1}{m'}{e'_1}$. Alors en appliquant \textsc{Ctx}
avec $C = +~\ctxEmpty$, on obtient $\mm{m}{e}{m'}{e'_1}$.
\item
  soit $\msi{m}{e_1} → Ω$.
  De \textsc{Eval-Err} avec $C = +~\ctxEmpty$ il vient$\msi{m}{e} → Ω$.
\end{itemize}

% }}}
\paragraph{\textsc{Unop-Plus-Float}:} % {{{
Ce cas est similaire à \textsc{Unop-Plus-Int}.
% }}}
\paragraph{\textsc{Unop-Minus-Int}:} % {{{
Ce cas est similaire à \textsc{Unop-Plus-Int}.
% }}}
\paragraph{\textsc{Unop-Minus-Float}:} % {{{
Ce cas est similaire à \textsc{Unop-Plus-Int}.
% }}}
\paragraph{\textsc{Unop-Not}:}%{{{
Ce cas est similaire à \textsc{Unop-Plus-Int}.
%}}}
\paragraph{\textsc{Ptr-Arith}:} % %{{{

Le schéma est similaire au cas \textsc{Op-Int}. Le seul cas intéressant arrive
lorsque $e_1$ et $e_2$ sont des valeurs. D'après le Lemme~\ref{lemma:canon}:

\begin{itemize}
\item $e_1 = \eNull$ ou $e_1 = φ$
\item $e_2 = n$
\end{itemize}

D'après \textsc{Exp-Binop}, $\mms{e}{e_1~\widehat{\opbin}~n}$.

On se réfère ensuite à la définition de $\widehat{\opbin}$
(page~\pageref{page:def-arith-ptr-error}): si $e_1$ est de la forme $φ[m]$,
alors $e_1~\widehat{\opbin}~n = φ[m+n]$. Donc $\mms{e}{φ[m+n]}$.

Dans les autres cas ($e_1 = \eNull$ ou $e_1 = φ$ avec $φ$ pas de la forme
$φ'[m]$), on a $e_1~\widehat{\opbin}~n = \serr{ptr}$. Donc d'après
\textsc{Exp-Err}, $\msi{m}{e} → \serr{ptr}$.

%}}}
\paragraph{\textsc{Addr}:} % {{{

On applique l'hypothèse de récurrence du Théorème~\ref{thm:progres-lv} à $lv$.

Les cas d'évaluation et d'erreur sont traités en appliquant respectivement
\textsc{Ctx} et \textsc{Eval-Err} avec $C = \&\ctxEmpty$. Dans le cas où $lv =
φ$, on peut appliquer \textsc{Exp-AddrOf}.

%}}}
\paragraph{\textsc{Set}:} % {{{

On applique l'hypothèse de récurrence à $lv$.

\begin{itemize}

\item $lv = φ$. On applique l'hypothèse de récurrence à $e$.

    \begin{itemize}
    \item $e = v$. Alors on peut appliquer \textsc{Exp-Set}. Puisque l'état
        mémoire est modifié en $m' = m{[φ ← v]}_Φ$, il reste à montrer que
        $\mcomp{Γ}{m'}$. % TODO

    \item $\mm{m}{e}{m'}{e'}$. On conclut avec $C = φ ← \ctxEmpty$.
    \item $\msi{m}{e} → Ω$. Idem.
    \end{itemize}

\item $\mm{m}{lv}{m'}{lv'}$. On conclut avec $C = \ctxEmpty ← e$.
\item $\msi{m}{lv} → Ω$. Idem.

\end{itemize}

%}}}
\paragraph{\textsc{Array}:} % {{{

On va appliquer l'hypothèse de récurrence à $e_1$, puis si $e_1 = v_1$, on
l'applique à $e_2$, etc. Alors on se retrouve dans un des cas suivants:

\begin{itemize}
\item $∃ p ∈ [1;n],e'_p,m: e_1 = v_1, …, e_{p-1} = v_{p-1}, \mm{m}{e_p}{m'}{e'_p}$.
  Alors on peut appliquer \textsc{Ctx} avec
  $C = [v_1, …, v_{p-1}, \ctxEmpty, e_{p+1}, …, e_n]$.
\item $∃ p ∈ [1;n],Ω :     e_1 = v_1, …, e_{p-1} = v_{p-1}, \msi{m}{e_p} → Ω$.
  Dans ce cas \textsc{Eval-Err} est applicable avec ce même $C$.
\item $e_1 = v_1, …, e_n = v_n$.
  Alors on peut appliquer \textsc{Exp-Array} en construisant un tableau.
\end{itemize}

%}}}
\paragraph{\textsc{Struct}:} % {{{

Le schéma de preuve est similaire au cas \textsc{Array}.
En cas de pas d'évaluation ou d'erreur, on utilise le contexte
$C = \eStruct{l_1: v_1, …, l_{p-1}: v_{p-1}, \ctxEmpty, l_{p+1}: e_{p+1}, …, l_n:
e_n}$; et dans le cas où toutes les expressions sont évaluées, on applique
\textsc{Exp-Struct}.

%}}}
\paragraph{\textsc{Call}:} % {{{

On commence par appliquer l'hypothèse de récurrence à $e$. Dans le cas d'un pas
d'évaluation ou d'erreur, on applique respectivement \textsc{Ctx} ou
\textsc{Eval-Err} avec $C = \ctxEmpty (e_1, …, e_n)$.
Reste le cas où $e$ est une valeur: d'après le Lemme~\ref{lemma:canon}, $e$ est
de la forme $f = \eFun{\vec{a}}{i}$.

Ensuite, appliquons le même schéma que pour \textsc{Array}.
En cas de pas d'évaluation ou d'erreur, on utilise
\textsc{Ctx} ou \textsc{Eval-Err} avec
$C = f (v_1, …, v_{p-1}, \ctxEmpty, e_{p+1}, …, e_n)$.
Le seul cas restant est celui où l'expression considérée a pour forme
$f (v_1, …, v_n)$
avec
$f = \eFun{\vec{a}}{i}$
:
\textsc{Exp-Call} s'applique alors. Cette règle modifie la mémoire :

\[ \semrule{Exp-Call} \]

Il reste à prouver que $\mcomp{Γ}{m_4}$.

% TODO

%}}}
\paragraph{\textsc{Fun}:} % {{{

Ce cas est direct: la règle \textsc{Exp-Fun} s'applique.

\[ \semrule{Exp-Fun} \]

\[ \disprule{Fun} \]

%}}}
% TODO {Pass}
% TODO {Seq}
% TODO {Exp}
% TODO {Decl}
% TODO {If}
% TODO {While}
% TODO {Return}
\end{proof}

\section{Préservation}
\label{proof:preservation}

(Théorème~\ref{thm:preservation})


\begin{proof}

On procède par induction sur la forme de l'expression $e$.

\paragraph{ Constante }             $c$ % {{{ TODO
On détaille par exemple le cas d'une constante entière.

D'une part, $Γ ⊢ e : t$ donc d'après le Lemme~\ref{lemma:canon}, $t = \tInt$.

D'autre part, la seule règle d'évaluation possible est \textsc{Exp-Cst} qui
évalue en une constante entière.

% TODO check si c'est vrai aussi pour Null
% TODO ajouter un typage des valeurs?

% }}}
\paragraph{ Accès mémoire }         $ lv $ % {{{ TODO

% }}}
\paragraph{ Opération unaire }      $ \opun~e $ % {{{ TODO


% }}}
\paragraph{ Opération binaire }     $ e~\opbin~e $ % {{{ TODO

% }}}
\paragraph{ Pointeur }              $ \& lv $ % {{{ TODO

% }}}
\paragraph{ Affectation }           $ lv ← e $ % {{{ TODO

% }}}
\paragraph{ Structure }             $ \eStruct{ l_1 : e_1 ; ; l_n : e_n } $ % {{{ TODO

% }}}
\paragraph{ Tableau }               $ \eArray{e_1 ;…; e_n} $ % {{{ TODO

% }}}
\paragraph{ Fonction }              $ f $ % {{{ TODO

% }}}
\paragraph{ Appel de fonction }     $ e (e_1, …, e_n) $ % {{{ TODO

% }}}

\end{proof}

%\section*{TODO}

%\begin{itemize}
%\item
  %titres dans les chapeaux buggés
%\end{itemize}

\section{Progrès pour les types qualifiés}
\label{proof:progres-qualif}

(Théorème~\ref{thm:progres-qual})

\section{Préservation pour les types qualifiés}
\label{proof:preservation-qualif}

(Théorème~\ref{thm:preservation-qualif})

  %% TODO


\backmatter{}

%\chapterstar{Liste des figures}

\clearpage{}

\renewcommand\listfigurename{Liste des figures}
\listoffigures{}

\clearpage{}

\renewcommand{\listtheoremname}{Liste des définitions}
\listoftheorems[ignoreall,show={definition}]

\renewcommand{\listtheoremname}{Liste des théorèmes et lemmes}
\listoftheorems[ignoreall,show={theorem,lemma}]

\chapter{Références web}

\insertlinks{}

\bibliographystyle{alpha}
\bibliography{typ}

\cleartoverso
\newgeometry{margin=2cm}

\begin{headingpage}
\thispagestyle{empty}

\begin{center}\Large \textbf{Résumé}\end{center}

Les noyaux de système d'exploitation manipulent des données fournies par les
programmes utilisateur via certaines interfaces, comme les appels système. Si
celles-ci sont manipulées sans prendre une attention particulière, une faille de
sécurité connue sous le nom de \emph{Confused Deputy Problem} peut amener à des
fuites de données confidentielles ou l'elévation de privilège d'un attaquant.

Le but de cette thèse est d'utiliser des techniques de typage statique afin de
détecter les manipulations dangereuses de pointeurs contrôlés par l'espace
utilisateur.

De nombreux systèmes d'exploitation, dont le noyau Linux sur lequel nous nous
basons, sont écrits dans le langage C. Afin de pouvoir appliquer les techniques
de typage statique à ce langage, on commence par isoler un sous-langage sûr
nommé \langname, pour lequel on définit une sémantique opérationnelle et un
système de types. Les propriétés classiques de sûreté du typage sont établies.

Dans un second temps on étend \langname pour intégrer la notion de valeur
provenant de l'espace utilisateur. On montre que l'extension triviale du système
de types brise la sûreté du typage, et que l'on peut réétablir celle-ci en
donnant un type particulier aux pointeurs dont la valeur est contrôlée par
l'espace utilisateur, forçant leur déférencement à se faire au sein d'un
ensemble limité de fonctions sûres. Un cas d'étude est déroulé pour montrer que
cette technique permet de détecter un bug qui a frappé un pilote de carte
graphique AMD dans le noyau Linux.

Ces travaux ouvrent des perspectives sur deux points. D'une part, il est
possible d'appliquer cette analyse à de plus grandes parties du noyau Linux. Les
diverses interfaces qui exposent des données contrôlées par l'espace utilisateur
peuvent bénéficier de cette vérification statique. Ceci est d'autant plus vrai
que les pilotes de périphériques dont elles sont issues sont en général moins
testées que le reste du code du noyau. D'autre part, cette technique revient à
introduire des types abstraits opaques, séparant leur représentation de leur
interface. Ce concept étant absent de C, de nombreuses erreurs de programmation
notamment liées aux manipulations d'entiers peuvent également être vérifiées à
l'aide de cette méthode.

\begin{center}\Large \textbf{Abstract}\end{center}

Lorem ipsum dolor sit amet, consectetur adipiscing elit. Donec gravida, diam
ullamcorper volutpat varius, est dui blandit nibh, et posuere nisi tortor ut
lectus. Aliquam in ipsum at nisl luctus pulvinar. Nulla pulvinar accumsan
interdum. Aliquam varius est et ligula commodo varius quis gravida justo. Ut
sit amet eros orci, at iaculis metus. Cras placerat condimentum magna ut
dictum. Mauris id sollicitudin quam. Ut pretium aliquet lacus, et elementum
odio blandit quis.

Class aptent taciti sociosqu ad litora torquent per conubia nostra, per
inceptos himenaeos. Morbi tincidunt, justo sed venenatis imperdiet, diam dui
blandit nisl, et sodales urna mauris sit amet felis. Sed congue sodales
sollicitudin. Integer imperdiet tellus et ligula consequat quis commodo metus
posuere. Donec id vehicula mi. Nullam id tincidunt orci. Nunc luctus ligula ut
urna laoreet consequat. In hac habitasse platea dictumst. Sed imperdiet
ultricies diam id cursus. Quisque metus orci, tempor et dapibus eget,
ullamcorper in augue. Nullam non urna metus, id accumsan lacus. Duis erat
arcu, porttitor non laoreet vitae, ultrices vitae tortor. Pellentesque
interdum faucibus tellus eu semper. Nulla ac aliquet mauris.

In in metus purus. Nunc ullamcorper posuere turpis, vel laoreet tortor posuere
at. Ut euismod arcu sit amet dui varius vel varius neque pellentesque. Lorem
ipsum dolor sit amet, consectetur adipiscing elit. Aliquam nec mi interdum
ipsum eleifend tincidunt. Duis bibendum sem vitae risus facilisis mattis. Sed
varius eleifend nulla vel commodo. Morbi risus metus, egestas non commodo id,
ullamcorper a enim. Integer et augue et orci feugiat rhoncus a in augue.

Nulla elementum felis eu eros pretium quis porttitor felis faucibus. Vivamus a
vulputate tellus. Suspendisse potenti. Morbi dapibus, quam ac consectetur
hendrerit, lorem urna faucibus dui, in vestibulum tellus nisl sit amet ligula.
Morbi ultricies venenatis nisi sit amet scelerisque. Nam pharetra bibendum
lacinia. Suspendisse potenti. Suspendisse elementum tellus et mi hendrerit
sagittis.

Ut dignissim fermentum cursus. Praesent feugiat venenatis rutrum. Cras eget
nibh id quam ultrices pellentesque eget non elit. Maecenas justo lorem,
fermentum vitae dictum eu, egestas eu urna. In hac habitasse platea dictumst.
Cras nec arcu at purus aliquet accumsan. Nam interdum malesuada orci et
sagittis

% TODO (PLAS suit)

%Manipulating user-provided pointers in the kernel of an operating system can
%lead to security flaws if done in an incautious manner. We present an efficient
%system to detect and prevent this class of erroneous memory manipulation.

%At the core of our approach is \langname, an imperative language that we equip
%with a qualified type system, where two kinds of pointers are distinguished:
%\emph{safe} pointers, whose value is statically proved to be controlled by the
%kernel, and \emph{unsafe} ones, whose value comes from userspace through
%run-time system calls. Dereferencing unsafe pointers is forbidden in a static
%manner by the means of a strong type system.

%A concrete case study is described based on a bug that affected a video driver
%in the Linux kernel. We also explain a technique to automatically translate GNU
%C code to our core language, which will enable us to analyze larger fractions of
%the kernel in order to find similar vulnerabilities.

\end{headingpage}
\restoregeometry


% TODx bib en français

\end{document}

% vim: spelllang=fr
